\chapter*{제III장 (계속)\\[1.2ex]
\Large 연접층의 코호몰로지적 연구}

\section{국소 및 전역 Tor 함자; 퀸네트 공식}
\label{section:III.6.fr}

\subsection{서론}
\label{subsection:III.6.1.fr}

\oldpage[III]{137}

\begin{env}[6.1.1]
\label{III.6.1.1-fr}
$f:X\to Y$를 준스킴의 사상, $\mathscr F$를 준연접
$\mathscr O_X$-가군이라 하자. ``고차 직상''
$\mathrm R^n f_*(\mathscr F)$을 연구할 때 다음의 일반 문제를 고려하게
된다. ``밑변경'' 사상 $g:Y'\to Y$가 주어졌다고 하고
\[
  X'=X_{(Y')}=X\times_Y Y',\qquad
  \mathscr F'=\mathscr F\otimes_Y\mathscr O_{Y'},\qquad
  f'=f_{(Y')}:X'\to Y',
\]
로 놓자. $\mathrm R^n f_*(\mathscr F)$들이 알려져 있다고 가정하고,
고차 직상 $\mathrm R^n f'_*(\mathscr F')$에 관한 정보를 얻고자 한다.
쉽게 알 수 있듯이(예컨대 (7.7.2) 참조), 이 문제는 변하는 준연접
$\mathscr O_Y$-가군 $\mathscr G$에 대하여
$\mathrm R^n f_*(\mathscr F\otimes_{\mathscr O_Y}\mathscr G)$의 변화를
연구하는 것으로 귀착된다. 다시 말해 다음 함자를 연구해야 한다.
\[
  \mathscr G\longmapsto
  \mathrm R^n f_*(\mathscr F\otimes_{\mathscr O_Y}\mathscr G).
\]
$\mathscr F$가 $Y$ 위에서 평탄이면 함자
$\mathscr G\mapsto\mathscr F\otimes_{\mathscr O_Y}\mathscr G$는 완전하고
(\hyperref[0.6.7.4-ko]{0\textsubscript{I}, 6.7.4}), 따라서 합성함자
$\mathscr G\mapsto\mathrm R^n f_*(\mathscr F\otimes_{\mathscr O_Y}\mathscr G)$
도 여전히 코호몰로지 함자이다. 그러나 일반적인 경우에는 더 이상 그렇지
않다. 코호몰로지적 방법을 적용하기 위해서는
$\mathscr G\mapsto\mathrm R^n f_*(\mathscr F\otimes_{\mathscr O_Y}\mathscr G)$
를 언제나 코호몰로지 함자인 다른 함자들로 바꾸게 된다. 가군론의 ``Tor''
함자를 일반화하는 이 함자들은 6.3부터 6.7까지에서 정의된다. 이러한
일반화에는 하나는 ``국소'', 다른 하나는 ``전역''인 두 가지가 있으며,
6.7에서 논의할 스펙트럼 열들로 서로 연결된다. 이 스펙트럼 열들을 응용하면,
특히 어떤 조건 아래에서
\[
  \mathrm R^n(f_1\times f_2)_*(\mathscr F_1\otimes_Y\mathscr F_2)
\]
를 고차 직상 $\mathrm R^p f_{1*}(\mathscr F_1)$과
$\mathrm R^q f_{2*}(\mathscr F_2)$로 나타내는 ``퀸네트 공식''을 얻는다.
또 다른 스펙트럼 열들(6.8)은 가군의 ``Tor'' 함자에 대한 결합성 스펙트럼
열을 일반화한다. 마지막으로 밑변경 문제 자체도 스펙트럼 열들(6.9)로
이어진다.
\end{env}

\begin{env}[6.1.2]
\label{III.6.1.2-fr}
또한 (6.10)에서, 이렇게 정의한 코호몰로지 함자
$\mathscr G\mapsto\mathscr T_\bullet(\mathscr G)$가 $Y$ 위에서 국소적으로
다음 꼴임을 알게 된다.
\[
  \mathscr G\longmapsto
  \mathscr H_\bullet(\mathscr L_\bullet\otimes_Y\mathscr G).
\]
여기서 $\mathscr L_\bullet$은 국소 자유 $\mathscr O_Y$-가군들의 복합체로서
호모토피를 제외하고 정의되며, $\mathscr H_\bullet$은 호몰로지를 나타낸다.
그렇다면 $\mathscr T_\bullet$을 낳은 특수한 상황에서 벗어나 앞의 꼴인 함자
$\mathscr G\mapsto\mathscr H_\bullet(\mathscr L_\bullet\otimes_Y\mathscr G)$
를 일반적으로 연구하는 것이 유익하다.

\oldpage[III]{138}

필요한 경우 $\mathscr L_i$ 또는
$\mathscr H_i(\mathscr L_\bullet)$에 적절한 유한성 가정도 부과한다. 이것이
제7절의 대상이며, 그 내용은 본질적인 부분에서 제6절과 독립적으로 읽을 수
있다. 이 함자들의 가장 중요한 성질들은 $\mathscr T_\bullet$의 한 성분
$\mathscr T_i$의 완전성에 관한 것이다. 그러한 성질들을 확립하는 여러
판정법을 제시할 것이다. 그 응용으로서, (6.1.1)의 표기 아래 함자
$\mathscr G\mapsto\mathrm R^n f_*(\mathscr F\otimes_{\mathscr O_Y}\mathscr G)$
가 완전임을 보장하는 조건들을 얻게 된다. 이를 $\mathscr F$가 차원 $n$에서
$Y$ 위에 코호몰로지적으로 평탄하다고 표현할 것이다. 또한
$\mathscr T_\bullet(\mathscr G)=
\mathscr H_\bullet(\mathscr L_\bullet\otimes_Y\mathscr G)$의 성분
$\mathscr T_i$에서 중요한 또 하나의 성질은 함수
$y\mapsto\dim_{k(y)}(T_i(k(y)))$의 상반연속성이다. $\mathscr T_i$가
완전하면 이 성질은 연속성으로 강화되며, $Y$가 축약인 경우 그 역도
Grauert에 의해 참이다 (7.8.4).
\end{env}

\begin{env}[6.1.3]
\label{III.6.1.3-fr}
제6절과 제7절에서는 이 관점의 필요성이 뒤의 장들에서야 드러남에도,
어디서나 층 대신 층의 복합체를 인자로 삼아 초코호몰로지를 체계적으로
사용하였다. 이때 전개한 코호몰로지 형식체계는 쌍대성 형식체계까지
포함하여 연접층의 코호몰로지 함자들의 대수학을 정립할 이 논고의 한
장에서 더 투명해질 것이다. 그러나 여기에는 이 장의 범위를 벗어나는
전개가 필요하다.
\end{env}

\begin{env}[6.1.4]
\label{III.6.1.4-fr}
용어를 줄이기 위하여 아벨 범주 $C$의 두 복합체 $K^\bullet$,
$K'^\bullet$이 주어졌을 때, 복합체 사상
$f:K^\bullet\to K'^\bullet$에 대하여 합성사상 $f\circ g$와
$g\circ f$가 모두 항등사상과 호모토픽이 되게 하는 사상
$g:K'^\bullet\to K^\bullet$가 존재하면 $f$를 \emph{호모토피
동치사상}이라 하자. 용어를 남용하여 그러한 호모토피 동치사상이 존재할
때 $K^\bullet$과 $K'^\bullet$도 서로 호모토피 동치라고 할 것이다.
$C$의 복합체에 대하여 $C$에서 아벨 범주 $C'$로 가는 가법적 공변 함자
$T$의 초코호몰로지를 정의할 수 있을 때
(\hyperref[0.11.4.3-fr]{0, 11.4.3}), $C$의 복합체들의 호모토피
동치사상 $K^\bullet\to K'^\bullet$은 초코호몰로지에 대한 동형
\[
  \mathrm R^\bullet T(K^\bullet)\xrightarrow{\sim}
  \mathrm R^\bullet T(K'^\bullet)
\]
을 표준적으로 정의함은 자명하다(같은 곳).
\end{env}

\subsection{준스킴 위의 가군 복합체의 초코호몰로지}
\label{subsection:III.6.2.fr}

\begin{env}[6.2.1]
\label{III.6.2.1-fr}
$X$를 준스킴,
$\mathscr K^\bullet=(\mathscr K^i)_{i\in\mathbf Z}$를 미분 연산자의
차수가 $+1$인 $\mathscr O_X$-가군의 복합체라 하자. 모든 준스킴 사상
$f:X\to Y$와 모든 $n\in\mathbf Z$에 대하여 $\mathscr O_Y$-가군인
초코호몰로지 $\mathscr H^n(f,\mathscr K^\bullet)$를 정의했음을 상기하자
(\hyperref[0.12.4.1-fr]{0, 12.4.1}). 이것은
$\mathscr H_f^n(\mathscr K^\bullet)$ 또는
$\mathrm R^n f_*(\mathscr K^\bullet)$로도 쓴다. 초코호몰로지
$\mathscr H^\bullet(f,\mathscr K^\bullet)$는 두 스펙트럼 함자
$'\mathscr E(f,\mathscr K^\bullet)$와
$''\mathscr E(f,\mathscr K^\bullet)$의 수렴대상이며, 그
$\mathscr E_2$-항은 다음과 같다.
\begin{align}
  '\mathscr E_2^{pq}
    &=\mathscr H^p\bigl(\mathscr H^q(f,\mathscr K^\bullet)\bigr),
    \tag{6.2.1.1}\label{III.6.2.1.1.fr}\\
  ''\mathscr E_2^{pq}
    &=\mathscr H^p\bigl(f,\mathscr H^q(\mathscr K^\bullet)\bigr)
      =\mathrm R^p f_*\bigl(\mathscr H^q(\mathscr K^\bullet)\bigr).
    \tag{6.2.1.2}\label{III.6.2.1.2.fr}
\end{align}

\oldpage[III]{139}

여기서 $\mathscr H^q(f,\mathscr K^\bullet)$은 차수 $i$인 성분이
$\mathscr H^q(f,\mathscr K^i)=\mathrm R^q f_*(\mathscr K^i)$인
복합체이다(같은 곳). 또한 $Y$가 한 점으로 이루어져 있을 때 그에
대응하는 초코호몰로지는 $\mathbf H^\bullet(X,\mathscr K^\bullet)$로
쓰며, 이는 고려하는 한 점짜리 준스킴 $Y$와 무관한
$\Gamma(X,\mathscr O_X)$-가군들로 이루어져 있음을 상기하자. $Y=X$이고
$f=1_X$일 때는
$\mathscr H^n(f,\mathscr K^\bullet)=\mathscr H^n(\mathscr K^\bullet)$,
즉 복합체 $\mathscr K^\bullet$의 코호몰로지이고,
$i=i_0$이 아닌 모든 $i$에 대하여 $\mathscr K^i=0$일 때는
\[
  \mathscr H^n(f,\mathscr K^\bullet)
    =\mathrm R^{n-i_0}f_*(\mathscr K^{i_0}).
\]
스펙트럼 열 $'\mathscr E(f,\mathscr K^\bullet)$은 언제나 정칙이다.
$\mathscr K^\bullet$이 아래로 유계이면 두 스펙트럼 열은 쌍정칙이다
(\hyperref[0.12.4.1-fr]{0, 12.4.1}).

$\mathscr O_X$-가군의 복합체들의 모든 호모토피 동치사상
$h:\mathscr K^\bullet\to\mathscr K'^\bullet$은 초코호몰로지에 대한 동형
$\mathscr H^\bullet(f,\mathscr K^\bullet)\xrightarrow{\sim}
\mathscr H^\bullet(f,\mathscr K'^\bullet)$을 준다 (6.1.4).
$\mathscr H^\bullet(h):\mathscr H^\bullet(\mathscr K^\bullet)\to
\mathscr H^\bullet(\mathscr K'^\bullet)$이 동형이고
$\mathscr K^\bullet$과 $\mathscr K'^\bullet$이 아래로 유계라고만
가정해도 마찬가지이다. 실제로 이는 스펙트럼 열 (6.2.1.2)와
$\mathscr K'^\bullet$에 대한 그 유사물에
(\hyperref[0.11.1.5-fr]{0, 11.1.5})를 적용하면 곧바로 따른다. 끝으로
$X$의 모든 열린 덮개 $\mathfrak U=(U_\alpha)$에 대하여 초코호몰로지
$\mathbf H^\bullet(\mathfrak U,\mathscr K^\bullet)$를 이중복합체
$C^\bullet(\mathfrak U,\mathscr K^\bullet)$의 코호몰로지로도 정의하였다
(\hyperref[0.12.4.5-fr]{0, 12.4.5}). 이 이중복합체의 $(i,j)$-성분은
정의상 $C^i(\mathfrak U,\mathscr K^j)$이다. 가군
$\mathbf H^n(\mathfrak U,\mathscr K^\bullet)$도 역시
$\Gamma(X,\mathscr O_X)$ 위의 가군이다.
\end{env}

\begin{proposition}[6.2.2]
\label{III.6.2.2-fr}
$X$를 스킴, $\mathfrak U=(U_\alpha)$를 $X$의 아핀 열린집합들로
이루어진 덮개라 하자. 준연접 $\mathscr O_X$-가군의 모든 복합체
$\mathscr K^\bullet$에 대하여 초코호몰로지 가군
$\mathbf H^\bullet(X,\mathscr K^\bullet)$과
$\mathbf H^\bullet(\mathfrak U,\mathscr K^\bullet)$은 표준적으로
동형이다.
\end{proposition}

\begin{proof}
실제로 덮개 $\mathfrak U$의 열린집합들의 모든 유한 교집합 $V$는
아핀이다 (\hyperref[I.5.5.6-ko]{I, 5.5.6}). 따라서 모든 $i$와 모든
$q>0$에 대하여 $\mathrm H^q(V,\mathscr K^i)=0$이다
(\hyperref[III.1.3.1-fr]{1.3.1}). 그러므로 이 명제는
(\hyperref[0.12.4.7-fr]{0, 12.4.7})의 특수한 경우이다.
\end{proof}


\begin{proposition}[6.2.3]
\label{III.6.2.3-fr}
$f:X\to Y$를 준스킴의 분리 준콤팩트 사상이라 하자. 준연접
$\mathscr O_X$-가군의 모든 복합체 $\mathscr K^\bullet$에 대하여
$\mathscr O_Y$-가군 $\mathscr H^n(f,\mathscr K^\bullet)$은
준연접이다.
\end{proposition}

\begin{proof}
$\mathscr H^q(f,\mathscr K^i)=\mathrm R^q f_*(\mathscr K^i)$가 준연접
$\mathscr O_Y$-가군이므로
(\hyperref[III.1.4.10-fr]{1.4.10}), (6.2.1.1)에 따라
$'\mathscr E_2^{pq}$도 준연접 가군의 준동형의 핵을 그러한 준동형의
상으로 나눈 몫이어서 준연접이다
(\hyperref[I.4.1.1-ko]{I, 4.1.1}). 같은 이유로 첫째 스펙트럼 열의
모든 $\mathscr O_Y$-가군 $'\mathscr E_r^{pq}$,
$B_k('\mathscr E_2^{pq})$, $Z_k('\mathscr E_2^{pq})$도 준연접이다.
스펙트럼 열 $'\mathscr E(f,\mathscr K^\bullet)$의 정칙성으로부터
$Z_\infty('\mathscr E_2^{pq})$는 어떤
$Z_k('\mathscr E_2^{pq})$와 같으므로 준연접이고,
$B_\infty('\mathscr E_2^{pq})
=\varinjlim_k B_k('\mathscr E_2^{pq})$도 마찬가지로 준연접이다
(\hyperref[0.11.2.4-fr]{0, 11.2.4} 및
\hyperref[I.4.1.1-ko]{I, 4.1.1}). 따라서
$'\mathscr E_\infty^{pq}$도 준연접이다.

앞의 스펙트럼 열이 정칙이므로
$\mathscr H^n(f,\mathscr K^\bullet)$의 여과
$(F^p(\mathscr H^n(f,\mathscr K^\bullet)))_p$는 이산이고 소진적이다.
다시 말해 $\mathscr O_Y$-가군
$\mathscr H^n(f,\mathscr K^\bullet)$은 다음을 만족하는
$\mathscr O_Y$-가군의 증가열 $(\mathscr G_k)_{k\geq0}$의 합집합이다.
$\mathscr G_0=0$이고, 각 $\mathscr G_k/\mathscr G_{k-1}$은 어떤
$\mathscr O_Y$-가군 $'\mathscr E_\infty^{pq}$와 같으므로 준연접이다.
$k$에 대한 귀납법으로 $\mathscr G_k$들이 준연접임을 얻고
(\hyperref[III.1.4.17-fr]{1.4.17}),
$\mathscr H^n(f,\mathscr K^\bullet)=\varinjlim_k\mathscr G_k$이므로
명제가 증명된다
(\hyperref[I.4.1.1-ko]{I, 4.1.1}).
\end{proof}

\oldpage[III]{140}

\begin{corollary}[6.2.4]
\label{III.6.2.4-fr}
(\hyperref[III.6.2.3-fr]{6.2.3})의 가정 아래에서 $Y$의 모든 아핀
열린집합 $V$와 모든 $n\in\mathbf Z$에 대하여 표준 준동형
\[
  \mathbf H^n(f^{-1}(V),\mathscr K^\bullet)
  \longrightarrow
  \Gamma(V,\mathscr H^n(f,\mathscr K^\bullet))
  \tag{6.2.4.1}\label{III.6.2.4.1.fr}
\]
은 전단사이다.
\end{corollary}

\begin{proof}
증명은 (\hyperref[III.1.4.11-fr]{1.4.11})의 증명과 같다. 여기서
(\hyperref[III.6.2.2-fr]{6.2.2})를 사용하고,
$\mathscr F$를 $\mathscr K^\bullet$으로,
$\mathscr K^\bullet$을
$f_*(C^\bullet(\mathfrak U,\mathscr K^\bullet))$로,
$\mathscr H^\bullet(\mathscr K^\bullet)$을
$\mathscr H^\bullet(f,\mathscr K^\bullet)$으로 바꾼다. 마지막 것이
(\hyperref[III.6.2.3-fr]{6.2.3})에 따라 준연접
$\mathscr O_Y$-가군임도 유의한다.
\end{proof}

\begin{proposition}[6.2.5]
\label{III.6.2.5-fr}
$Y$를 국소 뇌터 준스킴, $f:X\to Y$를 고유 사상,
$\mathscr K^\bullet$을 $\mathscr O_X$-가군의 복합체라 하고,
$\mathscr O_X$-가군 $\mathscr H^q(\mathscr K^\bullet)$들이
연접이라고 하자. 그러면 $\mathscr O_Y$-가군
$\mathscr H^n(f,\mathscr K^\bullet)$들은 연접이다.
\end{proposition}

\begin{proof}
문제는 $Y$ 위에서 국소적이므로 $Y$가 뇌터 아핀이라고 가정해도 된다.
따라서 (\hyperref[III.6.2.4-fr]{6.2.4})에 의해
$\mathbf H^n(X,\mathscr K^\bullet)$이
$\Gamma(Y,\mathscr O_Y)$ 위의 유한 생성 가군임을 증명하면 된다.
한편
$\mathbf H^\bullet(X,\mathscr K^\bullet)
=\mathbf H^\bullet(\mathfrak U,\mathscr K^\bullet)$이다
(\hyperref[III.6.2.2-fr]{6.2.2}). $X$가 준콤팩트이므로
$\mathfrak U$는 유한이라고 가정할 수 있다. 각 복합체
$C^\bullet(\mathfrak U,\mathscr K^j)$의 코사슬은 정의상 교대적이므로
어떤 정수 $r>0$이 존재하여 $i<0$ 또는 $i>r$이면
$C^i(\mathfrak U,\mathscr K^j)=0$이다. 따라서
(\hyperref[0.11.3.3-fr]{0, 11.3.3})에 의해 이중복합체
$C^\bullet(\mathfrak U,\mathscr K^\bullet)$의 두 스펙트럼 열은
쌍정칙이다.

$\mathfrak U$에 속하는 집합들의 교집합은 아핀 열린집합이므로
(\hyperref[I.5.5.6-ko]{I, 5.5.6}), 각 함자
$\mathscr F\mapsto C^i(\mathfrak U,\mathscr F)$는 준연접
$\mathscr O_X$-가군의 범주에서 완전하다. 그러므로
$\mathrm H_I^q(C^i(\mathfrak U,\mathscr K^\bullet))
=C^i(\mathfrak U,\mathscr H^q(\mathscr K^\bullet))$이고,
$C^\bullet(\mathfrak U,\mathscr K^\bullet)$의 둘째 스펙트럼 열의
$E_2$-항은 (\hyperref[0.11.3.2-fr]{0, 11.3.2})에 따라 다음과 같다.
\[
  ''E_2^{pq}
  =\mathrm H^p(C^\bullet(\mathfrak U,\mathscr H^q(\mathscr K^\bullet)))
  =\mathbf H^p(\mathfrak U,\mathscr H^q(\mathscr K^\bullet))
  =\mathrm H^p(X,\mathscr H^q(\mathscr K^\bullet))
\]
(\hyperref[III.1.4.1-fr]{1.4.1})에 따른 마지막 가군들은 $f$가
고유이므로 유한 생성 $\Gamma(Y,\mathscr O_Y)$-가군이다
(\hyperref[III.3.2.1-fr]{3.2.1}). 스펙트럼 열
$''E(C^\bullet(\mathfrak U,\mathscr K^\bullet))$이 쌍정칙이므로
$\mathbf H^n(X,\mathscr K^\bullet)$이 유한 생성
$\Gamma(Y,\mathscr O_Y)$-가군임이 따른다
(\hyperref[0.11.1.8-fr]{0, 11.1.8}).

더욱이 $\mathscr K^\bullet$이 연접 $\mathscr O_X$-가군의 복합체이면,
$Y$와 $f$에 관한 (\hyperref[III.6.2.5-fr]{6.2.5})의 가정 아래
$\mathscr O_Y$-가군 $\mathscr H^n(f,\mathscr K^\bullet)$은 연접이다
(\hyperref[0.5.3.4-ko]{0\textsubscript{I}, 5.3.4}).
\end{proof}

\begin{env}[6.2.6]
\label{III.6.2.6-fr}
초코호몰로지 $\mathscr H^\bullet(f,\mathscr K^\bullet)$은 아래로 유계인
$\mathscr O_X$-가군의 복합체들의 범주에서 코호몰로지 함자이다
(\hyperref[0.12.4.4-fr]{0, 12.4.4}). 사상 $f$가 준콤팩트이고
$X$의 바탕공간이 국소 뇌터이면, 이는 모든 $\mathscr O_X$-가군의
복합체들의 범주에서 코호몰로지 함자이다. 실제로 이 경우 $f_*$가
귀납극한과 가환함이 (문제가 $Y$ 위에서 국소적이라는 점과 함께)
(G, II, 3.10.1)에서 따르므로
(\hyperref[0.11.5.2-ko]{0, 11.5.2})를 적용할 수 있다.

끝으로 $f$가 분리 사상이면
$\mathscr H^\bullet(f,\mathscr K^\bullet)$은 준연접
$\mathscr O_X$-가군의 복합체들의 범주에서 코호몰로지 함자이다.
$Y$가 아핀일 때 이는 즉시 성립한다. 이 경우 $X$가 스킴이므로 표준 동형
(\hyperref[III.6.2.2-fr]{6.2.2})에 의해
$\mathscr K^\bullet\mapsto
\mathbf H^\bullet(\mathfrak U,\mathscr K^\bullet)$이 준연접
$\mathscr O_X$-가군의 복합체들의 범주에서 코호몰로지 함자임을 보이는
문제로 귀착된다. 그런데 함자
$\mathscr K^\bullet\mapsto
C^\bullet(\mathfrak U,\mathscr K^\bullet)$가 이 범주에서 완전하므로
이는 자명하다 (\hyperref[I.1.3.7-ko]{I, 1.3.7}).

일반적인 경우에는 $Y$의 모든 아핀 열린집합 $V$에 대하여

\oldpage[III]{141}

$f^{-1}(V)$가 스킴이다. 따라서 앞의 결과를 적용하기 위하여 다음
사실만 확인하면 된다. 준연접 $\mathscr O_X$-가군의 복합체들의 완전열
\[
  0\longrightarrow\mathscr K'^\bullet
  \longrightarrow\mathscr K^\bullet
  \longrightarrow\mathscr K''^\bullet
  \longrightarrow0
\]
에 대하여 준동형
\[
  \partial:
  \mathscr H^n(f,\mathscr K''^\bullet)|V
  \longrightarrow
  \mathscr H^{n+1}(f,\mathscr K'^\bullet)|V
\]
은 이를 정의할 때 사용한 $f^{-1}(V)$의 아핀 열린 덮개
$\mathfrak U$에 의존하지 않는다. 실제로 $\mathfrak U'$이
$\mathfrak U$보다 더 세분된 아핀 열린 덮개이면 표준 동형들의 도식
\[
  \xymatrix@C=34pt@R=11pt{
    \mathbf H^\bullet(\mathfrak U,\mathscr K^\bullet)
      \ar[rrd]^-{\sim}\ar[dd]^{\iota}
    && \\
    &&
    \mathbf H^\bullet(f^{-1}(V),\mathscr K^\bullet)
    \\
    \mathbf H^\bullet(\mathfrak U',\mathscr K^\bullet)
      \ar[rru]_-{\sim}
    &&
  }
\]
은 가환하고, 다음 도식도 가환하기 때문이다.
\[
  \xymatrix@C=46pt@R=25pt{
    \mathbf H^n(\mathfrak U,\mathscr K''^\bullet)
      \ar[r]^-{\partial}\ar[d]^{\iota}
    &
    \mathbf H^{n+1}(\mathfrak U,\mathscr K'^\bullet)
      \ar[d]^{\iota}
    \\
    \mathbf H^n(\mathfrak U',\mathscr K''^\bullet)
      \ar[r]_-{\partial}
    &
    \mathbf H^{n+1}(\mathfrak U',\mathscr K'^\bullet).
  }
\]

앞의 조건들 가운데 하나가 성립하고
$\mathscr K^\bullet\to\mathscr K'^\bullet$이 호모토피 동치사상이면
(6.1.4), 이에 대응하는 동형
$\mathscr H^\bullet(f,\mathscr K^\bullet)\xrightarrow{\sim}
\mathscr H^\bullet(f,\mathscr K'^\bullet)$은
$\partial$-함자의 동형사상이다
(\hyperref[0.11.4.4-fr]{0, 11.4.4}).
\end{env}

\begin{env}[6.2.7]
\label{III.6.2.7-fr}
앞의 모든 내용은 미분 연산자의 차수가 $-1$인 준연접
$\mathscr O_X$-가군의 복합체 $\mathscr K_\bullet$에도 표기만 바꾸면
그대로 자연스럽게 적용된다. 모든 $i\in\mathbf Z$에 대하여
$\mathscr K^i=\mathscr K_{-i}$인 복합체
$\mathscr K^\bullet=(\mathscr K^i)$를 생각하면 충분하다.
\end{env}

\subsection{두 가군 복합체의 초Tor}
\label{subsection:III.6.3.fr}

\begin{env}[6.3.1]
\label{III.6.3.1-fr}
$A$를 가환환, $P_\bullet$, $Q_\bullet$을 미분 연산자의 차수가 $-1$인
두 $A$-가군 복합체라 하자. $L_{\bullet\bullet}$ (각각
$M_{\bullet\bullet}$)을 $P_\bullet$ (각각 $Q_\bullet$)의 사영
카르탕--아일렌베르크 분해라 하자
(\hyperref[0.11.6.1-fr]{0, 11.6.1}). 그러면 첫째 차수들의 합과 둘째
차수들의 합에 관하여
$L_{\bullet\bullet}\otimes_A M_{\bullet\bullet}$은 미분 연산자의
차수가 $-1$인 이중복합체이다. 그 호몰로지
$H_\bullet(L_{\bullet\bullet}\otimes_A M_{\bullet\bullet})$는 선택한
카르탕--아일렌베르크 분해 $L_{\bullet\bullet}$,
$M_{\bullet\bullet}$에 의존하지 않으며, 정의상 $P_\bullet$과
$Q_\bullet$에서 쌍함자 $P_\bullet\otimes_A Q_\bullet$의
초호몰로지이다 (\hyperref[0.11.6.5-fr]{0, 11.6.5}). 정의에 따라
\[
  \mathbf{Tor}_n^A(P_\bullet,Q_\bullet)
  =H_n(L_{\bullet\bullet}\otimes_A M_{\bullet\bullet})
  \tag{6.3.1.1}\label{III.6.3.1.1.fr}
\]

\oldpage[III]{142}

로 놓고, 이 $A$-가군을 두 복합체 $P_\bullet$, $Q_\bullet$의 지표
$n$인 \emph{초Tor}라 하자. 아래로 유계인 $A$-가군 복합체들의
범주에서 $\mathbf{Tor}_n^A(P_\bullet,Q_\bullet)$들은
$P_\bullet$, $Q_\bullet$에 관한 호몰로지 쌍함자를 이룸이 알려져
있다 (\hyperref[0.11.6.5-fr]{0, 11.6.5}). 더 나아가 다음이 성립한다.
\end{env}

\begin{proposition}[6.3.2]
\label{III.6.3.2-fr}
쌍함자 $\mathbf{Tor}_\bullet^A(P_\bullet,Q_\bullet)$은 두 스펙트럼
쌍함자 $'E(P_\bullet,Q_\bullet)$,
$''E(P_\bullet,Q_\bullet)$의 공통 수렴대상이며, 그 $E_2$-항은
다음과 같다.
\begin{align}
  'E^2_{pq}
    &=H_p\bigl(\operatorname{Tor}_q^A(P_\bullet,Q_\bullet)\bigr),
    \tag{6.3.2.1}\label{III.6.3.2.1.fr}\\
  ''E^2_{pq}
    &=\bigoplus_{q'+q''=q}
      \operatorname{Tor}_p^A\bigl(H_{q'}(P_\bullet),H_{q''}(Q_\bullet)\bigr).
    \tag{6.3.2.2}\label{III.6.3.2.2.fr}
\end{align}
여기서 (6.3.2.1)의
$\operatorname{Tor}_q^A(P_\bullet,Q_\bullet)$은 $A$-가군
$\operatorname{Tor}_q^A(P_i,Q_j)$들로 이루어진 이중복합체를 뜻한다.
스펙트럼 열 (6.3.2.2)는 언제나 정칙이다. $P_\bullet$과
$Q_\bullet$이 아래로 유계이거나 $A$의 코호몰로지 차원이 유한이면,
두 스펙트럼 열 (6.3.2.1)과 (6.3.2.2)는 쌍정칙이다.
\end{proposition}

\begin{proof}
이는 (\hyperref[0.11.6.5-fr]{0, 11.6.5})에서 따른다. 실제로 $A$의
코호몰로지 차원이 유한한 $n$이면 모든 $A$-가군은 길이 $n$인 사영
분해를 갖는다 (M, VI, 2.1).
\end{proof}

\begin{corollary}[6.3.3]
\label{III.6.3.3-fr}
$P'_\bullet$, $Q'_\bullet$을 두 $A$-가군 복합체,
$u:P_\bullet\to P'_\bullet$, $v:Q_\bullet\to Q'_\bullet$을 두
복합체 준동형이라 하자. 각각 $u$와 $v$에서 얻는 준동형
$H_\bullet(u):H_\bullet(P_\bullet)\to H_\bullet(P'_\bullet)$,
$H_\bullet(v):H_\bullet(Q_\bullet)\to H_\bullet(Q'_\bullet)$이
전단사이면, $u$와 $v$에서 얻는 준동형
$\mathbf{Tor}_\bullet^A(P_\bullet,Q_\bullet)\to
\mathbf{Tor}_\bullet^A(P'_\bullet,Q'_\bullet)$도 전단사이다.
\end{corollary}

\begin{proof}
실제로 $u$와 $v$에서 얻는 스펙트럼 열의 준동형
$''E(P_\bullet,Q_\bullet)\to''E(P'_\bullet,Q'_\bullet)$은 이 경우
$E_2$-항에서 동형이다. 이 스펙트럼 열들이
(\hyperref[III.6.3.2-fr]{6.3.2})에 따라 정칙이므로 결론은
(\hyperref[0.11.1.5-fr]{0, 11.1.5})에서 따른다.
\end{proof}

\begin{proposition}[6.3.4]
\label{III.6.3.4-fr}
$P_\bullet$, $Q_\bullet$을 아래로 유계인 두 $A$-가군 복합체라 하자.
$L_{\bullet\bullet}$ (각각 $M_{\bullet\bullet}$)을 평탄 $A$-가군들로
이루어진 이중복합체라 하고, 모든 $i$에 대하여
$L_{i,\bullet}$ (각각 $M_{i,\bullet}$)이 $P_i$ (각각 $Q_i$)의
분해라고 하자. 그러면 다음 표준 동형들이 성립한다.
\[
  \mathbf{Tor}_\bullet^A(P_\bullet,Q_\bullet)
  \xrightarrow{\sim}H_\bullet(L_{\bullet\bullet}\otimes_AQ_\bullet)
  \xrightarrow{\sim}H_\bullet(P_\bullet\otimes_AM_{\bullet\bullet})
  \xrightarrow{\sim}H_\bullet(L_{\bullet\bullet}\otimes_AM_{\bullet\bullet})
  \tag{6.3.4.1}\label{III.6.3.4.1.fr}
\]
\end{proposition}

\begin{proof}
이는 (\hyperref[0.11.6.5-fr]{0, 11.6.5, (ii), (iii)})와 평탄
$A$-가군의 정의에서 따른다.
\end{proof}

\begin{namedtheorem}[6.3.5]{주석}
\label{III.6.3.5-fr}
(i) (\hyperref[III.6.3.1-fr]{6.3.1})의 표기에서 이중복합체
$L_{\bullet\bullet}\otimes_A M_{\bullet\bullet}$와
$M_{\bullet\bullet}\otimes_A L_{\bullet\bullet}$는 표준적으로
동형이다. 따라서 표준 동형
$\mathbf{Tor}_\bullet^A(P_\bullet,Q_\bullet)\xrightarrow{\sim}
\mathbf{Tor}_\bullet^A(Q_\bullet,P_\bullet)$을 얻는다.

(ii) $F$와 $G$를 두 $A$-가군이라 하고, $P_\bullet$과 $Q_\bullet$을
각각 차수 0에서 $F$와 $G$이고 다른 차수에서는 0인 $A$-가군
복합체라 하자. $F$와 $G$의 사영 분해 $L_\bullet$, $M_\bullet$은
각각 0들을 덧붙여 $P_\bullet$과 $Q_\bullet$의
카르탕--아일렌베르크 분해로 볼 수 있다. 따라서 이 경우
$\mathbf{Tor}_\bullet^A(P_\bullet,Q_\bullet)
=\operatorname{Tor}_\bullet^A(F,G)$이다.
\end{namedtheorem}

\begin{proposition}[6.3.6]
\label{III.6.3.6-fr}
$(P_\bullet^\lambda)$, $(Q_\bullet^\mu)$를 $A$-가군 복합체들의 두
여과 귀납계라 하자. 그러면 표준 동형
\[
  \varinjlim_{\lambda,\mu}
  \mathbf{Tor}_\bullet^A(P_\bullet^\lambda,Q_\bullet^\mu)
  \xrightarrow{\sim}
  \mathbf{Tor}_\bullet^A
  \left(\varinjlim_\lambda P_\bullet^\lambda,
        \varinjlim_\mu Q_\bullet^\mu\right).
  \tag{6.3.6.1}\label{III.6.3.6.1.fr}
\]
이 성립한다.
\end{proposition}

\begin{proof}
$P_\bullet=\varinjlim_\lambda P_\bullet^\lambda$,
$Q_\bullet=\varinjlim_\mu Q_\bullet^\mu$라 두자. 함자성에 의하여
$\mathbf{Tor}_\bullet^A(P_\bullet^\lambda,Q_\bullet^\mu)$들은 분명히
귀납계를 이루며, 표준 사상 $P_\bullet^\lambda\to P_\bullet$,
$Q_\bullet^\mu\to Q_\bullet$에서 얻는 사상
$\mathbf{Tor}_\bullet^A(P_\bullet^\lambda,Q_\bullet^\mu)\to
\mathbf{Tor}_\bullet^A(P_\bullet,Q_\bullet)$들은 준동형의 귀납계를
이룬다.

\oldpage[III]{143}

따라서 표준 준동형 (6.3.6.1)을 얻고, 더 일반적으로 표준 준동형
\[
  \varinjlim_{\lambda,\mu}''E(P_\bullet^\lambda,Q_\bullet^\mu)
  \longrightarrow ''E(P_\bullet,Q_\bullet)
\]
을 얻는데, 그 수렴대상 사이의 준동형이 (6.3.6.1)이다. 또한 스펙트럼
열 $''E(P_\bullet,Q_\bullet)$은
(\hyperref[III.6.3.2-fr]{6.3.2})에 따라 정칙이고, 스펙트럼 열
$\varinjlim_{\lambda,\mu}''E(P_\bullet^\lambda,Q_\bullet^\mu)$도
정의 (\hyperref[0.11.1.7-fr]{0, 11.1.7})와
(\hyperref[0.11.3.3-fr]{0, 11.3.3})의 증명에 따라 정칙이다. 그러므로
(6.3.6.1)이 전단사임을 보이려면
(\hyperref[0.11.1.5-fr]{0, 11.1.5})에 따라 준동형
\[
  \varinjlim_{\lambda,\mu}''E(P_\bullet^\lambda,Q_\bullet^\mu)
  \longrightarrow ''E(P_\bullet,Q_\bullet)
  \tag{6.3.6.2}\label{III.6.3.6.2.fr}
\]
이 $E^2$-항들에 대하여 전단사임을 보이면 충분하다. 함자
$H_\bullet$이 가군 복합체들의 귀납극한과 가환하므로, 결국
$A$-가군들의 두 여과 귀납계 $(F^\lambda)$, $(G^\mu)$에 대하여 표준
준동형
\[
  \varinjlim_{\lambda,\mu}\operatorname{Tor}_\bullet^A(F^\lambda,G^\mu)
  \longrightarrow
  \operatorname{Tor}_\bullet^A
  \left(\varinjlim_\lambda F^\lambda,
        \varinjlim_\mu G^\mu\right)
\]
이 전단사임을 보이는 문제로 귀착된다. 이를 위하여 각 $F^\lambda$에
대하여 표준 자유 분해
\[
  L_\bullet^\lambda:\quad
  \cdots\longrightarrow L_{i+1}^\lambda
  \longrightarrow L_i^\lambda\longrightarrow\cdots
  \longrightarrow L_1^\lambda\longrightarrow L_0^\lambda\longrightarrow0
\]
를 생각하자. 여기서 $L_0^\lambda$는 $F^\lambda$의 원소들의 형식적
선형결합으로 이루어진 $A$-가군이고, $L_{i+1}^\lambda$는
$\operatorname{Ker}(L_i^\lambda\to L_{i-1}^\lambda)$의 원소들의
형식적 선형결합으로 이루어진 $A$-가군이다. $L_\bullet^\lambda$들이
복합체의 귀납계를 이룸은 곧 확인된다. 이제
$F=\varinjlim_\lambda F^\lambda$, $L_i=\varinjlim_\lambda L_i^\lambda$라
두면, 함자 $\varinjlim$이 완전하므로 $L_i$들은 $F$의 분해
$L_\bullet$을 이룬다. 또한 $L_i$들은 자유 $A$-가군들의 귀납극한이므로
평탄하다 (\hyperref[0.6.1.2-ko]{0\textsubscript{I}, 6.1.2}).
$G^\mu$마다 마찬가지로 표준 자유 분해 $M_\bullet^\mu$를 생각하면,
$M_\bullet=\varinjlim_\mu M_\bullet^\mu$는
$G=\varinjlim_\mu G^\mu$의 평탄 분해이다. 따라서
(\hyperref[III.6.3.5-fr]{6.3.5})와
(\hyperref[III.6.3.4-fr]{6.3.4})에 의하여
$\operatorname{Tor}_\bullet^A(\varinjlim_\lambda F^\lambda,
\varinjlim_\mu G^\mu)=H_\bullet(L_\bullet\otimes_A M_\bullet)$이다.
그런데 $H_\bullet$은 가군 복합체들의 귀납극한과 가환하므로
\[
  H_\bullet(L_\bullet\otimes_A M_\bullet)
  =\varinjlim_{\lambda,\mu}
   H_\bullet(L_\bullet^\lambda\otimes_A M_\bullet^\mu)
\]
이다. 한편
$H_\bullet(L_\bullet^\lambda\otimes_A M_\bullet^\mu)
=\operatorname{Tor}_\bullet^A(F^\lambda,G^\mu)$이므로 증명이 끝난다.

어떤 $i_0$가 존재하여 모든 $\lambda$, $\mu$와 $i<i_0$에 대하여
$P_i^\lambda=Q_i^\mu=0$이라고 가정하는 경우에도 같은 방식으로 표준
준동형
\[
  \varinjlim_{\lambda,\mu}'E(P_\bullet^\lambda,Q_\bullet^\mu)
  \longrightarrow'E(P_\bullet,Q_\bullet)
  \tag{6.3.6.3}\label{III.6.3.6.3.fr}
\]
이 전단사임을 보인다.
\end{proof}

\begin{proposition}[6.3.7]
\label{III.6.3.7-fr}
$P_\bullet$과 $Q_\bullet$이 아래로 유계라고 가정하자. 복합체
$P_\bullet$이 평탄 $A$-가군들로 이루어져 있으면, $Q_\bullet$에 관한
$\partial$-함자들의 표준 $A$-동형
\[
  \mathbf{Tor}_\bullet^A(P_\bullet,Q_\bullet)
  \xrightarrow{\sim}H_\bullet(P_\bullet\otimes_A Q_\bullet).
  \tag{6.3.7.1}\label{III.6.3.7.1.fr}
\]
이 성립한다.
\end{proposition}

\begin{proof}
실제로 스펙트럼 열 (\hyperref[III.6.3.2.1.fr]{6.3.2.1})은 쌍정칙이고
퇴화하며, 동형 (6.3.7.1)의 존재는
(\hyperref[0.11.1.6-fr]{0, 11.1.6})에서 따른다. 한편
$P_\bullet$의 사영 카르탕--아일렌베르크 분해를 써서 초Tor를 계산하면
(\hyperref[III.6.3.4-fr]{6.3.4}), 이렇게 정의한 동형이
$Q_\bullet$에 관한 $\partial$-함자들의 동형임을 곧 알 수 있다.
\end{proof}

\begin{env}[6.3.8]
\label{III.6.3.8-fr}
\oldpage[III]{144}
$\rho:A\to A'$을 환 준동형이라 하자. $\rho$에 표준적으로 결부된 차수
0의 함자적 $A$-준동형
\[
  \rho_{P_\bullet,Q_\bullet}:
  \mathbf{Tor}_\bullet^A(P_\bullet,Q_\bullet)
  \longrightarrow
  \mathbf{Tor}_\bullet^{A'}
  (P_\bullet\otimes_A A',Q_\bullet\otimes_A A').
  \tag{6.3.8.1}\label{III.6.3.8.1.fr}
\]
을 정의하고자 한다.

이를 위하여 $P_\bullet$의 사영 카르탕--아일렌베르크 분해
$L_{\bullet\bullet}$을 생각하자. 또 $P_\bullet\otimes_A A'$의 사영
카르탕--아일렌베르크 분해 $L'_{\bullet\bullet}$을 생각하자.
호모토피를 제외하고 정해지는 복합체의 $A'$-준동형
$L_{\bullet\bullet}\otimes_A A'\to L'_{\bullet\bullet}$을 정의할 수
있음을 보이겠다. 실제로 $L_{\bullet\bullet}$의 구성은 각 $i$에
대하여 $\mathrm B_i(P_\bullet)$의 사영 분해
$(X_{ij}^{\mathrm B})_{j\geqslant0}$와
$\mathrm H_i(P_\bullet)$의 사영 분해
$(X_{ij}^{\mathrm H})_{j\geqslant0}$를 임의로 주면 완전히 정해진다.
이들은 각각 $\mathrm B_i^{\mathrm I}(L_{\bullet\bullet})$ 및
$\mathrm H_i^{\mathrm I}(L_{\bullet\bullet})$와 같으며, 이로부터
차례로
\[
  \mathrm Z_i^{\mathrm I}(L_{\bullet\bullet})
  =\mathrm H_i^{\mathrm I}(L_{\bullet\bullet})
   \oplus\mathrm B_i^{\mathrm I}(L_{\bullet\bullet}),
  \qquad
  L_{i,\bullet}=\mathrm Z_i^{\mathrm I}(L_{\bullet\bullet})
   \oplus\mathrm B_{i-1}^{\mathrm I}(L_{\bullet\bullet}).
\]
을 얻는다. 그런데 $X_{i,\bullet}^{\mathrm B}\otimes_A A'$은 일반적으로
더 이상 $\mathrm B_i(P_\bullet)\otimes_A A'$의 분해가 아니지만,
여전히 사영 $A'$-가군들로 이루어진 복합체이다. 따라서 증대와
양립하고 호모토피를 제외하고 정해지는 $A'$-준동형
\[
  X_{i,\bullet}^{\mathrm B}\otimes_A A'
  \longrightarrow \mathrm B_i^{\mathrm I}(L'_{\bullet\bullet})
\]
이 있다 (M, V, 1.1). 마찬가지로 호모토피를 제외하고 정해지는
$A'$-준동형
\[
  X_{i,\bullet}^{\mathrm H}\otimes_A A'
  \longrightarrow \mathrm H_i^{\mathrm I}(L'_{\bullet\bullet})
\]
이 있다. 그러므로 위에서 상기한 구성에 의하여 모든 $i$에 대해
$A'$-준동형 $L_{i,\bullet}\otimes_A A'\to L'_{i,\bullet}$을 얻는다.
같은 구성에 의하여 이 준동형들은 ($i\in\mathbf Z$에 대하여)
미분작용소 $L_{i,\bullet}\to L_{i-1,\bullet}$ 및
$L'_{\bullet\bullet}$에 대한 유사한 미분작용소들과 양립한다. 따라서
이들은 구하고자 한 $A'$-준동형
$L_{\bullet\bullet}\otimes_A A'\to L'_{\bullet\bullet}$을 이룬다.

(6.3.8.1)을 정의하기 위하여 마찬가지로 $Q_\bullet$ (각각
$Q_\bullet\otimes_A A'$)의 사영 카르탕--아일렌베르크 분해
$M_{\bullet\bullet}$ (각각 $M'_{\bullet\bullet}$)과 $A'$-준동형
$M_{\bullet\bullet}\otimes_A A'\to M'_{\bullet\bullet}$을 생각하면
충분하다. 이 준동형들로부터 $A'$-준동형
\[
  (L_{\bullet\bullet}\otimes_A A')\otimes_{A'}
  (M_{\bullet\bullet}\otimes_A A')
  \longrightarrow
  L'_{\bullet\bullet}\otimes_{A'}M'_{\bullet\bullet},
\]
을 얻고, 이어서 합성하여 이중복합체의 $A$-준동형
\[
  L_{\bullet\bullet}\otimes_A M_{\bullet\bullet}
  \longrightarrow
  L'_{\bullet\bullet}\otimes_{A'}M'_{\bullet\bullet},
\]
을 얻는다. 호몰로지로 이행하면 (6.3.8.1)을 얻는다. 이는 호모토피를
제외하고 정해지는 복합체의 사상에서 나오므로 잘 정의된다.

$\rho':A'\to A''$을 두 번째 환 준동형이라 하고,
$\rho'':A\to A''$을 합성 준동형 $\rho'\circ\rho$라 하면, 분명히
\[
  \rho''_{P_\bullet,Q_\bullet}
  =\rho'_{P'_\bullet,Q'_\bullet}\circ\rho_{P_\bullet,Q_\bullet},
  \qquad
  P'_\bullet=P_\bullet\otimes_A A',
  \quad Q'_\bullet=Q_\bullet\otimes_A A'.
\]
끝으로 위에서 생각한 이중복합체의 사상
$L_{\bullet\bullet}\otimes_A M_{\bullet\bullet}
\to L'_{\bullet\bullet}\otimes_{A'}M'_{\bullet\bullet}$은 스펙트럼
열들의 $P_\bullet$과 $Q_\bullet$에 관한 함자적 사상
\[
  'E_{pq}^r(P_\bullet,Q_\bullet)
  \longrightarrow
  'E_{pq}^r(P_\bullet\otimes_A A',Q_\bullet\otimes_A A')
\]
과
\[
  ''E_{pq}^r(P_\bullet,Q_\bullet)
  \longrightarrow
  ''E_{pq}^r(P_\bullet\otimes_A A',Q_\bullet\otimes_A A'),
\]
을 정의함을 유의하자. 이들은 선택한 카르탕--아일렌베르크 분해와
무관하고, 앞의 추이성도 만족한다.
\end{env}

\begin{proposition}[6.3.9]
\label{III.6.3.9-fr}
$\rho:A\to A'$을 $A'$이 평탄 $A$-가군이 되게 하는 환 준동형이라
하자. 그러면 함자적인 표준 동형들
\[
  \mathbf{Tor}_\bullet^{A'}
  (P_\bullet\otimes_A A',Q_\bullet\otimes_A A')
  \xrightarrow{\sim}
  \mathbf{Tor}_\bullet^A(P_\bullet,Q_\bullet)\otimes_A A'.
  \tag{6.3.9.1}\label{III.6.3.9.1.fr}
\]
\[
  \left\{
  \begin{aligned}
  'E(P_\bullet\otimes_A A',Q_\bullet\otimes_A A')
    &\xrightarrow{\sim}'E(P_\bullet,Q_\bullet)\otimes_A A',\\
  ''E(P_\bullet\otimes_A A',Q_\bullet\otimes_A A')
    &\xrightarrow{\sim}''E(P_\bullet,Q_\bullet)\otimes_A A'.
  \end{aligned}
  \right.
  \tag{6.3.9.2}\label{III.6.3.9.2.fr}
\]
이 성립한다.
\end{proposition}

\begin{proof}
\oldpage[III]{145}
실제로 $M\otimes_A A'$이라는 $M$에 대한 함자의 완전성에 의하여,
$L_{\bullet\bullet}\otimes_A A'$과
$M_{\bullet\bullet}\otimes_A A'$은 각각
$P_\bullet\otimes_A A'$과 $Q_\bullet\otimes_A A'$의 사영
카르탕--아일렌베르크 분해이다. 따라서 결론이 따른다.
\end{proof}

\begin{env}[6.3.10]
\label{III.6.3.10-fr}
$\rho:A\to A'$을 환 준동형이라 하자. $A'$-가군 복합체
$P'_\bullet$마다 $P'_{\bullet[\rho]}$는 $A$-가군 복합체이다. 또한
동일한 대응 $P'_{\bullet[\rho]}\to P'_\bullet$은 표준 사상들의 합성
\[
  P'_{\bullet[\rho]}\longrightarrow
  P'_{\bullet[\rho]}\otimes_A A'
  \xrightarrow{\mu}P'_\bullet,
\]
으로 볼 수 있다. 여기서 $\mu$는 $A'$-준동형
$\mu(x\otimes a')=a'x$이다. $Q'_\bullet$을 두 번째 $A'$-가군
복합체라 하면, 차수 0의 함자적인 표준 준동형들
\[
  \mathbf{Tor}_\bullet^A(P'_{\bullet[\rho]},Q'_{\bullet[\rho]})
  \longrightarrow
  \mathbf{Tor}_\bullet^{A'}
  (P'_{\bullet[\rho]}\otimes_A A',Q'_{\bullet[\rho]}\otimes_A A')
  \longrightarrow
  \mathbf{Tor}_\bullet^{A'}(P'_\bullet,Q'_\bullet),
  \tag{6.3.10.1}\label{III.6.3.10.1.fr}
\]
을 얻는다. 첫째 화살표는
(\hyperref[III.6.3.8-fr]{6.3.8})에서 정의한 $A$-준동형이고, 둘째
화살표는 $A'$-준동형
$P'_{\bullet[\rho]}\otimes_A A'\to P'_\bullet$ 및
$Q'_{\bullet[\rho]}\otimes_A A'\to Q'_\bullet$로부터 함자성에 의하여
얻는다. (\hyperref[III.6.3.2-fr]{6.3.2})의 스펙트럼 열들에 대해서도
유사한 준동형들과 명백한 추이성들이 성립한다. 그 서술은 독자에게
맡긴다.
\end{env}

\begin{proposition}[6.3.11]
\label{III.6.3.11-fr}
$\rho:A\to A'$을 $A'$이 평탄 $A$-가군이 되게 하는 환 준동형이라
하자. 아래로 유계인 $A'$-가군 복합체 $P'_\bullet$과 아래로 유계인
$A$-가군 복합체 $Q_\bullet$에 대하여 함자적인 표준 동형
\[
  \mathbf{Tor}_\bullet^A(P'_{\bullet[\rho]},Q_\bullet)
  \xrightarrow{\sim}
  \mathbf{Tor}_\bullet^{A'}(P'_\bullet,Q_\bullet\otimes_A A').
  \tag{6.3.11.1}\label{III.6.3.11.1.fr}
\]
이 성립한다.
\end{proposition}

\begin{proof}
실제로 $M_{\bullet\bullet}$이 $Q_\bullet$의 사영
카르탕--아일렌베르크 분해이면,
$M_{\bullet\bullet}\otimes_A A'$은 $Q_\bullet\otimes_A A'$의 사영
카르탕--아일렌베르크 분해이고, 표준 동형을 제외하면
\[
  P'_{\bullet[\rho]}\otimes_A M_{\bullet\bullet}
  =P'_\bullet\otimes_{A'}(M_{\bullet\bullet}\otimes_A A');
\]
이다. 결론은 (\hyperref[III.6.3.4-fr]{6.3.4})에서 따른다.
\end{proof}

\begin{namedtheorem}[6.3.12]{주석}
\label{III.6.3.12-fr}
$(A^\lambda)$를 환들의 여과 귀납계라 하고,
$(P_\bullet^\lambda)$, $(Q_\bullet^\lambda)$를
$(A^\lambda)$-가군 복합체들의 두 귀납계라 하자. 그러면
(6.3.6.1)을 일반화하는 표준 동형
\[
  \varinjlim_\lambda
  \mathbf{Tor}_\bullet^{A^\lambda}
  (P_\bullet^\lambda,Q_\bullet^\lambda)
  \xrightarrow{\sim}
  \mathbf{Tor}_\bullet^A(P_\bullet,Q_\bullet),
  \tag{6.3.12.1}\label{III.6.3.12.1.fr}
\]
이 성립한다. 여기서 $A=\varinjlim A^\lambda$,
$P_\bullet=\varinjlim P_\bullet^\lambda$,
$Q_\bullet=\varinjlim Q_\bullet^\lambda$이다. $\lambda\leqslant\mu$일
때 (\hyperref[III.6.3.10-fr]{6.3.10})을 써서 준동형
\[
  \mathbf{Tor}_\bullet^{A^\lambda}
  (P_\bullet^\lambda,Q_\bullet^\lambda)
  \longrightarrow
  \mathbf{Tor}_\bullet^{A^\mu}
  (P_\bullet^\mu,Q_\bullet^\mu)
\]
을 정의하고 나면, 증명은
(\hyperref[III.6.3.6-fr]{6.3.6})의 증명과 같다.
\end{namedtheorem}

\begin{proposition}[6.3.13]
\label{III.6.3.13-fr}
$S$를 $A$의 곱셈적 부분집합, $P_\bullet$과 $Q_\bullet$을 $A$-가군
복합체라 하자. 이 복합체들에서 $S$의 원소들이 정하는 스칼라 곱셈
사상들이 전단사라고 가정하자. 따라서 $A'=S^{-1}A$라 하면
$P_\bullet$과 $Q_\bullet$은 $A'$-가군들로 이루어진다. 그러면 표준
동형
\[
  \mathbf{Tor}_\bullet^A(P_\bullet,Q_\bullet)
  \xrightarrow{\sim}
  \mathbf{Tor}_\bullet^{A'}(P_\bullet,Q_\bullet).
\]
이 성립한다.
\end{proposition}

\begin{proof}
실제로 가정에 따라 표준 준동형
$P_\bullet\to P_\bullet\otimes_A A'$,
$Q_\bullet\to Q_\bullet\otimes_A A'$은 전단사이다. 한편 초Tor의
함자성에 의하여 각 $s\in S$는
$\mathbf{Tor}_\bullet^A(P_\bullet,Q_\bullet)$에서 전단사인 스칼라 곱셈
사상을 정의한다. 따라서
\[
  \mathbf{Tor}_\bullet^A(P_\bullet,Q_\bullet)
  \longrightarrow
  \mathbf{Tor}_\bullet^A(P_\bullet,Q_\bullet)\otimes_A A'
\]

\oldpage[III]{146}

도 전단사인 준동형이다. $A'$은 평탄 $A$-가군이므로 결론은
(\hyperref[III.6.3.9-fr]{6.3.9})에서 따른다. 스펙트럼 열들에 대해서도
같은 방식으로 표준 동형들을 얻는다.
\end{proof}

\subsection{준연접 가군 복합체의 국소 초Tor 함자: 아핀 스킴의 경우}
\label{subsection:III.6.4-ko}

\begin{env}[6.4.1]
\label{III.6.4.1-fr}
$S$를 환이 $A$인 아핀 스킴, $X$, $Y$를 각각 환이 $B$, $C$인 두
아핀 $S$-스킴이라 하자. 따라서 $B$와 $C$는 $A$-대수이다. 준연접
$\mathscr O_X$-가군 (각각 준연접 $\mathscr O_Y$-가군)의 복합체
$\mathscr P_\bullet$ (각각 $\mathscr Q_\bullet$)은 모두
$\widetilde{P_\bullet}$ (각각 $\widetilde{Q_\bullet}$)의 꼴이다. 여기서
$P_\bullet$ (각각 $Q_\bullet$)은 $B$-가군 (각각 $C$-가군)의 복합체이다
(I, 1.3.7 및 1.3.8). 물론 $P_\bullet$과 $Q_\bullet$을 $A$-가군의
복합체로 보아 $\mathbf{Tor}_n^A(P_\bullet,Q_\bullet)$을 만들 수 있다.
더욱이 $\mathbf{Tor}_n^A(P_\bullet,Q_\bullet)$의 쌍함자성에 의하여
$A$-대수 $B$와 $C$가 이 $A$-가군에 작용하고, 이 작용들은 그것을
$(B,C)$-쌍가군, 같은 말로 $B\otimes_A C=A(X\times_S Y)$ 위의 가군으로
만든다. 따라서 이와 같이 준연접 $\mathscr O_{X\times_S Y}$-가군
\[
  \mathscr{T}\!or_n^{\mathscr O_S}(\mathscr P_\bullet,\mathscr Q_\bullet)
  =\bigl(\mathbf{Tor}_n^A(P_\bullet,Q_\bullet)\bigr)^{\sim}
  \tag{6.4.1.1}\label{III.6.4.1.1.fr}
\]
을 정의하였다. 이를 복합체 $\mathscr P_\bullet$과
$\mathscr Q_\bullet$의 \emph{지표 $n$인 국소 초Tor}라 하며,
$\mathscr{T}\!or_n^S(\mathscr P_\bullet,\mathscr Q_\bullet)$로도 쓴다.
\end{env}

\begin{namedtheorem}[6.4.2]{보조정리}
\label{III.6.4.2-fr}
(\hyperref[III.6.4.1-fr]{6.4.1})의 표기에서, 환 $A$가
$A=R^{-1}A'$의 꼴이라고 가정하자. 여기서 $A'$은 환이고 $R$은
$A'$의 곱셈적 부분집합이다. $S'=\operatorname{Spec}(A')$라 하자.
그러면 $X$와 $Y$를 $S'$-준스킴으로 볼 수 있고,
$X\times_{S'}Y=X\times_S Y$이다 (I, 1.6.2 및 3.2.4). 이때
\[
  \mathscr{T}\!or_\bullet^{S'}(\mathscr P_\bullet,\mathscr Q_\bullet)
  =\mathscr{T}\!or_\bullet^S(\mathscr P_\bullet,\mathscr Q_\bullet).
\]
\end{namedtheorem}

\begin{proof}
이는 공식 (\hyperref[III.6.4.1.1.fr]{6.4.1.1})과
(\hyperref[III.6.3.13-fr]{6.3.13})에서 따른다.
\end{proof}

\begin{env}[6.4.3]
\label{III.6.4.3-fr}
(\hyperref[III.6.4.1-fr]{6.4.1})의 표기와 가정 아래,
$\mathscr F=\widetilde F$를 준연접 $\mathscr O_X$-가군,
$\mathscr G=\widetilde G$를 준연접 $\mathscr O_Y$-가군이라 하자.
$\mathscr F$와 $\mathscr G$를 가군 복합체로 보아 그 지표 $n$인 초Tor를
$\mathscr{T}\!or_n^{\mathscr O_S}(\mathscr F,\mathscr G)$ 또는
$\mathscr{T}\!or_n^S(\mathscr F,\mathscr G)$로 쓴다.
(\hyperref[III.6.3.5-fr]{6.3.5 (ii)})에 의하여
\[
  \mathscr{T}\!or_n^{\mathscr O_S}(\mathscr F,\mathscr G)
  =\bigl(\operatorname{Tor}_n^A(F,G)\bigr)^{\sim}.
  \tag{6.4.3.1}\label{III.6.4.3.1.fr}
\]

이제 준연접 가군 복합체 $\mathscr P_\bullet$,
$\mathscr Q_\bullet$의 일반적인 경우로 돌아가자. 공식
(\hyperref[III.6.4.1.1.fr]{6.4.1.1})과 (6.4.3.1), 그리고 명제
(\hyperref[III.6.3.2-fr]{6.3.2})에 의하여
$\mathscr{T}\!or_\bullet^S(\mathscr P_\bullet,\mathscr Q_\bullet)$은 두
스펙트럼 열 $'\mathscr E(\mathscr P_\bullet,\mathscr Q_\bullet)$,
$''\mathscr E(\mathscr P_\bullet,\mathscr Q_\bullet)$의 수렴대상이고,
그 $E_2$-항들은 다음과 같다.
\[
  '\mathscr E_{pq}^2
  =\mathscr H_p\bigl(\mathscr{T}\!or_q^S
    (\mathscr P_\bullet,\mathscr Q_\bullet)\bigr)
  \tag{6.4.3.2}\label{III.6.4.3.2.fr}
\]
\[
  ''\mathscr E_{pq}^2
  =\bigoplus_{q'+q''=q}
  \mathscr{T}\!or_p^S
  \bigl(\mathscr H_{q'}(\mathscr P_\bullet),
        \mathscr H_{q''}(\mathscr Q_\bullet)\bigr)
  \tag{6.4.3.3}\label{III.6.4.3.3.fr}
\]
여기서 $\mathscr{T}\!or_q^S(\mathscr P_\bullet,\mathscr Q_\bullet)$은
성분이 $\mathscr{T}\!or_q^S(\mathscr P_i,\mathscr Q_j)$인 준연접
$\mathscr O_{X\times_S Y}$-가군의 이중복합체이다.
\end{env}

\begin{env}[6.4.4]
\label{III.6.4.4-fr}
이제 다른 두 아핀 스킴
$X^{(1)}=\operatorname{Spec}(B^{(1)})$,
$Y^{(1)}=\operatorname{Spec}(C^{(1)})$을 생각하자. 여기서
$B^{(1)}$과 $C^{(1)}$은 $A$-대수이다. 또 $A$-준동형
$\varphi:B\to B^{(1)}$, $\psi:C\to C^{(1)}$에 대응하는 두

\oldpage[III]{147}

$S$-사상 $u:X^{(1)}\to X$, $v:Y^{(1)}\to Y$가 주어졌다고 하자.
$\mathscr O_{X^{(1)}}$-가군의 복합체
$u^*(\mathscr P_\bullet)=(P_\bullet\otimes_B B^{(1)})^{\sim}$과
$\mathscr O_{Y^{(1)}}$-가군의 복합체
$v^*(\mathscr Q_\bullet)=(Q_\bullet\otimes_C C^{(1)})^{\sim}$을
생각하자. 표준 $A$-준동형
\[
  P_\bullet\longrightarrow P_\bullet\otimes_B B^{(1)},
  \qquad
  Q_\bullet\longrightarrow Q_\bullet\otimes_C C^{(1)}
\]
은 함자성에 의하여 $A$-준동형
\[
  \mathbf{Tor}_\bullet^A(P_\bullet,Q_\bullet)
  \longrightarrow
  \mathbf{Tor}_\bullet^A
  (P_\bullet\otimes_B B^{(1)},Q_\bullet\otimes_C C^{(1)});
\]
을 준다. 더욱이 다시 함자성에 의하여 이 준동형은 실제로
$(B\otimes_A C)$-가군의 준동형이다. 따라서 이와 같이
$(u\times_S v)$-준동형
\[
  \theta:\mathscr{T}\!or_\bullet^S(\mathscr P_\bullet,\mathscr Q_\bullet)
  \longrightarrow
  \mathscr{T}\!or_\bullet^S
  (u^*(\mathscr P_\bullet),v^*(\mathscr Q_\bullet))
  \tag{6.4.4.1}\label{III.6.4.4.1.fr}
\]
을 정의하였고, 따라서 $\mathscr O_{X^{(1)}\times_S Y^{(1)}}$-가군의
준동형
\[
  \theta^\sharp:(u\times_S v)^*
  \bigl(\mathscr{T}\!or_\bullet^S(\mathscr P_\bullet,\mathscr Q_\bullet)\bigr)
  \longrightarrow
  \mathscr{T}\!or_\bullet^S
  (u^*(\mathscr P_\bullet),v^*(\mathscr Q_\bullet))
  \tag{6.4.4.2}\label{III.6.4.4.2.fr}
\]
을 정의하였다. 이는 아래로 유계인 준연접 가군 복합체의 범주들에서
두 변수에 대한 $\partial$-함자의 사상임이 분명하다.

준동형 (6.4.4.2)는 반드시 전단사인 것은 아니다. 그러나 다음이
성립한다.
\end{env}

\begin{namedtheorem}[6.4.5]{보조정리}
\label{III.6.4.5-fr}
(\hyperref[III.6.4.4-fr]{6.4.4})의 표기에서 $u$와 $v$가 열린 몰입이라고
가정하자. 그러면 준동형 (6.4.4.2)는 전단사이다.
\end{namedtheorem}

\begin{proof}
$X^{(1)}$ (각각 $Y^{(1)}$)을 $X$ (각각 $Y$)의 열린집합과 동일시하자.
그러면 $X^{(1)}$ (각각 $Y^{(1)}$)은 $f\in B$ (각각 $g\in C$)인
$\mathrm D(f)$ (각각 $\mathrm D(g)$) 꼴의 열린집합들의 합집합이고,
그 위에 유도되는 준스킴 $\mathrm D(\varphi(f))$와 $\mathrm D(f)$
(각각 $\mathrm D(\psi(g))$와 $\mathrm D(g)$)은 동형이다. 따라서
$X^{(1)}$ (각각 $Y^{(1)}$)이 $\mathrm D(f)$ (각각 $\mathrm D(g)$)의
꼴인 경우에 보조정리를 증명하면 충분하다. 실제로 이 점이 확립되었다고
하고 일반적인 경우로 돌아가면, $\theta^\sharp$의 각 열린집합
$\mathrm D(f)\times_S\mathrm D(g)$에 대한 제한이 동형임을 보이면
충분하다. 그런데 $u_1:\mathrm D(f)\to X^{(1)}$,
$v_1:\mathrm D(g)\to Y^{(1)}$이 표준 포함사상이면 앞의 제한은 바로
$(u_1\times_S v_1)^*(\theta^\sharp)$이다. 정의
(\hyperref[III.6.4.4-fr]{6.4.4})와
($0_{\mathrm I}$, 4.4.8)에 의하여, 이를 표준 준동형
\[
  (u_1\times_S v_1)^*
  \bigl(\mathscr{T}\!or_\bullet^S
    (u^*(\mathscr P_\bullet),v^*(\mathscr Q_\bullet))\bigr)
  \longrightarrow
  \mathscr{T}\!or_\bullet^S
  ((u')^*(\mathscr P_\bullet),(v')^*(\mathscr Q_\bullet))
  \tag{6.4.5.1}\label{III.6.4.5.1.fr}
\]
과 합성하면 표준 준동형
\[
  (u'\times_S v')^*
  \bigl(\mathscr{T}\!or_\bullet^S(\mathscr P_\bullet,\mathscr Q_\bullet)\bigr)
  \longrightarrow
  \mathscr{T}\!or_\bullet^S
  ((u')^*(\mathscr P_\bullet),(v')^*(\mathscr Q_\bullet))
  \tag{6.4.5.2}\label{III.6.4.5.2.fr}
\]
을 얻음은 곧 알 수 있다. 여기서 $u'=u\circ u_1$이고
$v'=v\circ v_1$이다. (6.4.5.1)과 (6.4.5.2)가 동형임을 안다면,
$(u_1\times_S v_1)^*(\theta^\sharp)$도 동형임이 따른다.

이제 $X^{(1)}=\mathrm D(f)$이고 $Y^{(1)}=\mathrm D(g)$라고 가정하자.
따라서 $B^{(1)}=B_f$, $C^{(1)}=C_g$이다.
$u^*(\mathscr P_\bullet)$ (각각 $v^*(\mathscr Q_\bullet)$)은
$((P_\bullet)_f)^{\sim}$ (각각 $((Q_\bullet)_g)^{\sim}$)와 동일시된다.
한편 $X^{(1)}\times_S Y^{(1)}$은
$X\times_S Y=\operatorname{Spec}(B\otimes_A C)$의 열린 부분스킴
$\mathrm D(f\otimes g)$와 동일시된다 (II, 4.3.2.4). 따라서 함자성에
의하여 표준 준동형 $P_\bullet\to(P_\bullet)_f$,
$Q_\bullet\to(Q_\bullet)_g$에서 얻는 준동형
\[
  \bigl(\mathbf{Tor}_\bullet^A(P_\bullet,Q_\bullet)\bigr)_{f\otimes g}
  \longrightarrow
  \mathbf{Tor}_\bullet^A((P_\bullet)_f,(Q_\bullet)_g)
  \tag{6.4.5.3}\label{III.6.4.5.3.fr}
\]

\oldpage[III]{148}

이 전단사임을 보이면 된다. 그런데 ($0_{\mathrm I}$, 1.6.1)에 따라
$(P_\bullet)_f=\varinjlim P_\bullet^{(n)}$으로 쓸 수 있다. 여기서
$P_\bullet^{(n)}$들은 모두 $P_\bullet$과 동일한 $B$-가군 복합체이고,
$m\leqslant n$일 때 사상
$P_\bullet^{(m)}\to P_\bullet^{(n)}$은 $f^{n-m}$을 곱하는 사상이다.
$f$를 $g$로 바꾸면 $Q_\bullet$에 대해서도 유사한 결과를 얻는다.
한편 준동형 $P_\bullet^{(m)}\to P_\bullet^{(n)}$과
$Q_\bullet^{(m)}\to Q_\bullet^{(n)}$에 대응하는 준동형
\[
  \mathbf{Tor}_\bullet^A(P_\bullet^{(m)},Q_\bullet^{(m)})
  \longrightarrow
  \mathbf{Tor}_\bullet^A(P_\bullet^{(n)},Q_\bullet^{(n)})
\]
은 정의에 따라 $(f\otimes g)^{n-m}$을 곱하는 사상이다. 이제
(6.4.5.3)의 왼쪽 항에 ($0_{\mathrm I}$, 1.6.1)을 적용하고
(\hyperref[III.6.3.6-fr]{6.3.6})을 적용하면 결론이 따른다.
\end{proof}

\begin{env}[6.4.6]
\label{III.6.4.6-fr}
(\hyperref[III.6.4.4-fr]{6.4.4})의 표기에서 마찬가지로 스펙트럼
함자들의 표준 준동형
\[
  \left\{
  \begin{aligned}
  (u\times_S v)^*\bigl('\mathscr E(\mathscr P_\bullet,\mathscr Q_\bullet)\bigr)
    &\longrightarrow
      '\mathscr E(u^*(\mathscr P_\bullet),v^*(\mathscr Q_\bullet)),\\
  (u\times_S v)^*\bigl(''\mathscr E(\mathscr P_\bullet,\mathscr Q_\bullet)\bigr)
    &\longrightarrow
      ''\mathscr E(u^*(\mathscr P_\bullet),v^*(\mathscr Q_\bullet)).
  \end{aligned}
  \right.
  \tag{6.4.6.1}\label{III.6.4.6.1.fr}
\]
을 정의한다. (\hyperref[III.6.4.5-fr]{6.4.5})의 논증에 따르면 $u$와
$v$가 열린 몰입일 때 준동형 (6.4.6.1)은 전단사이다. 실제로
(6.3.6.2)와 (6.3.6.3)을 고려하면 그 논증은 $E^2$-항들에 대하여
동형임을 보이고, (\hyperref[III.6.4.5-fr]{6.4.5})는 수렴대상에 대하여
동형임을 보인다. 따라서 (\hyperref[0.11.1.2-fr]{0, 11.1.2})와
(\hyperref[0.11.2.4-fr]{0, 11.2.4})를 써서 결론을 얻는다.
\end{env}

\subsection{준연접 가군 복합체의 국소 초Tor 함자: 일반적인 경우}

\begin{env}[6.5.1]
\label{III.6.5.1-fr}
이제 임의의 준스킴 $S$와 임의의 두 $S$-준스킴 $X$, $Y$를 생각하자.
$\mathscr P_\bullet$ (각각 $\mathscr Q_\bullet$)을 준연접
$\mathscr O_X$-가군 (각각 $\mathscr O_Y$-가군)의 복합체라 하자.
$Z=X\times_S Y$라 두자. 준연접 $\mathscr O_Z$-가군
$\mathscr{T}\!or_n^S(\mathscr P_\bullet,\mathscr Q_\bullet)$을 정의하여
$\mathscr P_\bullet$과 $\mathscr Q_\bullet$의 \emph{국소 초Tor}라
부르겠다. $S$, $X$, $Y$가 아핀일 때 이는
(\hyperref[subsection:III.6.4-ko]{6.4})에서 이미 정의한 것과 일치한다.

$\mathscr P_\bullet$과 $\mathscr Q_\bullet$이 각각 차수 0의 항
$\mathscr F$와 $\mathscr G$만 남고 다른 항들은 0인 경우에는
$\mathscr{T}\!or_n^S(\mathscr P_\bullet,\mathscr Q_\bullet)$ 대신
$\mathscr{T}\!or_n^S(\mathscr F,\mathscr G)$라고 쓰겠다.
\end{env}

\begin{env}[6.5.2]
\label{III.6.5.2-fr}
먼저 $S$가 아핀이라고 가정하고, $(X_\lambda)$, $(Y_\mu)$를 각각
$X$, $Y$의 아핀 열린집합들로 이루어진 덮개라 하자. 그러면
$Z_{\lambda\mu}=X_\lambda\times_S Y_\mu$들은 $Z$의 아핀 열린 덮개를
이룬다. $\mathscr P_{\lambda\bullet}=\mathscr P_\bullet|X_\lambda$,
$\mathscr Q_{\mu\bullet}=\mathscr Q_\bullet|Y_\mu$라 두자. 따라서 각
쌍 $(\lambda,\mu)$에 대하여 준연접 $\mathscr O_{Z_{\lambda\mu}}$-가군
\[
  \mathscr F_{\lambda\mu}
  =\mathscr{T}\!or_n^S
   (\mathscr P_{\lambda\bullet},\mathscr Q_{\mu\bullet}),
\]
을 얻는다. $\mathscr F_{\lambda\mu}$들이 붙이기 조건
($0_{\mathrm I}$, 3.3.1)을 만족함을 보여야 한다. 이를 위해서는 임의의
아핀 열린집합 $U\subset X_\lambda\cap X_{\lambda'}$ (각각
$V\subset Y_\mu\cap Y_{\mu'}$)에 대하여 $\mathscr F_{\lambda\mu}$와
$\mathscr F_{\lambda'\mu'}$의 $U\times_S V$에 대한 제한들이 표준적으로
동형임을 확인하면 충분하다. 그런데 이는 두 제한이 모두
$\mathscr{T}\!or_n^S(\mathscr P_\bullet|U,\mathscr Q_\bullet|V)$와
표준적으로 동형이라는 사실에서 곧 따른다
(\hyperref[III.6.4.5-fr]{6.4.5}). 더 나아가 이 정의와
(\hyperref[III.6.4.5-fr]{6.4.5})에 따라 이렇게 정의한
$\mathscr O_Z$-가군은 동형을 제외하면

\oldpage[III]{149}

생각한 열린 덮개들에 의존하지 않는다. 따라서 이를
$\mathscr{T}\!or_n^S(\mathscr P_\bullet,\mathscr Q_\bullet)$로 쓴다.
끝으로 (\hyperref[III.6.4.5-fr]{6.4.5})에 의하여 $X$ (각각 $Y$)의
임의의 열린부분 $U$ (각각 $V$)에 대하여
$\mathscr{T}\!or_n^S(\mathscr P_\bullet,\mathscr Q_\bullet)$의
$U\times_S V$에 대한 제한은
$\mathscr{T}\!or_n^S(\mathscr P_\bullet|U,\mathscr Q_\bullet|V)$와
표준적으로 동형이다.
\end{env}

\begin{env}[6.5.3]
\label{III.6.5.3-fr}
이제 $S$가 임의인 일반적인 경우로 넘어가자. $(S_\alpha)$를 $S$의
아핀 열린 덮개라 하고, $X_\alpha$ (각각 $Y_\alpha$)를 $X$ (각각
$Y$)에서 $S_\alpha$의 역상이라 하자. 다시 층
\[
  \mathscr{T}\!or_n^{S_\alpha}
  (\mathscr P_\bullet|X_\alpha,\mathscr Q_\bullet|Y_\alpha)
  =\mathscr G_\alpha
\]
들이 붙이기 조건을 만족함을 보여야 한다. 이를 위해서는
$S_\alpha\cap S_\beta$에 포함되는 임의의 아핀 열린집합 $T$에 대하여
$\mathscr G_\alpha$와 $\mathscr G_\beta$의 $U\times_S V$에 대한 제한을
$\mathscr{T}\!or_n^T(\mathscr P_\bullet|U,\mathscr Q_\bullet|V)$ 위에서
표준적으로 동일시하면 충분하다. 여기서 $U$, $V$는 각각 $X$, $Y$에서
$T$의 역상이다. 또한 $T$가 동시에 $\mathrm D(f_\alpha)$와
$\mathrm D(f_\beta)$로 쓰이는 경우로 한정해도 된다. 여기서
$f_\alpha$ (각각 $f_\beta$)는 $S_\alpha$ (각각 $S_\beta$) 위의
$\mathscr O_S$의 단면이다. 그런데
$\mathscr{T}\!or_n^T(\mathscr P_\bullet|U,\mathscr Q_\bullet|V)$는
(\hyperref[III.6.4.2-fr]{6.4.2})에 의하여 한편으로
$\mathscr{T}\!or_n^{S_\alpha}(\mathscr P_\bullet|U,
\mathscr Q_\bullet|V)$와, 다른 한편으로
$\mathscr{T}\!or_n^{S_\beta}(\mathscr P_\bullet|U,
\mathscr Q_\bullet|V)$와 표준적으로 동형이다. 또 방금
(\hyperref[III.6.5.2-fr]{6.5.2})에서 $\mathscr G_\alpha$와
$\mathscr G_\beta$를 각각 이 두 가군과 표준적으로 동일시하였다.
이로써 $\mathscr O_Z$-가군
$\mathscr{T}\!or_n^S(\mathscr P_\bullet,\mathscr Q_\bullet)$의 정의가
완성된다. 더 나아가 $X$ (각각 $Y$)의 임의의 열린부분 $U$ (각각
$V$)에 대하여
$\mathscr{T}\!or_n^S(\mathscr P_\bullet|U,\mathscr Q_\bullet|V)$는
$\mathscr{T}\!or_n^S(\mathscr P_\bullet,\mathscr Q_\bullet)$의
$U\times_S V$에 대한 제한과 표준적으로 동형이다.

이와 같이 아래로 유계인 준연접 가군 복합체의 범주들에서
$\mathscr O_Z$-가군의 범주에 값을 갖는 두 변수에 대한 $\partial$-함자
$\mathscr{T}\!or_\bullet^S(\mathscr P_\bullet,\mathscr Q_\bullet)$을
정의했음은 분명하다. 실제로
(\hyperref[III.6.4.5-fr]{6.4.5})와 (6.4.4.2)가 두 변수에 대한
$\partial$-함자의 사상이라는 주석에 의하여 이 문제는 $X$, $Y$, $S$
위에서 국소적임이 분명하다. $\mathscr P_\bullet$과
$\mathscr Q_\bullet$이 각각 차수 0의 항 $\mathscr F$, $\mathscr G$만
남는 경우에는, (\hyperref[III.6.4.1.1.fr]{6.4.1.1})에 의하여
$\mathscr{T}\!or_0^S(\mathscr F,\mathscr G)$는
(I, 9.1.2)에서 정의한 외부 텐서곱
$\mathscr F\otimes_S\mathscr G$와 다르지 않다. 이는 실제로
(I, 9.1.3)에서 따른다.
\end{env}

\begin{env}[6.5.4]
\label{III.6.5.4-fr}
앞의 구성과 (\hyperref[III.6.4.6-fr]{6.4.6})의 주석에 의하여
$\mathscr{T}\!or_\bullet^S(\mathscr P_\bullet,\mathscr Q_\bullet)$은 두
스펙트럼 함자 $'\mathscr E(\mathscr P_\bullet,\mathscr Q_\bullet)$,
$''\mathscr E(\mathscr P_\bullet,\mathscr Q_\bullet)$의 수렴대상이고,
그 $E_2$-항들은 다음과 같다.
\[
  '\mathscr E_{pq}^2
  =\mathscr H_p\bigl(\mathscr{T}\!or_q^S
    (\mathscr P_\bullet,\mathscr Q_\bullet)\bigr)
  \tag{6.5.4.1}\label{III.6.5.4.1.fr}
\]
\[
  ''\mathscr E_{pq}^2
  =\bigoplus_{q'+q''=q}
   \mathscr{T}\!or_p^S
   \bigl(\mathscr H_{q'}(\mathscr P_\bullet),
         \mathscr H_{q''}(\mathscr Q_\bullet)\bigr).
  \tag{6.5.4.2}\label{III.6.5.4.2.fr}
\]
스펙트럼 열 (6.5.4.2)는 항상 정칙이다. $\mathscr P_\bullet$과
$\mathscr Q_\bullet$이 아래로 유계이면 두 스펙트럼 열은 쌍정칙이다.
앞의 두 스펙트럼 열이 쌍정칙인 또 다른 경우는 다음과 같다.
\end{env}

\begin{env}[6.5.5]
\label{III.6.5.5-fr}
위상공간 $T$ 위의 환층 $\mathscr A$가 \emph{코호몰로지 차원
$\leqslant n$}이라는 것은 모든 $t\in T$에 대하여 환 $\mathscr A_t$의
코호몰로지 차원이 $\leqslant n$이라는 뜻이다. 이때 환 달린 공간
$(T,\mathscr A)$도 \emph{코호몰로지 차원 $\leqslant n$}이라고 한다.
어떤 정수 $n$에 대하여 코호몰로지 차원이 $\leqslant n$인 환층
(각각 환 달린 공간)을 코호몰로지 차원이 \emph{유한}하다고 한다.
$\mathscr A_t$들이 뇌터 (가환) 국소환인 경우, 그 코호몰로지 차원이
$\leqslant n$이라는 것은 그 환들이 정칙이고 (크룰) 차원이
$\leqslant n$이라는 뜻임을 유의하자

\oldpage[III]{150}

($0_{\mathrm{IV}}$, 17.3.1). 제IV장에서 도입할 차원론의 용어로는,
국소 뇌터 준스킴 $T$의 코호몰로지 차원이 $\leqslant n$이라는 것과
$T$가 정칙이고 ($0_{\mathrm I}$, 4.1.4) 차원이 $\leqslant n$이라는
것은 동치이다. 이는 $T$의 모든 아핀 열린집합 $U$에 대하여 환
$\Gamma(U,\mathscr O_T)$의 코호몰로지 차원이 $\leqslant n$이라는
뜻이다 ($0_{\mathrm{IV}}$, 17.2.6). 이제 이 마지막 주석과
(\hyperref[III.6.3.2-fr]{6.3.2})에 의하여 $S$가 국소 뇌터이고
코호몰로지 차원이 유한이면 스펙트럼 열
$'\mathscr E(\mathscr P_\bullet,\mathscr Q_\bullet)$과
$''\mathscr E(\mathscr P_\bullet,\mathscr Q_\bullet)$은 쌍정칙이다.

$X\times_S Y$에서 $Y\times_S X$로 가는 표준 동형에 의하여
$\mathscr{T}\!or_\bullet^S(\mathscr P_\bullet,\mathscr Q_\bullet)$은
동형을 제외하면
$\mathscr{T}\!or_\bullet^S(\mathscr Q_\bullet,\mathscr P_\bullet)$로
변환됨이 분명하다.
\end{env}

\begin{proposition}[6.5.6]
\label{III.6.5.6-fr}
$(\mathscr P_{\alpha\bullet})$를 준연접 $\mathscr O_X$-가군 복합체의
여과 귀납계라 하자. 그러면 표준 동형
\[
  \varinjlim_\alpha
  \mathscr{T}\!or_\bullet^S
  (\mathscr P_{\alpha\bullet},\mathscr Q_\bullet)
  \xrightarrow{\sim}
  \mathscr{T}\!or_\bullet^S
  \left(\varinjlim_\alpha\mathscr P_{\alpha\bullet},\mathscr Q_\bullet\right).
  \tag{6.5.6.1}\label{III.6.5.6.1.fr}
\]
이 성립한다.
\end{proposition}

\begin{proof}
문제는 $S$, $X$, $Y$ 위에서 국소적이므로 이들이 아핀이라고 가정할
수 있고, 그러면 명제는 (\hyperref[III.6.3.6-fr]{6.3.6})으로 환원된다.
\end{proof}

\begin{namedtheorem}[6.5.7]{주석}
\label{III.6.5.7-fr}
\textnormal{(i)} 특히 $S=X=Y$인 경우를 생각하자. 이때
$\mathscr P_\bullet$과 $\mathscr Q_\bullet$은 준연접
$\mathscr O_S$-가군의 두 복합체이다. 그러면
$\mathscr{T}\!or_n^S(\mathscr P_\bullet,\mathscr Q_\bullet)$은 준연접
$\mathscr O_S$-가군이다. 더 나아가 각 점 $z\in S$에 대하여
(\hyperref[III.6.5.6-fr]{6.5.6})에 의해 표준 동형
\[
  \bigl(\mathscr{T}\!or_n^S
    (\mathscr P_\bullet,\mathscr Q_\bullet)\bigr)_z
  \xrightarrow{\sim}
  \mathbf{Tor}_n^{\mathscr O_z}
  \bigl((\mathscr P_\bullet)_z,(\mathscr Q_\bullet)_z\bigr)
  \tag{6.5.7.1}\label{III.6.5.7.1.fr}
\]
이 성립한다. 실제로 문제는 국소적이고,
(\hyperref[III.6.4.1.1.fr]{6.4.1.1})에 의하여 가군의 경우로 환원된다.

\textnormal{(ii)} 같은 환 달린 공간 $(X,\mathscr O_X)$ 위의 두
$\mathscr O_X$-가군 복합체 $\mathscr P_\bullet$,
$\mathscr Q_\bullet$의 경우로 초Tor의 정의를 일반화할 수 있다. 실제로
$X$의 각 열린집합 $U$에 대하여
$A(U)=\Gamma(U,\mathscr O_X)$,
$P_\bullet(U)=\Gamma(U,\mathscr P_\bullet)$,
$Q_\bullet(U)=\Gamma(U,\mathscr Q_\bullet)$라 두자. 그러면 $A(U)$-가군
\[
  \mathbf{Tor}_n^{A(U)}(P_\bullet(U),Q_\bullet(U))
\]
들은 $X$ 위의 준층을 이루고, 이 준층을 층화하여 얻는
$\mathscr O_X$-가군을
$\mathscr{T}\!or_n^X(\mathscr P_\bullet,\mathscr Q_\bullet)$로 쓴다.
$X$가 준스킴이면 (\hyperref[III.6.3.12-fr]{6.3.12})에 의하여 이
$\mathscr O_X$-가군은 위에서 정의한 초Tor와 표준적으로 동형이다.
이 일반화는 더 전개하지 않겠다.
\end{namedtheorem}

\begin{proposition}[6.5.8]
\label{III.6.5.8-fr}
$X$, $Y$를 두 $S$-준스킴, $\mathscr F$ (각각 $\mathscr G$)를 준연접
$\mathscr O_X$-가군 (각각 $\mathscr O_Y$-가군)이라 하자.
$\mathscr F$ 또는 $\mathscr G$가 $S$-평탄이면 $n\neq0$에 대하여
$\mathscr{T}\!or_n^S(\mathscr F,\mathscr G)=0$이다.
\end{proposition}

\begin{proof}
문제는 $X$, $Y$ 위에서 국소적이므로 $X$, $Y$, $S$가 각각 환이
$B$, $C$, $A$인 아핀이라고 가정할 수 있다. 또
$\mathscr F=\widetilde M$, $\mathscr G=\widetilde N$이라 하자. 여기서
$M$ (각각 $N$)은 $B$-가군 (각각 $C$-가군)이다. 예를 들어
$\mathscr F$가 $S$-평탄이라고 하자. 이는 모든 $s\in S$에 대하여
$M_s$가 평탄 $A_s$-가군이라는 뜻이다 ($0_{\mathrm I}$, 6.7.1).
따라서 $M$은 평탄 $A$-가군이고 ($0_{\mathrm I}$, 6.3.3), 모든
$C$-가군 $N$과 $n>0$에 대하여
$\operatorname{Tor}_n^A(M,N)=0$임을 안다 ($0_{\mathrm I}$, 6.1.1).
그러므로 (\hyperref[III.6.4.1.1.fr]{6.4.1.1})에서 결론이 따른다.
\end{proof}

\begin{namedtheorem}[6.5.9]{따름정리}
\label{III.6.5.9-fr}
$X$, $Y$를 두 $S$-준스킴, $\mathscr P_\bullet$ (각각
$\mathscr Q_\bullet$)을 아래로 유계인 준연접 $\mathscr O_X$-가군

\oldpage[III]{151}

(각각 $\mathscr O_Y$-가군)의 복합체라 하자. 모든
$\mathscr P_i$가 $S$-평탄이라고 가정하자. 그러면
$\mathscr Q_\bullet$에 관한 $\partial$-함자들의 표준 동형
\[
  \mathscr{T}\!or_\bullet^S(\mathscr P_\bullet,\mathscr Q_\bullet)
  \xrightarrow{\sim}
  \mathscr H_\bullet(\mathscr P_\bullet\otimes_S\mathscr Q_\bullet).
  \tag{6.5.9.1}\label{III.6.5.9.1.fr}
\]
이 성립한다.
\end{namedtheorem}

\begin{proof}
$S$, $X$, $Y$가 아핀일 때 이는 바로
(\hyperref[III.6.3.7-fr]{6.3.7})이다. 일반적인 경우는
(\hyperref[III.6.5.2-fr]{6.5.2})와
(\hyperref[III.6.5.3-fr]{6.5.3})의 논증으로 이행한다.
\end{proof}

\begin{namedtheorem}[6.5.10]{따름정리}
\label{III.6.5.10-fr}
$X$가 $S$ 위에서 평탄하고 ($0_{\mathrm I}$, 6.7.1),
$\mathscr P_\bullet$과 $\mathscr Q_\bullet$이 아래로 유계이며, 모든
$\mathscr P_i$가 국소 자유 $\mathscr O_X$-가군이라고 가정하자. 이들은
유한형일 필요가 없다. 그러면 준동형 (6.5.9.1)은 전단사이다.
\end{namedtheorem}

\begin{proof}
실제로 평탄성은 정의에 따라 $X$ 위에서 점별 성질이고, 평탄 가군의
임의의 직합은 평탄하므로 ($0_{\mathrm I}$, 6.1.2),
(\hyperref[III.6.5.9-fr]{6.5.9})의 가정이 만족된다.
\end{proof}

\begin{proposition}[6.5.11]
\label{III.6.5.11-fr}
$X'$, $Y'$을 두 $S$-준스킴, $f:X\to X'$, $g:Y\to Y'$을 두 아핀
$S$-사상이라 하자. $\mathscr P_\bullet$ (각각
$\mathscr Q_\bullet$)을 준연접 $\mathscr O_X$-가군 (각각
$\mathscr O_Y$-가군)의 복합체라 하자. 그러면 함자적인 표준 동형
\[
  (f\times_S g)_*
  \bigl(\mathscr{T}\!or_\bullet^S
    (\mathscr P_\bullet,\mathscr Q_\bullet)\bigr)
  \xrightarrow{\sim}
  \mathscr{T}\!or_\bullet^S
  (f_*(\mathscr P_\bullet),g_*(\mathscr Q_\bullet)).
  \tag{6.5.11.1}\label{III.6.5.11.1.fr}
\]
이 성립한다.
\end{proposition}

\begin{proof}
$f$와 $g$가 아핀이므로 $f_*(\mathscr P_\bullet)$과
$g_*(\mathscr Q_\bullet)$은 준연접 가군의 복합체이다 (II, 1.2.6).
$Z'=X'\times_S Y'$라 두면 (6.5.11.1)의 두 항은 준연접
$\mathscr O_{Z'}$-가군이다 (\hyperref[III.6.5.1-fr]{6.5.1}). 문제는
$S$, $X'$, $Y'$이 아핀인 경우로 쉽게 환원된다. 그러면 가정에 따라
$X$와 $Y$도 아핀이고,
(\hyperref[III.6.4.1.1.fr]{6.4.1.1})과 (I, 1.6.3)에서 곧 결론이
따른다.
\end{proof}

\begin{namedtheorem}[6.5.12]{주석}
\label{III.6.5.12-fr}
$X'$, $Y'$을 두 $S$-준스킴이라 하고 $X'$이 $S$-평탄이라고 가정하자.
$X$를 $X'$의 닫힌 부분준스킴, $i:X\to X'$을 표준 단사사상이라 하자.
$\mathscr P_\bullet$ (각각 $\mathscr Q_\bullet$)을 아래로 유계인
준연접 $\mathscr O_X$-가군 (각각 $\mathscr O_{Y'}$-가군)의 복합체라
하자. 끝으로 $\mathscr L'_{\bullet\bullet}$을
$i_*(\mathscr P_\bullet)$의 분해로서 국소 자유 $\mathscr O_{X'}$-가군들로
이루어지고 다음 성질을 만족하는 것이라 하자. $X'$의 모든 점은 모든
$j$에 대하여 $\mathscr L'_{j,\bullet}|U$가
$i_*(\mathscr P_j)|U$의 자유 분해가 되는 아핀 열린근방 $U$를 갖는다.
그러면 표준 동형
\[
  (i\times_S 1)_*
  \bigl(\mathscr{T}\!or_\bullet^S
    (\mathscr P_\bullet,\mathscr Q_\bullet)\bigr)
  \xrightarrow{\sim}
  \mathscr H_\bullet
  (\mathscr L'_{\bullet\bullet}\otimes_S\mathscr Q_\bullet).
  \tag{6.5.12.1}\label{III.6.5.12.1.fr}
\]
이 성립한다.

$S$, $X'$, $Y'$이 아핀이고 모든 $j$에 대하여
$\mathscr L'_{j,\bullet}$이 $i_*(\mathscr P_j)$의 자유 분해이면,
(\hyperref[III.6.5.11-fr]{6.5.11})에 의하여 $X'=X$가 $S$-평탄인
경우로 환원되고, (\hyperref[III.6.3.4-fr]{6.3.4})를 적용하면 충분하다.
일반적인 경우에는 동형 (6.5.12.1)을 국소적으로 정의하고, 이 정의가
실제로 전역 동형을 주는지 확인해야 한다. 이를 위하여 첫째 동형
(6.3.4.1)의 정의로 돌아가야 한다. 그 동형은
(\hyperref[0.11.6.5-fr]{0, 11.6.5})와
(\hyperref[0.11.5.3-ko]{0, 11.5.3})의 스펙트럼 열 동형에서 나온다.
그 동형은 다시 이중복합체의 사상
$L_{\bullet\bullet}\to L''_{\bullet\bullet}$에서 얻는다. 여기서
$L_{\bullet\bullet}$은 $P_\bullet$의 주어진 분해이고,
$L''_{\bullet\bullet}$은 아래로 유계인 $A$-가군 복합체의 범주에서
$P_\bullet$의 사영 분해이다

\oldpage[III]{152}

(\hyperref[0.11.5.2.2-ko]{0, 11.5.2.2} 참조). 두 그러한 분해 사이에는
호모토피 동치사상이 존재하므로 (M, V, 1.2), 동형 (6.3.4.1)은 선택한
사영 분해 $L''_{\bullet\bullet}$에 의존하지 않는다. 이로부터 우리의
주장이 따른다.
\end{namedtheorem}

\begin{proposition}[6.5.13]
\label{III.6.5.13-fr}
$X$, $Y$를 두 $S$-준스킴이라 하고, 다음 조건 가운데 하나가
성립한다고 가정하자.
\begin{enumerate}
  \item[(i)] $X$와 $Z=X\times_S Y$는 국소 뇌터이고 $X$는 $S$ 위에서
    평탄하다.
  \item[(ii)] $S$와 $X$는 국소 뇌터이고 $Y$는 $S$ 위에서 유한형이다.
\end{enumerate}
$\mathscr P_\bullet$ (각각 $\mathscr Q_\bullet$)을 아래로 유계인
준연접 $\mathscr O_X$-가군 (각각 $\mathscr O_Y$-가군)의 복합체라
하자. 또한 모든 $n$에 대하여
$\mathscr H_n(\mathscr P_\bullet)$ (각각
$\mathscr H_n(\mathscr Q_\bullet)$)이 유한형 $\mathscr O_X$-가군
(각각 $\mathscr O_Y$-가군)이라고 가정하자. 그러면
$\mathscr{T}\!or_n^S(\mathscr P_\bullet,\mathscr Q_\bullet)$은
연접 $\mathscr O_Z$-가군이다.
\end{proposition}

\begin{proof}
$\mathscr P_\bullet$과 $\mathscr Q_\bullet$이 아래로 유계이므로
스펙트럼 열
$''\mathscr E(\mathscr P_\bullet,\mathscr Q_\bullet)$은 쌍정칙이다
(\hyperref[III.6.5.4-fr]{6.5.4}). 두 경우 (i), (ii) 모두에서 $Z$가
국소 뇌터이므로, (0, 11.1.8)에 의하여 $''\mathscr E_{pq}^2$의 항들이
연접임을 보이면 충분하다. 따라서
$\mathscr H_n(\mathscr P_\bullet)$과
$\mathscr H_n(\mathscr Q_\bullet)$에 대한 가정 및
$''\mathscr E_{pq}^2$의 식 (6.5.4.2)는 이 명제가
$\mathscr P_\bullet$과 $\mathscr Q_\bullet$이 차수 0의 항으로 축약된
특수한 경우와 동치임을 보여 준다. 달리 말해 이 명제는 그
% R71 authority ends "autrement dit à son" here. The bracketed referent is
% an explicit reversible editorial completion pending printed-witness adjudication.
[따름정리 (\hyperref[III.6.5.14-fr]{6.5.14})]와 동치이다.
\end{proof}

\begin{corollary}[6.5.14]
\label{III.6.5.14-fr}
(\hyperref[III.6.5.13-fr]{6.5.13})의 조건 (i), (ii) 가운데 하나가
성립한다고 가정하고, $\mathscr F$ (각각 $\mathscr G$)를 유한형
준연접 $\mathscr O_X$-가군 (각각 $\mathscr O_Y$-가군)이라 하자.
그러면 $\mathscr{T}\!or_n^S(\mathscr F,\mathscr G)$는 연접
$\mathscr O_Z$-가군이다.
\end{corollary}

\begin{proof}
문제는 $X$와 $Y$ 위에서 국소적이므로 $S$, $X$, $Y$가 아핀이라고
가정할 수 있다.

\textnormal{(i)} (i)의 가정 아래에서 $S$, $X$, $Z$는 뇌터이다.
% Authority note: admitted R71 has S, X, and Z; the unadmitted live leaf omits S.
% This translation remains bound to R71 and the locus is deliberately reversible.
따라서 유한형 $\mathscr O_X$-가군들로 이루어진 $\mathscr F$의
국소 자유 분해 $\mathscr L_\bullet$이 존재한다 (I, 1.3.7). $X$가
$S$ 위에서 평탄하므로 (\hyperref[III.6.3.4-fr]{6.3.4})에 의하여
\[
  \mathscr{T}\!or_n^S(\mathscr F,\mathscr G)
  =\mathscr H_n(\mathscr L_\bullet\otimes_S\mathscr G);
\]
이다. 그런데 $\mathscr L_i\otimes_S\mathscr G$는 유한형 준연접
$\mathscr O_Z$-가군이므로 (I, 9.1.1) 연접이다. 따라서
$\mathscr H_n(\mathscr L_\bullet\otimes_S\mathscr G)$는 연접이다
($0_{\mathrm I}$, 5.3.4).

\textnormal{(ii)} 이제 조건 (ii)가 성립한다고 가정하자. 환 $A(Y)$는
유한 개의 미지수를 갖는 $A(S)$-다항식 대수의 몫이므로 (I, 6.3.3),
$Y$는 $S$ 위에서 평탄하고 유한형인 어떤 아핀 $S$-준스킴 $Y'$의
닫힌 $S$-부분준스킴이다. $Y'$은 뇌터이므로 (I, 6.3.7),
$j:Y\to Y'$을 표준 단사사상이라 할 때 유한형
$\mathscr O_{Y'}$-가군들로 이루어진 $j_*(\mathscr G)$의 국소 자유
분해 $\mathscr M_\bullet$이 존재한다. (\hyperref[III.6.5.12-fr]{6.5.12})에
의하여 $\mathscr{T}\!or_n^S(\mathscr F,\mathscr G)$는
$\mathscr O_{Z'}$-가군
$\mathscr H_n(\mathscr F\otimes_S\mathscr M_\bullet)$의
$1\times j$에 의한 역상이다. 여기서 $Z'=X\times_S Y'$이다. (i)에서와
마찬가지로 $\mathscr F\otimes_S\mathscr M_i$는 연접
$\mathscr O_{Z'}$-가군임을 알 수 있고, 다시
($0_{\mathrm I}$, 5.3.4)에 의하여 결론을 얻는다.
\end{proof}

\begin{env}[6.5.15]
\label{III.6.5.15-fr}
두 $S$-준스킴 $X$, $Y$ 위의 준연접 가군층 복합체
$\mathscr P_\bullet$, $\mathscr Q_\bullet$에 관하여 위에서 전개한
이론은 다음 경우에도 어렵지 않게 일반화된다. 임의의 유한 개의
$S$-준스킴 $X^{(i)}$ ($1\leqslant i\leqslant m$)와 각 $X^{(i)}$ 위의
준연접 $\mathscr O_{X^{(i)}}$-가군 복합체
$\mathscr P_\bullet^{(i)}$를 생각하자. 이때
$Z=X^{(1)}\times_S X^{(2)}\times_S\cdots\times_S X^{(m)}$라 두면,
준연접 $\mathscr O_Z$-가군
$\mathscr{T}\!or_q^S(\mathscr P_\bullet^{(1)},\mathscr P_\bullet^{(2)},
\ldots,\mathscr P_\bullet^{(m)})$를 이와 같이 정의한다. 이 일반적인
경우의 이론 전개는 독자에게 맡기고, 뒤에서 참조할 수 있도록
$\mathscr{T}\!or_\bullet^S(\mathscr P_\bullet^{(1)},
\mathscr P_\bullet^{(2)},\ldots,\mathscr P_\bullet^{(m)})$로 수렴하는
둘째 (정칙) 스펙트럼 열의 $E_2$-항만 적어 두겠다.

\oldpage[III]{153}

\[
  ''\mathscr E_{pq}^2
  =\bigoplus_{q_1+q_2+\cdots+q_m=q}
   \mathscr{T}\!or_p^S
   \bigl(\mathscr H_{q_1}(\mathscr P_\bullet^{(1)}),\ldots,
         \mathscr H_{q_m}(\mathscr P_\bullet^{(m)})\bigr).
  \tag{6.5.15.1}\label{III.6.5.15.1.fr}
\]
임의 개수의 복합체에 대한 이 초Tor 함자들이 낳는 결합법칙 스펙트럼
열은 (6.8)에서 연구하겠다.
\end{env}

\subsection{준연접 가군 복합체의 전역 초Tor 함자와 퀸네트 스펙트럼 열:
아핀 밑의 경우}
\label{subsection:III.6.6.fr}

\begin{env}[6.6.1]
\label{III.6.6.1-fr}
아핀 스킴 $S=\operatorname{Spec}(A)$와 두 준콤팩트 $S$-스킴
$X^{(i)}$ ($i=1,2$)를 생각하자. $\mathscr P_\bullet^{(i)}$를 아래로
유계이고 미분 연산자의 차수가 $-1$인 준연접
$\mathscr O_{X^{(i)}}$-가군의 복합체라 하자 ($i=1,2$). 한편
$X^{(i)}$를 아핀 열린집합들로 덮는 유한 덮개
$\mathfrak U^{(i)}=(U_\alpha^{(i)})$를 생각하자. 다음과 같이 두자.
\[
  X=X^{(1)}\times_S X^{(2)},
\]
이는 준콤팩트 $S$-스킴이다 (I, 5.5.1 및 6.6.4). 또한
$\mathfrak U=\mathfrak U^{(1)}\times_S\mathfrak U^{(2)}$를 아핀
열린집합 $U_\alpha^{(1)}\times_S U_\beta^{(2)}$들로 이루어진 $X$의
덮개라 하자. 정수쌍 $p\leqslant0$, $q\in\mathbf Z$마다 층
$\mathscr P_q^{(i)}$를 계수로 하는 덮개 $\mathfrak U^{(i)}$의 교대
$(-p)$-코체인군
\[
  C^{-p}(\mathfrak U^{(i)},\mathscr P_q^{(i)})
\]
은 $A$-가군이다 (G, II, 5.1). $p>0$이면
$C^p(\mathfrak U^{(i)},\mathscr P_q^{(i)})=0$으로 둔다. 이로써 두
미분 연산자의 차수가 모두 $-1$인 $A$-가군의 이중복합체
\[
  C^\bullet(\mathfrak U^{(i)},\mathscr P_\bullet^{(i)})
\]
를 정의하였다. 정의들에 의하여 ($0_{\mathrm I}$, 12.1.2), 이
이중복합체를 통상 전체 차수에 대한 단일 복합체로 볼 때 그 호몰로지
$A$-가군
\[
  \mathrm H_n\bigl(C^\bullet
  (\mathfrak U^{(i)},\mathscr P_\bullet^{(i)})\bigr)
\]
은 초코호몰로지 $A$-가군
\[
  \mathbf H^{-n}
  (\mathfrak U^{(i)},\mathscr P^{(i)\bullet}),
\]
과 다름없다. 여기서 $\mathscr P^{(i)\bullet}$은
$\mathscr P_{-q}^{(i)}$를 차수 $q$의 성분으로 취하여 얻는, 미분
연산자의 차수가 $+1$인 복합체이다. 표기의 남용으로 이를
$\mathbf H^{-n}(\mathfrak U^{(i)},\mathscr P_\bullet^{(i)})$라고
쓰겠다. 그러면 (\hyperref[III.6.2.2-fr]{6.2.2})에 의하여 이
$A$-가군은 초코호몰로지 $A$-가군
$\mathbf H^{-n}(X^{(i)},\mathscr P^{(i)\bullet})$과 표준적으로
동형이다. 후자도 마찬가지로
$\mathbf H^{-n}(X^{(i)},\mathscr P_\bullet^{(i)})$라고 쓰겠다. 따라서
이는 선택한 유한 덮개 $\mathfrak U^{(i)}$에 의존하지 않는다.
\end{env}

\begin{env}[6.6.2]
\label{III.6.6.2-fr}
두 $A$-가군 이중복합체
\[
  L_{\bullet\bullet}^{(i)}
  =C^\bullet(\mathfrak U^{(i)},\mathscr P_\bullet^{(i)})
  \qquad(i=1,2)
\]
와 이 두 이중복합체에 대한 공변 쌍함자
$L_{\bullet\bullet}^{(1)}\otimes_A L_{\bullet\bullet}^{(2)}$에,
이중복합체에 관한 함자의 초호몰로지 일반 이론
($0_{\mathrm I}$, 11.7.4)을 적용하겠다. 여기서 생각하는 코체인들은
교대이고 덮개 $\mathfrak U^{(i)}$들은 유한이므로,
$C^{-p}(\mathfrak U^{(i)},\mathscr P_q^{(i)})$는 $p$의 유한 개의
값에 대해서만 $\ne0$임을 주목하자. 이 유한 개수는 $q$에 의존하지
않는다. 특히 각 $L_{\bullet\bullet}^{(i)}$의 두 차수는 아래로
유계이다. 다음을
\[
  \mathbf{Tor}_n^A(L_{\bullet\bullet}^{(1)},L_{\bullet\bullet}^{(2)})
\]
또는
\[
  \mathbf{Tor}_n^S
  (\mathfrak U^{(1)},\mathfrak U^{(2)};
   \mathscr P_\bullet^{(1)},\mathscr P_\bullet^{(2)})
\]
$L_{\bullet\bullet}^{(1)}\otimes_A L_{\bullet\bullet}^{(2)}$의 $n$째
초호몰로지 가군이라 하겠다. 이를 덮개 $\mathfrak U^{(1)}$과
$\mathfrak U^{(2)}$에 상대적인 $\mathscr P_\bullet^{(1)}$과
$\mathscr P_\bullet^{(2)}$의 지표 $n$인 \emph{초Tor}라고 부른다.
$\mathscr P_\bullet^{(1)}$과 $\mathscr P_\bullet^{(2)}$이 각각 차수
0의 항 $\mathscr F^{(1)}$, $\mathscr F^{(2)}$로 축약되면 그
초Tor를
\[
  \mathbf{Tor}_n^S
  (\mathfrak U^{(1)},\mathfrak U^{(2)};
   \mathscr F^{(1)},\mathscr F^{(2)})
\]
라고 쓴다. 일반적인 규약에 따라
\[
  \mathbf{Tor}_n^S
  (\mathfrak U^{(1)},\mathfrak U^{(2)};
   \mathscr P_\bullet^{(1)},\mathscr P_\bullet^{(2)}),
\]
로 그 지표 $(j,k)$인 성분이
\[
  \mathbf{Tor}_n^S
  (\mathfrak U^{(1)},\mathfrak U^{(2)};
   \mathscr P_j^{(1)},\mathscr P_k^{(2)}).
\]
인 이중복합체도 나타내기로 한다.

$\mathfrak U^{(i)}$에 속하는 집합들의 교집합은 아핀이므로
(I, 5.5.6 및 1.3.11), $L_{\bullet\bullet}^{(i)}$는
$\mathscr P_\bullet^{(i)}$에 대한 완전 함자이다.

\oldpage[III]{154}

따라서
$\mathbf{Tor}_\bullet^S
(\mathfrak U^{(1)},\mathfrak U^{(2)};
 \mathscr P_\bullet^{(1)},\mathscr P_\bullet^{(2)})$은
$\mathscr P_\bullet^{(1)}$, $\mathscr P_\bullet^{(2)}$에 대한 공변
두 변수 $\partial$-함자이며 그 값은 $A$-가군의 범주에 놓인다
($0_{\mathrm I}$, 11.7.3). 또한 이 쌍함자는 여섯 쌍정칙 스펙트럼
함자의 공통 수렴대상임이 알려져 있다 ($0_{\mathrm I}$, 11.7.2).
이들을
\[
  {}^{(t)}\mathscr E^S
  (\mathfrak U^{(1)},\mathfrak U^{(2)};
   \mathscr P_\bullet^{(1)},\mathscr P_\bullet^{(2)})
\]
또는
\[
  {}^{(t)}\mathscr E
  (\mathfrak U^{(1)},\mathfrak U^{(2)};
   \mathscr P_\bullet^{(1)},\mathscr P_\bullet^{(2)}),
\]
라고 나타내겠다. 여기서 $t$는 $a$, $b$, $a'$, $b'$, $c$, $d$
가운데 하나로 바꾸어야 하며, 그 $E_2$-항들은 다음과 같다.
\[
\begin{aligned}
  {}^{(a)}\mathscr E_{pq}^2
  &=\bigoplus_{q_1+q_2=q}
    \mathbf{Tor}_p^A
    \bigl(\mathrm H_{q_1}(L_{\bullet\bullet}^{(1)}),
          \mathrm H_{q_2}(L_{\bullet\bullet}^{(2)})\bigr),\\
  {}^{(b)}\mathscr E_{pq}^2
  &=\mathrm H_p\bigl(
    \mathbf{Tor}_q^{A,\mathrm{II}}
    (L_{\bullet\bullet}^{(1)},L_{\bullet\bullet}^{(2)})\bigr),\\
  {}^{(a')}\mathscr E_{pq}^2
  &=\bigoplus_{q_1+q_2=q}
    \mathbf{Tor}_p^A
    \bigl(\mathrm H_{q_1}^{\mathrm I}(L_{\bullet\bullet}^{(1)}),
          \mathrm H_{q_2}^{\mathrm I}(L_{\bullet\bullet}^{(2)})\bigr),\\
  {}^{(b')}\mathscr E_{pq}^2
  &=\mathrm H_p\bigl(
    \mathbf{Tor}_q^A
    (L_{\bullet\bullet}^{(1)},L_{\bullet\bullet}^{(2)})\bigr),\\
  {}^{(c)}\mathscr E_{pq}^2
  &=\bigoplus_{q_1+q_2=q}
    \mathbf{Tor}_p^A
    \bigl(\mathrm H_{q_1}^{\mathrm{II}}(L_{\bullet\bullet}^{(1)}),
          \mathrm H_{q_2}^{\mathrm{II}}(L_{\bullet\bullet}^{(2)})\bigr),\\
  {}^{(d)}\mathscr E_{pq}^2
  &=\mathrm H_p\bigl(
    \mathbf{Tor}_q^{A,\mathrm I}
    (L_{\bullet\bullet}^{(1)},L_{\bullet\bullet}^{(2)})\bigr),
\end{aligned}
\]
여기서 표기들은 초호몰로지 일반 이론의 표기들과 일치한다. 이제
이하에서 이 초항들을 더 구체적으로 설명하겠다.
\end{env}

\begin{env}[6.6.3]
\label{III.6.6.3-fr}
\emph{스펙트럼 열 $(a)$와 $(a')$.}
(\hyperref[III.6.6.1-fr]{6.6.1})에서 이중복합체
$L_{\bullet\bullet}^{(i)}$의 호몰로지 가군
$\mathrm H_n(L_{\bullet\bullet}^{(i)})$이
$\mathbf H^{-n}(X^{(i)},\mathscr P_\bullet^{(i)})$와 같음을 보았다.
따라서
\[
  {}^{(a)}\mathscr E_{pq}^2
  =\bigoplus_{q_1+q_2=q}
   \mathbf{Tor}_p^A
   \bigl(\mathbf H^{-q_1}(X^{(1)},\mathscr P_\bullet^{(1)}),
         \mathbf H^{-q_2}(X^{(2)},\mathscr P_\bullet^{(2)})\bigr).
\]
정의에 따라 복합체
$\mathrm H_n^{\mathrm I}(L_{\bullet\bullet}^{(i)})$의 차수 $k$인 항은
호몰로지 가군
$\mathrm H_n(C^\bullet(\mathfrak U^{(i)},\mathscr P_k^{(i)}))$, 즉
정의상 코호몰로지 가군
$\mathbf H^{-n}(\mathfrak U^{(i)},\mathscr P_k^{(i)})$이다. 이 가군은
$\mathbf H^{-n}(X^{(i)},\mathscr P_k^{(i)})$와 표준적으로 동형임이
알려져 있다 (1.4.1). 따라서
\[
  {}^{(a')}\mathscr E_{pq}^2
  =\bigoplus_{q_1+q_2=q}
   \mathbf{Tor}_p^A
   \bigl(\mathbf H^{-q_1}(X^{(1)},\mathscr P_\bullet^{(1)}),
         \mathbf H^{-q_2}(X^{(2)},\mathscr P_\bullet^{(2)})\bigr).
\]
\end{env}

\begin{env}[6.6.4]
\label{III.6.6.4-fr}
\emph{스펙트럼 열 $(b)$와 $(b')$.}
정의에 따라
$\mathbf{Tor}_q^{A,\mathrm{II}}
(L_{\bullet\bullet}^{(1)},L_{\bullet\bullet}^{(2)})$는 차수 $(h,k)$인
항이 다음 $A$-가군인 이중복합체이다.
\[
  \mathbf{Tor}_q^A
  \bigl(C^{-h}(\mathfrak U^{(1)},\mathscr P_\bullet^{(1)}),
        C^{-k}(\mathfrak U^{(2)},\mathscr P_\bullet^{(2)})\bigr).
\]
$\Phi^{(i)}$를 $\mathfrak U^{(i)}$의 지표집합이라 하자. 정의에 따라
가군 복합체
$C^r(\mathfrak U^{(i)},\mathscr P_\bullet^{(i)})$ ($r\geqslant0$)는
복합체 $\Gamma(U_\rho^{(i)},\mathscr P_\bullet^{(i)})$들의 직합이다.
여기서 $U_\rho^{(i)}$는 $\xi\in\rho$에 대한
$U_\xi^{(i)}$들의 교집합이고, $\rho$는
$\mathfrak P(\Phi^{(i)})$를 달린다. 따라서 $A$-가군
\[
  \mathbf{Tor}_q^A
  \bigl(C^{-h}(\mathfrak U^{(1)},\mathscr P_\bullet^{(1)}),
        C^{-k}(\mathfrak U^{(2)},\mathscr P_\bullet^{(2)})\bigr)
\]
은 $A$-가군
\[
  \mathbf{Tor}_q^A
  \bigl(\Gamma(U_\sigma^{(1)},\mathscr P_\bullet^{(1)}),
        \Gamma(U_\tau^{(2)},\mathscr P_\bullet^{(2)})\bigr),
\]
들의 직합이다. 여기서 $\sigma$ (각각 $\tau$)는
$\operatorname{Card}(\sigma)=-(h+1)$ (각각
$\operatorname{Card}(\tau)=-(k+1)$)을 만족하는
$\mathfrak P(\Phi^{(1)})$ (각각 $\mathfrak P(\Phi^{(2)})$)의 원소들을
달린다.
% Authority note: admitted R71 has -(h+1), -(k+1); the unadmitted live leaf
% has 1-h, 1-k. This translation remains bound to R71 and is reversible.
$X^{(1)}$과 $X^{(2)}$은 스킴이므로 $U_\rho^{(i)}$들은 아핀이다.
따라서 (\hyperref[III.6.4.1.1.fr]{6.4.1.1})에 의하여
\[
  \mathbf{Tor}_q^A
  \bigl(\Gamma(U_\sigma^{(1)},\mathscr P_\bullet^{(1)}),
        \Gamma(U_\tau^{(2)},\mathscr P_\bullet^{(2)})\bigr)
  =\Gamma\bigl(U_\sigma^{(1)}\times_S U_\tau^{(2)},
    \mathscr{T}\!or_q^S
    (\mathscr P_\bullet^{(1)},\mathscr P_\bullet^{(2)})\bigr).
\]

\oldpage[III]{155}

그러므로 ${}^{(b)}\mathscr E_{pq}^2$는
$\Phi^{(1)}$과 $\Phi^{(2)}$ 위의 이중 교대 코체인 복합체
\[
  L^\bullet(\Phi^{(1)},\Phi^{(2)};\mathscr S)
\]
의 $(-p)$째 코호몰로지 가군이다. 그 계수계는
\[
  \mathscr S:(\sigma,\tau)\leadsto
  \Gamma\bigl(U_\sigma^{(1)}\times_S U_\tau^{(2)},
    \mathscr{T}\!or_q^S
    (\mathscr P_\bullet^{(1)},\mathscr P_\bullet^{(2)})\bigr)
\]
이다 ($0_{\mathrm I}$, 11.8.4). 이 복합체의 코호몰로지는
$\Phi^{(1)}$과 $\Phi^{(2)}$ 위에서 $\mathscr S$에 값을 갖는 모든
코체인의 복합체
$C^\bullet(\Phi^{(1)},\Phi^{(2)};\mathscr S)$의 코호몰로지와 같고,
또한 다음 원소들의 선형결합으로 이루어진 복합체
$P^\bullet(\Phi^{(1)},\Phi^{(2)};\mathscr S)$의 코호몰로지와도 같음이
알려져 있다 ($0_{\mathrm I}$, 11.8.5 및 11.8.6).
\[
  \lambda(\sigma,\tau)\in
  \Gamma\bigl(U_\sigma^{(1)}\times_S U_\tau^{(2)},
    \mathscr{T}\!or_q^S
    (\mathscr P_\bullet^{(1)},\mathscr P_\bullet^{(2)})\bigr),
\]
여기서 $\sigma=(\alpha_0,\ldots,\alpha_h)$와
$\tau=(\beta_0,\ldots,\beta_h)$는 같은 개수의 원소를 갖는 열이다.
그런데
\[
  U_\sigma^{(1)}\times_S U_\tau^{(2)}
  =(U_{\alpha_0}^{(1)}\times_S U_{\beta_0}^{(2)})\cap\cdots\cap
    (U_{\alpha_h}^{(1)}\times_S U_{\beta_h}^{(2)})
  \qquad\text{(I, 3.2.7)}.
\]
$\mathfrak U$를 아핀 열린집합
$U_\alpha^{(1)}\times_S U_\beta^{(2)}$들로
$Z=X^{(1)}\times_S X^{(2)}$를 덮는 덮개라 하자.
$X^{(1)}\times_S X^{(2)}$이 스킴이라는 점을 고려하면, 끝으로
(1.3.1)에 의하여
\[
  {}^{(b)}\mathscr E_{pq}^2
  =\mathbf H^{-p}\bigl(X^{(1)}\times_S X^{(2)},
    \mathscr{T}\!or_q^S
    (\mathscr P_\bullet^{(1)},\mathscr P_\bullet^{(2)})\bigr).
\]

둘째로
$\mathbf{Tor}_q^A(L_{\bullet\bullet}^{(1)},L_{\bullet\bullet}^{(2)})$는
차수 $(h,k)$인 항이, $h_1+h_2=h$ 및 $k_1+k_2=k$를 만족하는
$A$-가군
\[
  \mathbf{Tor}_q^A
  \bigl(C^{-h_1}(\mathfrak U^{(1)},\mathscr P_{k_1}^{(1)}),
        C^{-h_2}(\mathfrak U^{(2)},\mathscr P_{k_2}^{(2)})\bigr)
\]
들의 직합인 이중복합체이다.
$C^r(\mathfrak U^{(i)},\mathscr P_s^{(i)})$를 위와 같이 풀어 쓰면,
이 항도 $A$-가군
\[
  \Gamma\bigl(U_\sigma^{(1)}\times_S U_\tau^{(2)},
    \mathscr{T}\!or_q^S
    (\mathscr P_{k_1}^{(1)},\mathscr P_{k_2}^{(2)})\bigr),
\]
들의 직합임을 알 수 있다. 여기서 $k_1+k_2=k$이고,
$\sigma$ (각각 $\tau$)는
$\operatorname{Card}(\sigma)+\operatorname{Card}(\tau)=-h-2$를
만족하는 $\mathfrak P(\Phi^{(1)})$ (각각
$\mathfrak P(\Phi^{(2)})$)의 원소들을 달린다.
% Authority note: admitted R71 has -h-2; the unadmitted live leaf has 2-h.
% This translation remains bound to R71 and the locus is deliberately reversible.
지금 계산하는 ${}^{(b')}\mathscr E_{pq}^2$는 이중복합체
$N^{\bullet\bullet}=(N^{hk})$의 $(-p)$째 코호몰로지 가군이다.
여기서 단일 복합체 $N^{\bullet k}$는 $\Phi^{(1)}$과
$\Phi^{(2)}$ 위에서 계수계
\[
  \mathscr S_k:(\sigma,\tau)\leadsto
  \Gamma\left(U_\sigma^{(1)}\times_S U_\tau^{(2)},
    \bigoplus_{k_1+k_2=k}
    \mathscr{T}\!or_q^S
    (\mathscr P_{-k_1}^{(1)},\mathscr P_{-k_2}^{(2)})\right),
\]
에 값을 갖는 이중 교대 코체인의 복합체이다. 이 계수계들은 복합체
$\mathscr S^\bullet$을 이루며, 그 미분은 이중복합체
$\mathscr{T}\!or_q^S
(\mathscr P_\bullet^{(1)},\mathscr P_\bullet^{(2)})$에 결부된 단일
복합체의 미분에서 온다. $N^{\bullet\bullet}$의 코호몰로지는
이중복합체
$C^\bullet(\Phi^{(1)},\Phi^{(2)};\mathscr S^\bullet)$의 코호몰로지와
같고 ($0_{\mathrm I}$, 11.8.9), 또한 이중복합체
$P^\bullet(\Phi^{(1)},\Phi^{(2)};\mathscr S^\bullet)$의 코호몰로지와도
같다. 후자의 차수 $(h,k)$인 원소들은 다음 원소들의 선형결합이다.
\[
  \lambda(\sigma,\tau)\in
  \Gamma\left(U_\sigma^{(1)}\times_S U_\tau^{(2)},
    \bigoplus_{k_1+k_2=k}
    \mathscr{T}\!or_q^S
    (\mathscr P_{-k_1}^{(1)},\mathscr P_{-k_2}^{(2)})\right),
\]
여기서 $\sigma=(\alpha_0,\ldots,\alpha_h)$와
$\tau=(\beta_0,\ldots,\beta_h)$는 같은 개수의 원소를 갖는 열이다
($0_{\mathrm I}$, 11.8.10). 그러면 위와 마찬가지로
${}^{(b')}\mathscr E_{pq}^2$는
$C^\bullet(\mathfrak U,\mathscr Q^\bullet)$의 $(-p)$째 코호몰로지
가군임을 알 수 있다. 여기서 $\mathscr Q^\bullet$은
$\mathscr O_Z$-가군의 이중복합체
$\mathscr{T}\!or_q^S
(\mathscr P_\bullet^{(1)},\mathscr P_\bullet^{(2)})$에 결부된 단일
복합체이다. (\hyperref[III.6.6.1-fr]{6.6.1})의 규약에 따라
\[
  {}^{(b')}\mathscr E_{pq}^2
  =\mathbf H^{-p}\bigl(X^{(1)}\times_S X^{(2)},
    \mathscr{T}\!or_q^S
    (\mathscr P_\bullet^{(1)},\mathscr P_\bullet^{(2)})\bigr).
\]
\end{env}

\oldpage[III]{156}
\begin{env}[6.6.5]
\label{III.6.6.5-fr}
\emph{스펙트럼 열 $(c)$와 $(d)$.}
정의에 따라 복합체
$\mathrm H_n^{\mathrm{II}}(L_{\bullet\bullet}^{(i)})$의 차수 $h$인
항은 $A$-가군
$C^{-h}(\mathfrak U^{(i)},\mathscr H_n(\mathscr P_\bullet^{(i)}))$이다.
이는 함자 $C^{-h}$의 완전성에서 따른다. 따라서 두 덮개에 상대적인
두 가군의 초Tor 정의 (\hyperref[III.6.6.2-fr]{6.6.2})에 의하여
\[
  {}^{(c)}\mathscr E_{pq}^2
  =\bigoplus_{q_1+q_2=q}
   \mathbf{Tor}_p^S
   \bigl(\mathfrak U^{(1)},\mathfrak U^{(2)};
         \mathscr H_{q_1}(\mathscr P_\bullet^{(1)}),
         \mathscr H_{q_2}(\mathscr P_\bullet^{(2)})\bigr).
\]
끝으로 정의에 따라
$\mathbf{Tor}_q^{A,\mathrm I}
(L_{\bullet\bullet}^{(1)},L_{\bullet\bullet}^{(2)})$는 차수 $(h,k)$인
항이 $A$-가군
$\mathbf{Tor}_q^S
(\mathfrak U^{(1)},\mathfrak U^{(2)};
 \mathscr P_h^{(1)},\mathscr P_k^{(2)})$인 이중복합체이다. 따라서
\[
  {}^{(d)}\mathscr E_{pq}^2
  =\mathrm H_p\bigl(
    \mathbf{Tor}_q^S
    (\mathfrak U^{(1)},\mathfrak U^{(2)};
     \mathscr P_\bullet^{(1)},\mathscr P_\bullet^{(2)})\bigr).
\]
\end{env}

\begin{env}[6.6.6]
\label{III.6.6.6-fr}
이중복합체의 함자에 대한 초호몰로지 이론
($0_{\mathrm I}$, 11.7.3)은 (\hyperref[III.6.3.4-fr]{6.3.4})에서와
마찬가지로 다음을 보여 준다. 아래로 유계인 가군 복합체의 범주에서
$L_{\bullet\bullet}^{(i)}$의 임의의 \emph{평탄}
카르탕--아일렌베르크 분해 $M_{\bullet\bullet}^{(i)}$를 택하면
($i=1,2$), 두 변수 $\partial$-함자의 표준 동형
\[
\begin{aligned}
  \mathbf{Tor}_\bullet^S
  (\mathfrak U^{(1)},\mathfrak U^{(2)};
   \mathscr P_\bullet^{(1)},\mathscr P_\bullet^{(2)})
  &\xrightarrow{\sim}
   \mathrm H_\bullet
   (M_{\bullet\bullet}^{(1)}\otimes_A M_{\bullet\bullet}^{(2)})\\
  &\xrightarrow{\sim}
   \mathrm H_\bullet
   (M_{\bullet\bullet}^{(1)}\otimes_A L_{\bullet\bullet}^{(2)})
   \xrightarrow{\sim}
   \mathrm H_\bullet
   (L_{\bullet\bullet}^{(1)}\otimes_A M_{\bullet\bullet}^{(2)}).
\end{aligned}
\tag{6.6.6.1}\label{III.6.6.6.1.fr}
\]
을 얻는다.
\end{env}

\begin{env}[6.6.7]
\label{III.6.6.7-fr}
이제 (\hyperref[III.6.6.2-fr]{6.6.2})에서 정의한 전역 초Tor와 이에
대응하는 여섯 스펙트럼 열이, 그것들을 정의하는 데 사용한 유한 아핀
열린 덮개 $\mathfrak U^{(i)}$에 표준 동형을 제외하면 의존하지 않음을
보이겠다. 이를 위해 같은 종류의 다른 두 덮개 $\mathfrak V^{(i)}$가
$i=1,2$ 각각에 대하여 $\mathfrak U^{(i)}$보다 \emph{더 세분되어}
있을 때, 스펙트럼 함자의 \emph{표준 동형}
\[
  {}^{(t)}\mathscr E
  (\mathfrak U^{(1)},\mathfrak U^{(2)};
   \mathscr P_\bullet^{(1)},\mathscr P_\bullet^{(2)})
  \xrightarrow{\sim}
  {}^{(t)}\mathscr E
  (\mathfrak V^{(1)},\mathfrak V^{(2)};
   \mathscr P_\bullet^{(1)},\mathscr P_\bullet^{(2)})
  \tag{6.6.7, $t$}\label{III.6.6.7.t.fr}
\]
이 존재함을 보이면 충분하다. 여기서 $t$는 $a$, $b$, $a'$, $b'$,
$c$, $d$ 가운데 하나로 바꾼다.

그런데 $i=1,2$에 대하여 호모토피를 제외하면 잘 정의되는 이중복합체
준동형
\[
  C^\bullet(\mathfrak U^{(i)},\mathscr P_\bullet^{(i)})
  \longrightarrow
  C^\bullet(\mathfrak V^{(i)},\mathscr P_\bullet^{(i)})
\]
이 존재한다 (G, II, 5.7.1). 이로부터 이미 표준적으로 정의되고
수렴대상의 경계 연산자와 가환하는 준동형 (6.6.7, $t$)를 얻는다
($0_{\mathrm I}$, 11.3.2). 또한 스펙트럼 열 $(a)$, $(b)$, $(a')$,
$(b')$의 $E_2$-항 계산은 이 네 스펙트럼 열에 대하여 준동형
(6.6.7, $t$)가 $E_2$-항에서 전단사임을 보여 준다. 이 스펙트럼
열들은 쌍정칙이므로, (6.6.7, $t$)는 이 네 스펙트럼 함자의 동형이며,
따라서 그 공통 수렴대상에 대한 \emph{두 변수 $\partial$-함자의
동형}이다 ($0_{\mathrm I}$, 11.1.5).

특히 준연접 $\mathscr O_{X^{(i)}}$-가군 $\mathscr F^{(i)}$
($i=1,2$)에 대하여 표준 준동형
\[
  \mathbf{Tor}_\bullet^S
  (\mathfrak U^{(1)},\mathfrak U^{(2)};
   \mathscr F^{(1)},\mathscr F^{(2)})
  \longrightarrow
  \mathbf{Tor}_\bullet^S
  (\mathfrak V^{(1)},\mathfrak V^{(2)};
   \mathscr F^{(1)},\mathscr F^{(2)})
\]
은 전단사이다. (\hyperref[III.6.6.5-fr]{6.6.5})의 계산을 고려하면
(6.6.7, $t$)는 $t=c$, $t=d$에 대해서도 $E_2$-항의 동형이다.
위와 마찬가지로 (6.6.7, $t$)가 $t=c$, $t=d$에 대해서도 스펙트럼
열의 동형임을 얻는다.

\oldpage[III]{157}

동형 (6.6.7, $t$)들은 $X^{(1)}$과 $X^{(2)}$의 유한 아핀 열린 덮개
쌍 $(\mathfrak U^{(1)},\mathfrak U^{(2)})$들의 여과집합 위에서
스펙트럼 함자의 귀납계를 정의한다고 볼 수 있다. 이 계의 귀납극한을
\[
  {}^{(t)}\mathscr E
  (X^{(1)},X^{(2)};
   \mathscr P_\bullet^{(1)},\mathscr P_\bullet^{(2)})
  \quad\text{ou}\quad
  {}^{(t)}\mathscr E^S
  (X^{(1)},X^{(2)};
   \mathscr P_\bullet^{(1)},\mathscr P_\bullet^{(2)})
\]
로 나타내겠다. 이 스펙트럼 함자의 수렴대상을
$\mathbf{Tor}_\bullet^S
(X^{(1)},X^{(2)};\mathscr P_\bullet^{(1)},\mathscr P_\bullet^{(2)})$로
나타내고, 이를 두 복합체 $\mathscr P_\bullet^{(1)}$과
$\mathscr P_\bullet^{(2)}$의 \emph{전역 초Tor}라고 부르겠다.
$\mathscr P_\bullet^{(1)}$과 $\mathscr P_\bullet^{(2)}$이 각각 차수
0의 항 $\mathscr F^{(1)}$과 $\mathscr F^{(2)}$로 축약되면
\[
  \mathbf{Tor}_\bullet^S
  (X^{(1)},X^{(2)};\mathscr F^{(1)},\mathscr F^{(2)}),
\]
라고 쓰겠다. 일반적인 규약에 따라
\[
  \mathbf{Tor}_q^S
  (X^{(1)},X^{(2)};\mathscr P_\bullet^{(1)},\mathscr P_\bullet^{(2)})
\]
는 다음 가군들로 이루어진 이중복합체이다.
\[
  \mathbf{Tor}_q^S
  (X^{(1)},X^{(2)};\mathscr P_h^{(1)},\mathscr P_k^{(2)}).
\]
\end{env}

\begin{env}[6.6.8]
\label{III.6.6.8-fr}
가정은 (\hyperref[III.6.6.1-fr]{6.6.1})과 같다고 하자. 이제 두
$S$-사상 $f_i:X^{(i)}\to Y^{(i)}$를 생각하자. 여기서
$Y^{(i)}=\operatorname{Spec}(B_i)$는 \emph{아핀} $S$-스킴이고,
$B_i$는 $A$-대수이다 ($i=1,2$). 이는 $A$-준동형
$B_i\to\Gamma(X^{(i)},\mathscr O_{X^{(i)}})$를 정의한다 (I, 2.2.4).
따라서 (\hyperref[III.6.6.2-fr]{6.6.2})에서 정의한 각
$L_{\bullet\bullet}^{(i)}$는 $B_i$-가군의 이중복합체이고,
$L_{\bullet\bullet}^{(1)}\otimes_A L_{\bullet\bullet}^{(2)}$는
$(B_1\otimes_A B_2)$-가군의 사중복합체이다. 그러므로 그 여섯
초호몰로지 스펙트럼 함자는 $(B_1\otimes_A B_2)$-가군의 스펙트럼
열 범주에 값을 갖는다고 볼 수 있다. 다음과 같이 두자.
$Y=Y^{(1)}\times_S Y^{(2)}=\operatorname{Spec}(B_1\otimes_A B_2)$.
이 가군들에 결부된 준연접 $\mathscr O_Y$-가군들을 생각할 수 있다
(I, 1.3.4). $t=a$, $b$, $a'$, $b'$, $c$, $d$에 대하여 여섯
$\mathscr O_Y$-가군 스펙트럼 열
\[
  \bigl({}^{(t)}\mathscr E
  (X^{(1)},X^{(2)};
   \mathscr P_\bullet^{(1)},\mathscr P_\bullet^{(2)})\bigr)^{\sim}
\]
을
${}^{(t)}\mathscr E
(f_1,f_2;\mathscr P_\bullet^{(1)},\mathscr P_\bullet^{(2)})$로
나타내고, 그 공통 수렴대상
\[
  \bigl(\mathbf{Tor}_\bullet^S
  (X^{(1)},X^{(2)};
   \mathscr P_\bullet^{(1)},\mathscr P_\bullet^{(2)})\bigr)^{\sim}.
\]
을
$\mathfrak{Tor}_\bullet^S
(f_1,f_2;\mathscr P_\bullet^{(1)},\mathscr P_\bullet^{(2)})$로
나타내겠다. $\mathscr P_\bullet^{(i)}$가 차수 0의 항
$\mathscr F^{(i)}$로 축약되면 ($i=1,2$), 이를
$\mathfrak{Tor}_\bullet^S(f_1,f_2;\mathscr F^{(1)},\mathscr F^{(2)})$
로 나타내겠다.
\end{env}

\subsection{준연접 가군 복합체의 전역 초Tor 함자와 퀸네트 스펙트럼 열:
일반적인 경우}

\begin{env}[6.7.1]
\label{III.6.7.1-fr}
이제 (\hyperref[III.6.6.8-fr]{6.6.8})의 정의들을 다음 경우로
일반화하겠다. $S$는 임의의 준스킴, $X^{(i)}$, $Y^{(i)}$는
$S$-준스킴이고, $f_i:X^{(i)}\to Y^{(i)}$는 \emph{분리이고
준콤팩트인} 사상이다. 이때 아래로 유계인 준연접
$\mathscr O_{X^{(i)}}$-가군의 복합체 $\mathscr P_\bullet^{(i)}$
($i=1,2$)의 각 쌍에 대하여, 모든 $n$에 대해 준연접
$\mathscr O_Y$-가군
$\mathfrak{Tor}_n^S
(f_1,f_2;\mathscr P_\bullet^{(1)},\mathscr P_\bullet^{(2)})$와 여섯
스펙트럼 함자를 정의해야 한다. $S$, $Y^{(1)}$, $Y^{(2)}$이
아핀이면 이들은 (\hyperref[III.6.6.8-fr]{6.6.8})의 정의로
환원되어야 한다. 여기서 $Y=Y^{(1)}\times_S Y^{(2)}$로 둔다.

먼저 $S=\operatorname{Spec}(A)$가 \emph{아핀}이라고 가정하되,
$Y^{(1)}$과 $Y^{(2)}$은 임의라 하자. $W^{(i)}$를 $Y^{(i)}$의
아핀 열린집합이라 하자. 그러면 $f_i^{-1}(W^{(i)})$는 준콤팩트
$S$-스킴이고, $W=W^{(1)}\times_S W^{(2)}$는 $Y$의 아핀
열린집합이다. $f'_i:f_i^{-1}(W^{(i)})\to W^{(i)}$를 $f_i$의
제한이라 하고, $\mathscr P'_\bullet{}^{(i)}$를 제한
$\mathscr P_\bullet^{(i)}|f_i^{-1}(W^{(i)})$라 하자 ($i=1,2$).
(\hyperref[III.6.6.8-fr]{6.6.8})에 따라 $W$ 위에는 준연접
$(\mathscr O_Y|W)$-가군의 스펙트럼 열
${}^{(t)}\mathscr E
(f'_1,f'_2;\mathscr P'_\bullet{}^{(1)},\mathscr P'_\bullet{}^{(2)})$가
정의되어 있다. 이들이 붙이기 조건을 만족함을 확인해야 한다
($0_{\mathrm I}$, 3.3.1).

이는 곧 $Y^{(i)}=\operatorname{Spec}(B_i)$가 아핀이고
$W^{(i)}=D(g_i)$, $g_i\in B_i$인 경우로 환원된다. 그러면
$Y=\operatorname{Spec}(B_1\otimes_A B_2)$ 안에서
$W=D(g_1\otimes g_2)$이다 (II, 4.3.2.1).

\oldpage[III]{158}

$X'^{(i)}=f_i^{-1}(W^{(i)})$라 두면, 스펙트럼 함자의 표준 동형
\[
  {}^{(t)}\mathscr E
  (X'^{(1)},X'^{(2)};
   \mathscr P'_\bullet{}^{(1)},\mathscr P'_\bullet{}^{(2)})
  \xrightarrow{\sim}
  {}^{(t)}\mathscr E
  (X^{(1)},X^{(2)};
   \mathscr P_\bullet^{(1)},\mathscr P_\bullet^{(2)})
  \otimes_B B_g
  \tag{6.7.1.1}\label{III.6.7.1.1.fr}
\]
을 확립하면 된다. 여기서 $B=B_1\otimes_A B_2$ 및
$g=g_1\otimes g_2$로 두었다. 이를 위해 $X^{(i)}$의 유한 아핀
열린 덮개 $\mathfrak U^{(i)}$에서 시작하자 ($i=1,2$).
$\mathfrak U'^{(i)}$를 $X'^{(i)}$ 위에 유도된 덮개라 하면, 이것도
아핀 열린집합들로 이루어진다 (I, 5.5.10). 정확히는
\[
  C^\bullet(\mathfrak U'^{(i)},\mathscr P'_\bullet{}^{(i)})
  =C^\bullet(\mathfrak U^{(i)},\mathscr P_\bullet^{(i)})
   \otimes_{B_i}(B_i)_{g_i}.
\]
$L'_{\bullet\bullet}{}^{(i)}
:=C^\bullet(\mathfrak U'^{(i)},\mathscr P'_\bullet{}^{(i)})$로 두면
\[
  L'_{\bullet\bullet}{}^{(1)}\otimes_A L'_{\bullet\bullet}{}^{(2)}
  =\bigl(L_{\bullet\bullet}^{(1)}\otimes_{B_1}(B_1)_{g_1}\bigr)
   \otimes_A
   \bigl(L_{\bullet\bullet}^{(2)}\otimes_{B_2}(B_2)_{g_2}\bigr);
\]
이다. 표준 동형을 제외하면
$B_g=(B_1)_{g_1}\otimes_A(B_2)_{g_2}$이므로, 표준 동형을 제외하면
\[
  L'_{\bullet\bullet}{}^{(1)}\otimes_A L'_{\bullet\bullet}{}^{(2)}
  =(L_{\bullet\bullet}^{(1)}\otimes_A L_{\bullet\bullet}^{(2)})
    \otimes_B B_g.
\]
$M_{\bullet\bullet\bullet}^{(i)}$가
$L_{\bullet\bullet}^{(i)}$의 사영 카르탕--아일렌베르크 분해이고,
이를 $B_i$-가군들로 이루어진 것으로 택했다고 하자.
$(B_i)_{g_i}$가 $B_i$ 위에서 \emph{평탄}하므로
$M'_{\bullet\bullet\bullet}{}^{(i)}
:=M_{\bullet\bullet\bullet}^{(i)}\otimes_{B_i}(B_i)_{g_i}$는 이중복합체
$L'_{\bullet\bullet}{}^{(i)}$의 사영 카르탕--아일렌베르크 분해이다.
더욱이
\[
  M'_{\bullet\bullet\bullet}{}^{(1)}\otimes_A M'_{\bullet\bullet\bullet}{}^{(2)}
  =(M_{\bullet\bullet\bullet}^{(1)}\otimes_A M_{\bullet\bullet\bullet}^{(2)})
    \otimes_B B_g.
\]
따라서 원하는 동형 (6.7.1.1)은 이중복합체의 초호몰로지 정의와
$B$-가군 $G$에 대한 함자 $G\mapsto G\otimes_B B_g$의 완전성에서
곧바로 따른다.
\end{env}

\begin{env}[6.7.2]
\label{III.6.7.2-fr}
이제 $S$가 임의라고 가정하고 $u_i:Y^{(i)}\to S$를 구조 사상이라
하자 ($i=1,2$). $(S_\alpha)$를 $S$의 아핀 열린 덮개라 하고,
$Y_\alpha^{(i)}=u_i^{-1}(S_\alpha)$,
$X_\alpha^{(i)}=f_i^{-1}(Y_\alpha^{(i)})$로 두자. 또한
$f_{i\alpha}:X_\alpha^{(i)}\to Y_\alpha^{(i)}$를 $f_i$의 제한이라
하자. 이는 \emph{분리이고 준콤팩트인} 사상이다.
$Y_\alpha=Y_\alpha^{(1)}\times_{S_\alpha}Y_\alpha^{(2)}$들은 $Y$의
열린 덮개를 이루며, 각 $Y_\alpha$ 위에는
(\hyperref[III.6.7.1-fr]{6.7.1})에 의하여 스펙트럼 함자
\[
  {}^{(t)}\mathscr E^{S_\alpha}
  \bigl(f_{1\alpha},f_{2\alpha};
    \mathscr P_\bullet^{(1)}|X_\alpha^{(1)},
    \mathscr P_\bullet^{(2)}|X_\alpha^{(2)}\bigr);
\]
가 정의되어 있다. 이 함자들이 붙이기 조건을 만족함을 다시 보여야
한다.

이는 곧 다음 상황으로 환원된다. $S=\operatorname{Spec}(A)$는
아핀이고, $S'=D(h)$이며 $h\in A$이고,
$u_i(Y^{(i)})\subset S'$이다. 또한
$Y^{(i)}=\operatorname{Spec}(B_i)$가 아핀이라고 가정할 수 있다.
표준 동형
\[
  {}^{(t)}\mathscr E^S
  (X^{(1)},X^{(2)};
   \mathscr P_\bullet^{(1)},\mathscr P_\bullet^{(2)})
  \xrightarrow{\sim}
  {}^{(t)}\mathscr E^{S'}
  (X^{(1)},X^{(2)};
   \mathscr P_\bullet^{(1)},\mathscr P_\bullet^{(2)}).
  \tag{6.7.2.1}\label{III.6.7.2.1.fr}
\]
을 정의하면 된다. 그런데 (\hyperref[III.6.6.2-fr]{6.6.2})의
표기에서 $L_{\bullet\bullet}^{(i)}$들은 $A_h$-가군들로 이루어져
있으므로, 표준 동형을 제외하면
$L_{\bullet\bullet}^{(1)}\otimes_{A_h}L_{\bullet\bullet}^{(2)}
:=L_{\bullet\bullet}^{(1)}\otimes_A L_{\bullet\bullet}^{(2)}$이다.
$L_{\bullet\bullet}^{(i)}$의 사영 카르탕--아일렌베르크 분해
$M_{\bullet\bullet\bullet}^{(i)}$를 $A_h$-가군들로 이루어진 것으로
택할 수 있으므로, 이로부터 원하는 표준 동형을 곧 얻는다.
\end{env}

요약하면 다음 정리를 증명하였다.

\begin{theorem}[6.7.3]
\label{III.6.7.3-fr}
$S$를 준스킴이라 하고, $f_i:X^{(i)}\to Y^{(i)}$를 $S$-준스킴의
분리이고 준콤팩트인 $S$-사상, $\mathscr P_\bullet^{(i)}$를 아래로
유계인 준연접 $\mathscr O_{X^{(i)}}$-가군의 복합체라 하자
($i=1,2$). $u_i:Y^{(i)}\to S$를 구조 사상이라 하고
$Y=Y^{(1)}\times_S Y^{(2)}$로 두자. 준연접
$\mathscr O_Y$-가군의 범주에 값을 갖는 두 변수 $\partial$-함자
$\mathfrak{Tor}_\bullet^S
(f_1,f_2;\mathscr P_\bullet^{(1)},\mathscr P_\bullet^{(2)})$가 존재하며
다음 성질을 만족한다. $S_\alpha$가 $S$의 아핀 열린집합이고,
$V^{(i)}$가 $Y_\alpha^{(i)}=u_i^{-1}(S_\alpha)$의 아핀 열린집합이며
($i=1,2$), $V=V^{(1)}\times_{S_\alpha}V^{(2)}$이면
\[
  \mathfrak{Tor}_\bullet^S
  (f_1,f_2;\mathscr P_\bullet^{(1)},\mathscr P_\bullet^{(2)})|V
  =\Bigl(\mathbf{Tor}_\bullet^{S_\alpha}
    \bigl(f_1^{-1}(V^{(1)}),f_2^{-1}(V^{(2)});
      \mathscr P_\bullet^{(1)}|f_1^{-1}(V^{(1)}),
      \mathscr P_\bullet^{(2)}|f_2^{-1}(V^{(2)})\bigr)\Bigr)^{\sim}.
\]

\oldpage[III]{159}

이 쌍함자는 여섯 쌍정칙 스펙트럼 함자
\[
  {}^{(t)}\mathscr E
  (f_1,f_2;\mathscr P_\bullet^{(1)},\mathscr P_\bullet^{(2)})
  \qquad(t=a,b,a',b',c,d)
\]
의 수렴대상이며, 그 $E_2$-항들은 다음과 같다.
\[
\begin{aligned}
  {}^{(a)}\mathscr E_{pq}^2
  &=\bigoplus_{q_1+q_2=q}
    \mathscr{T}\!or_p^S
    \bigl(\mathscr H^{-q_1}(f_1,\mathscr P_\bullet^{(1)}),
          \mathscr H^{-q_2}(f_2,\mathscr P_\bullet^{(2)})\bigr),\\
  {}^{(b)}\mathscr E_{pq}^2
  &=\mathscr H^{-p}\bigl(f_1\times_S f_2,
    \mathfrak{Tor}_q^S
    (\mathscr P_\bullet^{(1)},\mathscr P_\bullet^{(2)})\bigr),\\
  {}^{(a')}\mathscr E_{pq}^2
  &=\bigoplus_{q_1+q_2=q}
    \mathfrak{Tor}_p^S
    \bigl(\mathscr H^{-q_1}(f_1,\mathscr P_\bullet^{(1)}),
          \mathscr H^{-q_2}(f_2,\mathscr P_\bullet^{(2)})\bigr),\\
  {}^{(b')}\mathscr E_{pq}^2
  &=\mathscr H^{-p}\bigl(f_1\times_S f_2,
    \mathscr{T}\!or_q^S
    (\mathscr P_\bullet^{(1)},\mathscr P_\bullet^{(2)})\bigr),\\
  {}^{(c)}\mathscr E_{pq}^2
  &=\bigoplus_{q_1+q_2=q}
    \mathfrak{Tor}_p^S
    \bigl(f_1,f_2;
      \mathscr H_{q_1}(\mathscr P_\bullet^{(1)}),
      \mathscr H_{q_2}(\mathscr P_\bullet^{(2)})\bigr),\\
  {}^{(d)}\mathscr E_{pq}^2
  &=\mathscr H_p\bigl(\mathfrak{Tor}_q^S
    (f_1,f_2;\mathscr P_\bullet^{(1)},\mathscr P_\bullet^{(2)})\bigr).
\end{aligned}
\]
스펙트럼 열 $(a)$와 $(b)$를 \emph{퀸네트 스펙트럼 열}이라고 한다.

$\mathscr P_\bullet^{(1)}$과 $\mathscr P_\bullet^{(2)}$이 각각 차수
0의 항으로 축약되면 스펙트럼 열 $(a)$와 $(a')$ (각각 $(b)$와
$(b')$)는 동일하다. 이 경우 스펙트럼 열 $(c)$와 $(d)$는 퇴화하므로
관심 대상이 아니다.
\end{theorem}

\begin{remark}[6.7.4]
\label{III.6.7.4-fr}
위에서 정의한 전역 초Tor는
(\hyperref[III.6.2.1-fr]{6.2.1})에서 정의한
\emph{초코호몰로지 가군}과
(\hyperref[III.6.5.3-fr]{6.5.3})에서 정의한 \emph{국소 초Tor}를
모두 특수한 경우로 포함한다. 임의의 \emph{준콤팩트이고 분리인}
사상 $f:X\to Y$와 아래로 유계인 준연접 $\mathscr O_X$-가군의
복합체 $\mathscr P_\bullet$에 대하여, $\mathscr P_\bullet$에 대한
$\partial$-함자의 표준 동형
\[
  \mathfrak{Tor}_n^Y(f,1_Y;\mathscr P_\bullet,\mathscr O_Y)
  \xrightarrow{\sim}
  \mathscr H^{-n}(f,\mathscr P_\bullet)
  \qquad\text{(모든 $n\in\mathbf Z$에 대하여)}.
  \tag{6.7.4.1}\label{III.6.7.4.1.fr}
\]
이 성립함을 보이자. 실제로
(\hyperref[III.6.7.2-fr]{6.7.2})의 붙이기 방법에 의하여 곧 $Y$가
아핀인 경우로 환원된다. 그러면 (\hyperref[III.6.2.2-fr]{6.2.2})에
의하여 $X$의 같은 유한 아핀 열린 덮개 $\mathfrak U$를 사용하여
(6.7.4.1)의 양변을 계산할 수 있다. 첫째 항의 계산에서는 $Y$
자신만으로 이루어진 $Y$의 덮개도 사용한다.
(\hyperref[III.6.6.2-fr]{6.6.2})의 표기에서 이때
$L_{\bullet\bullet}^{(2)}$는 차수 $(0,0)$의 항 $A$로 축약되고,
결론은 ($0_{\mathrm I}$, 11.7.5)에서 따른다. 이 결과의 일반화는
(6.7.7)을 보라. 그러나 (6.7.4.1)의 첫째 항에서 $\mathscr O_Y$를
임의의 준연접 $\mathscr O_Y$-가군 $\mathscr F$로 바꾸면, 일반적으로
더 이상
$\mathscr H^{-n}(f,\mathscr P_\bullet\otimes_{\mathscr O_Y}\mathscr F)$와
동형이 아님을 주목하자. 이는 앞의 계산에서 이중복합체
$L_{\bullet\bullet}^{(1)}\otimes_A L_{\bullet\bullet}^{(2)}$가 여전히
$C^\bullet(\mathfrak U,\mathscr P_\bullet\otimes_Y\mathscr F)$와
동일시된다는 사실과 양립한다.

한편 두 변수 $\partial$-함자의 표준 동형
\[
  \mathfrak{Tor}_\bullet^S
  (1_{X^{(1)}},1_{X^{(2)}};
   \mathscr P_\bullet^{(1)},\mathscr P_\bullet^{(2)})
  \xrightarrow{\sim}
  \mathscr{T}\!or_\bullet^S
  (\mathscr P_\bullet^{(1)},\mathscr P_\bullet^{(2)}).
  \tag{6.7.4.2}\label{III.6.7.4.2.fr}
\]
이 성립한다. 실제로 (\hyperref[III.6.7.1-fr]{6.7.1})과
(\hyperref[III.6.7.2-fr]{6.7.2})에 의하여 다시 $S$와
$X^{(i)}$들이 아핀인 경우로 환원된다. 그러면 (6.7.4.2)의 첫째 항을
계산할 때 $\mathfrak U^{(i)}$로 $X^{(i)}$ 하나만으로 이루어진
덮개를 택할 수 있다.

\oldpage[III]{160}

따라서 (\hyperref[III.6.6.2-fr]{6.6.2})의 표기에서
$L_{\bullet\bullet}^{(i)}$는
$\Gamma(X^{(i)},\mathscr P_\bullet^{(i)})$로 축약된다. 여기서는 이를
첫째 차수가 $\ne0$인 항들이 0인 이중복합체로 본다. 그러면
(6.7.4.2)의 양변이 같다는 것은
(\hyperref[III.6.4.1.1.fr]{6.4.1.1})과
(\hyperref[III.6.3.1-fr]{6.3.1})에서 따른다.
\end{remark}

\begin{proposition}[6.7.5]
\label{III.6.7.5-fr}
$u:\mathscr P_\bullet^{(1)}\to\mathscr Q_\bullet^{(1)}$를 아래로 유계인
준연접 $\mathscr O_{X^{(1)}}$-가군 복합체의 준동형이라 하고, $u$가
유도하는 준동형
\[
  \mathscr H_\bullet(u):
  \mathscr H_\bullet(\mathscr P_\bullet^{(1)})
  \longrightarrow
  \mathscr H_\bullet(\mathscr Q_\bullet^{(1)})
\]
이 전단사라고 가정하자. 그러면 $u$가 유도하는 준동형
\[
  {}^{(t)}\mathscr E
  (f_1,f_2;\mathscr P_\bullet^{(1)},\mathscr P_\bullet^{(2)})
  \longrightarrow
  {}^{(t)}\mathscr E
  (f_1,f_2;\mathscr Q_\bullet^{(1)},\mathscr P_\bullet^{(2)})
\]
은 $t=a$, $t=b$, $t=c$에 대하여 전단사이다.
\end{proposition}

\begin{proof}
스펙트럼 열 $(c)$에 관한 주장은 이 열이 쌍정칙이고, 가정에 의하여
고려하는 준동형이 $E_2$-항에서 전단사라는 사실에서 따른다
($0_{\mathrm I}$, 11.1.5). 이는 이미
\[
  \mathfrak{Tor}_\bullet^S
  (f_1,f_2;\mathscr P_\bullet^{(1)},\mathscr P_\bullet^{(2)})
  \longrightarrow
  \mathfrak{Tor}_\bullet^S
  (f_1,f_2;\mathscr Q_\bullet^{(1)},\mathscr P_\bullet^{(2)})
\]
이 전단사임을 보여 준다. 관계식
(\hyperref[III.6.7.4.1.fr]{6.7.4.1})과
(\hyperref[III.6.7.4.2.fr]{6.7.4.2})를 적용하면 먼저 $u$가 유도하는
준동형
\[
  \mathscr H^{-n}(f_1,\mathscr P_\bullet^{(1)})
  \longrightarrow
  \mathscr H^{-n}(f_1,\mathscr Q_\bullet^{(1)})
  \quad\text{및}\quad
  \mathscr{T}\!or_n^S
  (\mathscr P_\bullet^{(1)},\mathscr P_\bullet^{(2)})
  \longrightarrow
  \mathscr{T}\!or_n^S
  (\mathscr Q_\bullet^{(1)},\mathscr P_\bullet^{(2)})
\]
이 전단사임을 알 수 있다. 그러면 스펙트럼 열 $(a)$와 $(b)$에 관한
주장은 이 열들이 쌍정칙이고
(\hyperref[III.6.7.3-fr]{6.7.3}), 고려하는 준동형이 $E_2$-항에서
전단사라는 사실에서 따른다 ($0_{\mathrm I}$, 11.1.5).

한편
$u:\mathscr P_\bullet^{(1)}\to\mathscr Q_\bullet^{(1)}$가
\emph{호모토피 동치사상}이면, 여섯 스펙트럼 열 모두에 대하여 표준
동형
\[
  {}^{(t)}\mathscr E
  (f_1,f_2;\mathscr P_\bullet^{(1)},\mathscr P_\bullet^{(2)})
  \longrightarrow
  {}^{(t)}\mathscr E
  (f_1,f_2;\mathscr Q_\bullet^{(1)},\mathscr P_\bullet^{(2)})
\]
을 얻는다. 실제로 $S$와 $Y^{(i)}$들이 아핀이면 $u$로부터
이중복합체의 호모토피 동치사상
\[
  C^\bullet(\mathfrak U^{(1)},\mathscr P_\bullet^{(1)})
  \longrightarrow
  C^\bullet(\mathfrak U^{(1)},\mathscr Q_\bullet^{(1)}),
\]
을 얻고, 주장은 초호몰로지의 일반 이론에서 따른다
($0_{\mathrm I}$, 11.3.2). 일반적인 경우로의 이행은 붙이기로 한다.
이때 복합체의 호모토피 동치사상에서 그 복합체들의 사영
카르탕--아일렌베르크 분해 사이의 호모토피 동치사상이 유도된다는
사실을 사용한다 (M, XVII, 1.2).
\end{proof}

\begin{proposition}[6.7.6]
\label{III.6.7.6-fr}
복합체 $\mathscr P_\bullet^{(1)}$ 또는 복합체
$\mathscr P_\bullet^{(2)}$가 $S$-평탄 가군들로 이루어져 있다고
가정하자. 두 복합체는 아래로 유계이다. 그러면 두 변수
$\partial$-함자의 표준 동형
\[
  \mathfrak{Tor}_n^S
  (f_1,f_2;\mathscr P_\bullet^{(1)},\mathscr P_\bullet^{(2)})
  \xrightarrow{\sim}
  \mathscr H^{-n}\bigl(f_1\times_S f_2,
    \mathscr P_\bullet^{(1)}\otimes_S\mathscr P_\bullet^{(2)}\bigr).
  \tag{6.7.6.1}\label{III.6.7.6.1.fr}
\]
이 성립한다.
\end{proposition}

\begin{proof}
먼저 $S$, $Y^{(1)}$, $Y^{(2)}$이 아핀이라고 가정하자. 그러면
(\hyperref[III.6.6.2-fr]{6.6.2})의 상황에 놓이며 그 표기를
유지한다. 예를 들어 $\mathscr P_\bullet^{(1)}$이 $S$-평탄 가군들로
이루어져 있다고 가정하고,
(\hyperref[III.6.6.6-fr]{6.6.6})의 주석을 사용하여 초Tor를 계산하자.
따라서 이것은
$L_{\bullet\bullet}^{(1)}\otimes_A M_{\bullet\bullet\bullet}^{(2)}$의
호몰로지이다. 여기서 $M_{\bullet\bullet\bullet}^{(2)}$는
($0_{\mathrm I}$, 11.7.1)의 의미에서
$L_{\bullet\bullet}^{(2)}$의 사영 카르탕--아일렌베르크 분해이다.
한편 가정에 의하여 가군 $L_{ij}^{(1)}$들은 $A$ 위에서
\emph{평탄}하다 (I, 4.15.1). 그러므로 ($0_{\mathrm I}$, 11.7.5)에서
표준 동형
\[
  \mathbf{Tor}_\bullet^S
  (\mathfrak U^{(1)},\mathfrak U^{(2)};
   \mathscr P_\bullet^{(1)},\mathscr P_\bullet^{(2)})
  \xrightarrow{\sim}
  \mathscr H_\bullet
  (L_{\bullet\bullet}^{(1)}\otimes_A L_{\bullet\bullet}^{(2)}).
  \tag{6.7.6.2}\label{III.6.7.6.2.fr}
\]
을 얻는다.

한편 $L_{\bullet\bullet}^{(1)}\otimes_A L_{\bullet\bullet}^{(2)}$에서
$C^\bullet(\mathfrak U,\mathscr Q_\bullet)$로 가는 이중복합체의 자연
준동형이 있다. 여기서 $\mathfrak U$는 아핀 열린집합
$U_\alpha^{(1)}\times_S U_\beta^{(2)}$들로
$Z=X^{(1)}\times_S X^{(2)}$를 덮는 덮개이고,

\oldpage[III]{161}

$\mathscr Q_\bullet
:=\mathscr P_\bullet^{(1)}\otimes_S\mathscr P_\bullet^{(2)}$는 전체
차수에 대한 단일 복합체로 본다. 실제로 이 준동형의 정의는
(\hyperref[III.6.6.4-fr]{6.6.4})에서 $q=0$일 때 스펙트럼 열
$(b')$를 계산하는 과정에서 본질적으로 주어졌다. 그 표기를 유지하면
다음 세 자연 준동형을 고려하면 된다. 첫째,
$N^{\bullet k}$에서
$C^\bullet(\Phi^{(1)},\Phi^{(2)};\mathscr S_k)$로 가는 준동형
($0_{\mathrm I}$, 11.8.5), 둘째, 이 복합체에서
$P^\bullet(\Phi^{(1)},\Phi^{(2)};\mathscr S_k)$로 가는 준동형
($0_{\mathrm I}$, 11.8.6), 끝으로 이 코체인 복합체에서 교대 코체인의
부분복합체로 가는 준동형이다 ($0_{\mathrm I}$, 11.8.7). 이와 같이
정의한 이중복합체의 준동형은
(\hyperref[III.6.6.4-fr]{6.6.4})에서 본 바와 같이 호몰로지의
동형을 준다. 따라서 (6.7.6.2)와 합성하여 동형
\[
  \mathbf{Tor}_n^S
  (\mathfrak U^{(1)},\mathfrak U^{(2)};
   \mathscr P_\bullet^{(1)},\mathscr P_\bullet^{(2)})
  \xrightarrow{\sim}
  \mathbf H^{-n}
  (\mathfrak U,\mathscr P_\bullet^{(1)}\otimes_S\mathscr P_\bullet^{(2)}).
  \tag{6.7.6.3}\label{III.6.7.6.3.fr}
\]
을 얻었다.

이제 이와 같이 정의한 동형이 선택한 열린 덮개들에 의존하지 않음을
증명해야 한다. (6.7.6.3)의 둘째 항은
(\hyperref[III.6.2.2-fr]{6.2.2})에 의하여
$\mathbf H^{-n}
(X^{(1)}\times_S X^{(2)},
 \mathscr P_\bullet^{(1)}\otimes_S\mathscr P_\bullet^{(2)})$와
표준적으로 동형이다. 이 독립성은
(\hyperref[III.6.6.7-fr]{6.6.7})을 사용하여 증명한다.
그 표기에서 호모토피 동치사상을 제외하면 가환하는 도식
\[
\xymatrix@C=18pt@R=28pt{
  C^\bullet(\mathfrak U^{(1)},\mathscr P_\bullet^{(1)})
  \otimes_A
  C^\bullet(\mathfrak U^{(2)},\mathscr P_\bullet^{(2)})
  \ar[r]\ar[d]
  & C^\bullet(\mathfrak U,\mathscr Q_\bullet)\ar[d]
  \\
  C^\bullet(\mathfrak V^{(1)},\mathscr P_\bullet^{(1)})
  \otimes_A
  C^\bullet(\mathfrak V^{(2)},\mathscr P_\bullet^{(2)})
  \ar[r]
  & C^\bullet(\mathfrak V,\mathscr Q_\bullet)
}
\]
이 있으며, 가로 화살표들은 위에서 정의한 준동형들이다. 끝으로
붙이기에 의하여 일반적인 경우로 이행해야 한다. 이는
(\hyperref[III.6.7.1-fr]{6.7.1})과
(\hyperref[III.6.7.2-fr]{6.7.2})에서와 같이 어렵지 않게 이루어지며,
자세한 내용은 독자에게 맡긴다.
\end{proof}

\begin{proposition}[6.7.7]
\label{III.6.7.7-fr}
복합체 $\mathscr P_\bullet^{(1)}$과
$\mathscr P_\bullet^{(2)}$이 아래로 유계이고, 가군
$\mathscr H^{-n}(f_1,\mathscr P_\bullet^{(1)})$들이 모두
$S$-평탄이거나 가군
$\mathscr H^{-n}(f_2,\mathscr P_\bullet^{(2)})$들이 모두
$S$-평탄이라고 가정하자. 그러면 모든 $n\in\mathbf Z$에 대하여
두 변수 $\partial$-함자의 표준 동형
\[
  \mathfrak{Tor}_n^S
  (f_1,f_2;\mathscr P_\bullet^{(1)},\mathscr P_\bullet^{(2)})
  \xrightarrow{\sim}
  \bigoplus_{q_1+q_2=n}
  \mathscr H^{-q_1}(f_1,\mathscr P_\bullet^{(1)})
  \otimes_S
  \mathscr H^{-q_2}(f_2,\mathscr P_\bullet^{(2)}).
  \tag{6.7.7.1}\label{III.6.7.7.1.fr}
\]
이 성립한다.
\end{proposition}

\begin{proof}
실제로 (\hyperref[III.6.5.8-fr]{6.5.8})에 의하여
(\hyperref[III.6.7.3-fr]{6.7.3})의 스펙트럼 열 $(a)$는 퇴화한다.
또한 이 열은 쌍정칙이므로
(\hyperref[0.11.1.6-fr]{$0_{\mathrm I}$, 11.1.6})에서 명제가 곧바로
따른다 (\hyperref[III.6.7.3-fr]{6.7.3}).
\end{proof}

\begin{theorem}[6.7.8]
\label{III.6.7.8-fr}
다음 세 조건을 가정하자.
$1^\circ$ 복합체 $\mathscr P_\bullet^{(1)}$과
$\mathscr P_\bullet^{(2)}$은 아래로 유계이다.
$2^\circ$ 복합체 $\mathscr P_\bullet^{(1)}$ 또는 복합체
$\mathscr P_\bullet^{(2)}$는 $S$-평탄 가군들로 이루어져 있다.
$3^\circ$
\oldpage[III]{162}
가군 $\mathscr H^{-n}(f_1,\mathscr P_\bullet^{(1)})$들이 모두
$S$-평탄이거나 가군
$\mathscr H^{-n}(f_2,\mathscr P_\bullet^{(2)})$들이 모두
$S$-평탄이다. 그러면 모든 $n\in\mathbf Z$에 대하여 두 변수
$\partial$-함자의 표준 동형
\[
  \mathscr H^n\bigl(f_1\times_S f_2,
    \mathscr P_\bullet^{(1)}\otimes_S\mathscr P_\bullet^{(2)}\bigr)
  \xrightarrow{\sim}
  \bigoplus_{n_1+n_2=n}
  \mathscr H^{n_1}(f_1,\mathscr P_\bullet^{(1)})
  \otimes_S
  \mathscr H^{n_2}(f_2,\mathscr P_\bullet^{(2)}).
  \tag{6.7.8.1}\label{III.6.7.8.1.fr}
\]
이 성립한다 (``퀸네트 공식'').
\end{theorem}

\begin{proof}
이는 (\hyperref[III.6.7.6-fr]{6.7.6})과
(\hyperref[III.6.7.7-fr]{6.7.7})에서 따른다.

$S$, $Y^{(1)}$, $Y^{(2)}$이 아핀이면, (6.7.8.1)의 역동형사상은
(\hyperref[III.6.7.6-fr]{6.7.6})의 표기 아래 이중복합체의 준동형
\[
  C^\bullet(\mathfrak U^{(1)},\mathscr P_\bullet^{(1)})
  \otimes_A
  C^\bullet(\mathfrak U^{(2)},\mathscr P_\bullet^{(2)})
  \longrightarrow
  C^\bullet(\mathfrak U,\mathscr Q_\bullet)
\]
으로부터 (G, I, 2.7)에 정의된 절차로 얻는다. 이는 (G, I, 5.5)에서
따른다.
\end{proof}

\begin{proposition}[6.7.9]
\label{III.6.7.9-fr}
다음 세 조건이 성립한다고 가정하자.
$1^\circ$ $S$, $Y^{(1)}$, $Y^{(2)}$은 국소 뇌터이고,
$f_1$, $f_2$는 고유하며, $Y^{(1)}$ 또는 $Y^{(2)}$은 $S$ 위에서
유한형이다.
$2^\circ$ $\mathscr P_\bullet^{(1)}$과
$\mathscr P_\bullet^{(2)}$은 아래로 유계이다.
$3^\circ$ 모든 $n\in\mathbf Z$와 $i=1,2$에 대하여
$\mathscr H_n(\mathscr P_\bullet^{(i)})$는 연접 가군이다.
그러면 $Y=Y^{(1)}\times_S Y^{(2)}$라 할 때
$\mathfrak{Tor}_n^S
(f_1,f_2;\mathscr P_\bullet^{(1)},\mathscr P_\bullet^{(2)})$는
연접 $\mathscr O_Y$-가군이다.
\end{proposition}

\begin{proof}
(\hyperref[III.6.5.13-fr]{6.5.13})에 의하여 국소 초Tor
$\mathscr{T}\!or_n^S(\mathscr P_\bullet^{(1)},\mathscr P_\bullet^{(2)})$는
연접 $\mathscr O_X$-가군이다. 여기서
$X=X^{(1)}\times_S X^{(2)}$는 국소 뇌터이다. 실제로 가정에 의하여
$X^{(i)}$ 가운데 하나가 $S$ 위에서 유한형이기 때문이다
(I, 6.3.4 및 6.3.8). 또한 $Y$는 국소 뇌터이고
$f_1\times_S f_2$는 고유하므로 (II, 5.4.2),
(\hyperref[III.6.2.5-fr]{6.2.5})에 의하여
(\hyperref[III.6.7.3-fr]{6.7.3})의 항
${}^{(b)}\mathscr E_{pq}^2$는 연접 $\mathscr O_Y$-가군이다.
가정 $2^\circ$에 의하여
(\hyperref[III.6.7.3-fr]{6.7.3})의 모든 스펙트럼 열이 쌍정칙이므로
(\hyperref[0.11.1.8-fr]{$0_{\mathrm I}$, 11.1.8})에서 결론이 따른다.
\end{proof}

\begin{env}[6.7.10]
\label{III.6.7.10-fr}
이제 $Y'^{(i)}$를 두 $S$-준스킴,
$v_i:Y'^{(i)}\to Y^{(i)}$를 두 $S$-사상이라 하자 ($i=1,2$).
$Y'=Y'^{(1)}\times_S Y'^{(2)}$라 두고, 그 곱
$v=v_1\times_S v_2$를 $S$-사상 $Y'\to Y$로 보자. 한편 각
$i=1,2$에 대하여 $S$-준스킴 $X'^{(i)}$와 두 $S$-사상
$u_i:X'^{(i)}\to X^{(i)}$, $f'_i:X'^{(i)}\to Y'^{(i)}$를 택하여
도식
\[
\xymatrix@C=36pt@R=28pt{
  X'^{(i)}\ar[r]^{u_i}\ar[d]_{f'_i}
  & X^{(i)}\ar[d]^{f_i}
  \\
  Y'^{(i)}\ar[r]_{v_i}
  & Y^{(i)}
}
\tag{6.7.10.1}\label{III.6.7.10.1.fr}
\]
이 가환한다고 하자. 여기서 사상 $f'_i$는 \emph{분리이고
준콤팩트}이다. 그러면 $t=a,a',b,b',c,d$에 대하여 스펙트럼 함자의
표준 $\mathscr O_{Y'}$-준동형
\[
  v^*\bigl({}^{(t)}\mathscr E
  (f_1,f_2;\mathscr P_\bullet^{(1)},\mathscr P_\bullet^{(2)})\bigr)
  \longrightarrow
  {}^{(t)}\mathscr E
  (f'_1,f'_2;u_1^*(\mathscr P_\bullet^{(1)}),
                 u_2^*(\mathscr P_\bullet^{(2)}))
  \tag{6.7.10.2}\label{III.6.7.10.2.fr}
\]
이 존재한다. 이를 정의하기 위하여 먼저
$S=\operatorname{Spec}(A)$,
$Y^{(i)}=\operatorname{Spec}(B_i)$,
$Y'^{(i)}=\operatorname{Spec}(B'_i)$가 아핀이라고 가정하자. 그러면
$X^{(i)}$와 $X'^{(i)}$는 준콤팩트 스킴이다. 스펙트럼 열
$ {}^{(t)}\mathscr E
(f_1,f_2;\mathscr P_\bullet^{(1)},\mathscr P_\bullet^{(2)})$를
계산하기 위하여 (\hyperref[III.6.6.1-fr]{6.6.1})에서와 같이
$X^{(i)}$의 아핀 열린집합들로 이루어진 유한 덮개
$\mathfrak U^{(i)}$를 택한다 ($i=1,2$). 그리고
$ {}^{(t)}\mathscr E
(f'_1,f'_2;u_1^*(\mathscr P_\bullet^{(1)}),
u_2^*(\mathscr P_\bullet^{(2)}))$를 계산하기 위하여
$X'^{(i)}$의 아핀 열린집합들로 이루어진 유한
덮개 $\mathfrak U'^{(i)}$를 택한다.

\oldpage[III]{163}

이 덮개는 각각 $u_i^{-1}(\mathfrak U^{(i)})$보다
\emph{더 세분되어} 있다 ($i=1,2$). 이중복합체
$C^\bullet(\mathfrak U^{(i)},\mathscr P_\bullet^{(i)})
=L_{\bullet\bullet}^{(i)}$에서
$C^\bullet(u_i^{-1}(\mathfrak U^{(i)}),
u_i^*(\mathscr P_\bullet^{(i)}))$로 가는 표준 준동형이 존재한다
($0_{\mathrm I}$, 4.4.3.2). 또한 $\mathfrak U'^{(i)}$에서
$u_i^{-1}(\mathfrak U^{(i)})$로 가는 단체 사상을 택하면
(G, II, 5.7), 이중복합체의 준동형
\[
  C^\bullet(u_i^{-1}(\mathfrak U^{(i)}),u_i^*(\mathscr P_\bullet^{(i)}))
  \longrightarrow
  C^\bullet(\mathfrak U'^{(i)},u_i^*(\mathscr P_\bullet^{(i)}))
\]
을 정의할 수 있다. 따라서 합성하여 이중복합체의 준동형
\[
  L_{\bullet\bullet}^{(i)}
  \longrightarrow
  L_{\bullet\bullet}^{\prime(i)}
  =C^\bullet(\mathfrak U'^{(i)},u_i^*(\mathscr P_\bullet^{(i)})).
\]
을 얻는다. 더욱이 단체 사상을 바꾸면 이 준동형은 그와 호모토픽인
준동형으로 바뀐다 (G, II, 5.7.1). 따라서 잘 정의된 스펙트럼 함자의
준동형
\[
  {}^{(t)}\mathscr E
  (\mathfrak U^{(1)},\mathfrak U^{(2)};
   \mathscr P_\bullet^{(1)},\mathscr P_\bullet^{(2)})
  \longrightarrow
  {}^{(t)}\mathscr E
  (\mathfrak U'^{(1)},\mathfrak U'^{(2)};
   u_1^*(\mathscr P_\bullet^{(1)}),u_2^*(\mathscr P_\bullet^{(2)})).
  \tag{6.7.10.3}\label{III.6.7.10.3.fr}
\]
을 얻는다.

$\mathfrak V^{(i)}$를 $X^{(i)}$의 아핀 유한 덮개로서
$\mathfrak U^{(i)}$보다 더 세분된 것이라 하고,
$\mathfrak V'^{(i)}$를 $X'^{(i)}$의 아핀 유한 덮개로서
$u_i^{-1}(\mathfrak V^{(i)})$와 $\mathfrak U'^{(i)}$보다 모두 더
세분된 것이라 하자. 그러면 도식
\[
\xymatrix@C=34pt@R=30pt{
  C^\bullet(\mathfrak U^{(i)},\mathscr P_\bullet^{(i)})
  \ar[r]\ar[d]
  & C^\bullet(\mathfrak V^{(i)},\mathscr P_\bullet^{(i)})\ar[d]
  \\
  C^\bullet(\mathfrak U'^{(i)},u_i^*(\mathscr P_\bullet^{(i)}))
  \ar[r]
  & C^\bullet(\mathfrak V'^{(i)},u_i^*(\mathscr P_\bullet^{(i)}))
}
\]
이 가환함을 곧바로 확인할 수 있다. 따라서 준동형 (6.7.10.3)은
택한 덮개 $\mathfrak U^{(i)}$와 $\mathfrak U'^{(i)}$에 본질적으로
의존하지 않는다. 그러므로 실제로 $A$-가군의 준동형
\[
  {}^{(t)}E
  (X^{(1)},X^{(2)};\mathscr P_\bullet^{(1)},\mathscr P_\bullet^{(2)})
  \longrightarrow
  {}^{(t)}E
  (X'^{(1)},X'^{(2)};
   u_1^*(\mathscr P_\bullet^{(1)}),u_2^*(\mathscr P_\bullet^{(2)})).
  \tag{6.7.10.4}\label{III.6.7.10.4.fr}
\]
을 정의한 것이다. 그런데 $u_i^*(\mathscr P_\bullet^{(i)})$의 정의와
(6.7.10.1)의 가환성에 의하여 이 준동형은
$(B_1\otimes_A B_2)$-가군의 준동형이기도 하다. 또한 (6.7.10.4)의
둘째 항은 $(B'_1\otimes_A B'_2)$-가군들로 이루어져 있으므로,
(6.7.10.4)에서 표준적인 $(B'_1\otimes_A B'_2)$-가군의 준동형
\[
  {}^{(t)}E
  (X^{(1)},X^{(2)};\mathscr P_\bullet^{(1)},\mathscr P_\bullet^{(2)})
  \otimes_{B_1\otimes_A B_2}(B'_1\otimes_A B'_2)
  \longrightarrow
  {}^{(t)}E
  (X'^{(1)},X'^{(2)};
   u_1^*(\mathscr P_\bullet^{(1)}),u_2^*(\mathscr P_\bullet^{(2)}))
  \tag{6.7.10.5}\label{III.6.7.10.5.fr}
\]
을 얻는다. (I, 1.6.5)를 고려하면 이것은 바로 지금 다루는 특수한
경우에 구한 준동형 (6.7.10.2)이다.

이제 (\hyperref[III.6.7.1-fr]{6.7.1})과
(\hyperref[III.6.7.2-fr]{6.7.2})의 붙이기를 따라 일반적인 경우로
이행하면 된다. 둘째 이행은 즉시 된다. 첫째 이행에 대해서는
(\hyperref[III.6.7.1-fr]{6.7.1})에서와 같이 원소
$g_i\in B_i$와 그 상 $g'_i\in B'_i$, 텐서곱
$B=B_1\otimes_A B_2$ 안의 $g=g_1\otimes g_2$, 그리고
$B'=B'_1\otimes_A B'_2$ 안에서 그 상
$g'=g'_1\otimes g'_2$를 생각한다. 그러면 모든 것은

\oldpage[III]{164}

표준 동형
\[
  (M\otimes_B B')_{g'}
  \xrightarrow{\sim}
  M_g\otimes_{B_g}B'_{g'}
\]
을 사용하는 것으로 귀착된다 ($0_{\mathrm I}$, 1.5.4). 자세한 내용은
독자에게 맡긴다.
\end{env}

\begin{env}[6.7.11]
\label{III.6.7.11-fr}
위에서 두 $S$-사상 $X^{(i)}\to Y^{(i)}$와 아래로 유계인 준연접
가군 복합체 $\mathscr P_\bullet^{(i)}$ 두 개에 대하여 전개한 전역
초Tor 이론은 다음 일반적인 경우로 곧바로 확장된다. 준스킴 $S$,
$S$-준스킴들의 유한 족 $Y^{(i)}$ ($i\in I$), 분리이고 준콤팩트인
$S$-사상들의 유한 족
$f_i:X^{(i)}\to Y^{(i)}$, 그리고 각 $i$에 대하여 아래로 유계인
준연접 $\mathscr O_{X^{(i)}}$-가군 복합체
$\mathscr P_\bullet^{(i)}$가 주어졌다고 하자. $Y$를 $S$-준스킴
$Y^{(i)}$들의 곱이라 하면, 각 정수 $n\in\mathbf Z$에 대하여 준연접
$\mathscr O_Y$-가군
$\mathfrak{Tor}_n^S((f_i)_{i\in I};
(\mathscr P_\bullet^{(i)})_{i\in I})$를 정의한다. 이 가군들은 각
복합체 $\mathscr P_\bullet^{(i)}$에 대하여 공변
$\partial$-함자를 이루며, 이 함자는 여섯 스펙트럼 함자
$ {}^{(t)}\mathscr E((f_i)_{i\in I};
(\mathscr P_\bullet^{(i)})_{i\in I})$의 공통 수렴대상이다.

$I=\{1,2\}$인 경우에 위에서 한 정의와 논증을 이 일반적인 경우에
되풀이하는 일은 독자에게 맡긴다. 다만 $I$가 한 원소만 가지면
(\hyperref[III.6.2.7-fr]{6.2.7})에서 정의한 초코호몰로지
$\mathscr H^\bullet(f,\mathscr P_\bullet)$를 다시 얻는다는 점만
기록하자. 이는 (\hyperref[III.6.7.4-fr]{6.7.4})에서도 이미
관찰하였다. $I$가 $\mathbf N$의 구간 $1\leq i\leq m$이면 다음과
같이 쓰겠다.
\[
  \mathfrak{Tor}_n^S
  (f_1,\ldots,f_m;
   \mathscr P_\bullet^{(1)},\ldots,\mathscr P_\bullet^{(m)})
  \quad\text{대신}\quad
  \mathfrak{Tor}_n^S
  ((f_i)_{i\in I};(\mathscr P_\bullet^{(i)})_{i\in I}).
\]
\end{env}

\begin{proposition}[6.7.12]
\label{III.6.7.12-fr}
표기는 (\hyperref[III.6.7.11-fr]{6.7.11})과 같다고 하자. $J$를
$I$의 부분집합이라 하고, $i\in I-J$에 대하여
$X^{(i)}=Y^{(i)}=S$, $f_i$는 항등사상이며
$\mathscr P_\bullet^{(i)}$는 차수 $0$의 항
$\mathscr O_S$에만 집중된 복합체라고 하자. 그러면
$\partial$-함자의 표준 동형
\[
  \mathfrak{Tor}_\bullet^S
  ((f_i)_{i\in I};(\mathscr P_\bullet^{(i)})_{i\in I})
  \xrightarrow{\sim}
  \mathfrak{Tor}_\bullet^S
  ((f_i)_{i\in J};(\mathscr P_\bullet^{(i)})_{i\in J}).
  \tag{6.7.12.1}\label{III.6.7.12.1.fr}
\]
이 존재한다.
\end{proposition}

\begin{proof}
붙이기는 평소와 같이 이루어지므로 $S$와 $Y^{(i)}$들이 아핀일 때만
이 동형을 정의하면 충분하다. $i\in I-J$에 대해서는 $S$ 하나만으로
이루어진 덮개 $\mathfrak U^{(i)}$를 택할 수 있다. 그러면
$L_{\bullet\bullet}^{(i)}
=C^\bullet(\mathfrak U^{(i)},\mathscr P_\bullet^{(i)})$는 쌍차수
$(0,0)$의 유일한 항
$\Gamma(S,\mathscr O_S)=A(S)$에만 집중된다. 따라서 동형
(6.7.12.1)은 자명하다.
\end{proof}

\begin{remark}[6.7.13]
\label{III.6.7.13-fr}
표기는 (\hyperref[III.6.7.3-fr]{6.7.3})과 같다고 하고, 표준
$S$-동형
$Y^{(1)}\times_S Y^{(2)}\to Y^{(2)}\times_S Y^{(1)}$을 생각하자
(I, 3.3.5). 이 동형에 의한
$\mathfrak{Tor}_\bullet^S(f_1,f_2;
\mathscr P_\bullet^{(1)},\mathscr P_\bullet^{(2)})$의 상은
$\mathfrak{Tor}_\bullet^S(f_2,f_1;
\mathscr P_\bullet^{(2)},\mathscr P_\bullet^{(1)})$이다. 이 문제는
국소적이므로 (\hyperref[III.6.6.2-fr]{6.6.2})에서 다룬 경우로
환원된다. $M_{\bullet\bullet\bullet}^{(i)}$를
$L_{\bullet\bullet}^{(i)}$의 사영 카르탕--아일렌베르크 분해라
하면 ($i=1,2$), 위 동형은
$M_{\bullet\bullet\bullet}^{(1)}\otimes_A
M_{\bullet\bullet\bullet}^{(2)}$를
$M_{\bullet\bullet\bullet}^{(2)}\otimes_A
M_{\bullet\bullet\bullet}^{(1)}$로 옮긴다. 이 세겹복합체들에 결부된
단일 복합체의 호몰로지를 취하면 주장이 따른다.
\end{remark}

\subsection{전역 초Tor의 결합법칙 스펙트럼 열}
\label{subsection:III.6.8.fr}

\begin{env}[6.8.1]
\label{III.6.8.1-fr}
가정과 표기는 (\hyperref[III.6.7.11-fr]{6.7.11})과 같다고 하자.
특히 $\mathscr P_\bullet^{(i)}$들은 아래로 유계라고 가정한다.
지표 집합 $I$의 분할 $(I_j)_{j\in J}$가 주어졌다고 하자. 전역 초Tor
$\mathfrak{Tor}_n^S((f_i)_{i\in I};
(\mathscr P_\bullet^{(i)})_{i\in I})$와 각각의 “부분” 초Tor
\[
  \mathcal T_{\bullet j}
  =\mathfrak{Tor}_\bullet^S
  ((f_i)_{i\in I_j};(\mathscr P_\bullet^{(i)})_{i\in I_j}).
\]
사이의 “결합법칙” 관계를 제시하고자 한다.

\oldpage[III]{165}

표기를 단순하게 하기 위하여 $I$가 구간 $1\leq i\leq m$이고, 분할
$(I_j)$가 두 구간 $\{1,2,\ldots,r\}$와
$\{r+1,\ldots,m\}$로 이루어진 경우만 다루겠다.
\end{env}

\begin{proposition}[6.8.2]
\label{III.6.8.2-fr}
다음과 같이 나타내는 표준 쌍정칙 스펙트럼 함자
\[
  {}^{(e)}\mathscr E^S
  (f_1,\ldots,f_m;
   \mathscr P_\bullet^{(1)},\ldots,\mathscr P_\bullet^{(m)})
\]
가 존재한다. 이를 “결합법칙 스펙트럼 함자”라 한다. 간단히
$ {}^{(e)}\mathscr E(f_1,\ldots,f_m;
\mathscr P_\bullet^{(1)},\ldots,\mathscr P_\bullet^{(m)})$로도 쓴다.
그 수렴대상은
$\mathfrak{Tor}_\bullet^S(f_1,\ldots,f_m;
\mathscr P_\bullet^{(1)},\ldots,\mathscr P_\bullet^{(m)})$이고,
$E_2$-항은 다음과 같다.
\[
\begin{split}
  {}^{(e)}\mathscr E_{pq}^2
  =\bigoplus_{q_1+q_2=q}
  \mathscr{T}\!or_p^S\Bigl(
  &\mathfrak{Tor}_{q_1}^S
  (f_1,\ldots,f_r;
   \mathscr P_\bullet^{(1)},\ldots,\mathscr P_\bullet^{(r)}),\\
  &\mathfrak{Tor}_{q_2}^S
  (f_{r+1},\ldots,f_m;
   \mathscr P_\bullet^{(r+1)},\ldots,\mathscr P_\bullet^{(m)})
  \Bigr).
\end{split}
\]
\end{proposition}

\begin{proof}
이 명제에서는 $Y$를 $Z^{(1)}\times_S Z^{(2)}$와 표준적으로
동일시하였다. 여기서
$Z^{(1)}=Y^{(1)}\times_S Y^{(2)}\times_S\cdots\times_S Y^{(r)}$이고
$Z^{(2)}=Y^{(r+1)}\times_S\cdots\times_S Y^{(m)}$이다. $S$와
$Y^{(i)}$들이 아핀인 경우만 다루겠다. 이 특수한 경우에서 일반적인
경우로는 (\hyperref[III.6.7.1-fr]{6.7.1})과
(\hyperref[III.6.7.2-fr]{6.7.2})에서 전개한 방법으로 이행하며,
어려움이 없는 자세한 논증은 독자에게 맡긴다. 따라서 다음 따름정리를
증명하겠다.
\end{proof}

\begin{corollary}[6.8.3]
\label{III.6.8.3-fr}
$A$를 환, $S=\operatorname{Spec}(A)$라 하고, $X^{(i)}$
($1\leq i\leq m$)를 준콤팩트 $S$-스킴이라 하자. 각 $i$에 대하여
$\mathscr P_\bullet^{(i)}$를 아래로 유계인 준연접
$\mathscr O_{X^{(i)}}$-가군 복합체라 하자. 수렴대상이
\[
  \mathbf{Tor}_\bullet^S
  (X^{(1)},\ldots,X^{(m)};
   \mathscr P_\bullet^{(1)},\ldots,\mathscr P_\bullet^{(m)})
\]
이고 $E_2$-항이 다음과 같은 표준 쌍정칙 스펙트럼 함자가 존재한다.
\[
\begin{split}
  {}^{(e)}E_{pq}^2
  =\bigoplus_{q_1+q_2=q}
  \operatorname{Tor}_p^A\Bigl(
  &\mathbf{Tor}_{q_1}^S
  (X^{(1)},\ldots,X^{(r)};
   \mathscr P_\bullet^{(1)},\ldots,\mathscr P_\bullet^{(r)}),\\
  &\mathbf{Tor}_{q_2}^S
  (X^{(r+1)},\ldots,X^{(m)};
   \mathscr P_\bullet^{(r+1)},\ldots,\mathscr P_\bullet^{(m)})
  \Bigr).
\end{split}
\]
\end{corollary}

\begin{proof}
(\hyperref[III.6.6.2-fr]{6.6.2})의 정의에 따라 이 초Tor를 계산하려면
각 $i$에 대하여 $X^{(i)}$의 유한 아핀 열린 덮개
$\mathfrak U^{(i)}$를 택하고, 이중복합체
$L_{\bullet\bullet}^{(i)}
=C^\bullet(\mathfrak U^{(i)},\mathscr P_\bullet^{(i)})$와, 각각에 대한
($0_{\mathrm I}$, 11.7.1)의 의미에서의 사영 카르탕--아일렌베르크
분해 $M_{\bullet\bullet\bullet}^{(i)}$를 생각한다. 이 세겹복합체들의 텐서곱
\[
  M_{\bullet\bullet\bullet}
  =\bigotimes_{i=1}^m M_{\bullet\bullet\bullet}^{(i)}
\]
을 취하고 $M_{\bullet\bullet\bullet}$의 호몰로지를 취한다.
$M_{\bullet\bullet\bullet}$을 두 단일 복합체
\[
  N'_\bullet=\bigotimes_{i=1}^r M_{\bullet\bullet\bullet}^{(i)},
  \qquad
  N''_\bullet=\bigotimes_{i=r+1}^m M_{\bullet\bullet\bullet}^{(i)},
\]
의 텐서곱인 단일 복합체 $N_\bullet$로 보자. 여기서
$N'_\bullet$과 $N''_\bullet$은
$M_{\bullet\bullet\bullet}^{(i)}$들의 전차수의 합으로 등급화된다.
더욱이 복합체 $N'_\bullet$과 $N''_\bullet$의 $A$-가군들은
사영이다. 따라서 (\hyperref[III.6.5.9-fr]{6.5.9})에 의하여
\[
  H_\bullet(M_{\bullet\bullet\bullet})
  =\mathbf{Tor}_\bullet^A(N'_\bullet,N''_\bullet);
\]
를 얻는다. 구하는 스펙트럼 열은 바로 복합체 $N'_\bullet$과
$N''_\bullet$에 적용한 (6.3.2.2)의 열이다. 이때 앞에서 얻은 이
복합체들의 호몰로지 가군에 대한 해석을 사용한다. 즉 처음의 주석을
각 부분곱
$X^{(1)}\times_S\cdots\times_S X^{(r)}$와
$X^{(r+1)}\times_S\cdots\times_S X^{(m)}$에 적용한다.
끝으로 정칙성은 (\hyperref[III.6.3.2-fr]{6.3.2})와 다음 사실에서
따른다. 복합체 $\mathscr P_\bullet^{(i)}$들이 아래로 유계이므로
$M_{\bullet\bullet\bullet}^{(i)}$들도 아래로 유계이고, 따라서
$N'_\bullet$과 $N''_\bullet$도 아래로 유계이다.
\end{proof}

\oldpage[III]{166}
\subsection{전역 초Tor의 밑변경 스펙트럼 열}
\label{subsection:III.6.9.fr}

\begin{env}[6.9.1]
\label{III.6.9.1-fr}
가정과 표기는 계속 (\hyperref[III.6.7.11-fr]{6.7.11})과 같다고
하자. 특히 $\mathscr P_\bullet^{(i)}$들은 아래로 유계라고 가정한다.
준스킴의 사상 $g:S'\to S$를 생각하고
$Y'^{(i)}=Y_{(S')}^{(i)}$, $X'^{(i)}=X_{(S')}^{(i)}$,
$\mathscr P_\bullet'^{(i)}
=\mathscr P_\bullet^{(i)}
\otimes_{\mathscr O_S}\mathscr O_{S'}$
라 두자. 따라서 $\mathscr P_\bullet'^{(i)}$는 준연접
$\mathscr O_{X'^{(i)}}$-가군의 복합체이다. 또한
$f'_i=(f_i)_{(S')}:X'^{(i)}\to Y'^{(i)}$라 두자. 이는 분리이고
준콤팩트인 사상이다 (I, 5.5.1 및 6.6.4). 이제
$Y'=Y\times_S S'=Y_{(S')}$라 할 때 다음 두 준연접
$\mathscr O_{Y'}$-가군 사이의 관계를 연구하고자 한다.
$\mathfrak{Tor}_n^{S'}((f'_i)_{i\in I};
(\mathscr P_\bullet'^{(i)})_{i\in I})$ 및
$\mathfrak{Tor}_n^S((f_i)_{i\in I};
(\mathscr P_\bullet^{(i)})_{i\in I})
\otimes_{\mathscr O_S}\mathscr O_{S'}$.
특히 단순한 경우는 다음과 같다. $I$가 한 원소만 가지고
$\mathscr P_\bullet$가 한 가군에만 집중되면
(\hyperref[III.1.4.15-fr]{1.4.15})로 환원된다.
\end{env}

\begin{proposition}[6.9.2]
\label{III.6.9.2-fr}
사상 $g:S'\to S$가 평탄이면, 복합체
$\mathscr P_\bullet^{(i)}$들에 대한 $\partial$-함자의 표준 동형
\[
  \mathfrak{Tor}_\bullet^S
  ((f_i)_{i\in I};(\mathscr P_\bullet^{(i)})_{i\in I})
  \otimes_{\mathscr O_S}\mathscr O_{S'}
  \xrightarrow{\sim}
  \mathfrak{Tor}_\bullet^{S'}
  ((f'_i)_{i\in I};(\mathscr P_\bullet'^{(i)})_{i\in I}).
  \tag{6.9.2.1}\label{III.6.9.2.1.fr}
\]
이 성립한다.
\end{proposition}

\begin{proof}
붙이기는 (\hyperref[III.6.7.1-fr]{6.7.1})과
(\hyperref[III.6.7.2-fr]{6.7.2})의 방법으로 이루어지므로, 다시
$S$, $S'$, $Y^{(i)}$들이 아핀인 경우만 다루면 된다.
$S=\operatorname{Spec}(A)$, $S'=\operatorname{Spec}(A')$라 하고,
각 $i$에 대하여 $X^{(i)}$의 아핀 열린 덮개
$\mathfrak U^{(i)}$를 택하자. $u_i:X'^{(i)}\to X^{(i)}$가 표준
사영이면 $u_i^{-1}(\mathfrak U^{(i)})$는 $X'^{(i)}$의 아핀 열린
덮개이다 (II, 1.5.5). 이를 $\mathfrak U'^{(i)}$로 나타내자.
그러면 분명히
\[
  C^\bullet(\mathfrak U'^{(i)},\mathscr P_\bullet'^{(i)})
  =C^\bullet(\mathfrak U^{(i)},\mathscr P_\bullet^{(i)})\otimes_A A',
\]
동형 (6.9.2.1)의 존재는 이제 곧바로 따른다. 실제로
$M_{\bullet\bullet\bullet}^{(i)}$가
$L_{\bullet\bullet}^{(i)}
=C^\bullet(\mathfrak U^{(i)},\mathscr P_\bullet^{(i)})$의
사영 카르탕--아일렌베르크 분해이고 그 성분들이 $A$-가군이라면
($0_{\mathrm I}$, 11.7.1의 의미에서), $A'$이 평탄 $A$-가군이라는 가정에 의하여
$M_{\bullet\bullet\bullet}^{(i)}\otimes_A A'$는
$L_{\bullet\bullet}^{(i)}\otimes_A A'$의 사영
카르탕--아일렌베르크 분해이며 그 성분들은 $A'$-가군이다. 같은 가정에
의하여 또한
\[
  H_\bullet\left(
    \bigotimes_{i=1}^m
    (M_{\bullet\bullet\bullet}^{(i)}\otimes_A A')
  \right)
  =H_\bullet\left(
    \bigotimes_{i=1}^m M_{\bullet\bullet\bullet}^{(i)}
  \right)\otimes_A A'.
\]
\end{proof}

\begin{env}[6.9.2.2]
\label{III.6.9.2.2-fr}
$I$가 원소 $1$ 하나만 가지면 공식 (6.9.2.1)은
(\hyperref[III.6.7.7-fr]{6.7.7})에서 직접 따른다는 점에 유의하자.
실제로 거기에서 $Y_2=X_2=S'$, $f_2=1_{S'}$로 두고,
$\mathscr P_\bullet^{(2)}$를 차수 $0$의 항
$\mathscr O_{S'}$에만 집중된 복합체로 택한다. 그러면
초코호몰로지 $\mathscr H^n(1_{S'},\mathscr O_{S'})$는
$n\ne0$이면 영이고, $n=0$이면 $\mathscr O_{S'}$이다
(\hyperref[III.6.2.1-fr]{6.2.1}).
\end{env}

일반적인 경우에는 $\mathscr O_{S'}$ 대신 아래로 유계인 준연접
$\mathscr O_{S'}$-가군 복합체 $\mathscr Q'_\bullet$을 도입한다.
표기를 단순하게 하기 위하여 $I=\{1,2,\ldots,m\}$이라 하면
$\partial$-함자
\[
  \mathfrak{Tor}_\bullet^S
  (f_1,\ldots,f_m,1_{S'};
   \mathscr P_\bullet^{(1)},\ldots,\mathscr P_\bullet^{(m)},
   \mathscr Q'_\bullet).
\]
를 생각할 수 있다.

\begin{proposition}[6.9.3]
\label{III.6.9.3-fr}
$t=e,f,f'$에 대하여
$ {}^{(t)}\mathscr E(f_1,\ldots,f_m;
\mathscr P_\bullet^{(1)},\ldots,\mathscr P_\bullet^{(m)},
\mathscr Q'_\bullet)$로 나타내는 세 표준 쌍정칙 스펙트럼 함자가
존재한다. 그 공통 수렴대상은
$\mathfrak{Tor}_\bullet^S(f_1,\ldots,f_m,1_{S'};
\mathscr P_\bullet^{(1)},\ldots,\mathscr P_\bullet^{(m)},
\mathscr Q'_\bullet)$이고, 그 $E_2$-항들은 각각 다음과 같다.

\oldpage[III]{167}
\[
\begin{split}
  {}^{(e)}\mathscr E_{pq}^2
  =\bigoplus_{q'+q''=q}
  \mathscr{T}\!or_p^S\Bigl(
  &\mathfrak{Tor}_{q'}^S
  (f_1,\ldots,f_m;
   \mathscr P_\bullet^{(1)},\ldots,\mathscr P_\bullet^{(m)}),\\
  &\mathscr H_{q''}(\mathscr Q'_\bullet)
  \Bigr),
\end{split}
\]
\[
\begin{split}
  {}^{(f)}\mathscr E_{pq}^2
  =\bigoplus_{q_1+q_2+\cdots+q_{m+1}=q}
  \mathscr{T}\!or_p^{S'}\Bigl(
  &\mathfrak{Tor}_{q_1}^S
  (f_1,1_{S'};\mathscr P_\bullet^{(1)},\mathscr O_{S'}),\ldots,\\
  &\mathfrak{Tor}_{q_m}^S
  (f_m,1_{S'};\mathscr P_\bullet^{(m)},\mathscr O_{S'}),
  \mathscr H_{q_{m+1}}(\mathscr Q'_\bullet)
  \Bigr),
\end{split}
\]
\[
\begin{split}
  {}^{(f')}\mathscr E_{pq}^2
  =\bigoplus_{q_1+\cdots+q_m=q}
  \mathfrak{Tor}_p^{S'}\Bigl(
  &f'_1,\ldots,f'_m,1_{S'};\\
  &\mathscr{T}\!or_{q_1}^S
  (\mathscr P_\bullet^{(1)},\mathscr O_{S'}),\ldots,
  \mathscr{T}\!or_{q_m}^S
  (\mathscr P_\bullet^{(m)},\mathscr O_{S'}),
  \mathscr Q'_\bullet
  \Bigr).
\end{split}
\]
\end{proposition}

\begin{proof}
열 $(e)$는 $r=m$으로 둔
(\hyperref[III.6.8.2-fr]{6.8.2})의 결합법칙 열이다. 다른 두
스펙트럼 열을 정의할 때에도 $S$, $S'$, $Y^{(i)}$들이 아핀인
경우만 다루면 된다. 일반적인 경우로의 이행은
(\hyperref[III.6.7.1-fr]{6.7.1})과
(\hyperref[III.6.7.2-fr]{6.7.2})의 방법으로 이루어지며, 그 내용은
독자에게 맡긴다. 따라서 다음 따름정리를 증명하겠다.
\end{proof}

\begin{corollary}[6.9.4]
\label{III.6.9.4-fr}
$A$를 환, $A'$을 $A$-대수라 하고,
$S=\operatorname{Spec}(A)$, $S'=\operatorname{Spec}(A')$라 하자.
$X^{(i)}$ ($1\leq i\leq m$)를 준콤팩트 $S$-스킴들이라 하고, 각
$i$에 대하여
$X'^{(i)}=X_{(S')}^{(i)}$라 두자. 이는 준콤팩트 $S'$-스킴이다.
각 $i$에 대하여 $\mathscr P_\bullet^{(i)}$를 준연접
$\mathscr O_{X^{(i)}}$-가군 복합체라 하고, 끝으로 $Q'_\bullet$을
$A'$-가군 복합체라 하자. 이 복합체들은 모두 아래로 유계라고
가정한다. $\mathscr P_\bullet^{(i)}$들과 $Q'_\bullet$에 관하여
쌍정칙인 스펙트럼 함자 세 개가 존재한다. 그 공통 수렴대상은
\[
  \mathbf{Tor}_\bullet^S
  (X^{(1)},\ldots,X^{(m)},S';
   \mathscr P_\bullet^{(1)},\ldots,\mathscr P_\bullet^{(m)},
   \widetilde{Q'_\bullet})
\]
이고, 그 $E_2$-항들은 각각 다음과 같다.
\[
\begin{split}
  {}^{(e)}E_{pq}^2
  =\bigoplus_{q'+q''=q}
  \operatorname{Tor}_p^A\Bigl(
  &\mathbf{Tor}_{q'}^S
  (X^{(1)},\ldots,X^{(m)};
   \mathscr P_\bullet^{(1)},\ldots,\mathscr P_\bullet^{(m)}),\\
  &H_{q''}(Q'_\bullet)
  \Bigr),
\end{split}
\]
\[
\begin{split}
  {}^{(f)}E_{pq}^2
  =\bigoplus_{q_1+q_2+\cdots+q_{m+1}=q}
  \operatorname{Tor}_p^{A'}\Bigl(
  &\mathbf{Tor}_{q_1}^S
  (X^{(1)},S';\mathscr P_\bullet^{(1)},\mathscr O_{S'}),\ldots,\\
  &\mathbf{Tor}_{q_m}^S
  (X^{(m)},S';\mathscr P_\bullet^{(m)},\mathscr O_{S'}),
  H_{q_{m+1}}(Q'_\bullet)
  \Bigr),
\end{split}
\]
\[
\begin{split}
  {}^{(f')}E_{pq}^2
  =\bigoplus_{q_1+\cdots+q_m=q}
  \mathbf{Tor}_p^{S'}\Bigl(
  &X'^{(1)},\ldots,X'^{(m)},S';\\
  &\mathscr{T}\!or_{q_1}^S
  (\mathscr P_\bullet^{(1)},\mathscr O_{S'}),\ldots,
  \mathscr{T}\!or_{q_m}^S
  (\mathscr P_\bullet^{(m)},\mathscr O_{S'}),
  \widetilde{Q'_\bullet}
  \Bigr).
\end{split}
\]
\end{corollary}

\begin{proof}
이 스펙트럼 함자들 가운데 첫 번째 것은
(\hyperref[III.6.8.3-fr]{6.8.3})에서 다루었으며, 여기에는
참고를 위하여 포함했을 뿐이므로 다시 논하지 않겠다. 나머지 두
함자를 정의하기 위하여 각 $i$에 대하여 $X^{(i)}$의 유한 아핀
열린 덮개 $\mathfrak U^{(i)}$를 생각하자. 또한
$u_i:X'^{(i)}\to X^{(i)}$가 표준 사영이면, 이에 대응하는 유한
아핀 열린 덮개를
$\mathfrak U'^{(i)}=u_i^{-1}(\mathfrak U^{(i)})$라 하자.

(\hyperref[III.6.6.6-fr]{6.6.6})에 의하여
$\mathbf{Tor}_\bullet^S(X^{(1)},\ldots,X^{(m)},S';
\mathscr P_\bullet^{(1)},\ldots,\mathscr P_\bullet^{(m)},
\widetilde{Q'_\bullet})$는 다음과 같이 얻어진다. 각
$1\leq i\leq m$에 대하여
$L_{\bullet\bullet}^{(i)}=C^\bullet(\mathfrak U^{(i)},
\mathscr P_\bullet^{(i)})$의 사영 카르탕--아일렌베르크 분해
$M_{\bullet\bullet\bullet}^{(i)}$를 택하고
($0_{\mathrm I}$, 11.7.1의 의미에서), 세겹복합체
\[
  M_{\bullet\bullet\bullet}
  =M_{\bullet\bullet\bullet}^{(1)}\otimes_A
   M_{\bullet\bullet\bullet}^{(2)}\otimes_A\cdots\otimes_A
   M_{\bullet\bullet\bullet}^{(m)}\otimes_A Q'_\bullet
\]
를 생각하여 그 호몰로지를 취한다. 여기서 $Q'_\bullet$은 마지막
두 차수가 $0$에만 놓이는 세겹복합체로 본다.
$M_{\bullet\bullet\bullet}'^{(i)}
=M_{\bullet\bullet\bullet}^{(i)}\otimes_A A'$라 두면,
$Q'_\bullet$이 $A'$-가군 복합체라는 점에 유의하여
\[
  M_{\bullet\bullet\bullet}
  =M_{\bullet\bullet\bullet}'^{(1)}\otimes_{A'}
   M_{\bullet\bullet\bullet}'^{(2)}\otimes_{A'}\cdots\otimes_{A'}
   M_{\bullet\bullet\bullet}'^{(m)}\otimes_{A'} Q'_\bullet.
\]
를 얻는다. 이제 각각의 복합체
$M_{\bullet\bullet\bullet}'^{(i)}$를 전차수에 관한 단일 복합체로
보자. 이 복합체는 사영 $A'$-가군들로 이루어져 있다.
(\hyperref[III.6.3.7-fr]{6.3.7})을 임의 개수의 복합체로 확장하여
적용하면 $H_\bullet(M_{\bullet\bullet\bullet})$은 또한
\[
  \operatorname{Tor}_\bullet^{A'}
  (M_{\bullet\bullet\bullet}'^{(1)},\ldots,
   M_{\bullet\bullet\bullet}'^{(m)},Q'_\bullet);
\]
과 같다. 따라서 이는 (\hyperref[III.6.5.15-fr]{6.5.15})에 의하여
원하는 정칙성을 가지며 $E_2$-항이
\[
  E_{pq}^2
  =\bigoplus_{q_1+\cdots+q_{m+1}=q}
  \operatorname{Tor}_p^{A'}
  (H_{q_1}(M_{\bullet\bullet\bullet}'^{(1)}),\ldots,
   H_{q_m}(M_{\bullet\bullet\bullet}'^{(m)}),
   H_{q_{m+1}}(Q'_\bullet)).
\]
인 스펙트럼 열의 수렴대상이다. 여기서
$\mathscr P_\bullet^{(i)}$가 아래로 유계이면
$M_{\bullet\bullet\bullet}^{(i)}$의 세 차수도 모두 아래로 유계이다.

\oldpage[III]{168}
한편 전역 초Tor의 정의에 의하여
\[
  H_{q_i}(M_{\bullet\bullet\bullet}'^{(i)})
  =H_{q_i}(M_{\bullet\bullet\bullet}^{(i)}\otimes_A A')
  =\mathbf{Tor}_{q_i}^S
  (X^{(i)},S';\mathscr P_\bullet^{(i)},\mathscr O_{S'})
\]
이므로, 이로부터 원하는 열 $(f)$를 얻는다. 다른 한편
$M_{\bullet\bullet\bullet}'^{(i)}$를 다음과 같은 이중복합체로 볼 수
있다. 그 첫째 차수는 세겹복합체
$M_{\bullet\bullet\bullet}^{(i)}$의 첫째 차수와 둘째 차수의 합이고,
둘째 차수는 그 세겹복합체의 셋째 차수이다.
$M_{\bullet\bullet\bullet}'^{(i)}$를 이루는 가군들은 사영
$A'$-가군들이므로, 초호몰로지의 일반 이론에 의하여 이중복합체
\[
  M_{\bullet\bullet\bullet}'^{(1)}\otimes_{A'}
  M_{\bullet\bullet\bullet}'^{(2)}\otimes_{A'}\cdots\otimes_{A'}
  M_{\bullet\bullet\bullet}'^{(m)}\otimes_{A'} Q'_\bullet
\]
의 호몰로지는 그 초호몰로지와 표준적으로 동형이다
(\hyperref[0.11.6.5-fr]{0, 11.6.5}). 따라서 이는 $E_2$-항이
\[
  E_{pq}^2
  =\bigoplus_{q_1+\cdots+q_{m+1}=q}
  \operatorname{Tor}_p^{A'}
  (H_{q_1}^{\mathrm{II}}(M_{\bullet\bullet\bullet}'^{(1)}),\ldots,
   H_{q_m}^{\mathrm{II}}(M_{\bullet\bullet\bullet}'^{(m)}),
   H_{q_{m+1}}^{\mathrm{II}}(Q'_\bullet)).
\]
인 스펙트럼 열의 수렴대상이다. 그런데 $Q'_\bullet$의 둘째 차수는
$0$에만 놓이므로 $n\ne0$이면
$H_n^{\mathrm{II}}(Q'_\bullet)=0$이고
$H_0^{\mathrm{II}}(Q'_\bullet)=Q'_\bullet$이다. 따라서 앞의 공식은
\[
  E_{pq}^2
  =\bigoplus_{q_1+\cdots+q_m=q}
  \operatorname{Tor}_p^{A'}
  (H_{q_1}^{\mathrm{II}}(M_{\bullet\bullet\bullet}'^{(1)}),\ldots,
   H_{q_m}^{\mathrm{II}}(M_{\bullet\bullet\bullet}'^{(m)}),Q'_\bullet).
\]
로도 쓸 수 있다. 또한
\[
  H_{q_i}^{\mathrm{II}}(M_{\bullet\bullet\bullet}'^{(i)})
  =H_{q_i}^{\mathrm{II}}(M_{\bullet\bullet\bullet}^{(i)}\otimes_A A')
  =\operatorname{Tor}_{q_i}^A(L_{\bullet\bullet}^{(i)},A')
\]
이 성립한다. 이는 (\hyperref[III.6.3.4-fr]{6.3.4})에서 따른다.
그런데
$L_{-j,k}^{(i)}=C^j(\mathfrak U^{(i)},\mathscr P_k^{(i)})$는
$V$가 덮개 $\mathfrak U^{(i)}$의 $j+1$개 집합의 아핀 교집합들을
달릴 때 $\Gamma(V,\mathscr P_k^{(i)})$들의 직합이다.
$V'=u_i^{-1}(V)$라 하면 $V'$은 $X'^{(i)}$에서 아핀이며,
(\hyperref[III.6.4.1.1.fr]{6.4.1.1})에 의하여
\[
  \Gamma\bigl(V',\mathscr{T}\!or_{q_i}^S
  (\mathscr P_k^{(i)},\mathscr O_{S'})\bigr)
  =\operatorname{Tor}_{q_i}^A(\Gamma(V,\mathscr P_k^{(i)}),A')
\]
을 얻는다. 따라서 이중복합체
$H_{q_i}^{\mathrm{II}}(M_{\bullet\bullet\bullet}'^{(i)})$는
\[
  C^\bullet\bigl(\mathfrak U'^{(i)},
  \mathscr{T}\!or_{q_i}^S(\mathscr P_\bullet^{(i)},\mathscr O_{S'})\bigr),
\]
로 주어진다. 이로써 마침내 열 $(f')$의 $E_2$-항에 대하여 원하는
표현을 얻는다. 이 열이 주어진 조건 아래에서 쌍정칙이라는 사실은
통상적인 방법으로 확인된다. 이때 $\mathscr P_\bullet^{(i)}$가
아래로 유계이면
$M_{\bullet\bullet\bullet}'^{(i)}$의 모든 차수도 아래로 유계라는
점을 고려하면 된다.
\end{proof}

\begin{remark}[6.9.5]
\label{III.6.9.5-fr}
(\hyperref[III.6.7.6-fr]{6.7.6})에서와 같이,
$\mathscr P_\bullet^{(i)}$들과 $\mathscr Q'_\bullet$을 각각 이들과
호모토피 동치인 복합체들로 바꾸어도 열 $(e)$, $(f)$, $(f')$는
표준 동형까지는 변하지 않음을 알 수 있다. 또한 열 $(f)$에 대해서는
호몰로지에서 동형을 유도하는 복합체의 준동형들
$\mathscr P_\bullet^{(i)}\to\mathscr R_\bullet^{(i)}$,
$\mathscr Q'_\bullet\to\mathscr T'_\bullet$, 즉
$\mathscr H_\bullet(\mathscr P_\bullet^{(i)})\xrightarrow{\sim}
\mathscr H_\bullet(\mathscr R_\bullet^{(i)})$,
$\mathscr H_\bullet(\mathscr Q'_\bullet)\xrightarrow{\sim}
\mathscr H_\bullet(\mathscr T'_\bullet)$
을 주는 준동형들이 스펙트럼 열의 동형
\[
\begin{split}
  {}^{(f)}\mathscr E
  (f_1,\ldots,f_m;
   \mathscr P_\bullet^{(1)},\ldots,\mathscr P_\bullet^{(m)},
   \mathscr Q'_\bullet)
  \xrightarrow{\sim}\;&
  {}^{(f)}\mathscr E
  (f_1,\ldots,f_m;\\
  &\mathscr R_\bullet^{(1)},\ldots,\mathscr R_\bullet^{(m)},
   \mathscr T'_\bullet);
\end{split}
\]
을 준다. 증명은 (\hyperref[III.6.7.6-fr]{6.7.6})의 결과와 열 $(f)$의
정칙성을 고려하면 (\hyperref[III.6.7.6-fr]{6.7.6})의 증명과 같다.
\end{remark}

\begin{corollary}[6.9.6]
\label{III.6.9.6-fr}
(\hyperref[subsection:III.6.9.fr]{6.9.1})의 조건 아래 다음을
가정하자.
\begin{enumerate}
  \item[$1^\circ$] 복합체 $\mathscr P_\bullet^{(i)}$들은 $S$ 위
  평탄한 가군들로 이루어지고, $\mathscr O_Y$-가군
  \[
    \mathfrak{Tor}_n^S
    (f_1,\ldots,f_m;
     \mathscr P_\bullet^{(1)},\ldots,\mathscr P_\bullet^{(m)})
  \]
  들은 $S$ 위 평탄하다.
  \item[$2^\circ$] $\mathscr P_\bullet^{(i)}$들과
  $\mathscr Q'_\bullet$은 아래로 유계이다.
\end{enumerate}

\oldpage[III]{169}
그러면
$\mathscr P_\bullet'^{(i)}
=\mathscr P_\bullet^{(i)}\otimes_{\mathscr O_S}\mathscr O_{S'}$라
두었을 때 다음 함자적인 표준 동형들이 성립한다.
\[
\begin{split}
  &\mathfrak{Tor}_n^{S'}
  (f'_1,\ldots,f'_m,1_{S'};
   \mathscr P_\bullet'^{(1)},\ldots,
   \mathscr P_\bullet'^{(m)},\mathscr Q'_\bullet)
  \xrightarrow{\sim}\\
  &\qquad
  \bigoplus_{n'+n''=n}
  \mathfrak{Tor}_{n'}^S
  (f_1,\ldots,f_m;
   \mathscr P_\bullet^{(1)},\ldots,
   \mathscr P_\bullet^{(m)})
  \otimes_{\mathscr O_S}\mathscr H_{n''}(\mathscr Q'_\bullet).
\end{split}
\tag{6.9.6.1}\label{III.6.9.6.1.fr}
\]
특히 $\mathscr Q'_\bullet$이 차수 $0$의 한 항
$\mathscr F'$에만 집중되어 있으면 다음 함자적인 표준 동형들이
성립한다.
\[
\begin{split}
  &\mathfrak{Tor}_n^{S'}
  (f'_1,\ldots,f'_m,1_{S'};
   \mathscr P_\bullet'^{(1)},\ldots,
   \mathscr P_\bullet'^{(m)},\mathscr F')
  \xrightarrow{\sim}\\
  &\qquad
  \mathfrak{Tor}_n^S
  (f_1,\ldots,f_m;
   \mathscr P_\bullet^{(1)},\ldots,
   \mathscr P_\bullet^{(m)})
  \otimes_{\mathscr O_S}\mathscr F'
\end{split}
\tag{6.9.6.2}\label{III.6.9.6.2.fr}
\]
더욱 특별히 $\mathscr F'=\mathscr O_{S'}$이면
\[
\begin{split}
  &\mathfrak{Tor}_n^{S'}
  (f'_1,\ldots,f'_m;
   \mathscr P_\bullet'^{(1)},\ldots,
   \mathscr P_\bullet'^{(m)})
  \xrightarrow{\sim}\\
  &\qquad
  \mathfrak{Tor}_n^S
  (f_1,\ldots,f_m;
   \mathscr P_\bullet^{(1)},\ldots,
   \mathscr P_\bullet^{(m)})
  \otimes_{\mathscr O_S}\mathscr O_{S'}.
\end{split}
\tag{6.9.6.3}\label{III.6.9.6.3.fr}
\]
\end{corollary}

\begin{proof}
$\mathscr P_\bullet^{(i)}$를 이루는 가군들에 대한 평탄성 가정에
의하여 $q_i\ne0$이면 복합체
$\mathscr{T}\!or_{q_i}^S
(\mathscr P_\bullet^{(i)},\mathscr O_{S'})$는 영이다
(\hyperref[III.6.5.8-fr]{6.5.8}). 따라서 열 $(f')$는 퇴화한다.
또한 가정 $2^\circ$에 의하여 이 열은 쌍정칙이므로
(\hyperref[III.6.9.3-fr]{6.9.3}), 가장자리 준동형
\[
\begin{split}
  &\mathfrak{Tor}_n^S
  (f_1,\ldots,f_m,1_{S'};
   \mathscr P_\bullet^{(1)},\ldots,
   \mathscr P_\bullet^{(m)},\mathscr Q'_\bullet)
  \longrightarrow {}^{(f')}\mathscr E_{n0}^2\\
  &\qquad=
  \mathfrak{Tor}_n^{S'}
  (f'_1,\ldots,f'_m,1_{S'};
   \mathscr P_\bullet'^{(1)},\ldots,
   \mathscr P_\bullet'^{(m)},\mathscr Q'_\bullet)
\end{split}
\tag{6.9.6.4}\label{III.6.9.6.4.fr}
\]
은 전단사이다 ($0_{\mathrm I}$, 11.1.6).
$\mathscr O_Y$-가군
\[
  \mathfrak{Tor}_n^S
  (f_1,\ldots,f_m;
   \mathscr P_\bullet^{(1)},\ldots,\mathscr P_\bullet^{(m)})
\]
들에 대한 평탄성 가정에 의하여 $p\ne0$이면
${}^{(e)}\mathscr E_{pq}^2=0$이다
(\hyperref[III.6.5.8-fr]{6.5.8}). 따라서 열 $(e)$도 퇴화하고,
쌍정칙이므로 가장자리 준동형
\[
\begin{split}
  &\mathfrak{Tor}_n^S
  (f_1,\ldots,f_m,1_{S'};
   \mathscr P_\bullet^{(1)},\ldots,
   \mathscr P_\bullet^{(m)},\mathscr Q'_\bullet)
  \longrightarrow {}^{(e)}\mathscr E_{0n}^2\\
  &\qquad=
  \bigoplus_{n'+n''=n}
  \mathfrak{Tor}_{n'}^S
  (f_1,\ldots,f_m;
   \mathscr P_\bullet^{(1)},\ldots,
   \mathscr P_\bullet^{(m)})
  \otimes_{\mathscr O_S}\mathscr H_{n''}(\mathscr Q'_\bullet)
\end{split}
\tag{6.9.6.5}\label{III.6.9.6.5.fr}
\]
은 전단사이다 ($0_{\mathrm I}$, 11.1.6). 앞의 두 동형을 합성하면
동형 (6.9.6.1)을 얻는다. 이제
$q\ne0$이면 $\mathscr H_q(\mathscr Q'_\bullet)=0$이고
$\mathscr H_0(\mathscr Q'_\bullet)=\mathscr F'$이므로 동형
(6.9.6.2)는 여기서 곧바로 따른다. 끝으로 (6.9.6.2)에서
$\mathscr F'=\mathscr O_{S'}$인 경우는
(\hyperref[III.6.7.12-fr]{6.7.12})를 고려하면 동형 (6.9.6.3)을
준다.
\end{proof}

\begin{corollary}[6.9.7]
\label{III.6.9.7-fr}
(\hyperref[subsection:III.6.9.fr]{6.9.1})의 조건 아래 $S$와 $S'$이
아핀이라고 하고, 각 $i$에 대하여 정수 $d_i$ ($1\leq i\leq m$)가
주어졌다고 하자. 그러면 오직 $S$, $X^{(i)}$들 및 $d_i$들에만 의존하는
정수 $N$이 존재하여 다음 성질을 가진다. 임의의 정수 $n_0$와 다음
조건을 만족하는 모든 복합체들 $\mathscr P_\bullet^{(i)}$에 대하여,
$n\leq n_0$이면 (\hyperref[III.6.9.6.3.fr]{6.9.6.3})의 표준 동형들이
성립한다.
\begin{enumerate}
  \item[$1^\circ$] $k<d_i$이면 $\mathscr P_k^{(i)}=0$이다.
  \item[$2^\circ$] $k<n_0+N$이면 $\mathscr P_k^{(i)}$는 $S$ 위
  평탄하다.
  \item[$3^\circ$]
  $\mathfrak{Tor}_q^S(f_1,\ldots,f_m;
  \mathscr P_\bullet^{(1)},\ldots,\mathscr P_\bullet^{(m)})$는
  $q<n_0+N$이면 $S$ 위 평탄하다.
\end{enumerate}
\end{corollary}

\begin{proof}
$k<r$이면 $\mathscr P_k^{(i)}$가 $S$ 위 평탄하다고 하자. 그러면
$k<r$이고 $q_i\ne0$이면
$\mathscr{T}\!or_{q_i}^S(\mathscr P_k^{(i)},\mathscr O_{S'})=0$이다.
이제 다음을 계산하자.
\[
\begin{split}
  \mathfrak{Tor}_p^{S'}\Bigl(
  &f'_1,\ldots,f'_m;
  \mathscr{T}\!or_{q_1}^S(\mathscr P_\bullet^{(1)},\mathscr O_{S'}),
  \ldots,\\
  &\mathscr{T}\!or_{q_m}^S(\mathscr P_\bullet^{(m)},\mathscr O_{S'})
  \Bigr).
\end{split}
\tag{6.9.7.1}\label{III.6.9.7.1.fr}
\]

\oldpage[III]{170}
이를 (\hyperref[III.6.6.2-fr]{6.6.2})의 방법으로, $X^{(i)}$의 고정된
아핀 덮개 $\mathfrak U^{(i)}$ ($1\leq i\leq m$)의 역상을 사용하여
계산한다. 이 덮개들은 $S'$과 $\mathscr P_\bullet^{(i)}$들에 무관하게
고정되어 있다. 항 $L_{jk}^{(i)}$는 $j<-N_i$이면 영이며, $N_i$는
$\mathfrak U^{(i)}$에만 의존한다. 어떤 $q_i$가 영이 아니면, 그
$p$차 호몰로지가 (\hyperref[III.6.9.7.1.fr]{6.9.7.1})인 단일 복합체의
항들은 다음보다 작은 모든 차수에서 영이다.
\[
  <r+\sum_{j\ne i}d_j-\sum_{i=1}^m N_i,
\]
따라서 (\hyperref[III.6.9.7.1.fr]{6.9.7.1})은 $p<r-N$이면 영이다.
여기서 $N$은 다음 수들 가운데 가장 큰 수를 뜻한다.
\[
  \sum_{i=1}^m N_i-\sum_{j\ne i}d_j.
\]
그러므로 $q\ne0$이고 $p<r-N$이면
${}^{(f')}\mathscr E_{pq}^2=0$이다. 다른 한편 $q<0$이면
${}^{(f')}\mathscr E_{pq}^2=0$이므로, 가장자리 준동형
(\hyperref[III.6.9.6.4.fr]{6.9.6.4})는 $n<r-N$이면 전단사이다
(M, XV, 5.6) (이때 $\mathscr Q'_\bullet=\mathscr O_{S'}$이다).
둘째로,
\[
  \mathfrak{Tor}_q^S
  (f_1,\ldots,f_m;
   \mathscr P_\bullet^{(1)},\ldots,\mathscr P_\bullet^{(m)})
\]
가 $q<r$이면 $S$ 위 평탄하다고 하자. 그러면 $p\ne0$이고 $q<r$이면
${}^{(e)}\mathscr E_{pq}^2=0$이다. 한편 $p<0$이면
${}^{(e)}\mathscr E_{pq}^2=0$이므로, 가장자리 준동형
(\hyperref[III.6.9.6.5.fr]{6.9.6.5})는 $n<r$이면 전단사이다. 이로써
증명이 끝난다.
\end{proof}


응용에서 (\hyperref[III.6.9.3-fr]{6.9.3})의 가장 중요한 경우는
$m=1$이고 $\mathscr Q'_\bullet$이 차수 $0$의 한 항 $\mathscr F'$에만
집중되어 있는 경우이다. 뒤의 참조를 위하여 이 경우를 다시
진술한다\footnote{(\hyperref[III.6.9.8-fr]{6.9.8})에서 다룬 경우,
특히 스펙트럼 열 (\hyperref[III.6.9.8.3.fr]{6.9.8.3})과
(\hyperref[III.6.9.8.4.fr]{6.9.8.4})는 J.-P. Serre가 1957년에 우리에게
알려 주었다.}:

\begin{proposition}[6.9.8]
\label{III.6.9.8-fr}
$S$를 준스킴, $g:S'\to S$를 사상,
$f:X\to Y$를 $S$-준스킴의 분리 준콤팩트 $S$-사상이라 하자.
$\mathscr P_\bullet$을 아래로 유계인 준연접 $\mathscr O_X$-가군의
복합체, $\mathscr F'$을 준연접 $\mathscr O_{S'}$-가군이라 하자.
$\mathscr P_\bullet$과 $\mathscr F'$에 관하여 쌍정칙이고 준연접
$\mathscr O_{Y_{(S')}}$-가군의 범주에 값을 가지며 공통 수렴대상
$\mathfrak{Tor}_\bullet^S(f,1_{S'};\mathscr P_\bullet,\mathscr F')$을
갖는 두 스펙트럼 함자가 존재한다. 그 $E_2$-항은 각각
\[
  {}'\mathscr E_{pq}^2
  =\mathscr{T}\!or_p^S
  (\mathscr H^{-q}(f,\mathscr P_\bullet),\mathscr F')
  \tag{6.9.8.1}\label{III.6.9.8.1.fr}
\]
과
\[
  {}''\mathscr E_{pq}^2
  =\mathscr H^{-p}
  (f',\mathscr{T}\!or_q^S(\mathscr P_\bullet,\mathscr F')),
  \tag{6.9.8.2}\label{III.6.9.8.2.fr}
\]
이다. 여기서 $f'=f_{(S')}:X_{(S')}\to Y_{(S')}$이다.
\end{proposition}

\begin{proof}
문제의 열들은 (\hyperref[III.6.9.3-fr]{6.9.3})에서 출발하지 않고도,
$X^{(1)}=X$, $Y^{(1)}=Y$, $X^{(2)}=Y^{(2)}=S'$, $f_1=f$,
$f_2=1_{S'}$로 두었을 때 (\hyperref[III.6.7.3-fr]{6.7.3})의 열
$(a)$와 $(b')$에서 얻을 수 있다. $S=S'=Y$이고 $Y$가 아핀이면,
$E_2$-항이 각각
\[
  {}'\mathscr E_{pq}^2
  =\mathscr{T}\!or_p^Y
  (\mathscr H^{-q}(f,\mathscr P_\bullet),\mathscr F)
  \tag{6.9.8.3}\label{III.6.9.8.3.fr}
\]
과
\[
  {}''\mathscr E_{pq}^2
  =\mathscr H^{-p}
  (f,\mathscr{T}\!or_q^Y(\mathscr P_\bullet,\mathscr F))
  \tag{6.9.8.4}\label{III.6.9.8.4.fr}
\]
인 두 스펙트럼 열을 얻는다. 준연접이고 $Y$ 위 평탄한 모든
$\mathscr O_Y$-가군 $\mathscr F$에 대하여, 이 두 열은
(\hyperref[III.6.7.6-fr]{6.7.6})에 의하여 $\mathscr O_X$-가군 복합체
$\mathscr P_\bullet\otimes_Y\mathscr F$에 관한 함자 $f_*$의
초코호몰로지
$\mathscr H^\bullet(f,\mathscr P_\bullet\otimes_Y\mathscr F)$를
수렴대상으로 한다. $\mathscr P_\bullet$이 $Y$ 위 평탄한
$\mathscr O_X$-가군들로 이루어져 있으면 여기서 $\mathscr F$은 임의의
준연접 $\mathscr O_Y$-가군이어도 된다. 이 열들은
(\hyperref[III.6.2.1-fr]{6.2.1})의 열들과는 다르다.
\end{proof}

\begin{corollary}[6.9.9]
\label{III.6.9.9-fr}
(\hyperref[III.6.9.8-fr]{6.9.8})의 조건 아래, 복합체
$\mathscr P_\bullet$이 아래로 유계이고 $S$ 위 평탄한 가군들로
이루어지며 $\mathscr O_Y$-가군
$\mathscr H^n(f,\mathscr P_\bullet)$이 $S$ 위 평탄하다고 가정하자.

\oldpage[III]{171}
그러면 다음 함자적인 표준 동형들이 성립한다.
\[
  \mathfrak{Tor}_n^{S'}
  (f',1_{S'};\mathscr P'_\bullet,\mathscr F')
  \xrightarrow{\sim}
  \mathscr H^{-n}(f,\mathscr P_\bullet)
  \otimes_{\mathscr O_S}\mathscr F'
  \tag{6.9.9.1}\label{III.6.9.9.1.fr}
\]
여기서
$\mathscr P'_\bullet=\mathscr P_\bullet\otimes_{\mathscr O_S}
\mathscr O_{S'}$이다. 특히 $\mathscr F'=\mathscr O_{S'}$이면 다음
함자적인 표준 동형들이 성립한다.
\[
  \mathscr H^n(f',\mathscr P'_\bullet)
  \xrightarrow{\sim}
  \mathscr H^n(f,\mathscr P_\bullet)
  \otimes_{\mathscr O_S}\mathscr O_{S'}.
  \tag{6.9.9.2}\label{III.6.9.9.2.fr}
\]
이는 (\hyperref[III.6.9.6-fr]{6.9.6})에서 $m=1$인 특수한 경우이다.
더 특별히 다음이 성립한다.
\end{corollary}

\begin{corollary}[6.9.10]
\label{III.6.9.10-fr}
$S$를 준스킴, $f:X\to Y$를 $S$-준스킴의 분리 준콤팩트
$S$-사상이라 하고, $\mathscr P_\bullet$을 $S$ 위 평탄한 준연접
$\mathscr O_X$-가군들로 이루어진 아래로 유계인 복합체라 하자.
또 $\mathscr O_Y$-가군 $\mathscr H^n(f,\mathscr P_\bullet)$이
$S$ 위 평탄하다고 가정하자. 각 $s\in S$에 대하여 $X_s$와 $Y_s$를
각각 섬유 $X\otimes_S k(s)$와 $Y\otimes_S k(s)$라 하고,
$f_s:X_s\to Y_s$를 사상 $f\times_S1$,
$\mathscr P_\bullet^s$를 $\mathscr O_{X_s}$-가군의 복합체
$\mathscr P_\bullet\otimes_{\mathscr O_S}k(s)$라 하자. 그러면 다음
함자적인 표준 동형들이 성립한다.
\[
  \mathscr H^n(f_s,\mathscr P_\bullet^s)
  \xrightarrow{\sim}
  \mathscr H^n(f,\mathscr P_\bullet)\otimes_{\mathscr O_S}k(s).
  \tag{6.9.10.1}\label{III.6.9.10.1.fr}
\]
\end{corollary}

따라서 적절한 평탄성 가정 아래 유도함자
$R^nf_*(\mathscr F)$의 형성이 “섬유로의 이행과 가환하는” 경우를
하나 얻는다. 이는 §~7에서 다른 방법으로 다시 보게 될 것이다.

\subsection{일부 코호몰로지 함자의 국소 구조.}

\begin{proposition}[6.10.1]
\label{III.6.10.1-fr}
$S=\operatorname{Spec}(A)$를 아핀 스킴이라 하고,
$Y^{(i)}$ ($1\leq i\leq n$)를 $S$ 위 평탄한 아핀 $S$-스킴의
유한족이라 하자. 각 $i$에 대하여
$f_i:X^{(i)}\to Y^{(i)}$를 분리 준콤팩트 $S$-사상이라 하고,
$\mathscr P_\bullet^{(i)}$를 아래로 유계인 준연접
$\mathscr O_{X^{(i)}}$-가군 복합체라 하자. $Y$를 $S$-스킴
$Y^{(i)}$들의 곱이라 하자. $S$ 위에서 평탄한 준연접
$\mathscr O_Y$-가군들로 이루어진 복합체 $\mathscr R_\bullet$이
존재하여 다음 성질을 갖는다. 모든 아핀 $S$-스킴 $S'$과 아래로
유계인 준연접 $\mathscr O_{S'}$-가군 복합체
$\mathscr Q'_\bullet$에 대하여 동형
\[
  \mathbf{Tor}_\bullet^S
  (f_1,\ldots,f_n,1_{S'};
   \mathscr P_\bullet^{(1)},\ldots,\mathscr P_\bullet^{(n)},
   \mathscr Q'_\bullet)
  \xrightarrow{\sim}
  \mathscr H_\bullet
  (\mathscr R_\bullet\otimes_{\mathscr O_S}\mathscr Q'_\bullet)
  \tag{6.10.1.1}\label{III.6.10.1.1.fr}
\]
이 성립하며, 이는 $\mathscr Q'_\bullet$에 관한 $\partial$-함자들의
동형이다. 또한 모든 아핀 $S$-스킴 사이의 $S$-사상
$u:S''\to S'$에 대하여 도식
\[
\xymatrix@C=12pt@R=30pt{
  {\begin{gathered}
    \mathbf{Tor}_\bullet^S
    (f_1,\ldots,f_n,1_{S'};\\
    \mathscr P_\bullet^{(1)},\ldots,\mathscr P_\bullet^{(n)},
    \mathscr Q'_\bullet)
  \end{gathered}}
  \ar[r]^-{\sim}\ar[d]
  &
  {\mathscr H_\bullet
    (\mathscr R_\bullet\otimes_{\mathscr O_S}\mathscr Q'_\bullet)}
  \ar[d]
  \\
  {\begin{gathered}
    \mathbf{Tor}_\bullet^S
    (f_1,\ldots,f_n,1_{S''};\\
    \mathscr P_\bullet^{(1)},\ldots,\mathscr P_\bullet^{(n)},
    u^*(\mathscr Q'_\bullet))
  \end{gathered}}
  \ar[r]^-{\sim}
  &
  {\mathscr H_\bullet
    (\mathscr R_\bullet\otimes_{\mathscr O_S}u^*(\mathscr Q'_\bullet))}
}
\tag{6.10.1.2}\label{III.6.10.1.2.fr}
\]
은 가환한다. 여기서 수직 화살표들은
(\hyperref[III.6.7.10-fr]{6.7.10})에서 정의한
$(1_Y\times_Su)$-표준 사상들이다.
\end{proposition}

\begin{proof}
\oldpage[III]{172}
(\hyperref[III.6.6.6-fr]{6.6.6})의 주석을 고려하여
(\hyperref[III.6.6.2-fr]{6.6.2})의 방법으로 초Tor를 계산하자.
(\hyperref[III.6.6.2-fr]{6.6.2})의 표기를 쓰면, 각
$L_{\bullet\bullet}^{(i)}$는 $A_i$-가군의 이중복합체이다. 여기서
$A_i$는 $Y^{(i)}$의 환이다. 따라서 이것은 $A_i$-가군들로 이루어진
사영 카르탕--아일렌베르크 분해
$M_{\bullet\bullet\bullet}^{(i)}$를 갖는다
($0_{\mathrm I}$, 11.7.1의 의미에서). 또한
(\hyperref[III.6.6.6-fr]{6.6.6})에 의하여
(\hyperref[III.6.10.1.1.fr]{6.10.1.1})의 좌변은 다음과
표준적으로 동형이다.
\[
  \mathscr H_\bullet
  (M_{\bullet\bullet\bullet}\otimes_A Q'_\bullet),
\]
여기서
\[
  M_{\bullet\bullet\bullet}
  =
  M_{\bullet\bullet\bullet}^{(1)}
  \otimes_A M_{\bullet\bullet\bullet}^{(2)}
  \otimes_A\cdots\otimes_A M_{\bullet\bullet\bullet}^{(n)}
\]
이고 $\mathscr Q'_\bullet=\widetilde{Q'_\bullet}$이다. 가정에 의하여
$A_i$들은 평탄 $A$-가군이므로,
$M_{\bullet\bullet\bullet}^{(i)}$들은 평탄 $A$-가군의
세겹복합체이고 ($0_{\mathrm I}$, 6.2.1),
$M_{\bullet\bullet\bullet}$도 마찬가지이다. 또한 $B$가 $Y$의
환, 곧 $A_i$들의 텐서곱이면, $M_{\bullet\bullet\bullet}$은
$B$-가군의 세겹복합체이다. 따라서 $\mathscr O_Y$-가군 복합체
\[
  \mathscr R_\bullet=(M_{\bullet\bullet\bullet})^\sim
\]
은 요구를 충족한다. 여기서 $M_{\bullet\bullet\bullet}$은 단일
복합체로 보며, 이 결론은
(\hyperref[III.6.7.10-fr]{6.7.10})에서 곧바로 따른다.
\end{proof}

\begin{corollary}[6.10.2]
\label{III.6.10.2-fr}
(\hyperref[III.6.10.1-fr]{6.10.1})의 진술에서
$\mathscr R_\bullet$이 아래로 유계라고 가정할 수 있다.
$\mathscr P_\bullet^{(i)}$들이 위로 유계이고 $Y^{(i)}$들의
코호몰로지 차원이 유한하면, $\mathscr R_\bullet$이 위로
유계라고 가정할 수 있다.
\end{corollary}

\begin{proof}
첫째 주장은 각 $M_{\bullet\bullet\bullet}^{(i)}$의 세 차수가 모두
아래로 유계라는 사실에서 따른다. 다른 한편 환 $A_i$들의
코호몰로지 차원이 유한하면 각 $M_{\bullet\bullet\bullet}^{(i)}$의
셋째 차수는 유한개의 값만 취하고, 그 첫째 차수도 구성에 의하여
마찬가지이다 (\hyperref[III.6.6.2-fr]{6.6.2}). 또한
$\mathscr P_\bullet^{(i)}$의 차수가 위로 유계이면 그 둘째 차수도
위로 유계이므로 (\hyperref[III.6.6.2-fr]{6.6.2}), 둘째 주장도
곧바로 따른다.
\end{proof}

\begin{remarks}[6.10.3]
\label{III.6.10.3-fr}
\textnormal{(i)} (\hyperref[III.6.10.1-fr]{6.10.1})의 표기에서
$\mathscr R_\bullet$은 $S$ 위 평탄한 $\mathscr O_Y$-가군들로
이루어져 있으므로 (\hyperref[III.6.5.9-fr]{6.5.9}),
$\mathscr H_\bullet(\mathscr R_\bullet\otimes_S\mathscr Q'_\bullet)$은
$\mathbf{Tor}_\bullet(\mathscr R_\bullet,\mathscr Q'_\bullet)$과
동형이다. 따라서 이는 (\hyperref[III.6.5.4-fr]{6.5.4})에 의하여
$E_2$-항이 다음과 같은 정칙 스펙트럼 열의 수렴대상이다.
\[
  \mathscr E_{pq}^2
  =
  \bigoplus_{q'+q''=q}
  \mathbf{Tor}_p^S
  \bigl(
    \mathscr H_{q'}(\mathscr R_\bullet),
    \mathscr H_{q''}(\mathscr Q'_\bullet)
  \bigr)
  \tag{6.10.3.1}\label{III.6.10.3.1.fr}
\]
이는 다름 아닌 밑변경의 스펙트럼 열
$(e)$ (\hyperref[III.6.9.3-fr]{6.9.3})이다.

\medskip
\textnormal{(ii)} $\mathscr R'_\bullet$을 $S$ 위 평탄한 준연접
$\mathscr O_Y$-가군들로 이루어진 둘째 복합체라 하고,
$g:\mathscr R'_\bullet\to\mathscr R_\bullet$을 다음 사상이 동형이
되게 하는 복합체 준동형이라 하자.
\[
  \mathscr H_\bullet(g):
  \mathscr H_\bullet(\mathscr R'_\bullet)
  \longrightarrow
  \mathscr H_\bullet(\mathscr R_\bullet)
\]
그러면 (\hyperref[III.6.3.3-fr]{6.3.3})과
(\hyperref[III.6.5.9-fr]{6.5.9})에 의하여 $g$에서
$\mathscr Q'_\bullet$에 관한 $\partial$-함자들의 동형
\[
  \mathscr H_\bullet
  (\mathscr R'_\bullet\otimes_S\mathscr Q'_\bullet)
  \xrightarrow{\sim}
  \mathscr H_\bullet
  (\mathscr R_\bullet\otimes_S\mathscr Q'_\bullet)
\]
을 얻으며, 도식
\[
\xymatrix@C=28pt@R=30pt{
  \mathscr H_\bullet
  (\mathscr R'_\bullet\otimes_S\mathscr Q'_\bullet)
  \ar[r]^-{\sim}\ar[d]
  &
  \mathscr H_\bullet
  (\mathscr R_\bullet\otimes_S\mathscr Q'_\bullet)
  \ar[d]
  \\
  \mathscr H_\bullet
  (\mathscr R'_\bullet\otimes_Su^*(\mathscr Q'_\bullet))
  \ar[r]^-{\sim}
  &
  \mathscr H_\bullet
  (\mathscr R_\bullet\otimes_Su^*(\mathscr Q'_\bullet))
}
\]
\par\medskip\noindent
은 가환한다. 따라서 복합체 $\mathscr R_\bullet$은
(\hyperref[III.6.10.1-fr]{6.10.1})의 성질들에 의하여 완전히
결정되는 것은 아니다.

\medskip
\oldpage[III]{173}
\textnormal{(iii)} (\hyperref[III.6.10.1-fr]{6.10.1})의 증명에서
$M_{\bullet\bullet\bullet}^{(i)}$들이 자유 $A_i$-가군들로
이루어졌다고 가정할 수 있다. 이는 ($0_{\mathrm I}$, 11.5.2.1)의
“쌍대화된” 증명에서 쉽게 따른다. 그러면
\[
  M_{\bullet\bullet\bullet}^{(i)}\otimes_{A_i}B
\]
는 자유 $B$-가군들로 이루어진다. 또한
$M_{\bullet\bullet\bullet}$은 이들의 $B$ 위 텐서곱과 같으므로,
(\hyperref[III.6.10.1-fr]{6.10.1})에서 더 나아가
$\mathscr R_\bullet$을 자유 $B$-가군 복합체에 결부된 것으로
택할 수 있다. 또한 (M, XVII, 1.2)에 의하여 세겹복합체
$M_{\bullet\bullet\bullet}$은 각 이중복합체
$L_{\bullet\bullet}^{(i)}$에 함자적으로 의존한다. 따라서 각
$X^{(i)}$의 유한 덮개를 고정하면, 각
$\mathscr P_\bullet^{(i)}$에도 함자적으로 의존한다. 다만 여기서
세겹복합체의 “사상”은 호모토피 관계에 따른 준동형들의 동치류로 이해해야
한다. 한편 $X^{(i)}$의 덮개를 더 세분된 덮개로 바꾸면
$L_{\bullet\bullet}^{(i)}$에 대하여 정확히 호모토피를 제외하고
정의되는 준동형들이 생기므로
(\hyperref[III.6.6.8-fr]{6.6.8}), 앞의 사상에 관한 약속 아래
세겹복합체 $M_{\bullet\bullet\bullet}$은 마침내 각
$\mathscr P_\bullet^{(i)}$에 관한 함자가 된다. 우리는 이 함자적
의존성, 특히 복합체의 완전열
\[
  0\longrightarrow\mathscr P_\bullet^{\prime(i)}
  \longrightarrow\mathscr P_\bullet^{(i)}
  \longrightarrow\mathscr P_\bullet^{\prime\prime(i)}
  \longrightarrow0
\]
에 대한 $\mathscr R_\bullet$의 거동을
(\hyperref[III.6.1.3-fr]{6.1.3})에서 언급한 코호몰로지 함자들의
일반 대수학을 다루는 장에서 명확히 할 것이다.
\end{remarks}

\begin{scholium}[6.10.4]
\label{III.6.10.4-fr}
$\mathscr R_\bullet$이 $S$ 위 평탄한 $\mathscr O_Y$-가군들로
이루어져 있다는 사실에서
\[
  \mathscr H_\bullet
  (\mathscr R_\bullet\otimes_S\mathscr Q'_\bullet)
\]
이 $\mathscr Q'_\bullet$에 관한 호몰로지 함자임이 쉽게 따른다
((7.7.1)의 논증을 보라). 바로 이 성질이
(\hyperref[III.6.1.1-fr]{6.1.1})에서 언급했듯이 초Tor를 도입한
동기이다. 실제로 다음과 같이 두자.
\[
  X=X^{(1)}\times_SX^{(2)}\times_S\cdots\times_SX^{(n)},
  \qquad
  f=f_1\times_Sf_2\times_S\cdots\times_Sf_n,
\]
\[
  \mathscr P_\bullet
  =
  \mathscr P_\bullet^{(1)}
  \otimes_S\mathscr P_\bullet^{(2)}
  \otimes_S\cdots\otimes_S\mathscr P_\bullet^{(n)},
\]
\[
  X'=X\times_SS',
  \qquad
  Y'=Y\times_SS',
  \qquad
  f'=f\times_S1_{S'}.
\]
밑변경 문제는 준연접 $\mathscr O_{S'}$-가군 $\mathscr N'$에 관한
함자로서 초코호몰로지
\[
  \mathscr H^\bullet_{f'}
  (\mathscr P_\bullet\otimes_S\mathscr N')
\]
를 연구하게 하고, 또는 준연접 $\mathscr O_S$-가군 $\mathscr N$에
관한 함자로서 초코호몰로지
\[
  \mathscr H^\bullet_f
  (\mathscr P_\bullet\otimes_S\mathscr N)
\]
를 연구하게 한다. $\mathscr P_\bullet^{(i)}$들이, 따라서
$\mathscr P_\bullet$도, $S$ 위 평탄하면 앞의 결과와
(\hyperref[III.6.7.6-fr]{6.7.6})에 의하여 이 함자는 실제로
$\mathscr N$에 관한 코호몰로지 함자이다. 그러나
$\mathscr P_\bullet^{(i)}$들에 관한 평탄성 가정을 더 이상 하지
않으면 그렇지 않으며, 이때는
\[
  \mathscr H^\bullet_f
  (\mathscr P_\bullet\otimes_S\mathscr N)
\]
의 연구에 통상적인 호몰로지 대수학의 방법을 더 이상 적용할 수
없다.

다만 우리는 주로 $n=1$, $Y=S$이고 $\mathscr P_\bullet$이 $Y$ 위
평탄한 $\mathscr O_X$-가군들로 이루어진 경우를 사용한다. 이 경우
다음 정리를 얻는다.
\end{scholium}

\begin{theorem}[6.10.5]
\label{III.6.10.5-fr}
$Y=\operatorname{Spec}(A)$를 아핀 뇌터 스킴,
$f:X\to Y$를 고유 사상, $\mathscr P_\bullet$을 $Y$ 위 평탄한 연접
$\mathscr O_X$-가군들로 이루어진 아래로 유계인 복합체라 하자.
그러면 아래로 유계인 $\mathscr O_Y$-가군 복합체
$\mathscr L_\bullet$이 존재한다. 그 각 항 $\mathscr L_i$는
$\mathscr O_Y^{m_i}$ 꼴이며, 아래로 유계인 준연접
$\mathscr O_Y$-가군 복합체 $\mathscr Q_\bullet$에 관한
$\partial$-함자들의 동형
\[
  \mathscr H^\bullet
  (f,\mathscr P_\bullet\otimes_Y\mathscr Q_\bullet)
  \xrightarrow{\sim}
  \mathscr H_\bullet
  (\mathscr L_\bullet\otimes_Y\mathscr Q_\bullet)
  \tag{6.10.5.1}\label{III.6.10.5.1.fr}
\]
이 성립한다. 또한 모든 사상 $u:Y'\to Y$에 대하여 다음과 같이
두자.
\[
  X'=X_{(Y')},
  \qquad
  f'=f_{(Y')},
  \qquad
  \mathscr P'_\bullet
  =
  \mathscr P_\bullet\otimes_Y\mathscr O_{Y'},
  \qquad
  \mathscr L'_\bullet=u^*(\mathscr L_\bullet),
\]
\oldpage[III]{174}
여기서 $\mathscr L'_\bullet$은 유한 계수의 국소 자유
$\mathscr O_{Y'}$-가군들로 이루어진 복합체이다. 그러면 아래로
유계인 준연접 $\mathscr O_{Y'}$-가군 복합체
$\mathscr Q'_\bullet$에 관한 $\partial$-함자들의 동형
\[
  \mathscr H^\bullet
  (f',\mathscr P'_\bullet\otimes_{Y'}\mathscr Q'_\bullet)
  \xrightarrow{\sim}
  \mathscr H_\bullet
  (\mathscr L'_\bullet\otimes_{Y'}\mathscr Q'_\bullet)
  \tag{6.10.5.2}\label{III.6.10.5.2.fr}
\]
이 성립하며, 도식
\[
\xymatrix@C=25pt@R=30pt{
  \mathscr H^\bullet
  (f,\mathscr P_\bullet\otimes_Y\mathscr Q_\bullet)
  \ar[r]^-{\sim}\ar[d]
  &
  \mathscr H_\bullet
  (\mathscr L_\bullet\otimes_Y\mathscr Q_\bullet)
  \ar[d]
  \\
  \mathscr H^\bullet
  (f',\mathscr P'_\bullet\otimes_{Y'}u^*(\mathscr Q_\bullet))
  \ar[r]^-{\sim}
  &
  \mathscr H_\bullet
  (\mathscr L'_\bullet\otimes_{Y'}u^*(\mathscr Q_\bullet))
}
\tag{6.10.5.3}\label{III.6.10.5.3.fr}
\]
은 가환한다.
\end{theorem}

\begin{proof}
(\hyperref[III.6.10.1-fr]{6.10.1})을 적용하면 먼저 준연접
$\mathscr O_Y$-가군 복합체 $\mathscr R_\bullet$을 얻는다. 이 복합체는
(\hyperref[III.6.10.2-fr]{6.10.2})에 따라 아래로 유계이고,
(\hyperref[III.6.10.3-fr]{6.10.3, (iii)})에 따라 $Y$-자유이다. 또한
$\mathscr Q_\bullet$에 관한 $\partial$-함자들의 동형
\[
  \mathscr H^\bullet
  (f,\mathscr P_\bullet\otimes_Y\mathscr Q_\bullet)
  \xrightarrow{\sim}
  \mathscr H_\bullet
  (\mathscr R_\bullet\otimes_Y\mathscr Q_\bullet)
  \tag{6.10.5.4}\label{III.6.10.5.4.fr}
\]
을 얻는다. 다만 \emph{선험적으로} $\mathscr R_\bullet$의 각 항이
유한형 $\mathscr O_Y$-가군일 필요는 없다. 그러나
$\mathscr Q_\bullet$이 한 항 $\mathscr O_Y$만으로 이루어진 복합체인
경우에 (6.10.5.4)를 적용하면
$\mathscr H_\bullet(\mathscr R_\bullet)$이
$\mathscr H^\bullet(f,\mathscr P_\bullet)$과 동형임을 알 수 있다.
따라서 이는 연접 $\mathscr O_Y$-가군들로 이루어진다
(\hyperref[III.6.2.5-fr]{6.2.5}). 이제 (0, 11.9.2)에 의하여,
유한 생성 자유 $A$-가군들에 결부된 $\mathscr O_Y$-가군들로
이루어진 아래로 유계인 복합체 $\mathscr L_\bullet$과 복합체 준동형
\[
  \mathscr L_\bullet\longrightarrow\mathscr R_\bullet,
\]
이 존재하며, 이에 따라 호몰로지에 유도되는 준동형
\[
  \mathscr H_\bullet(\mathscr L_\bullet)
  \longrightarrow
  \mathscr H_\bullet(\mathscr R_\bullet)
\]
은 전단사이다. 그러므로
(\hyperref[III.6.10.3-fr]{6.10.3, (ii)})에 의하여 동형
(6.10.5.1)을 얻는다. $Y'$이 아핀일 때
(\hyperref[III.6.10.5-fr]{6.10.5})의 나머지 주장들은
(\hyperref[III.6.10.1-fr]{6.10.1})과
(\hyperref[III.6.10.3-fr]{6.10.3, (ii)})에서 따른다. 일반적인
경우에는 $Y'$의 아핀 열린집합 덮개 $(V_\alpha)$와 각
$V_\alpha$에 대한 동형 (6.10.5.2)를 생각하자. 아핀 열린집합
$W\subset V_\alpha\cap V_\beta$에 대하여 $V_\alpha$와
$V_\beta$에 대응하는 동형들의 $W$ 위 제한이 모두 $W$에
대응하는 동형과 일치함을 확인하면 충분하다. 이는 표준 포함 사상
$W\to V_\alpha$와 $W\to V_\beta$에
(6.10.1.2)의 도식을 적용했을 때 그 도식이 가환한다는 사실에서
따른다.
\end{proof}

\begin{remark}[6.10.6]
뒤의 장들에서는 주로 $\mathscr P_\bullet$이 $Y$ 위 평탄한 연접
$\mathscr O_X$-가군 하나 $\mathscr F$만으로 이루어진 경우에
(\hyperref[III.6.10.5-fr]{6.10.5})를 적용할 것이다. 이때
\[
  \mathscr H_n(\mathscr L_\bullet)
  =
  \mathscr H^{-n}(f,\mathscr F)
  =
  R^{-n}f_*(\mathscr F)
  \qquad(6.2.1),
\]
이므로 $n>0$이면 $\mathscr H_n(\mathscr L_\bullet)=0$이다. 뒤에서
(7.7.12, (i))에 의하여 $\mathscr L_\bullet$이 차수
$\leqslant0$인 항들만 갖는다고 가정할 수 있음을 보일 것이다.
따라서 그 항의 수는 유한하다. 다만 각 $\mathscr L_i$가 유한 생성
자유 $A$-가군에 결부된다는 가정을, 유한 생성 국소 자유
$\mathscr O_Y$-가군이라는 가정으로 바꾸어야 한다.

이러한 $\mathscr O_X$-가군 $\mathscr F$에 대응하는 복합체
$\mathscr L_\bullet$은
\oldpage[III]{175}
앞의 차수 제한을 제외하면 어떠한 특별한 성질도 갖는 것 같지 않다.
그러면 역으로 다음을 물을 수 있다. 유한 생성 사영 $A$-가군들에
결부된 $\mathscr O_Y$-가군들로 이루어지고, 아래로 유계이며, 차수
$>0$인 항들이 영인 복합체 $\mathscr L_\bullet$이 주어졌다고 하자.
$Y$ 위 사영이고 평탄한 $Y$-스킴 $f:X\to Y$와 국소 자유
$\mathscr O_X$-가군 $\mathscr F$가 존재하여
\[
  \mathscr H^\bullet
  (f,\mathscr F\otimes_Y\mathscr Q_\bullet)
  \xrightarrow{\sim}
  \mathscr H_\bullet
  (\mathscr L_\bullet\otimes_Y\mathscr Q_\bullet)
\]
이 $\mathscr Q_\bullet$에 함자적인 동형이 되는가? 그러한 결과의
의의는 고유 $Y$-스킴 위에서 연접이고 $Y$-평탄한 가군들의
코호몰로지 이론을 뇌터 환 $A$ 위의 유한 생성 사영 $A$-가군
복합체들을 ``호모토피를 제외하고'' 보는 이론으로 완전히 환원하는
데 있다.
\end{remark}

\section{가군의 공변 호몰로지 함자에서의 밑변경 연구}
\label{section:III.7.fr}

\subsection{$A$-가군의 함자}
\label{subsection:III.7.1.fr}

\begin{env}[7.1.1]
\label{III.7.1.1-fr}
$A$를 반드시 가환일 필요는 없는 환이라 하자. 왼쪽 $A$-가군의
범주를 $\mathbf{Ab}_A$로 쓰고, $\mathbf Z$-가군, 즉 아벨군의
범주를 단순히 $\mathbf{Ab}$로 쓰겠다. 공변 가법 함자
\[
  T:\mathbf{Ab}_A\longrightarrow\mathbf{Ab}
\]
와 $(A,A)$-쌍가군 $M$이 주어지면, $T(M)$에는 오른쪽 $A$-가군
구조가 자연스럽게 주어진다. 실제로 모든 $a\in A$에 대하여 왼쪽
$A$-가군 $M$의 자기준동형 $x\mapsto xa$를 $h_{a,M}$ 또는 간단히
$h_a$로 쓰자. 가정에 따라 $T(h_a)$는 $\mathbf Z$-가군 $T(M)$의
자기준동형이다. 또한 $T$가 공변 가법 함자이므로 $a,b\in A$에 대하여
\[
  T(h_{ab})
  =
  T(h_b\circ h_a)
  =
  T(h_b)\circ T(h_a),
  \qquad\text{그리고}\qquad
  T(h_{a+b})
  =
  T(h_a+h_b)
  =
  T(h_a)+T(h_b);
\]
따라서 사상
\[
  (y,a)\longmapsto T(h_a)(y)
\]
은 $T(M)$ 위에 오른쪽 $A$-가군 구조를 주는 외부 연산이다. 특히
$T(A_s)$는 오른쪽 $A$-가군이다.
\end{env}

\begin{env}[7.1.2]
\label{III.7.1.2-fr}
$A$가 가환환이면 모든 $A$-가군 $M$에 $T(M)$의 $A$-가군 구조가
자연스럽게 주어짐이 (\hyperref[III.7.1.1-fr]{7.1.1})에서 따른다.
또한 $u:M\to N$이 $A$-가군의 준동형이면 모든 $a\in A$에 대하여
\[
  u\circ h_{a,M}=h_{a,N}\circ u,
\]
이고, 따라서
\[
  T(u)\circ T(h_{a,M})
  =
  T(h_{a,N})\circ T(u),
\]
그러므로 $T(u):T(M)\to T(N)$은 $A$-가군의 준동형이다. 따라서
$T$를 $\mathbf{Ab}_A$에서 자기 자신으로 가는 공변 가법 함자로 볼
수 있다. 더 정확히 말하면, 이렇게 하여 공변 가법 함자
$\mathbf{Ab}_A\to\mathbf{Ab}$의 범주와 공변 $A$-선형 함자
\[
  T:\mathbf{Ab}_A\longrightarrow\mathbf{Ab}_A,
\]
의 범주 사이에 표준 동치가 정의된다. 여기서 $A$-선형이라는 것은
\[
  T(h_{a,M})=h_{a,T(M)}
  \qquad\text{모든 }a\in A\text{에 대하여}.
\]
가 성립한다는 뜻이다. 각 $A$-가군에 그 밑에 놓인 $\mathbf Z$-가군을
대응시키는 포함 함자
\[
  I:\mathbf{Ab}_A\longrightarrow\mathbf{Ab},
\]
가 완전하고 충실하므로, 앞의 동치로 서로 대응하는 두 함자의
완전성에 관한 성질은 같다.
\end{env}

\begin{env}[7.1.3]
\label{III.7.1.3-fr}
계속해서 $A$가 가환환이라고 가정하자. $B$를 반드시 가환일 필요는
없는 $A$-대수라 하고, $\rho:A\to B$를 이 대수 구조에 대응하는 환
준동형이라 하자. 이 준동형은 스칼라 제한에 의한 공변
\oldpage[III]{176}
가법 함자
\[
  \rho_*:M\longmapsto M_{[\rho]}
\]
를 왼쪽 $B$-가군의 범주 $\mathbf{Ab}_B$에서 왼쪽 $A$-가군의 범주
$\mathbf{Ab}_A$로 정의한다. 이를 $T$와 합성하여 공변 가법 함자
\[
  T_{(B)}:
  \mathbf{Ab}_B
  \xrightarrow{\rho_*}
  \mathbf{Ab}_A
  \xrightarrow{T}
  \mathbf{Ab},
\]
를 얻는다. 활자상의 이유로 이를 $T^{(B)}$ 또는 $T\otimes_A B$로도
쓰며, $T$에서 $A$의 스칼라를 $B$로 확장하여 얻은 함자라고 한다.
물론 $B$가 가환이면 $T_{(B)}$를 $\mathbf{Ab}_B$에서 자기 자신으로
가는 함자로 볼 수 있다
(\hyperref[III.7.1.2-fr]{7.1.2}). $B$가 가환이고 $C$가
$B$-대수이면 곧
\[
  T_{(C)}=(T_{(B)})_{(C)};
\]
임을 알 수 있다. 스칼라 확장은 $T$에 관하여 함자적이고 가법적이다.
또한 $T$가 귀납극한 또는 직합과 가환하면 $T_{(B)}$도 그러하며,
$T$가 왼쪽 완전(각각 오른쪽 완전, 완전)이면 $T_{(B)}$도 같은
성질을 갖는다. 실제로 $\rho_*$는 완전하며 귀납극한과 직합에
가환한다.
\end{env}

\begin{env}[7.1.4]
\label{III.7.1.4-fr}
계속해서 $A$가 가환환이라고 하고,
$T:\mathbf{Ab}_A\to\mathbf{Ab}_A$를 귀납극한과 가환하는 공변 가법
$A$-선형 함자라 하자. 그러면 $A$의 모든 곱셈적 부분집합 $S$와 모든
$A$-가군 $M$에 대하여 $M$에 함자적인 $A$-가군의 표준 동형
\[
  T(S^{-1}M)\xrightarrow{\sim}S^{-1}T(M).
  \tag{7.1.4.1}\label{III.7.1.4.1.fr}
\]
이 성립한다. 실제로 먼저 $S$가 어떤 $f\in A$의 거듭제곱
$f^n$ $(n\geqslant0)$들의 집합이라고 하자. 이때
\[
  M_f=\varinjlim M_n,
\]
임이 알려져 있다. 여기서 $(M_n,\varphi_{nm})$은 $A$-가군
$M_n=M$들로 이루어진 귀납계이고,
\[
  \varphi_{nm}:z\longmapsto f^{\,n-m}z
  \qquad(0_{\mathrm I}, 1.6.1);
\]
이다. 따라서 $T$에 관한 가정으로부터 이 경우의 동형
(7.1.4.1)을 얻는다. 이제 $S$가 임의이면 $S^{-1}M$은 $f\in S$에
대한 $M_f$들의 귀납극한이므로 ($0_{\mathrm I}$, 1.4.5), 같은 방식으로
결론을 얻는다. 또한 동형 (7.1.4.1)의 함자성에 의하여 이는
$S^{-1}A$-가군의 동형이다. 따라서 표준 동형에 따라 다음과 같이
쓸 수 있다.
\[
  T_{(S^{-1}A)}(S^{-1}M)
  =
  S^{-1}T(M)
  =
  T(S^{-1}M).
  \tag{7.1.4.2}\label{III.7.1.4.2.fr}
\]
$S=A-\mathfrak p$가 $A$의 소 아이디얼 $\mathfrak p$의 여집합이면
$T_{(A_{\mathfrak p})}$ 대신 $T_{\mathfrak p}$라고 쓴다.
\end{env}

\begin{proposition}[7.1.5]
\label{III.7.1.5-fr}
(\hyperref[III.7.1.4-fr]{7.1.4})의 가정 아래에서 $A$의 모든 극대
아이디얼 $\mathfrak m$에 대하여 $T_{\mathfrak m}$이 왼쪽 완전
(각각 오른쪽 완전, 완전)이면 $T$도 왼쪽 완전(각각 오른쪽 완전,
완전)이다.
\end{proposition}

\begin{proof}
실제로 $A$-가군 $M$의 두 부분가군 $N$, $P$가 $A$의 모든 극대
아이디얼 $\mathfrak m$에 대하여
$N_{\mathfrak m}=P_{\mathfrak m}$을 만족하면 $N=P$임이 알려져 있다
(Bourbaki, \emph{Alg. comm.}, chap. II, \S~3, no~3, th.~1).
\end{proof}

\subsection{텐서곱 함자의 특징짓기}
\label{subsection:III.7.2.fr}

\begin{env}[7.2.1]
\label{III.7.2.1-fr}
$A$를 반드시 가환일 필요는 없는 환, $M$ (각각 $N$)을 왼쪽
(각각 오른쪽) $A$-가군, $P$를 $\mathbf Z$-가군이라 하자.
$\mathbf Z$-준동형
\[
  v:N\otimes_A M\longrightarrow P
\]
을 주는 것은 다음 조건을 만족하는 $\mathbf Z$-쌍선형 사상
\[
  u:N\times M\longrightarrow P
\]
을 주는 것과 동치임을 상기하자.
\[
  u(ta,x)=u(t,ax)
  \qquad\text{모든 }a\in A,\ t\in N,\ x\in M\text{에 대하여},
\]
두 사상의 관계는 $v(t\otimes x)=u(t,x)$이다. 한편 $u$를 주는 것은
$M$에서 $\operatorname{Hom}_{\mathbf Z}(N,P)$로 가는
\oldpage[III]{177}
$\mathbf Z$-준동형 $x\mapsto f_x$로서
\[
  f_{ax}(t)=f_x(ta)
  \qquad\text{모든 }a\in A,\ t\in N,\ x\in M\text{에 대하여},
\]
를 만족하는 것을 주는 것과 동치이다. 이때 두 사상의 관계는
$u(t,x)=f_x(t)$이다.
\end{env}

\begin{env}[7.2.2]
\label{III.7.2.2-fr}
$T:\mathbf{Ab}_A\to\mathbf{Ab}$를 공변 가법 함자라 하자. 모든 왼쪽
$A$-가군 $M$에 대하여 $M$에 함자적인 $\mathbf Z$-가군의 표준 준동형
\[
  t_M:T(A_s)\otimes_A M\longrightarrow T(M).
  \tag{7.2.2.1}\label{III.7.2.2.1.fr}
\]
을 정의하겠다. (\hyperref[III.7.2.1-fr]{7.2.1})에 의하여 이를 위해서는
$M$에서 $\operatorname{Hom}_{\mathbf Z}(T(A_s),T(M))$로 가는
$\mathbf Z$-준동형 $x\mapsto t'_M(x)$로서
\[
  t'_M(ax)(y)=t'_M(x)(ya)
  \qquad\text{모든 }a\in A,\ x\in M,\ y\in T(A_s)\text{에 대하여}.
\]
를 정의하면 충분하다. 이를 위하여
$\operatorname{Hom}_{\mathbf Z}(T(A_s),T(M))$에는 $T(A_s)$의 오른쪽
$A$-가군 구조에서 유도되는 왼쪽 $A$-가군 구조가 표준적으로
주어짐에 유의하자. $a\in A$,
$v\in\operatorname{Hom}_{\mathbf Z}(T(A_s),T(M))$일 때 그 외부 연산은
\[
  (a\cdot v)(y)=v(ya)
  \qquad\text{모든 }y\in T(A_s)\text{에 대하여}.
\]
로 주어진다. 이제 $t'_M$을 두 표준 준동형의 합성
\[
  M\xrightarrow{\sim}\operatorname{Hom}_A(A_s,M)
  \xrightarrow{T}
  \operatorname{Hom}_{\mathbf Z}(T(A_s),T(M)),
\]
으로 정의한다. 둘째 화살표는 사상 $u\mapsto T(u)$이고, 첫째 화살표는
$\xi\in A$, $x\in M$에 대하여 $\theta_x(\xi)=\xi x$를 만족하는
$A$-가군의 표준 동형 $x\mapsto\theta_x$이다. 이때
$\theta_{ax}=\theta_x\circ h_a$이므로
\[
  T(\theta_{ax})
  =
  T(\theta_x\circ h_a)
  =
  T(\theta_x)\circ T(h_a)
\]
이고, 따라서 $y\in T(A_s)$에 대하여
\[
  T(\theta_{ax})(y)
  =
  T(\theta_x)(T(h_a)(y))
  =
  T(\theta_x)(ya)
\]
이다. 마지막 등식은 $T(A_s)$ 위의 외부 연산의 정의에서 따른다.
이로써 $t_M$의 존재가 증명된다. 이 준동형이 $M$에 함자적임도
곧 확인된다. 즉 왼쪽 $A$-가군의 모든 준동형 $w:M\to M'$에 대하여
도식
\[
\xymatrix@C=30pt@R=30pt{
  T(A_s)\otimes_A M
  \ar[r]^-{t_M}\ar[d]_{1\otimes w}
  &
  T(M)\ar[d]^{T(w)}
  \\
  T(A_s)\otimes_A M'
  \ar[r]_-{t_{M'}}
  &
  T(M')
}
\tag{7.2.2.2}\label{III.7.2.2.2.fr}
\]
은 가환한다.

준동형 (7.2.2.1)의 함자성에 의하여 $A$가 가환이면 이는 $A$-가군의
준동형이다 (\hyperref[III.7.1.2-fr]{7.1.2} 참조).
\end{env}

\begin{env}[7.2.3]
\label{III.7.2.3-fr}
$A$가 가환일 때에는 더 일반적으로 모든 $A$-가군 $N$에 대하여
$A$-가군의 표준 준동형
\[
  T(N)\otimes_A M\longrightarrow T(N\otimes_A M)
  \tag{7.2.3.1}\label{III.7.2.3.1.fr}
\]
을 정의할 수 있다. 이를 위해서는 (\hyperref[III.7.2.2-fr]{7.2.2})의
구성에서 준동형 $\theta_x$를 모든 $y\in N$을 $y\otimes x$에 대응시키는
$A$-가군의 준동형
\[
  N\longrightarrow N\otimes_A M
\]
으로 바꾸면 된다. 이 준동형이 $M$과 $N$에 함자적임은 곧 알 수 있다.

\oldpage[III]{178}
특히 $B$가 $A$-대수일 때($B$가 가환일 필요는 없다), $M$에 함자적인
준동형
\[
  (T(M))_{(B)}
  =
  T(M)\otimes_A B
  \longrightarrow
  T(M\otimes_A B)
  =
  T_{(B)}(M_{(B)})
  \tag{7.2.3.2}\label{III.7.2.3.2.fr}
\]
을 얻는다. (7.2.3.1)이 $M$에 함자적이므로 이는 $B$-가군의
준동형이다.

또한 다음 가환도식을 얻는다.
\[
\xymatrix@C=30pt@R=30pt{
  T(A)\otimes_A M
  \ar[r]^-{t_M}\ar[d]
  &
  T(M)\ar[d]
  \\
  T_{(B)}(B_s)\otimes_B M_{(B)}
  \ar[r]_-{t_{M_{(B)}}}
  &
  T_{(B)}(M_{(B)})
}
\tag{7.2.3.3}\label{III.7.2.3.3.fr}
\]
여기서 오른쪽 세로 화살표는 표준 준동형 다음에 (7.2.3.2)를 취한 합성
준동형
\[
  T(M)
  \longrightarrow
  T(M)\otimes_A B
  \longrightarrow
  T(M\otimes_A B)
  =
  T_{(B)}(M_{(B)})
\]
이다. 한편 (7.2.3.3)의 왼쪽 세로 화살표는 준동형
\[
  T(A)\otimes_A M
  \longrightarrow
  T_{(B)}(B_s)\otimes_B(B\otimes_A M)
  =
  T_{(B)}(B_s)\otimes_A M
\]
이다. 여기서
\[
  T(A)\longrightarrow T_{(B)}(B_s)=T(B)
\]
는 $T(\rho)$이고, $\rho$는 $A$-가군의 준동형
$A\to B_{[\rho]}$로 본 것이다.
\end{env}

\begin{lemma}[7.2.4]
\label{III.7.2.4-fr}
$T$가 $\mathbf{Ab}_A$에서 $\mathbf{Ab}$로 가는 공변 가법 함자이고
직합과 가환하면, 모든 자유 $A$-가군 $L$에 대하여 표준 준동형
$t_L$ (7.2.2.1)은 동형이다.
\end{lemma}

\begin{proof}
실제로
\[
  L=\bigoplus_{\alpha\in I}L_\alpha
\]
이고 모든 $\alpha\in I$에 대하여 $L_\alpha$는 $A_s$와 동형이다.
(\hyperref[III.7.2.2-fr]{7.2.2})에서 준 $t_M$의 정의에 따르면
\[
  t_L=\bigoplus_{\alpha\in I}t_{L_\alpha},
\]
이다. 실제로 $T$에 관한 가정에 의하여
\[
  T:
  \operatorname{Hom}_A(A_s,L)
  \longrightarrow
  \operatorname{Hom}_{\mathbf Z}(T(A_s),T(L))
\]
는 $\mathbf Z$-선형 사상
\[
  T_\alpha:
  \operatorname{Hom}_A(A_s,L_\alpha)
  \longrightarrow
  \operatorname{Hom}_{\mathbf Z}(T(A_s),T(L_\alpha))
\]
들의 직합이다. 따라서 $L=A_s$인 경우에 보조정리를 증명하면
충분하다. 그런데 이때 $t_L$은 모든 오른쪽 $A$-가군에 대하여
성립하는 표준 동형
\[
  T(A_s)\otimes_A A_s\xrightarrow{\sim}T(A_s)
\]
과 다르지 않다.
\end{proof}

\begin{proposition}[7.2.5]
\label{III.7.2.5-fr}
$T$가 $\mathbf{Ab}_A$에서 $\mathbf{Ab}$로 가고 직합과 가환하는 공변
가법 함자라 하자. 다음 조건들은 동치이다.
\begin{enumerate}
  \item[a)] $T$는 오른쪽 완전하다.
  \item[b)] 모든 왼쪽 $A$-가군 $M$에 대하여 표준 준동형 $t_M$
    (7.2.2.1)은 동형이다.
  \item[$b'$)] $T$는 반완전하고, 모든 왼쪽 $A$-가군 $M$에 대하여
    준동형 $t_M$은 전사이다.
  \item[c)] $T$는 $M$에 관하여 $N\otimes_A M$ 꼴인 함자와 동형이다.
    여기서 $N$은 오른쪽 $A$-가군이다.
\end{enumerate}
\end{proposition}

\begin{proof}
$b)$가 $c)$를 함의하고 $c)$가 $a)$를 함의함은 명백하다. $a)$가
$b)$를 함의함을 보이자.
\oldpage[III]{179}
모든 왼쪽 $A$-가군 $M$에 대하여
\[
  T'(M)=T(A_s)\otimes_A M
\]
으로 놓자. $L$과 $L'$이 자유 왼쪽 $A$-가군인 완전열
\[
  L'\longrightarrow L\longrightarrow M\longrightarrow0,
\]
이 존재한다. $T$와 $T'$이 오른쪽 완전하므로 다음 가환도식을 얻는다.
\[
\xymatrix@C=22pt@R=25pt{
  T'(L')\ar[r]\ar[d]^{t_{L'}}
  &
  T'(L)\ar[r]\ar[d]^{t_L}
  &
  T'(M)\ar[r]\ar[d]^{t_M}
  &
  0
  \\
  T(L')\ar[r]
  &
  T(L)\ar[r]
  &
  T(M)\ar[r]
  &
  0
}
\]
여기서 두 행은 완전하다. (\hyperref[III.7.2.4-fr]{7.2.4})에 의하여
$t_L$과 $t_{L'}$은 동형이므로, 5-보조정리에 의하여 $t_M$도
동형이다. 끝으로 $b)$가 $b')$을 함의함은 명백하다. $b')$이 $a)$를
함의함을 보이려면 다음 보조정리를 증명하면 충분하다.
\end{proof}

\begin{lemma}[7.2.5.1]
\label{III.7.2.5.1-fr}
$\mathcal K$와 $\mathcal K'$을 두 아벨 범주, $F$와 $G$를
$\mathcal K$에서 $\mathcal K'$으로 가는 두 공변 가법 함자라 하자.
$f:F\to G$를 $\mathcal K$의 모든 대상 $E$에 대하여
\[
  f_E:F(E)\longrightarrow G(E)
\]
가 에피사상이 되게 하는 함자적 사상 (T, I, 1.2)이라 하자. $F$가
오른쪽 완전하고 $G$가 반완전하면 $G$는 오른쪽 완전하다.
\end{lemma}

\begin{proof}
실제로 $\mathcal K$ 안의 모든 에피사상 $v:E'\to E$에 대하여
\[
  G(v):G(E')\longrightarrow G(E)
\]
가 에피사상임을 보이면 충분하다. 그런데 가환도식
\[
\xymatrix@C=35pt@R=30pt{
  F(E')\ar[r]^{F(v)}\ar[d]_{f_{E'}}
  &
  F(E)\ar[d]^{f_E}
  \\
  G(E')\ar[r]_{G(v)}
  &
  G(E)
}
\]
에서 $F(v)$, $f_{E'}$, $f_E$는 에피사상이므로 $G(v)$도
에피사상이다.
\end{proof}

\begin{remark}[7.2.6]
\label{III.7.2.6-fr}
모든 오른쪽 $A$-가군 $N$에 대하여, 모든 왼쪽 $A$-가군 $M$에 대해
\[
  T_N(M)=N\otimes_A M
\]
으로 놓자. 그러면 $T_N$은 $\mathbf{Ab}_A$에서 $\mathbf{Ab}$로 가며
오른쪽 완전하고 직합과 가환하는 공변 가법 함자이다. $T_N(A_s)$를
$N$과 표준적으로 동일시하면, 이에 대응하는 준동형 (7.2.2.1)은 항등
준동형이 됨을 곧 확인할 수 있다. 따라서 (\hyperref[III.7.2.5-fr]{7.2.5},
$c)$)의 오른쪽 $A$-가군 $N$은 유일한 동형까지 결정되며
$T(A_s)$와 표준적으로 동형이다. 달리 말하면 함자적 대응
\[
  T\longmapsto T(A_s),
  \qquad
  N\longmapsto T_N
\]
은 오른쪽 $A$-가군의 범주와, $\mathbf{Ab}_A\to\mathbf{Ab}$인 공변
가법 함자 가운데 오른쪽 완전하고 직합과 가환하는 것들의 범주 사이의
동치 (T, I, 1.2)를 이룬다.
\end{remark}

\begin{proposition}[7.2.7]
\label{III.7.2.7-fr}
$A$를 그 근기 $\mathfrak m$으로 나눈 몫이 체 $k$인 왼쪽 아르틴
환이라 하자. $T$를 $\mathbf{Ab}_A$에서 $\mathbf{Ab}$로 가고 직합과
가환하는 공변 가법 함자라 하자. 그러면
(\hyperref[III.7.2.5-fr]{7.2.5})의 조건들은 다음 조건과도 동치이다.
\begin{enumerate}
  \item[d)] $T$는 반완전하고, 표준 준동형
    $\varepsilon:A_s\to k$에서 유도되는 준동형
    \[
      T(\varepsilon):T(A_s)\longrightarrow T(k)
    \]
    은 전사이다.
\end{enumerate}
\end{proposition}

\begin{proof}
\oldpage[III]{180}
(\hyperref[III.7.2.5-fr]{7.2.5})의 조건 $b')$이 $d)$를 함의함은
명백하다. $d)$가 $b')$을 함의함을 증명하자.

$\mathfrak m^n=0$인 정수 $n$이 존재한다. 모든 $A$-가군 $M$에 대하여
$M_h=\mathfrak m^hM$으로 놓자. $t_{M_h}$가 전사임을 $h$에 관한
내림귀납법으로 증명하겠다. $h=n$일 때 주장은 명백하다. $h<n$일 때에는
완전열
\[
  0\longrightarrow M_{h+1}
  \longrightarrow M_h
  \longrightarrow M_h/M_{h+1}
  \longrightarrow 0
\]
이 있고, 귀납가정에 의하여 $t_{M_{h+1}}$은 전사이다. 한편
$M_h/M_{h+1}$은 $\mathfrak m$에 의하여 소거되므로
$(A/\mathfrak m)$-가군이며, 다시 말해 $k$와 동형인 $A$-가군들의
직합이다. $T$가 직합과 가환하므로 $t_{M_h/M_{h+1}}$이 전사임을
보이기 위해서는 $t_k$가 전사임을 보이면 충분하다. 그런데 가환도식
\[
\xymatrix@C=35pt@R=30pt{
  T(A_s)\otimes_A A_s
  \ar[r]^-{t_{A_s}}\ar[d]_{1\otimes\varepsilon}
  &
  T(A_s)\ar[d]^{T(\varepsilon)}
  \\
  T(A_s)\otimes_A k
  \ar[r]_-{t_k}
  &
  T(k)
}
\]
의 가환성과 (\hyperref[III.7.2.4-fr]{7.2.4})에 의하여, 가정 $d)$는
$t_k$가 실제로 전사임을 함의한다. 증명을 끝내기 위해서는 $A$-가군의
완전열
\[
  0\longrightarrow M'
  \longrightarrow M
  \xrightarrow{v}M''
  \longrightarrow 0
\]
에서 $t_{M'}$과 $t_{M''}$이 전사이면 $t_M$도 전사임을 증명하면
충분하다. 그런데 가환도식
\[
\xymatrix@C=27pt@R=30pt{
  T'(M')\ar[r]\ar[d]_{t_{M'}}
  &
  T'(M)\ar[r]\ar[d]_{t_M}
  &
  T'(M'')\ar[r]\ar[d]_{t_{M''}}
  &
  0\ar[d]
  \\
  T(M')\ar[r]
  &
  T(M)\ar[r]_-{T(v)}
  &
  T(M'')\ar[r]
  &
  \operatorname{Coker}(T(v))
}
\]
에서 $T$가 반완전이라는 가정에 의하여 두 행은 완전하다. 귀납가정에
의하여 $t_{M'}$과 $t_{M''}$은 에피사상이고 마지막 세로 화살표는
모노사상이므로, 5-보조정리 (M, I, 1.1)에 의하여 $t_M$은
에피사상이다.
\end{proof}

\subsection{가군의 호몰로지 함자에 대한 완전성 판정}
\label{subsection:III.7.3.fr}

\begin{proposition}[7.3.1]
\label{III.7.3.1-fr}
$A$를 환(가환일 필요는 없다), $T_\bullet$을 $\mathbf{Ab}_A$에서
$\mathbf{Ab}$로 가고 직합과 가환하는 공변 호몰로지 함자 (T, II,
2.1)라 하자. $T_p$와 $T_{p-1}$이 정의되는 정수 $p$를 잡자. 다음
조건들은 동치이다.
\begin{enumerate}
  \item[a)] $T_p$는 오른쪽 완전하다.
  \item[b)] $T_{p-1}$은 왼쪽 완전하다.
  \oldpage[III]{181}
  \item[c)] 모든 왼쪽 $A$-가군 $M$에 대하여 표준 함자적 준동형
    (7.2.2.1)
    \[
      T_p(A_s)\otimes_A M\longrightarrow T_p(M)
      \tag{7.3.1.1}\label{III.7.3.1.1.fr}
    \]
    은 동형이다.
  \item[d)] 모든 왼쪽 $A$-가군 $M$에 대하여 준동형 (7.3.1.1)은
    에피사상이다.
  \item[e)] $T_p$는 $N$이 오른쪽 $A$-가군인 함자
    $M\mapsto N\otimes_A M$과 동형이다.
\end{enumerate}
더 나아가 $A$와 $\mathfrak m$이
(\hyperref[III.7.2.7-fr]{7.2.7})의 조건을 만족하면, 앞의 조건들은
다음 조건과도 동치이다.
\begin{enumerate}
  \item[f)] 표준 준동형
    \[
      T_p(\varepsilon):T_p(A_s)\longrightarrow T_p(k)
    \]
    은 에피사상이다.
\end{enumerate}
\end{proposition}

\begin{proof}
호몰로지 함자의 정의에 의하여 $T_i$가 정의되는 모든 $i$에 대하여
$T_i$는 반완전하고, 모든 완전열
\[
  0\longrightarrow M'\xrightarrow{u}M\xrightarrow{v}M''\longrightarrow0,
\]
에 대하여
\[
  \operatorname{Ker}(T_{i-1}(u))
  =
  \operatorname{Coker}(T_i(v)),
\]
이므로 $a)$와 $b)$는 동치임이 명백하다. 나머지 주장들은
(\hyperref[III.7.2.5-fr]{7.2.5})와
(\hyperref[III.7.2.7-fr]{7.2.7})에서 바로 따른다.
\end{proof}

\begin{corollary}[7.3.2]
\label{III.7.3.2-fr}
$A$를 가환환이라 하자. (\hyperref[III.7.3.1-fr]{7.3.1})의 표기 아래
$T_p$가 오른쪽 완전하다고 가정하자. $f\in A$가 $A$-가군 $M$의 어떤
0이 아닌 원소의 소멸자에도 속하지 않으면, $f$는
$T_{p-1}(M)$의 어떤 0이 아닌 원소의 소멸자에도 속하지 않는다. 특히
$A$가 정역이면 $A$-가군 $T_{p-1}(A)$는 꼬임이 없다.
\end{corollary}

\begin{proof}
실제로 $h_f$가 $M$의 곱셈 준동형 $x\mapsto fx$를 나타낸다면, 가정은
$h_f$가 단사임을 뜻한다. 따라서
(\hyperref[III.7.3.1-fr]{7.3.1})의 조건 $b)$에 의하여
$T_{p-1}(h_f)$도 단사이다.
\end{proof}

\begin{proposition}[7.3.3]
\label{III.7.3.3-fr}
$A$를 환, $T_\bullet$을 $\mathbf{Ab}_A$에서 $\mathbf{Ab}$로 가고
직합과 가환하는 공변 호몰로지 함자라 하자. $T_{p-1}$, $T_p$,
$T_{p+1}$이 정의되는 정수 $p$를 잡자. 다음 조건들은 동치이다.
\begin{enumerate}
  \item[a)] $T_p$는 완전하다.
  \item[b)] $T_{p+1}$과 $T_p$는 오른쪽 완전하다.
  \item[c)] $T_p$와 $T_{p-1}$은 왼쪽 완전하다.
  \item[d)] $T_{p+1}$은 오른쪽 완전하고 $T_{p-1}$은 왼쪽 완전하다.
  \item[e)] 모든 $A$-가군 $M$에 대하여 표준 준동형
    \[
      T_i(A_s)\otimes_A M\longrightarrow T_i(M)
      \tag{7.3.3.1}\label{III.7.3.3.1.fr}
    \]
    은 $i=p$와 $i=p+1$일 때 동형이다.
  \item[$e'$)] 모든 $A$-가군 $M$에 대하여 표준 준동형 (7.3.3.1)은
    $i=p$와 $i=p+1$일 때 에피사상이다.
  \item[f)] 모든 $A$-가군 $M$에 대하여 준동형 (7.3.3.1)은 $i=p$일
    때 동형이고, $T_p(A_s)$는 평탄한 오른쪽 $A$-가군이다.
  \item[$f'$)] 모든 $A$-가군 $M$에 대하여 준동형 (7.3.3.1)은
    $i=p$일 때 에피사상이고, $T_p(A_s)$는 평탄한 오른쪽
    $A$-가군이다.
\end{enumerate}
\end{proposition}

\begin{proof}
조건 $a)$, $b)$, $c)$, $d)$의 동치는
(\hyperref[III.7.3.1-fr]{7.3.1})의 조건 $a)$와 $b)$의 동치에서
따른다. $b)$, $e)$, $e')$의 동치는 (\hyperref[III.7.3.1-fr]{7.3.1})의
조건 $a)$, $c)$, $d)$의 동치에서 따른다. 끝으로 $T_p(A_s)$가
평탄하다는 것은 함자
\[
  M\mapsto T_p(A_s)\otimes_A M
\]
가 왼쪽 완전하다는 뜻이므로, $a)$, $f)$, $f')$의 동치도 다시
(\hyperref[III.7.3.1-fr]{7.3.1})의 조건 $a)$, $c)$, $d)$의 동치에서
따른다.
\end{proof}

\oldpage[III]{182}

\begin{corollary}[7.3.4]
\label{III.7.3.4-fr}
$A$가 가환이고 $T_p$가 완전하며, 더 나아가 $T_p(A)$가 유한 표시를
갖는 $A$-가군이라고 가정하자. 그러면 함수
\[
  x\mapsto\operatorname{rang}_{\kappa(x)}(T_p(\kappa(x)))
\]
는 $X=\operatorname{Spec}(A)$에서 국소적으로 일정하고, 따라서
$\operatorname{Spec}(A)$가 연결이면 일정하다.
\end{corollary}

\begin{proof}
실제로 (\hyperref[III.7.3.3-fr]{7.3.3}, $f)$)에 의하여 $T_p(A)$는
평탄한 $A$-가군이므로 유한 생성 사영 가군이고, 따라서
$\bigl(T_p(A)\bigr)^{\sim}$은 국소 자유 $\mathscr O_X$-가군이다
(Bourbaki, \emph{Alg. comm.}, chap. II, \S 5,
n\textsuperscript{o} 2, th. 1). 또한
\[
  T_p(\kappa(x))=T_p(A)\otimes_A\kappa(x)
\]
이고 (\hyperref[III.7.3.3-fr]{7.3.3}, $e)$), $A$-가군 $T_p(A)$의
점 $x$에서의 계수(階數)가 국소적으로 일정함이 알려져 있으므로
(\emph{loc. cit.}) 따름정리가 성립한다.
\end{proof}

\begin{proposition}[7.3.5]
\label{III.7.3.5-fr}
$A$를 그 근기 $\mathfrak m$으로 나눈 몫이 체 $k$인 왼쪽 아르틴
환이라고 가정하자. 그러면 (\hyperref[III.7.3.3-fr]{7.3.3})의 조건들은
다음 각 조건과도 동치이다.
\begin{enumerate}
  \item[g)] 표준 준동형
    \[
      T_i(\varepsilon):T_i(A_s)\longrightarrow T_i(k)
    \]
    은 $i=p$와 $i=p+1$일 때 에피사상이다.
  \item[h)] $T_p(\varepsilon)$은 에피사상이고 $T_p(A_s)$는 평탄한
    오른쪽 $A$-가군이다. 이는 같은 말로 (Bourbaki,
    \emph{Alg. comm.}, chap. II, \S 3, n\textsuperscript{o} 2,
    prop. 5의 cor. 2) 자유 $A$-가군이다.
\end{enumerate}
더 나아가 $A$가 가환이고 $A$-가군 $T_p(k)$의 길이가 유한한 $d$라고
가정하자. 그러면 앞의 조건들은 다음 각 조건과도 동치이다.
\begin{enumerate}
  \item[i)] 모든 유한 길이 $A$-가군 $M$에 대하여
    \[
      \operatorname{long}(T_p(M))=d\mathbin{.}\operatorname{long}(M).
      \tag{7.3.5.1}\label{III.7.3.5.1.fr}
    \]
  \item[j)] 다음 등식이 성립한다.
    \[
      \operatorname{long}(T_p(A))=d\mathbin{.}\operatorname{long}(A).
      \tag{7.3.5.2}\label{III.7.3.5.2.fr}
    \]
\end{enumerate}
\end{proposition}

$g)$와 $h)$가 (\hyperref[III.7.3.3-fr]{7.3.3})의 조건들과 동치임은
(\hyperref[III.7.2.7-fr]{7.2.7})에서 곧 따른다. 나머지 주장들을
증명하기 위하여 다음 보조정리를 사용하겠다.

\begin{lemma}[7.3.5.3]
\label{III.7.3.5.3-fr}
$\mathcal K$, $\mathcal K'$을 두 아벨 범주,
$F:\mathcal K\to\mathcal K'$을 공변 가법 함자라 하자. $F$가
반완전이고, $\mathcal K$의 모든 단순 대상 $S$에 대하여 $F(S)$가
$\mathcal K'$의 유한 길이 대상이라고 가정하자. 그러면
$\mathcal K$의 모든 유한 길이 대상 $E$에 대하여 $F(E)$는
$\mathcal K'$의 유한 길이 대상이다. $\mathcal K$의 유한 길이
대상들로 이루어진 모든 완전열
\[
  0\longrightarrow E'\xrightarrow{u}E\xrightarrow{v}E''\longrightarrow0
\]
에 대하여
\[
  \operatorname{long}F(E)
  \leq
  \operatorname{long}F(E')+\operatorname{long}F(E'')
  \tag{7.3.5.4}\label{III.7.3.5.4.fr}
\]
가 성립한다. (7.3.5.4)의 양변이 같기 위한 필요충분조건은 열
\[
  0\longrightarrow F(E')\longrightarrow F(E)\longrightarrow F(E'')
  \longrightarrow0
\]
이 완전한 것이다.
\end{lemma}

\begin{proof}
실제로 가정에 의하여 열
\[
  F(E')\xrightarrow{F(u)}F(E)\xrightarrow{F(v)}F(E'')
\]
은 완전하다. $F(E')$과 $F(E'')$이 유한 길이라고 가정하면
$\operatorname{Im}(F(u))$와 $\operatorname{Im}(F(v))$도 유한 길이이고,
$\operatorname{Ker}(F(v))=\operatorname{Im}(F(u))$이므로 $F(E)$도
유한 길이이며
\[
  \operatorname{long}F(E)
  =
  \operatorname{long}\operatorname{Im}(F(u))
  +
  \operatorname{long}\operatorname{Im}(F(v))
  \leq
  \operatorname{long}F(E')+\operatorname{long}F(E'').
  \tag{7.3.5.5}\label{III.7.3.5.5.fr}
\]
이다.
\oldpage[III]{183}
$E$의 길이에 관한 귀납법으로 첫째 주장이 이미 증명된다. 더욱이
(7.3.5.5)의 양변이 같을 수 있으려면
\[
  \operatorname{long}\operatorname{Im}(F(u))=\operatorname{long}F(E')
\]
(이는 $\operatorname{long}\operatorname{Ker}(F(u))=0$, 또는
$\operatorname{Ker}(F(u))=0$과 동치이다)이고
\[
  \operatorname{long}\operatorname{Im}(F(v))=\operatorname{long}F(E'')
\]
(이는 $\operatorname{long}\operatorname{Coker}(F(v))=0$, 또는
$\operatorname{Coker}(F(v))=0$과 동치이다)이어야 한다.
\end{proof}

이제 $M$이 유한 길이 $A$-가군이면($A$는 가환), $M$의 한
조르당--횔더 합성열의 몫들은 반드시 $A$-가군 $k$와 동형이다. 따라서
(\hyperref[III.7.3.5.4.fr]{7.3.5.4})로부터 $M$의 길이에 관한 귀납법으로
다음을 얻는다.
\[
  \operatorname{long}T_p(M)\leq d\mathbin{.}\operatorname{long}(M).
  \tag{7.3.5.6}\label{III.7.3.5.6.fr}
\]
또한 (\hyperref[III.7.3.5.3-fr]{7.3.5.3})에 의하여 $T_p$가 완전하면
(\hyperref[III.7.3.5.1.fr]{7.3.5.1})의 등식이 성립한다. 따라서
(\hyperref[III.7.3.3-fr]{7.3.3})의 조건 $a)$는 $i)$를 함의한다.
$i)$가 $j)$를 함의함은 명백하므로, 다음 보조정리만 증명하면 된다.

\begin{lemma}[7.3.5.7]
\label{III.7.3.5.7-fr}
관계식
\[
  \operatorname{long}T_p(A)=d\mathbin{.}\operatorname{long}A
\]
이 성립하면 $T_p(\varepsilon)$은 에피사상이고 $T_p(A)$는 평탄한
$A$-가군이다.
\end{lemma}

\begin{proof}
실제로 완전열
\[
  0\longrightarrow\mathfrak m\longrightarrow A\longrightarrow k
  \longrightarrow0,
\]
에서 출발하면, (\hyperref[III.7.3.5.4.fr]{7.3.5.4})와
(\hyperref[III.7.3.5.6.fr]{7.3.5.6})에 의하여
\[
  \operatorname{long}T_p(A)
  \leq
  \operatorname{long}T_p(\mathfrak m)+\operatorname{long}T_p(k)
  \leq
  d\bigl(\operatorname{long}\mathfrak m+\operatorname{long}k\bigr)
  =d\mathbin{.}\operatorname{long}A
\]
이고, 등식이 성립하려면 (\hyperref[III.7.3.5.3-fr]{7.3.5.3})에 의하여
열
\[
  0\longrightarrow T_p(\mathfrak m)\longrightarrow T_p(A)
  \longrightarrow T_p(k)\longrightarrow0
  \tag{7.3.5.8}\label{III.7.3.5.8.fr}
\]
이 완전해야 한다. (\hyperref[III.7.2.7-fr]{7.2.7})과
(\hyperref[III.7.2.5-fr]{7.2.5})에 의하여 $T_p$는 함자
$M\mapsto N\otimes_A M$과 동형이고, 열 (7.3.5.8)의 완전성으로부터
$\operatorname{Tor}$ 완전열을 사용하여
$\operatorname{Tor}_1^A(N,k)=0$을 얻는다. 따라서
$N=T_p(A)$는 평탄한 $A$-가군이다 (0, 10.1.3).
\end{proof}

\begin{lemma}[7.3.6]
\label{III.7.3.6-fr}
$A$를 환, $T_\bullet$을 $\mathbf{Ab}_A$에서 $\mathbf{Ab}$로 가고
직합과 가환하는 공변 호몰로지 함자라 하자. $T_p$와 $T_{p+1}$이
정의되고 $T_p$가 왼쪽 완전하다고 가정하자. $T_{p+1}$이 완전하기 위한
필요충분조건은 $T_{p+1}(A_s)$가 평탄한 오른쪽 $A$-가군인 것이다.
\end{lemma}

\begin{proof}
실제로 (\hyperref[III.7.3.1-fr]{7.3.1})에 의하여 표준 준동형
\[
  T_{p+1}(A_s)\otimes_A M\longrightarrow T_{p+1}(M)
\]
은 함자들의 동형임을 알고 있다. 이제 평탄한 $A$-가군의 정의를
적용하면 충분하다.
\end{proof}

\begin{proposition}[7.3.7]
\label{III.7.3.7-fr}
$A$를 환, $T_\bullet$을 $\mathbf{Ab}_A$에서 $\mathbf{Ab}$로 가고
직합과 가환하는 공변 호몰로지 함자라 하자. 어떤 $i_0$가 존재하여
$i\leq i_0$이면 $T_i$가 완전하다고 가정한다. 그러면 모든 정수
$p>i_0$에 대하여 다음 조건들은 동치이다.
\begin{enumerate}
  \item[a)] $q\leq p$이면 $T_q$는 완전하다.
  \item[b)] $q\leq p$이면 $T_q(A_s)$는 평탄한 오른쪽 $A$-가군이다.
  \item[c)] 모든 $A$-가군 $M$에 대하여 표준 준동형
    \[
      T_q(A_s)\otimes_A M\longrightarrow T_q(M)
    \]
    은 $q\leq p+1$이면 전사이다.
\end{enumerate}
\end{proposition}

\begin{proof}
가정에 의하여 $T_{i_0}$가 완전하므로 $a)$와 $b)$의 동치는
(\hyperref[III.7.3.6-fr]{7.3.6})과 $q$에 관한 귀납법에서 따른다.
$a)$와 $c)$의 동치는 (\hyperref[III.7.3.3-fr]{7.3.3})의 조건 $a)$와
$e')$의 동치에서 따른다.
\end{proof}

\oldpage[III]{184}

\begin{env}[7.3.8]
\label{III.7.3.8-fr}
$A$가 가환환이고, $B$가 반드시 가환일 필요는 없는 $A$-대수이며,
$T_\bullet$이 $\mathbf{Ab}_A$에서 $\mathbf{Ab}$로 가는 공변 호몰로지
함자이면, 정의 (\hyperref[III.7.1.3-fr]{7.1.3})에 의하여 $A$에서
$B$로 스칼라를 확장하여 얻은 $\mathbf{Ab}_B$에서 $\mathbf{Ab}$로 가는
함자, 즉
\[
  T_\bullet^{(B)}=(T_i^{(B)}),
\]
로 나타낼 함자도 호몰로지 함자이다.
\end{env}

\begin{corollary}[7.3.9]
\label{III.7.3.9-fr}
$T_\bullet$이 (\hyperref[III.7.3.7-fr]{7.3.7})의 일반 조건들을 만족하고
귀납극한과 가환한다고 가정하자. 더 나아가 $A$가 정역이고 모든
$T_n(A)$가 유한 표시를 갖는 $A$-가군이라고 하자. 그러면 모든 정수
$N$에 대하여 어떤 $f\in A-\{0\}$가 존재하여 함자
\[
  T_p^{(A_f)}:\mathbf{Ab}_{A_f}\longrightarrow\mathbf{Ab}
\]
가 $p\leq N$이면 완전하다.
\end{corollary}

\begin{proof}
가정에 의하여 $i\leq i_0$이면 $T_i$는 완전하므로, 그러한 $i$에
대하여 $T_i(A)$는 평탄하다. (\hyperref[III.7.3.7-fr]{7.3.7}, $b)$)에
의하여
\[
  T_p^{(A_f)}(A_f)=T_p(A_f)
\]
가 $i_0<p\leq N$이면 자유 $A_f$-가군이 되도록 $f$를 잡으면 충분하다.
그런데 $T_p$가 귀납극한과 가환하므로
(\hyperref[III.7.1.4-fr]{7.1.4})
$T_p(A_f)=(T_p(A))_f$이다. $x$가 $\operatorname{Spec}(A)$의
일반점이면 $(T_p(A))_x$는 $A$의 분수체 위의 유한차원 벡터공간이다.
각 $T_p(A)$가 유한 표시를 가지므로, 요구한 성질을 갖는 $f$가 실제로
존재한다 (Bourbaki, \emph{Alg. comm.}, chap. II, \S 5,
n\textsuperscript{o} 1, cor. de la prop. 2).

$T_i\ne0$인 지표 $i$가 유한개뿐이면, 어떤 $f\in A-\{0\}$가 존재하여
모든 $T_p^{(A_f)}$가 완전하다는 점에 유의하자.
\end{proof}

\begin{corollary}[7.3.10]
\label{III.7.3.10-fr}
$T_\bullet$이 (\hyperref[III.7.3.7-fr]{7.3.7})의 일반 조건들을 만족하고
귀납극한과 가환하며, $A$가 가환 뇌터 환이고 $T_n(A)$들이 유한 생성
$A$-가군이라고 가정하자. 그러면 모든 정수 $N$에 대하여
$\operatorname{Spec}(A)$의 한 조밀 열린집합 $U$가 존재하여, 모든
$p\leq N$에 대하여 함수
\[
  x\mapsto\operatorname{rang}_{\kappa(x)}(T_p(\kappa(x)))
\]
가 $U$에서 일정하다.
\end{corollary}

\begin{proof}
$\mathfrak p$를 $A$의 극소 소 아이디얼이라 하자. 가정에 의하여
$B=A/\mathfrak p$는 정역이고 $\operatorname{Spec}(B)$는 위상공간
$\operatorname{Spec}(A)$의 한 기약 성분과 동일시된다. $p\leq N$에
관한 귀납법으로, $B'=B_{f_p}$라 놓으면 $i\leq p$일 때
$T_i^{(B')}$가 완전하고 $T_i(B')$가 유한 생성 $B'$-가군이 되게 하는
어떤 $f_p\in B-\{0\}$가 존재함을 보이겠다. $p\leq i_0$이면 가정에
의하여 주장이 성립하며 $f_p=1$, 따라서 $B'=B=A/\mathfrak p$로 잡을
수 있다. 실제로 이때 $T_p$가 완전하므로 $T_p(B)$는
$T_p(A)/T_p(\mathfrak p)$와 동형이고, 따라서 유한 생성 $A$-가군이며
\emph{따라서 특히} 유한 생성 $B$-가군이다.

이제 $p$에 관하여 귀납하자. $f_p$는 어떤 $g_p\in A$의 $B$에서의
표준상이고, $A'=A_{g_p}$라 놓으면
$B'=A'/\mathfrak p'$이다. 여기서 $\mathfrak p'$은 $A'$의 극소 소
아이디얼이고 $\mathfrak p_{g_p}$와 같다. 또한
\[
  T_i(A_{g_p})=(T_i(A))_{g_p},
\]
이므로 $T_i(A')$들은 유한 생성 $A'$-가군이다. 따라서
$T_\bullet^{(A')}$은 $i_0$를 $p$로 바꾸었을 때 $T_\bullet$과 같은
가정들을 만족한다. 그러므로 $A'=A$이고 $T_p$가 완전한 경우로
한정해도 된다. 완전열
\[
  0\longrightarrow\mathfrak p\longrightarrow A\longrightarrow A/\mathfrak p
  \longrightarrow0
\]
은 완전열
\[
  T_{p+1}(A)\longrightarrow T_{p+1}(A/\mathfrak p)
  \xrightarrow{\partial}T_p(\mathfrak p)\longrightarrow T_p(A),
\]
을 준다. $T_p$가 완전하므로 맨 오른쪽 화살표는 단사이고, 따라서
$T_{p+1}(A/\mathfrak p)$는 $T_{p+1}(A)$의 몫이므로 유한 생성이다.
이제 (\hyperref[III.7.3.9-fr]{7.3.9})의 논증은 $T_p(A)$들이
$p\leq N$일 때 유한 생성이라는 사실만 사용했다는 점에 유의하자.
따라서 정역 $B$, 함자 $T_\bullet^{(B)}$, 그리고 $N=p+1$에 그 논증을
적용할 수 있고, 이로써 귀납 논증이 끝난다.

이제 어떤 $f_N\in B-\{0\}$가 존재하여
$\operatorname{Spec}(B_{f_N})$은 $\operatorname{Spec}(B)$에서
어디서나 조밀한 열린집합 $V$이고, $\operatorname{Spec}(A)$의 다른
어떤 기약 성분과도 만나지 않는다. 명제가 $B_{f_N}$에 대하여
\oldpage[III]{185}
증명되었다면, 명제의 함수들이 일정한 $V$의 어디서나 조밀한
열린집합 $W$가 존재한다. 실제로 모든 $x\in W$에 대하여
$A_x=(B_{f_N})_x$이다. $\operatorname{Spec}(A)$의 각 기약 성분에
대하여 같은 논증을 적용하면 따름정리가 증명된다. 따라서 $A$가
정역인 경우로 한정해도 된다. 그러면
(\hyperref[III.7.3.9-fr]{7.3.9})의 논증으로 어떤
$f\in A-\{0\}$가 존재하여 $p\leq N$일 때 $T_p(A_f)$가 유한 생성
자유 $A_f$-가군임을 얻는다. 이제
(\hyperref[III.7.3.4-fr]{7.3.4})에 의하여
(\hyperref[III.7.3.10-fr]{7.3.10})의 결론이 따른다.
\end{proof}

\begin{proposition}[7.3.11]
\label{III.7.3.11-fr}
$A$를 가환 국소환, $k$를 그 잉여체, $T_\bullet$을
$\mathbf{Ab}_A$에서 $\mathbf{Ab}$로 가고 직합과 가환하는 공변
호몰로지 함자라 하자. 어떤 $i_0$가 존재하여 $i\leq i_0$이면
$T_i$가 완전하고, 모든 $T_n(A)$가 유한 표시를 갖는 $A$-가군이라고
가정한다. 그러면 (\hyperref[III.7.3.7-fr]{7.3.7})의 서로 동치인
조건 $a)$, $b)$, $c)$는 다음 두 조건을 함의하며, 환이 더 나아가
축소환이면 이 두 조건과도 동치이다.
\begin{enumerate}
  \item[d)] 모든 $x\in\operatorname{Spec}(A)$에 대하여
    \[
      \operatorname{rang}_{\kappa(x)}T_q(\kappa(x))
      =
      \operatorname{rang}_kT_q(k)
      \quad q\leq p\text{일 때}.
    \]
  \item[$d'$)] $\operatorname{Spec}(A)$의 한 기약 성분의 모든 일반점
    $x_j$에 대하여
    \[
      \operatorname{rang}_{\kappa(x_j)}T_q(\kappa(x_j))
      =
      \operatorname{rang}_kT_q(k)
      \quad q\leq p\text{일 때}.
    \]
\end{enumerate}
\end{proposition}

\begin{proof}
$T_q(A)$가 유한 표시를 갖는 $A$-가군이므로,
(\hyperref[III.7.3.7-fr]{7.3.7})의 조건 $b)$는 $q\leq p$이면
$T_q(A)$가 자유 $A$-가군이라는 것과 동치이다 (Bourbaki,
\emph{Alg. comm.}, chap. II, \S 3, n\textsuperscript{o} 2, cor. 2
de la prop. 5). 조건 $c)$는
\[
  T_q(\kappa(x))=T_q(A)\otimes_A\kappa(x)
  \quad q\leq p\text{일 때},
\]
임을 함의한다. 따라서 (\hyperref[III.7.3.7-fr]{7.3.7})의 서로 동치인
조건들은 $d)$를 함의하고, $d)$가 $d')$을 함의함은 자명하다.

$A$가 축소환일 때 $d')$이 $a)$를 함의함을 증명하는 일만 남는다.
$q\leq i_0$이면 $T_q$가 완전하므로, $q\leq p$에 관하여 귀납하자.
$k\leq q<p$이면 $T_k$가 완전하다고 가정하고 $T_{q+1}(A)$가 자유
$A$-가군임을 보이겠다. 귀납 가정에 의하여 모든 $A$-가군 $M$에
대하여 $T_{q+1}(A)\otimes_A M$은 $T_{q+1}(M)$과 동형이다. 이는
(\hyperref[III.7.3.7-fr]{7.3.7})의 조건 $c)$와
(\hyperref[III.7.3.3-fr]{7.3.3})에서 따른다. 이 성질을
$M=\kappa(x_j)$와 $M=k$에 적용하면, 가정 $d')$에 의하여 모든 $j$에
대하여
\[
  \operatorname{rang}_{\kappa(x_j)}
  \bigl(T_{q+1}(A)\otimes_A\kappa(x_j)\bigr)
  =
  \operatorname{rang}_kT_{q+1}(k)
\]
를 얻는다. 그런데 이로부터 $T_{q+1}(A)$가 자유임이 따른다
(Bourbaki, \emph{Alg. comm.}, chap. II, \S 3,
n\textsuperscript{o} 2, prop. 7). 이로써 증명이 끝난다.
\end{proof}

앞의 결과들은 (\hyperref[subsection:III.7.4.fr]{7.4})에서 연구할
특별한 유형의 호몰로지 함자들에 대해서는 크게 강화된다. 실제로
$T_p$ 가운데 하나만을 사용하는 완전성 판정들을 얻게 될 것이다.

\subsection{함자 $H_\bullet(P_\bullet\otimes_A M)$의 완전성 판정}
\label{subsection:III.7.4.fr}

\begin{env}[7.4.1]
\label{III.7.4.1-fr}
$A$를 환(가환일 필요는 없다), $P_\bullet$을 평탄한 오른쪽
$A$-가군들로 이루어진 복합체라 하자. 그러면 모든 $k$에 대하여 함자
$M\mapsto P_k\otimes_A M$이 $\mathbf{Ab}_A$에서 완전하므로,
$\partial$-함자
\[
  T_\bullet(M)=H_\bullet(P_\bullet\otimes_A M)
  \tag{7.4.1.1}\label{III.7.4.1.1.fr}
\]
\oldpage[III]{186}
는 $\mathbf{Ab}_A$에서 $\mathbf{Ab}$로 가는 호몰로지 함자이다.
$A$가 가환일 때에는 이 함자가 분명히 $A$-선형이고
(\hyperref[III.7.1.2-fr]{7.1.2}), 귀납적 극한과 가환한다.

$A$가 가환이면 모든 $A$-대수 $B$에 대하여 호몰로지 함자
$T_\bullet^{(B)}$ (\hyperref[III.7.3.8-fr]{7.3.8})는 정의에 따라
다음으로 주어진다.
\[
  T_\bullet^{(B)}(N)
  =
  H_\bullet(P_\bullet\otimes_A N_{[\rho]})
  \tag{7.4.1.2}\label{III.7.4.1.2.fr}
\]
여기서 $\rho:A\to B$는 $B$의 대수 구조를 정하는 준동형이다. 또한
\[
  P_\bullet\otimes_A N_{[\rho]}
  =P_\bullet\otimes_A(B\otimes_B N)_{[\rho]}
  =(P_\bullet\otimes_A B)\otimes_B N,
\]
로 쓸 수 있으므로, 모든 $B$-가군 $N$에 대하여
\[
  T_\bullet^{(B)}(N)=H_\bullet(P'_\bullet\otimes_B N)
  \tag{7.4.1.3}\label{III.7.4.1.3.fr}
\]
임을 알 수 있다. 여기서 $P'_\bullet$은 평탄한 $B$-가군들로 이루어진
복합체 $P_\bullet\otimes_A B$이다 ($0_{\mathrm I}$, 6.2.1).
\end{env}

\begin{proposition}[7.4.2]
\label{III.7.4.2-fr}
(\hyperref[III.7.4.1-fr]{7.4.1})의 일반 조건 아래에서 정수
$p\in\mathbf Z$를 고정하면, 다음 조건들은 동치이다.
\begin{enumerate}
  \item[a)] $T_p$는 왼쪽 완전하다 (또는, 이와 동치로 $T_{p+1}$은
    오른쪽 완전하다).
  \item[b)]
    \[
      Z'_p(P_\bullet)
      =
      \operatorname{Coker}(P_{p+1}\longrightarrow P_p)
    \]
    는 평탄한 오른쪽 $A$-가군이다.
  \item[c)] 평탄한 오른쪽 $A$-가군들로 이루어진 복합체
    $P'_\bullet$이 존재하여 그 미분
    \[
      d'_{p+1}:P'_{p+1}\longrightarrow P'_p
    \]
    은 영이고, 호몰로지 함자
    $H_\bullet(P_\bullet\otimes_A M)$에서
    $H_\bullet(P'_\bullet\otimes_A M)$로의 동형사상이 존재한다.
\end{enumerate}
\end{proposition}

\begin{proof}
정의에 따라 $M$에 대하여 함자적인 완전열
\[
  0\longrightarrow T_p(M)
  \longrightarrow Z'_p(P_\bullet\otimes M)
  \longrightarrow P_{p-1}\otimes M
\]
이 있고, 여기서 텐서곱의 오른쪽 완전성에 의하여
\[
  Z'_p(P_\bullet\otimes M)
  =
  \operatorname{Coker}(P_{p+1}\otimes M\longrightarrow P_p\otimes M)
  =
  Z'_p(P_\bullet)\otimes M
\]
이다. 모든 준동형 $f:M\to N$에 대하여 따라서 다음 가환도식을 얻는다.
\[
\xymatrix@C=30pt@R=38pt{
  0\ar[r]
  &T_p(M)\ar[r]\ar[d]_u
  &Z'_p(P_\bullet)\otimes M\ar[r]\ar[d]_v
  &P_{p-1}\otimes M\ar[d]_w
  \\
  0\ar[r]
  &T_p(N)\ar[r]
  &Z'_p(P_\bullet)\otimes N\ar[r]
  &P_{p-1}\otimes N
}
\tag{7.4.2.1}\label{III.7.4.2.1.fr}
\]
두 행은 완전하다. $f$가 단사사상이면 $P_{p-1}$이 평탄하므로
$w$도 단사사상이다. $T_p$가 왼쪽 완전하면 $u$도 단사사상이므로
$v$가 단사사상이라고 결론지을 수 있고, 이로부터
$Z'_p(P_\bullet)$이 평탄함이 따른다. 역으로 이것이 평탄하면 모든
단사사상 $f:M\to N$에 대하여 $v$가 단사사상이므로, 도식
(\hyperref[III.7.4.2.1.fr]{7.4.2.1})은 $u$가 단사사상임을 보여 준다.
따라서 이미 반완전한 $T_p$는 왼쪽 완전하다. 이로써 $a)$와 $b)$는
동치이다.

$c)$가 $a)$를 함의함은 곧바로 알 수 있다. 실제로
$d'_{p+1}:P'_{p+1}\to P'_p$가 영이고
$0\to M'\to M\to M''\to0$이 $A$-가군들의 완전열이면, 완전열
\[
  H_{p+1}(P'_\bullet\otimes M'')
  \xrightarrow{\partial}H_p(P'_\bullet\otimes M')
  \longrightarrow H_p(P'_\bullet\otimes M)
\]
\oldpage[III]{187}
의 경계 연산자는 정의에 따라 영이다 (M, IV, 1). 따라서 $T_p$는
왼쪽 완전하다. 역으로 $b)$가 $c)$를 함의함을 보이자. 다음과 같이
놓는다.
\[
 C=Z'_p(P_\bullet),\qquad
 B=B_p(P_\bullet)=\operatorname{Im}(d_{p+1}),\qquad
 Z=Z_{p+1}(P_\bullet)=\operatorname{Ker}(d_{p+1}).
\]
다음 두 완전열이 있다.
\[
 0\longrightarrow B\longrightarrow P_p\longrightarrow C\longrightarrow0,
 \qquad
 0\longrightarrow Z\longrightarrow P_{p+1}\longrightarrow B\longrightarrow0.
\]
$C$와 $P_p$가 평탄하므로 $B$가 평탄하고, $B$와 $P_{p+1}$이
평탄하므로 $Z$가 평탄하다. 이 두 열은 임의의 왼쪽 $A$-가군 $M$을
텐서한 뒤에도 완전하다 ($0_{\mathrm I}$, 6.1.2). 특히 표준 준동형들은
$B\otimes_A M$을 $d_{p+1}\otimes1_M$의 상과 동일시하고,
$Z\otimes_A M$을 그 핵과 동일시한다.

복합체 $U_\bullet$을 $i\geq p+2$이면 $U_i=P_i$,
$U_{p+1}=Z$, $i\leq p$이면 $U_i=0$으로 정의하자. 차수 $p+2$의
미분은 $Z$에 값을 갖는 $d_{p+2}$이고, 0이 아닌 다른 미분들은
$P_\bullet$에서 물려받는다. 마찬가지로 $i\leq p-1$이면
$L_i=P_i$, $L_p=C$, $i\geq p+1$이면 $L_i=0$으로 $L_\bullet$을
정의하자. 그 차수 $p$의 미분은 몫 $C$ 위에 $d_p$가 유도하는
사상이다. 그러면
$P'_\bullet=U_\bullet\oplus L_\bullet$은 평탄 가군들의 복합체이고,
$i\ne p,p+1$이면 $P'_i=P_i$, $P'_p=C$, $P'_{p+1}=Z$이며
$d'_{p+1}=0$이다.

차수 $p+1$의 포함은 복합체의 사상 $a:U_\bullet\to P_\bullet$을
정의하고, 차수 $p$에서 몫으로 보내는 사상은
$b:P_\bullet\to L_\bullet$을 정의한다. $M$을 텐서한 뒤 첫 번째
사상은 모든 차수 $i\geq p+1$의 호몰로지에서 동형을 유도한다.
차수 $p+1$에서는 앞의 핵 식별에서, 차수 $p+2$에서는
$Z\otimes_A M\to P_{p+1}\otimes_A M$의 단사성에서 이 사실이
따르며, 더 높은 차수들은 바뀌지 않는다. 두 번째 사상은 모든 차수
$i\leq p$에서 동형을 유도한다. 차수 $p$에서는 상
$B\otimes_A M$으로 나눈 몫을 취하고, 차수 $p-1$에서는
$P_p\otimes_A M\to C\otimes_A M$의 전사성을 사용하며, 더 낮은
차수들은 바뀌지 않는다. 따라서 첫 번째 동형들은 역으로 취하고 두
번째 동형들은 그대로 취하면, 모든 $i$에 대하여 함자적인 동형
\[
 H_i(P_\bullet\otimes_A M)\xrightarrow{\sim}
 H_i(P'_\bullet\otimes_A M).
\]
들을 얻는다.

이 동형들은 임의의 완전열 $0\to M'\to M\to M''\to0$에 딸린 경계
연산자들과 가환한다. 차수 $i\geq p+2$와 $i\leq p$에서는 복합체의
사상 $a$와 $b$에 대한 함자성에서 이 사실이 따른다. 차수 $p+1$에서는
$P_\bullet\otimes_A M''$의 모든 순환이 $Z\otimes_A M''$에서 오고,
전사 $Z\otimes_A M\to Z\otimes_A M''$을 통하여
$P_\bullet\otimes_A M$의 순환으로 올려진다. 따라서 경계 연산자는
영이다. $P'_\bullet$의 경계 연산자도 $d'_{p+1}=0$이므로 영이다.
이로써 요구한 $\partial$-함자의 동형을 얻는다.
\end{proof}

(\hyperref[III.7.4.2-fr]{7.4.2})의 조건들은 $B_p(P_\bullet)$도
평탄함을 함의함에 유의하자. 실제로 완전열
\[
  0\longrightarrow B_p(P_\bullet)\longrightarrow P_p
  \longrightarrow Z'_p(P_\bullet)\longrightarrow0,
\]
에서 $P_p$와 $Z'_p(P_\bullet)$이 평탄하다 ($0_{\mathrm I}$, 6.1.2).

\begin{corollary}[7.4.3]
\label{III.7.4.3-fr}
$A$가 1차원 정칙 뇌터 환이라고 가정하자. 다시 말해 $A$는 데데킨트
환들의 곱이다 ($0_{\mathrm{IV}}$, 17.1.3과 17.3.7). 예를 들어
$A$는 주 아이디얼환일 수 있다. 그러면 $T_p$가 왼쪽 완전하기 위한
필요충분조건은 $T_p(A)$가 평탄한 $A$-가군인 것이다. $T_p$가
완전하기 위한 필요충분조건은 $T_p(A)$와 $T_{p-1}(A)$가 평탄한
$A$-가군인 것이다.
\end{corollary}

\begin{proof}
데데킨트 환 위의 가군 $M$에 대하여 $M$이 평탄하다는 것과 꼬임이
없다는 것은 동치임을 상기하자 ($0_{\mathrm I}$, 6.3.3과 6.3.4).
따라서 (\hyperref[III.7.4.3-fr]{7.4.3})의 가정 아래에서는 평탄한
$A$-가군의 모든 부분가군이 평탄하다.

(\hyperref[III.7.4.3-fr]{7.4.3})의 두 번째 주장은 첫 번째 주장으로부터
따른다. 실제로 $T_p$가 완전하다는 것은 $T_p$와 $T_{p-1}$이 왼쪽
완전하다는 뜻이다. 첫 번째 주장을 증명하기 위하여 완전열
\[
  0\longrightarrow H_p(P_\bullet)\longrightarrow Z'_p(P_\bullet)
  \longrightarrow B_{p-1}(P_\bullet)\longrightarrow0
\]
을 보자. 여기서 $B_{p-1}(P_\bullet)$은 평탄한 $A$-가군
$P_{p-1}$의 부분가군이므로 평탄하다. 따라서 $H_p(P_\bullet)$이
평탄하다는 것과 $Z'_p(P_\bullet)$이 평탄하다는 것은 동치이다
($0_{\mathrm I}$, 6.1.2).
\end{proof}

(\hyperref[III.7.4.2-fr]{7.4.2})의 가장 중요한 응용들은 다음과 같다.

\begin{proposition}[7.4.4]
\label{III.7.4.4-fr}
$A$를 뇌터 환, $P_\bullet$을 평탄한 $A$-가군들의 복합체라 하자.
각 $P_i$가 유한형이거나, 또는 각 $H_i(P_\bullet)$이 유한형
$A$-가군이고 어떤 $i_0$가 존재하여 $i<i_0$이면
$H_i(P_\bullet)=0$이라고 가정한다. $T$를 (7.4.1.1)로 정의되는
호몰로지 함자라 하자. 그러면 $(T_p)_y$
(\hyperref[III.7.1.4-fr]{7.1.4})가 오른쪽 완전(각각 왼쪽 완전,
완전)인 $y\in\operatorname{Spec}(A)$들의 집합 $U$는
$\operatorname{Spec}(A)$의 열린집합이다.
\end{proposition}

\begin{proof}
$P_\bullet$에 관한 두 번째 가정 아래에서는 $P_\bullet$을 유한형 자유
\oldpage[III]{188}
$A$-가군들로 이루어진 복합체 $P'_\bullet$로 바꿀 수 있고, 이때
$H_\bullet(P'_\bullet\otimes_A M)$은 $T_\bullet(M)$과
$\partial$-함자로서 동형이다 (0, 11.9.3). 따라서 언제나 첫 번째
가정으로 환원할 수 있으며, 이 경우 $Z'_i(P_\bullet)$들은 유한형이다.
또한 왼쪽 완전성에 관한 주장들만 증명하면 충분하다
((\hyperref[III.7.4.2-fr]{7.4.2}, $a)$ 참조).

이제 $x\in U$라 하자. 함자 $M\mapsto M_x$가 완전하므로
\[
  (Z'_p(P_\bullet))_x=Z'_p((P_\bullet)_x),
\]
이고, (7.4.1.3)을 고려하면 가정과
(\hyperref[III.7.4.2-fr]{7.4.2}, $b)$로부터
$(Z'_p(P_\bullet))_x$가 평탄한 $A_x$-가군임이 따른다. 이 가군은
유한형이고 $A_x$가 뇌터 국소환이므로 자유이다 (Bourbaki,
\emph{Alg. comm.}, chap. II, \S 3, n\textsuperscript{o} 2, cor. 2
de la prop. 5). 따라서 어떤 $f\in A$가 존재하여
$(Z'_p(P_\bullet))_f$는 $A_f$ 위에서 자유이다 (Bourbaki,
\emph{Alg. comm.}, chap. II, \S 5, n\textsuperscript{o} 1,
cor. de la prop. 2). 따라서 특히 모든 $y\in D(f)$에 대하여
$(Z'_p(P_\bullet))_y$가 $A_y$
위에서 자유이므로, (\hyperref[III.7.4.2-fr]{7.4.2}, $b)$에 의하여
증명이 끝난다.
\end{proof}

\begin{corollary}[7.4.5]
\label{III.7.4.5-fr}
(\hyperref[III.7.4.4-fr]{7.4.4})의 가정 아래에서 $A$가 정역이라고
추가로 가정하자. 그러면 $(T_p)_x$가 완전한
$x\in\operatorname{Spec}(A)$들의 집합 $U$는 공집합이 아닌
열린집합이다.
\end{corollary}

\begin{proof}
$\operatorname{Spec}(A)$의 일반점 $x$에서 $(T_p)_x$가 완전함을
보이면 충분하다. 그런데 $A_x$가 체이므로
$\mathbf{Ab}_{A_x}$ 위의 모든 가법 함자는 완전하며, 주장은 곧바로
따른다.
\end{proof}

\begin{proposition}[7.4.6]
\label{III.7.4.6-fr}
(\hyperref[III.7.4.4-fr]{7.4.4})의 일반 가정 아래에서
(\hyperref[III.7.4.2-fr]{7.4.2})의 조건 $a)$, $b)$, $c)$는 다음
조건과도 동치이다.
\begin{enumerate}
  \item[d)] 어떤 $A$-가군 $Q$와 함자적인 동형
    \[
      T_p(M)\xrightarrow{\sim}\operatorname{Hom}_A(Q,M).
      \tag{7.4.6.1}\label{III.7.4.6.1.fr}
    \]
    이 존재한다. 또한 이 성질에 의하여 $A$-가군 $Q$는 유일한
    동형을 제외하고 유일하게 결정되며, 유한형이다.
\end{enumerate}
\end{proposition}

\begin{proof}
$Q$의 유일성은 표현 가능한 함자의 표현 대상이 유일하다는 사실의
특별한 경우이다 (0, 8.1.5). (7.4.6.1)의 오른쪽 항이 왼쪽 완전함은
분명하다. 역으로 $T_p$가 왼쪽 완전할 때 $Q$의 존재를 증명하기
위해서는, 우선 (\hyperref[III.7.4.4-fr]{7.4.4})에서와 같이 각
$P_i$가 평탄하고 유한형인 경우로 환원할 수 있다. $A$가 뇌터
환이므로 이 $P_i$들은 유한형 사영 가군이다 (Bourbaki,
\emph{Alg. comm.}, chap. II, \S 5, n\textsuperscript{o} 2,
cor. du th. 1). 그러면 $P_i$의 쌍대 $\check P_i$도 유한형 사영
$A$-가군이고, $P_i$는 $\check P_i$의 쌍대와 표준적으로 동형이며,
표준 준동형
\[
  P_i\otimes_A M\longrightarrow\operatorname{Hom}_A(\check P_i,M)
\]
은 전단사이다 (Bourbaki, \emph{Alg.}, chap. II,
3\textsuperscript{e} éd., \S 4, n\textsuperscript{o} 2, prop. 2).
한편 (\hyperref[III.7.4.2-fr]{7.4.2}, $c)$에 의하여
$d_{p+1}:P_{p+1}\to P_p$가 영이라고 가정할 수 있다. 따라서
\[
  0\longrightarrow T_p(M)\xrightarrow{u}P_p\otimes M
  \xrightarrow{v}P_{p-1}\otimes M
\]
은 완전열이고, 여기서 $v=d_p\otimes1$이다. 이제
$Q'=\operatorname{Ker}(d_p)$라 놓으면 완전열
\[
  0\longrightarrow Q'\xrightarrow{w}P_p\xrightarrow{d_p}P_{p-1},
\]
을 얻고, 이를 전치하여 완전열
\[
  \check P_{p-1}\xrightarrow{{}^t d_p}\check P_p
  \xrightarrow{{}^t w}\check Q'\longrightarrow0.
\]
을 얻는다. 이제
\[
  Q=\check Q'=\operatorname{Coker}({}^t d_p)
\]
가 요구를 만족함을 보이자. 실제로 완전열
\[
  0\longrightarrow\operatorname{Hom}_A(Q,M)
  \longrightarrow\operatorname{Hom}_A(\check P_p,M)
  \xrightarrow{v'}\operatorname{Hom}_A(\check P_{p-1},M)
\]
\oldpage[III]{189}
이 있고, 여기서 $v'=\operatorname{Hom}({}^t d_p,1)$이다.
$P_i\otimes M$을 $\operatorname{Hom}(\check P_i,M)$과 표준적으로
동일시하면 $v'$는 $v=d_p\otimes1$과 동일시된다. 따라서 요구한
함자적인 동형
\[
  T_p(M)\xrightarrow{\sim}\operatorname{Hom}_A(Q,M)
\]
을 얻는다. 또한 $Q$는 $\check P_p$의 몫이므로 유한형이다.
\end{proof}

\begin{proposition}[7.4.7]
\label{III.7.4.7-fr}
(\hyperref[III.7.4.4-fr]{7.4.4})의 일반 조건들이 성립한다고 가정하자.
그러면 모든 유한형 $A$-가군 $M$에 대하여 다음이 성립한다.
\begin{enumerate}
  \item[(i)] $T_i(M)$들은 유한형 $A$-가군이다.
  \item[(ii)] $A$의 모든 아이디얼 $\mathfrak m$에 대하여 표준 준동형
    \[
      \widehat{T_i(M)}
      \longrightarrow
      \varprojlim_n T_i\bigl(M\otimes_A(A/\mathfrak m^{n+1})\bigr)
      \tag{7.4.7.1}\label{III.7.4.7.1.fr}
    \]
    은 전단사이다. 여기서 왼쪽 항은 $\mathfrak m$-준아딕 위상에
    관한 $T_i(M)$의 분리 완비화이다.
\end{enumerate}
\end{proposition}

\begin{proof}
(\hyperref[III.7.4.4-fr]{7.4.4})에서와 같이 우선 각 $P_i$가 유한형인
경우로 환원한다. $A$가 뇌터 환이므로 $P_i\otimes_A M$의 부분가군들은
유한형이고, 따라서 (i)은 자명하다. 주장 (ii)는 다음 보조정리의 더
일반적인 결과이다.
\end{proof}

\begin{lemma}[7.4.7.2]
\label{III.7.4.7.2-fr}
$A$를 뇌터 환, $u:E\to F$를 유한형 $A$-가군들의 준동형이라 하자.
모든 유한형 $A$-가군 $M$에 대하여
\[
  K(M)=\operatorname{Ker}(u\otimes1_M),
  \qquad
  C(M)=\operatorname{Coker}(u\otimes1_M);
\]
로 놓자. 그러면 $A$의 모든 아이디얼 $\mathfrak m$에 대하여 표준
준동형들
\[
  \widehat{K(M)}\longrightarrow\varprojlim_n K(M_n),
  \qquad
  \widehat{C(M)}\longrightarrow\varprojlim_n C(M_n)
  \tag{7.4.7.3}\label{III.7.4.7.3.fr}
\]
은 전단사이다. 여기서
$M_n=M\otimes_A(A/\mathfrak m^{n+1})=M/\mathfrak m^{n+1}M$으로
놓았다.
\end{lemma}

\begin{proof}
$E\otimes M$과 $F\otimes M$이 유한형이고, 함자
$M\mapsto\widehat M$이 유한형 $A$-가군의 범주에서 완전하므로
($0_{\mathrm I}$, 7.3.3), $\widehat{K(M)}$과 $\widehat{C(M)}$은 각각
\[
  \widehat{(u\otimes1)}:
  \widehat{(E\otimes M)}\longrightarrow\widehat{(F\otimes M)}.
\]
의 핵과 여핵이다. 따라서 함자 $\varprojlim$의 왼쪽 완전성은
$\widehat{K(M)}=\varprojlim_n K(M_n)$을 보여 준다. 한편 텐서곱의
오른쪽 완전성으로부터
$C(M_n)=C(M)\otimes_A(A/\mathfrak m^{n+1})$이고, 따라서 정의에 의해
$\widehat{C(M)}=\varprojlim_n C(M_n)$이다.
\end{proof}

\begin{remark}[7.4.8]
\label{III.7.4.8-fr}
(6.10.5)와 (6.10.6)을 고려하면, 평탄성 가정을 하나 더 둘 때
(\hyperref[III.7.4.7-fr]{7.4.7})은 (4.1.7.1)이 동형이라는 사실, 곧
고유 사상에 대한 “제1 비교 정리”의 핵심을 다시 준다. 또한 이 명제는
연접 $\mathscr O_X$-가군 하나뿐 아니라 그러한 가군들의 복합체에도
적용된다. (\hyperref[III.7.4.7-fr]{7.4.7})과 (4.1.7.1)을 모두 특수한
경우로 포함하는 명제를 얻는 일은 흥미로울 것이다.

각 $P_i$가 유한형이면 (\hyperref[III.7.4.7-fr]{7.4.7})의 증명은
그들이 평탄 가군이라는 사실을 사용하지 않음에 유의하자. 따라서
$P_i$들이 평탄하다고도 유한형이라고도 가정하지 않고, 모든 $i$에
대하여 $H_i(P_\bullet)$이 유한형이며 $i<i_0$이면 영이라고만 가정할
때에도 (\hyperref[III.7.4.7-fr]{7.4.7})의 결론이 여전히 성립하는지를
검토할 필요가 있다. 이때 $P_\bullet$을 유한형 $A$-가군들의 복합체
$P'_\bullet$로 바꾸어 함자 $H_\bullet(P_\bullet\otimes M)$과
$H_\bullet(P'_\bullet\otimes M)$이, 더 이상 호몰로지 함자는 아니지만,
여전히 동형이 되게 할 수 있는가?
\end{remark}

\oldpage[III]{190}
\subsection{뇌터 국소환의 경우}
\label{subsection:III.7.5.fr}

\begin{env}[7.5.1]
\label{III.7.5.1-fr}
$A$를 뇌터 국소환, $\mathfrak m$을 그 극대 아이디얼이라 하자. 모든
$A$-가군 $M$에 대하여 $\widehat M$은 $\mathfrak m$-준아딕 위상에
관한 그 분리 완비화를 나타내며, 이는 다음과 동형이다.
\[
  \varprojlim_n\bigl(M\otimes_A(A/\mathfrak m^{n+1})\bigr)
  =
  \varprojlim_n\bigl(M/\mathfrak m^{n+1}M\bigr).
\]
$T$를 $\mathbf{Ab}_A$에서 $\mathbf{Ab}$로 가는 공변 가법 함자라 하자.
표준 준동형들 (7.2.3.1)
\[
  T(M)\otimes_A(A/\mathfrak m^{n+1})
  \longrightarrow
  T\bigl(M\otimes_A(A/\mathfrak m^{n+1})\bigr)
\]
은 분명히 $A$-준동형들의 사영계를 이루며, 따라서 극한에서 $M$에
대해 함자적인 $\widehat A$-준동형
\[
  \widehat{T(M)}
  \longrightarrow
  \varprojlim_n T(M_n)
  \tag{7.5.1.1}\label{III.7.5.1.1.fr}
\]
을 준다. 여기서
\[
  M_n=M\otimes_A(A/\mathfrak m^{n+1}),
  \qquad
  A_n=A/\mathfrak m^{n+1}.
\]
로 놓았다.
\end{env}

\begin{proposition}[7.5.2]
\label{III.7.5.2-fr}
$A$를 극대 아이디얼이 $\mathfrak m$인 뇌터 국소환, $k=A/\mathfrak m$을
그 잉여체라 하고, $T$를 $\mathbf{Ab}_A$에서 $\mathbf{Ab}$로 가는
공변 가법 함자로서 반완전이고 귀납극한과 가환한다고 하자. 더 나아가
모든 유한형 $A$-가군 $M$에 대하여 $T(M)$이 유한형 $A$-가군이고
표준 준동형 (7.5.1.1)이 동형이라고 가정하자. 이 조건들 아래에서
다음 성질들은 동치이다.
\begin{enumerate}
  \item[a)] $T$는 오른쪽 완전하다.
  \item[b)] 모든 $n$에 대하여 함자 $N\mapsto T(N)$은 유한형
    $A_n$-가군의 범주에서 오른쪽 완전하다. 이는 $T$가 유한 길이
    $A$-가군의 범주에서 오른쪽 완전하다는 것과 같다.
  \item[c)] 표준 준동형
    \[
      T(\varepsilon):T(A)\longrightarrow T(k)
    \]
    은 전사이다.
  \item[d)] 충분히 큰 모든 $n$에 대하여 표준 준동형
    \[
      T(A_n)\longrightarrow T(k)
    \]
    은 전사이다.
\end{enumerate}
\end{proposition}

\begin{proof}
$a)$가 $b)$를 함의함은 명백하다. $b)$가 $a)$를 함의함, 곧
$A$-가군의 전사 준동형 $u:M\to N$에 대하여 $T(u)$가 전사임을
보이자. $T$는 귀납극한과 가환하고, 가군의 범주에서 여과 지표집합에
대한 함자 $\varinjlim$은 완전하므로, $M$과 $N$이 유한형인 경우로
한정할 수 있다. 그러면 $T(M)$과 $T(N)$도 유한형이고 $A$는 뇌터
국소환이므로
\[
  \widehat{T(u)}:
  \widehat{T(M)}
  \longrightarrow
  \widehat{T(N)}
\]
가 전사임을 보이면 충분하다($0_{\mathrm I}$, 7.3.5 및
$0_{\mathrm I}$, 6.4.1). 가정에 의하여 $\widehat{T(M)}$과
$\widehat{T(N)}$은 각각 $\varprojlim_n T(M_n)$과
$\varprojlim_n T(N_n)$이므로, $\widehat{T(u)}$는 준동형들의 사영계
\[
  T(u\otimes1_{A_n}):T(M_n)\longrightarrow T(N_n).
\]
의 극한이다. 그런데 $b)$는 이 준동형들이 전사라는 뜻이다. 또한
$T(M_n)$은 유한형 $A_n$-가군이고 $A_n$은 가정에 따라 아르틴 환이다.
따라서 $\widehat{T(u)}$는 전사이다($0_{\mathrm I}$, 13.1.2 및
13.2.2). $a)$가 $c)$를 함의함은 명백하다. 한편 $T(\varepsilon)$은
\[
  T(A)\longrightarrow T(A_n)\longrightarrow T(k),
\]
로 분해되므로 $c)$는 $d)$를 함의한다. 끝으로 $T$는
$\mathbf{Ab}_{A_n}$에서 반완전이므로, $b)$와 $d)$가 동치임은
(\hyperref[III.7.2.7-fr]{7.2.7})에서 따른다. 이로써 증명이 끝난다.
\end{proof}

\begin{corollary}[7.5.3]
\label{III.7.5.3-fr}
(\hyperref[III.7.5.2-fr]{7.5.2})의 일반 가정들이 성립한다고 하자.
$T(k)=0$이면 모든 $A$-가군 $M$에 대하여 $T(M)=0$이다.
\end{corollary}

\begin{proof}
\oldpage[III]{191}
$k$는 유일한 단순 $A$-가군이므로,
(\hyperref[III.7.3.5.4.fr]{7.3.5.4})에서 모든 유한 길이 $A$-가군
$E$에 대하여 $T(E)=0$임이 따른다. 이제 $M$이 유한형이면
$\widehat{T(M)}$은 $\varprojlim_n T(M_n)$과 동형이고 $M_n$들은 유한
길이이므로 $\widehat{T(M)}=0$이다. 가정에 따라 $T(M)$은 유한형이고
$\widehat{T(M)}$의 부분가군과 동형이므로($0_{\mathrm I}$, 7.3.5),
$T(M)=0$이다. 끝으로 임의의 $A$-가군 $M$에 대하여 $T(M)$은 $M$의
유한형 부분가군 $N_\alpha$들에 대한 $T(N_\alpha)$들의 귀납극한이다.
이로써 증명이 끝난다.
\end{proof}

\begin{proposition}[7.5.4]
\label{III.7.5.4-fr}
$A$를 극대 아이디얼이 $\mathfrak m$인 뇌터 국소환, $k=A/\mathfrak m$을
그 잉여체라 하고, $T_\bullet$을 $\mathbf{Ab}_A$에서
$\mathbf{Ab}$로 가며 귀납극한과 가환하는 호몰로지 함자라 하자.
더 나아가 모든 $i$와 모든 유한형 $A$-가군 $M$에 대하여 $T_i(M)$이
유한형이고 표준 준동형
\[
  \widehat{T_i(M)}
  \longrightarrow
  \varprojlim_n T_i(M_n)
\]
이 전단사라고 가정하자. 주어진 정수 $p$에 대하여 다음 조건들은
동치이다.
\begin{enumerate}
  \item[a)] $T_p$는 완전하다.
  \item[b)] $T_p$는 오른쪽 완전하고, $T_p(A)$는 자유 $A$-가군이다.
  \item[c)] 표준 준동형들
    \[
      T_{p+1}(A)\longrightarrow T_{p+1}(k)
      \quad\text{및}\quad
      T_p(A)\longrightarrow T_p(k)
    \]
    은 전사이다.
  \item[d)] 모든 $n$에 대하여 표준 준동형들
    \[
      T_{p+1}(A_n)\longrightarrow T_{p+1}(k)
      \quad\text{및}\quad
      T_p(A_n)\longrightarrow T_p(k)
    \]
    은 전사이다.
  \item[e)] 모든 $n$에 대하여 함자 $N\mapsto T_p(N)$은 유한형
    $A_n$-가군의 범주에서 완전하다.
\end{enumerate}
\end{proposition}

\begin{proof}
(\hyperref[III.7.3.3-fr]{7.3.3})에 따라 $a)$는 $T_{p+1}$과 $T_p$가
오른쪽 완전하다는 것과 동치이다. 한편 $T_\bullet$은
$\mathbf{Ab}_{A_n}$의 범주에서 호몰로지 함자이므로,
(\hyperref[III.7.3.1-fr]{7.3.1})과 같은 논증에 의하여 $e)$는
$T_p$와 $T_{p+1}$이 유한형 $A_n$-가군의 범주에서 오른쪽 완전하다는
것과 동치이다. 따라서 (\hyperref[III.7.5.2-fr]{7.5.2})에서 $a)$와
$e)$가 동치이고, $a)$, $c)$, $d)$의 동치도 같은 명제에서 따른다.
끝으로 모든 유한형 평탄 $A$-가군은 자유임이 알려져 있다(Bourbaki,
\emph{Alg. comm.}, chap.~II, \S~3, n\textsuperscript{o}~2, prop.~5의
cor.~2). 이제 $a)$와 $b)$의 동치는
(\hyperref[III.7.3.1-fr]{7.3.1})과
(\hyperref[III.7.3.3-fr]{7.3.3})에서 따른다.
\end{proof}

\begin{corollary}[7.5.5]
\label{III.7.5.5-fr}
(\hyperref[III.7.5.4-fr]{7.5.4})의 일반 조건들이 성립한다고 하자.
\begin{enumerate}
  \item[(i)] $T_p(k)=0$이면 $T_p=0$이고, $T_{p+1}$은 오른쪽 완전하며
    $T_{p-1}$은 왼쪽 완전하다.
  \item[(ii)] $T_{p-1}(k)=T_{p+1}(k)=0$이면 $T_p$는 완전하고, 표준
    준동형
    \[
      T_p(A)\otimes_A M\longrightarrow T_p(M)
    \]
    은 전단사이며, $T_p(A)$는 자유 $A$-가군이다.
\end{enumerate}
\end{corollary}

\begin{proof}
$T_p$가 반완전이므로 (i)는 (\hyperref[III.7.5.3-fr]{7.5.3})에서 곧
따르고, 마지막 주장은 호몰로지 함자의 정의에서 따른다.
(\hyperref[III.7.3.3-fr]{7.3.3})을 함께 고려하면 (i)에서 (ii)의 앞 두
주장이 곧 따르며, $T_p(A)$가 자유라는 사실은
(\hyperref[III.7.5.4-fr]{7.5.4})에서 따른다.
\end{proof}

\begin{corollary}[7.5.6]
\label{III.7.5.6-fr}
(\hyperref[III.7.5.4-fr]{7.5.4})의 일반 가정들이 성립하고, 더 나아가
$A$가 이산 값매김환이라고 하자.
\begin{enumerate}
  \item[(i)] $T_p$가 오른쪽 완전일 필요충분조건은 $T_{p-1}(A)$가
    자유 $A$-가군인 것이다.
  \item[(ii)] $T_p$가 완전할 필요충분조건은 $T_p(A)$와 $T_{p-1}(A)$가
    자유 $A$-가군인 것이다.
\end{enumerate}
\end{corollary}

\begin{proof}
\oldpage[III]{192}
(ii)가 (i)를 함의함은 명백하다
((\hyperref[III.7.3.3-fr]{7.3.3}) 참조). (i)를 증명하기 위하여,
$f$가 $A$의 극대 아이디얼의 생성원(곧 $A$의 ``균등화원'')이라고
하자. 유한형 $A$-가군 $M$이 자유일(또는, 이 경우 같은 말인 평탄할)
필요충분조건은 $M$의 곱셈 준동형 $h_f:x\mapsto fx$가 단사인 것이다.
실제로 이는 여기서 $M$에 꼬임이 없다는 것과 동치이다
($0_{\mathrm I}$, 6.3.4). 이제 완전열
\[
  0\longrightarrow A\xrightarrow{h_f}A\longrightarrow k\longrightarrow0,
\]
을 생각하자. 이 열은 호몰로지 완전열
\[
  T_p(A)
  \longrightarrow T_p(k)
  \longrightarrow T_{p-1}(A)
  \xrightarrow{h_f}T_{p-1}(A).
\]
을 준다. 따라서 $T_{p-1}(A)$가 자유일 필요충분조건은
$T_p(A)\to T_p(k)$가 전사인 것이며, 결론은
(\hyperref[III.7.5.2-fr]{7.5.2})에서 따른다.
\end{proof}

\begin{remark}[7.5.7]
\label{III.7.5.7-fr}
(\hyperref[III.7.4.7-fr]{7.4.7})에 의하여
(\hyperref[III.7.4.4-fr]{7.4.4})의 일반 가정들은 (7.4.1.1)로 정의된
호몰로지 함자 $T_\bullet$이 (\hyperref[III.7.5.4-fr]{7.5.4})의 일반
가정들을 만족함을 함의한다. 따라서 이 경우
(\hyperref[III.7.5.6-fr]{7.5.6})은
(\hyperref[III.7.4.3-fr]{7.4.3})에 포함되어 있음에 유의하자.
\end{remark}

\subsection{완전성 성질들의 내림. 반연속성 정리와 Grauert의 완전성 판정법}
\label{subsection:III.7.6.fr}

\begin{proposition}[7.6.1]
\label{III.7.6.1-fr}
(7.4.1)의 조건들 아래에서 $B$를 가환 $A$-대수라 하자. $T_p$가 오른쪽
완전(각각 왼쪽 완전, 완전)이면 $T_p^{(B)}$도 그러하다. $B$가 충실
평탄 $A$-가군이면 그 역도 참이다.
\end{proposition}

\begin{proof}
첫 주장은 (7.1.3)의 자명한 주장의 특수한 경우이다. 역으로 먼저
$B$가 평탄 $A$-가군이라고 하자. 그러면 모든 $A$-가군 $M$에 대하여
\[
  H_\bullet\bigl(P_\bullet\otimes_A(M\otimes_A B)\bigr)
  =
  \bigl(H_\bullet(P_\bullet\otimes_A M)\bigr)\otimes_A B,
\]
이고, 이를 모든 $p$에 대하여 다시 쓰면 표준 동형을 제외하고
\[
  T_p(M)\otimes_A B
  =
  T_p^{(B)}(M_{(B)})
  \tag{7.6.1.1}\label{III.7.6.1.1.fr}
\]
이다. $T_p^{(B)}$가 오른쪽 완전(각각 왼쪽 완전, 완전)이라고 하자.
$M\mapsto M_{(B)}$는 완전 함자이므로 (7.6.1.1)의 왼쪽 항은 $M$에
대하여 오른쪽 완전(각각 왼쪽 완전, 완전)한 함자이다. 이제 $B$가
$A$ 위에서 충실 평탄이면 $T_p$가 같은 완전성 성질을 가짐이 따른다
($0_{\mathrm I}$, 6.4.1).
\end{proof}

\begin{proposition}[7.6.2]
\label{III.7.6.2-fr}
(7.4.1)의 조건들 아래에서, 더 나아가 $A$가 축소 뇌터 환이고 $P_i$들이
유한형 $A$-가군들이라고 가정하자. $T_p$가 오른쪽 완전(각각 왼쪽
완전, 완전)일 필요충분조건은, 이산 값매김환인 모든 $A$-대수 $B$에
대하여 $T_p^{(B)}$가 같은 성질을 갖는 것이다.
\end{proposition}

\begin{proof}
(\hyperref[III.7.3.1-fr]{7.3.1})과
(\hyperref[III.7.3.3-fr]{7.3.3})에 의하여 오른쪽 완전성만 고려하면
되고, 당연히 조건 (7.6.1)의 충분성만 증명하면 된다.
(\hyperref[III.7.4.2-fr]{7.4.2})에 의하여
$Z'_{p-1}(P_\bullet)$이 평탄 $A$-가군임을 보이면 충분하다.
$P_{p-1}$이 유한형이므로 $Z'_{p-1}(P_\bullet)$도 유한형이다. 따라서
판정법 ($0$, 10.2.8)에 의하면, 이산 값매김환인 모든 $A$-대수 $B$에
대하여 $Z'_{p-1}(P_\bullet)\otimes_A B$가 평탄 $B$-가군임을 보이면
충분하다. 그런데 $P_\bullet$은 평탄 $A$-가군들의 복합체이므로
\[
  Z'_{p-1}(P_\bullet)\otimes_A B
  =
  Z'_{p-1}(P_\bullet\otimes_A B);
\]
\oldpage[III]{193}
$P_\bullet\otimes_A B$는 평탄 $B$-가군들의 복합체이고
($0_{\mathrm I}$, 6.2.1), 모든 $B$-가군 $N$에 대하여
\[
  H_\bullet(P_\bullet\otimes_A N)
  =
  H_\bullet\bigl((P_\bullet\otimes_A B)\otimes_B N\bigr),
\]
따라서
\[
  T_p^{(B)}(N)
  =
  H_p\bigl((P_\bullet\otimes_A B)\otimes_B N\bigr);
\]
이다. 이제 (\hyperref[III.7.4.2-fr]{7.4.2})를 $T_p^{(B)}$에 적용하면,
$T_p^{(B)}$가 오른쪽 완전하다는 가정은
$Z'_{p-1}(P_\bullet\otimes_A B)$가 평탄 $B$-가군이라는 것과 동치임을
알 수 있다.

앞의 판정법은 이산 값매김환의 경우를 더 자세히 연구하도록 이끈다.
\end{proof}

\begin{proposition}[7.6.3]
\label{III.7.6.3-fr}
(7.4.1)의 조건들 아래에서 $A$가 1차원 정칙 뇌터 환이라고 가정하자.
다시 말해 $A$는 뇌터 환이고 모든 $x\in\operatorname{Spec}(A)$에
대하여 $A_x$는 체 또는 이산 값매김환이라고 하자. 그러면 모든 정수
$p$와 모든 $A$-가군 $M$에 대하여 $M$에 함자적인 표준 완전열
\[
  0
  \longrightarrow T_p(A)\otimes_A M
  \xrightarrow{t_M}T_p(M)
  \longrightarrow \operatorname{Tor}_1^A\bigl(T_{p-1}(A),M\bigr)
  \longrightarrow0.
  \tag{7.6.3.1}\label{III.7.6.3.1.fr}
\]
이 있다.
\end{proposition}

\begin{proof}
이하에서는 표기를 간단히 하기 위하여 통상적인 호몰로지 표기
$H_p(P_\bullet)$, $B_p(P_\bullet)$, $Z_p(P_\bullet)$,
$Z'_p(P_\bullet)$에서 복합체 $P_\bullet$의 언급을 생략하겠다. 다음
세 완전열이 있다.
\begin{gather*}
  0\longrightarrow H_p\longrightarrow Z'_p\longrightarrow B_{p-1}
  \longrightarrow0,\\
  0\longrightarrow B_{p-1}\longrightarrow Z_{p-1}\longrightarrow H_{p-1}
  \longrightarrow0,\\
  0\longrightarrow Z_{p-1}\longrightarrow P_{p-1}\longrightarrow B_{p-2}
  \longrightarrow0.
\end{gather*}
$P_{p-1}$과 $P_{p-2}$가 평탄하고, 모든
$x\in\operatorname{Spec}(A)$에 대하여 평탄 $A_x$-가군과 꼬임 없는
$A_x$-가군은 같은 것이므로, 이들의 각 부분가군 $B_{p-1}$,
$Z_{p-1}$, $B_{p-2}$도 평탄하다. 따라서 $M$과 텐서곱을 취하면 다음
완전열들을 얻는다.
\[
  0=\operatorname{Tor}_1^A(B_{p-1},M)
  \longrightarrow H_p\otimes M
  \longrightarrow Z'_p\otimes M
  \xrightarrow{u}B_{p-1}\otimes M
  \longrightarrow0
  \tag{7.6.3.2}\label{III.7.6.3.2.fr}
\]
\[
  0=\operatorname{Tor}_1^A(Z_{p-1},M)
  \longrightarrow\operatorname{Tor}_1^A(H_{p-1},M)
  \longrightarrow B_{p-1}\otimes M
  \xrightarrow{v}Z_{p-1}\otimes M
  \tag{7.6.3.3}\label{III.7.6.3.3.fr}
\]
\[
  0=\operatorname{Tor}_1^A(B_{p-2},M)
  \longrightarrow Z_{p-1}\otimes M
  \xrightarrow{w}P_{p-1}\otimes M.
  \tag{7.6.3.4}\label{III.7.6.3.4.fr}
\]
정의에 따라
$T_p(M)=\operatorname{Ker}(d_p\otimes1)/\operatorname{Im}(d_{p+1}\otimes1)$
이다. 따라서 이는 $d_p\otimes1$에서 몫으로 내려가 얻는 준동형
\[
  (P_p\otimes M)/\operatorname{Im}(d_{p+1}\otimes1)
  \longrightarrow P_{p-1}\otimes M
\]
의 핵이다. $Z'_p=P_p/B_p$의 정의에 따라 이 준동형은
$Z'_p\otimes M\to P_{p-1}\otimes M$으로도 쓸 수 있으며, 다시 합성
\[
  Z'_p\otimes M
  \xrightarrow{u}B_{p-1}\otimes M
  \xrightarrow{v}Z_{p-1}\otimes M
  \xrightarrow{w}P_{p-1}\otimes M.
\]
으로 볼 수 있다. (7.6.3.4)에 따라 $w$가 단사이므로 완전열
\[
  0\longrightarrow\operatorname{Ker}u
  \longrightarrow T_p(M)
  \longrightarrow\operatorname{Ker}v
  \longrightarrow0,
\]
을 얻는다. (7.6.3.2), (7.6.3.3)과 정의에 따른 $H_p=T_p(A)$를
고려하면, 이 열은 바로 (7.6.3.1)이다.
\end{proof}

\begin{remarks}[7.6.4]
\label{III.7.6.4-fr}
\begin{enumerate}
  \item[(i)] $H_\bullet(P_\bullet\otimes_A M)$은 이중복합체
    $P_\bullet\otimes_A M$의 호몰로지이다. 여기서 $M$은 차수 0의 항만
    있는 복합체로 간주한다. 따라서 이는 (6.3.6과 6.3.2)에 따라 $E_2$항이
    다음인 정칙 스펙트럼 열의 수렴대상이다.
    \[
      E^2_{pq}
      =
      \operatorname{Tor}_p^A\bigl(H_q(P_\bullet),M\bigr)
      =
      \operatorname{Tor}_p^A\bigl(T_q(A),M\bigr).
    \]
    \oldpage[III]{194}
    그런데 환 $A$에 대한 가정으로부터 임의의 $A$-가군 $E,F$와
    $p\geq2$에 대하여
    $\operatorname{Tor}_p^A(E,F)=0$임이 따른다
    ($0_{\mathrm{IV}}$, 17.2.2). 이것으로부터 완전열
    \[
      0\longrightarrow E^2_{0,q}
      \longrightarrow H_q(P_\bullet\otimes_A M)
      \longrightarrow E^2_{1,q-1}
      \longrightarrow0
    \]
    을 얻는다는 것이 알려져 있으며 (M, XV), 이 열은 바로 (7.6.3.1)이다.
  \item[(ii)] (7.3.1)을 고려하면, 완전열 (7.6.3.1)은 특수한 경우로
    (7.4.3)의 결과를 다시 준다.
\end{enumerate}
\end{remarks}

\begin{corollary}[7.6.5]
\label{III.7.6.5-fr}
(7.4.1)의 조건들 아래에서 $A$가 분수체 $K$와 잉여체 $k$를 갖는 이산
값매김환이고, $T_i(A)$들이 유한 생성 $A$-가군이라고 가정하자. 그러면
\[
  \operatorname{rang}_k T_p(k)
  \geq
  \operatorname{rang}_k\bigl(T_p(A)\otimes_A k\bigr)
  \geq
  \operatorname{rang}_A T_p(A)
  =
  \operatorname{rang}_K T_p(K).
  \tag{7.6.5.1}\label{III.7.6.5.1.fr}
\]
더 나아가 이 부등식의 양 끝 항이 같기 위한 필요충분조건은 $T_p$가
완전한 것이며, 이는 다시 $T_p(A)$와 $T_{p-1}(A)$가 자유
$A$-가군인 것과 동치이다.
\end{corollary}

\begin{proof}
실제로 완전열 (7.6.3.1)에서 $M=k$로 놓으면, 관련 가군들이 $k$ 위의
벡터공간이므로
\[
  \operatorname{rang}_k T_p(k)
  =
  \operatorname{rang}_k\bigl(T_p(A)\otimes_A k\bigr)
  +
  \operatorname{rang}_k\bigl(\operatorname{Tor}_1^A(T_{p-1}(A),k)\bigr).
\]
한편 $T_p(A)$는 국소 정역 $A$ 위의 유한 생성 가군이므로
(Bourbaki, \emph{Alg. comm.}, chap. II, \S 3,
n\textsuperscript{o} 2, cor. 1 de la prop. 4)
\[
  \operatorname{rang}_k\bigl(T_p(A)\otimes_A k\bigr)
  \geq
  \operatorname{rang}_A T_p(A)
  =
  \operatorname{rang}_K\bigl(T_p(A)\otimes_A K\bigr)
  \tag{7.6.5.2}\label{III.7.6.5.2.fr}
\]
이고, 또한 (7.6.5.2)의 두 항이 같기 위한 필요충분조건은 $T_p(A)$가
자유 $A$-가군인 것이다 (\emph{loc. cit.}, prop. 7). 한편 $K$는 평탄한
$A$-가군이므로 정의에 따라
\[
  T_p(A)\otimes_A K
  =
  H_p(P_\bullet)\otimes_A K
  =
  H_p(P_\bullet\otimes_A K)
  =
  T_p(K).
\]
따라서 부등식 (7.6.5.1)을 얻는다. 또한 등식이 성립하려면
1\textsuperscript{o} $T_p(A)$가 자유이고,
2\textsuperscript{o} $\operatorname{Tor}_1^A(T_{p-1}(A),k)=0$이어야
한다. 두 번째 조건은 알려진 바와 같이 ($0$, 10.1.3)
$T_{p-1}(A)$가 자유 $A$-가군이라는 것과 동치이다. 마지막으로
$T_i(A)$들은 유한 생성 $A$-가군이므로, 이들이 평탄하다는 것과
자유라는 것은 같은 뜻이며 (Bourbaki, \emph{Alg. comm.}, chap. II,
\S 3, n\textsuperscript{o} 2, cor. 2 de la prop. 5), (7.4.3)에 따라
결론을 얻는다.
\end{proof}

\begin{env}[7.6.6]
\label{III.7.6.6-fr}
계속해서 (7.4.1)의 가정들 아래에서, 모든
$x\in\operatorname{Spec}(A)$에 대하여 다음과 같이 놓는다.
\[
  d_p(x)=d_p^T(x)=\operatorname{rang}_{\kappa(x)}T_p(\kappa(x)).
  \tag{7.6.6.1}\label{III.7.6.6.1.fr}
\]
\end{env}

\begin{lemma}[7.6.7]
\label{III.7.6.7-fr}
$\varphi:A\to A'$을 환의 준동형이라 하고
\[
  f={}^a\!\varphi:
  \operatorname{Spec}(A')\longrightarrow\operatorname{Spec}(A)
\]
를 이에 대응하는 사상이라 하자 (I, 1.2.1). 만일
$T'_\bullet=T_\bullet^{(A')}$ (7.1.3)로 놓으면
\[
  d_p^{T'}=d_p^T\circ f.
  \tag{7.6.7.1}\label{III.7.6.7.1.fr}
\]
이다.
\end{lemma}

\begin{proof}
\oldpage[III]{195}
실제로 모든 $x'\in\operatorname{Spec}(A')$에 대하여 $x=f(x')$로 놓으면
\[
  H_\bullet(P_\bullet\otimes_A\kappa(x'))
  =
  H_\bullet\bigl((P_\bullet\otimes_A\kappa(x))
    \otimes_{\kappa(x)}\kappa(x')\bigr)
  =
  H_\bullet(P_\bullet\otimes_A\kappa(x))
    \otimes_{\kappa(x)}\kappa(x'),
\]
이는 $\kappa(x')$가 $\kappa(x)$ 위에서 평탄하기 때문이며,
따라서 관계 (7.6.7.1)을 얻는다.
\end{proof}

\begin{lemma}[7.6.8]
\label{III.7.6.8-fr}
환 $A$가 뇌터이고 복합체 $P_\bullet$이 유한 생성 $A$-가군들로
이루어져 있으면, $\operatorname{Spec}(A)$ 위의 함수
$x\mapsto d_p^T(x)$는 구성가능이다.
\end{lemma}

\begin{proof}
$X=\operatorname{Spec}(A)$의 모든 닫힌 기약 부분집합 $Y$에 대하여
$d_p$가 상수인 $Y$의 공집합이 아닌 열린집합 $U$가 존재함을 보이면
된다 ($0$, 9.2.2). $Y=\operatorname{Spec}(A/\mathfrak a)$이고 여기서
$\mathfrak a$는 $A/\mathfrak a$가 축소환이 되게 하는 $A$의
아이디얼이므로, (7.6.7)에 따라 $Y=X$이고 $A$가 뇌터 정역인 경우로
환원할 수 있다. 이 경우 주장은 (7.4.5)에서 따른다.
\end{proof}

\begin{theorem}[7.6.9]
\label{III.7.6.9-fr}
$A$를 뇌터 환, $P_\bullet$을 유한 생성 평탄 $A$-가군들로 이루어진
복합체, $T_\bullet(M)=H_\bullet(P_\bullet\otimes_A M)$을
$P_\bullet$이 정의하는 호몰로지 함자라 하자. 모든
$x\in\operatorname{Spec}(A)$에 대하여
\[
  d_p(x)=\operatorname{rang}_{\kappa(x)}T_p(\kappa(x)).
\]
로 놓는다. 그러면 다음이 성립한다.
\begin{enumerate}
  \item[(i)] 함수 $d_p$는 $\operatorname{Spec}(A)$에서 구성가능이고
    상반연속이다.
  \item[(ii)] $T_p$가 완전하면 $d_p$는 $\operatorname{Spec}(A)$에서
    연속이고, 따라서 국소 상수이다. 환 $A$가 축소이면 그 역도 참이다.
\end{enumerate}
\end{theorem}

\begin{proof}
\begin{enumerate}
  \item[(i)] 첫 번째 주장은 (7.6.8)에서 증명했다. 따라서 두 번째 주장을
    보이려면 ($0$, 9.2.4), $x'\ne x$가
    $\operatorname{Spec}(A)$에서 $x$의 일반화일 때
    $d_p(x')\leq d_p(x)$임을 보이면 충분하다. 이때 이산 값매김환 $B$와
    사상
    \[
      f:\operatorname{Spec}(B)\longrightarrow\operatorname{Spec}(A)
    \]
    이 존재하여, $a$가 $\operatorname{Spec}(B)$의 닫힌점이고 $b$가 그
    일반점이면 $f(a)=x$, $f(b)=x'$이다 (II, 7.1.9). 공식 (7.6.7.1)에
    따라 $\operatorname{Spec}(B)$에서 부등식 $d_p(a)\geq d_p(b)$을
    보이는 것으로 환원되는데, 이는 바로 부등식
    (7.6.5.1)이다\footnote{(i)의 증명을 이산 값매김환의 경우로 환원하는
    원리는 Hironaka가 구두로 알려 주었다.}.
  \item[(ii)] 첫 번째 주장은 이미 (7.3.4)에서 증명했다. 역을 증명하기
    위하여 값매김 판정법 (7.6.2)을 사용하자. 공식 (7.6.7.1)을 고려하면
    $A$가 이산 값매김환인 경우로 환원된다. 이때
    $\operatorname{Spec}(A)$는 두 점만 가지므로, $d_p$가 상수라는 가정은
    (7.6.5)에 따라 $T_p$가 완전함을 뜻한다.
\end{enumerate}
\end{proof}

\begin{corollary}[7.6.10]
\label{III.7.6.10-fr}
$A$를 뇌터 환, $\mathfrak p_i$ $(1\leq i\leq r)$를 그 극소 소아이디얼들,
$k_i$를 $A_{\mathfrak p_i}$의 잉여체 $(1\leq i\leq r)$라 하자.
\begin{enumerate}
  \item[(i)] 모든 $x\in\operatorname{Spec}(A)$에 대하여 어떤 지표 $i$가
    존재하여
    \[
      d_p(x)\geq\operatorname{rang}_{k_i}T_p(k_i).
      \tag{7.6.10.1}\label{III.7.6.10.1.fr}
    \]
    이다.
\end{enumerate}
특히 $A$가 정역이고 $K$가 그 분수체이면, 모든
$x\in\operatorname{Spec}(A)$에 대하여
\[
  d_p(x)\geq\operatorname{rang}_K T_p(K)
  \tag{7.6.10.2}\label{III.7.6.10.2.fr}
\]
이다.
\oldpage[III]{196}
\begin{enumerate}
  \item[(ii)] 더 나아가 $A$가 국소 축소환이고 $k$가 그 잉여체라고
    가정하자. $T_p$가 완전하기 위한 필요충분조건은
    \[
      \operatorname{rang}_k T_p(k)
      =
      \operatorname{rang}_{k_i}T_p(k_i)
      \qquad\text{모든 }1\leq i\leq r.
      \tag{7.6.10.3}\label{III.7.6.10.3.fr}
    \]
    이 성립하는 것이다.
\end{enumerate}
\end{corollary}

\begin{proof}
모든 $x$의 근방은 $\mathfrak p_i$ 중 하나를 포함하므로 (i)는
상반연속성의 정의에서 곧바로 따른다. 한편 $A$가 국소이면 극대
아이디얼 $\mathfrak m$의 $\operatorname{Spec}(A)$ 안에서의 유일한 근방은
$\operatorname{Spec}(A)$ 전체이므로, 모든
$x\in\operatorname{Spec}(A)$에 대하여
\[
  d_p(x)\leq\operatorname{rang}_k T_p(k)
\]
이다. 이 관계와 (i)를 합치면 조건 (7.6.10.3)에 따라 $d_p(x)$가
$\operatorname{Spec}(A)$에서 상수임을 알 수 있고, 따라서
(7.6.9, (ii))에 의하여 $T_p$가 완전하다. 역은 (7.6.9, (ii))에서
명백하다.
\end{proof}

\begin{remark}[7.6.11]
\label{III.7.6.11-fr}
(7.6.9, (i))의 주장을 모든 $x\in\operatorname{Spec}(A)$에 대한 부등식
\[
  \operatorname{rang}_{\kappa(x)}T_p(\kappa(x))
  \geq
  \operatorname{rang}_{\kappa(x)}
    \bigl(T_p(A)\otimes_A\kappa(x)\bigr)
  \tag{7.6.11.1}\label{III.7.6.11.1.fr}
\]
으로 강화할 수 있는지 물을 수 있다. 이 부등식은 $A$가 이산
값매김환이고 $x$가 그 극대 아이디얼일 때에는 실제로 성립한다 (7.6.5).
$A$가 극대 아이디얼 $\mathfrak m$과 잉여체 $k$를 갖는 뇌터 국소환인
경우로 한정하자. 그러면 다음 조건들은 서로 동치이다.
\begin{enumerate}
  \item[a)] 유한 생성 평탄 $A$-가군들로 이루어진 모든 복합체
    $P_\bullet$에 대하여
    \[
      \operatorname{rang}_k T_p(k)
      \geq
      \operatorname{rang}_k\bigl(T_p(A)\otimes_A k\bigr)
      \quad\text{모든 정수 }p.
      \tag{7.6.11.2}\label{III.7.6.11.2.fr}
    \]
  \item[b)] 모든 유한 생성 $A$-가군 $M$에 대하여
    \[
      \operatorname{rang}_k(M\otimes_A k)
      \geq
      \operatorname{rang}_k(\check M\otimes_A k).
      \tag{7.6.11.3}\label{III.7.6.11.3.fr}
    \]
  \item[c)] 모든 유한 생성 $A$-가군 $N$에 대하여
    \[
      \operatorname{rang}_k\bigl(\operatorname{Tor}_1^A(N,k)\bigr)
      \geq
      \operatorname{rang}_k\bigl(\operatorname{Tor}_2^A(N,k)\bigr).
      \tag{7.6.11.4}\label{III.7.6.11.4.fr}
    \]
\end{enumerate}
차수 이동에 의하면 (M, V, 7.2), 이는 모든 $i\geq1$에 대하여
\[
  \operatorname{rang}_k\bigl(\operatorname{Tor}_i^A(N,k)\bigr)
  \geq
  \operatorname{rang}_k\bigl(\operatorname{Tor}_{i+1}^A(N,k)\bigr).
  \tag{7.6.11.5}\label{III.7.6.11.5.fr}
\]
이 성립한다고 말하는 것과 같다.

증명에 관한 몇 가지 설명을 간단히 제시하자. a)가 b)를 함의함을 보기
위하여 완전열
\[
  L_1\xrightarrow{d}L_0\longrightarrow M\longrightarrow0
\]
을 잡는다. 여기서 $L_0$와 $L_1$은 유한 생성 자유 가군이다. a)를
복합체
\[
  P_1\xrightarrow{{}^t d}P_0,
  \qquad
  P_0=\check L_1,
  \quad
  P_1=\check L_0,
\]
에 적용하고 다른 항들은 영으로 둔다. 그러면 $T_1(A)=\check M$이고
\[
  T_1(k)
  =
  \operatorname{Hom}_A(M,k)
  =
  \operatorname{Hom}_k(M/\mathfrak mM,k),
\]
즉 $T_1(k)$는 벡터공간 $M\otimes_A k$의 쌍대이므로 이 벡터공간과
계수가 같다. b)가 c)를 함의함을 증명하기 위하여 먼저 다음 보조정리를
세운다.

\begin{lemma}[7.6.11.6]
\label{III.7.6.11.6-fr}
평탄 $A$-가군들로 이루어진 복합체
\[
  \cdots\longrightarrow0\longrightarrow
  P_1\xrightarrow{d}P_0
  \longrightarrow0\longrightarrow\cdots
\]
가 주어지면 완전열
\[
  0\longrightarrow\operatorname{Tor}_2^A(Z'_0,k)
  \longrightarrow T_1(A)\otimes_A k
  \longrightarrow T_1(k)
  \longrightarrow\operatorname{Tor}_1^A(Z'_0,k)
  \longrightarrow0.
  \tag{7.6.11.7}\label{III.7.6.11.7.fr}
\]
이 있다.
\end{lemma}

\begin{proof}
\oldpage[III]{197}
실제로 완전열 $0\to Z_1\to P_1\to B_0\to0$에서 완전열
\[
  0\longrightarrow\operatorname{Tor}_1^A(B_0,k)
  \longrightarrow Z_1\otimes k
  \longrightarrow P_1\otimes k
  \xrightarrow{u}B_0\otimes k
  \longrightarrow0.
\]
을 얻는다. 완전열 $0\to B_0\to Z_0\to Z'_0\to0$에서
$Z_0=P_0$가 평탄하므로
$\operatorname{Tor}_1^A(B_0,k)=\operatorname{Tor}_2^A(Z'_0,k)$이다.
정의에 따라 $Z_1=T_1(A)$이고, 마지막으로
$T_1(k)=\operatorname{Ker}(d\otimes1)$이며 $d\otimes1$은 다음과 같이
분해된다.
\[
  P_1\otimes k
  \xrightarrow{u}B_0\otimes k
  \xrightarrow{v}Z_0\otimes k
  =P_0\otimes k;
\]
$R=\operatorname{Ker}v$라 하면 $T_1(k)=u^{-1}(R)$이다. $u$가
전사이므로 $R=u(T_1(k))$이고, 정의에 따라
$R=\operatorname{Tor}_1^A(Z'_0,k)$이다. 이로써 완전열 (7.6.11.7)을
얻는다.
\end{proof}

이제 b)에서 c)를 이끌어 내기 위하여 완전열
\[
  L_1\xrightarrow{d}L_0\longrightarrow N\longrightarrow0,
\]
을 잡는다. 여기서 $L_0$와 $L_1$은 유한 생성 자유 가군이다.
$L_1$과 $L_0$로 이루어진 복합체에 대응하는 함자 $T$를 생각하자.
$L_0$와 $L_1$은 자유이므로 각각 자신의 이중쌍대와 동일시된다. 따라서
\[
  M=\operatorname{Coker}({}^t d),
  \qquad
  T_1(A)=\operatorname{Ker}(d)=\check M;
\]
이고, 한편
$M\otimes_A k=\operatorname{Coker}({}^t d\otimes1_k)$는
$\operatorname{Ker}(d\otimes1_k)$와 $k$ 위에서 계수가 같다. 따라서
가정 b)에 의하여
\[
  \operatorname{rang}_k\bigl(T_1(A)\otimes_A k\bigr)
  \leq
  \operatorname{rang}_k\bigl(T_1(k)\bigr);
\]
$Z'_0=N$이므로, 이제 완전열 (7.6.11.7)에서 부등식 (7.6.11.4)이
따른다.

마지막으로 c)가 a)를 함의함을 증명하기 위하여 보조정리 (7.6.11.6)을
두 항 복합체 $P_p\longrightarrow P_{p-1}$에 적용하자. 가정 c)를
$Z'_{p-1}$에 적용하면 $\operatorname{rang}_kR\geq
\operatorname{rang}_kS$를 얻는다. 여기서
\[
  R=\operatorname{Ker}(P_p\otimes k\longrightarrow P_{p-1}\otimes k)
  \quad\text{이고}\quad
  S=Z_p\otimes k.
\]
한편 $d_{p+1}:P_{p+1}\to P_p$를
\[
  P_{p+1}\xrightarrow{v}Z_p\xrightarrow{j}P_p,
\]
로 분해하면
\[
  \operatorname{Im}(d_{p+1}\otimes1)
  =
  (j\otimes1)\bigl(\operatorname{Im}(v\otimes1)\bigr).
\]
또한
\[
  T_p(A)\otimes_A k
  =(Z_p/B_p)\otimes_A k
  =(Z_p\otimes k)/\operatorname{Im}(v\otimes1),
\]
이고
\[
  T_p(k)=R/\operatorname{Im}(d_{p+1}\otimes1),
\]
이므로 부등식 (7.6.11.2)을 얻는다.

이제 국소환 $A$가 차원 $n$의 정칙환이라고 가정하자. 이때 [17]에 따라
$A$-가군 $\operatorname{Tor}_i^A(k,k)$는 외대수 거듭제곱
\[
  \bigwedge\nolimits^i(\mathfrak m/\mathfrak m^2);
\]
과 동형이다. 따라서 $n\geq4$이면 $N=k$에 대하여 조건 (7.6.11.4)가
성립하지 않는다.

반대로 국소 정역 $A$에서 모든 유한 생성 반사적 $A$-가군이 자유라고
가정하자. 차원 2의 정칙환은 이 조건을 만족한다. 그러면 조건
(7.6.11.3)이 성립한다. 실제로 유한 생성 $A$-가군 $M$의 쌍대
$\check M$은 반사적이므로 자유이고, 따라서
\[
  \operatorname{rang}_k(\check M\otimes_A k)
  =
  \operatorname{rang}_K(\check M)
  =
  \operatorname{rang}_K(M),
\]
이다. 여기서 $K$는 $A$의 분수체이다. 한편 $M\otimes_A k$의 모든
$k$-기저는 $M$의 한 생성계의 상들로 이루어진다는 것이 알려져 있으므로
(Bourbaki, \emph{Alg. comm.}, chap. II, \S 3,
n\textsuperscript{o} 2, cor. 2 de la prop. 4),
$\operatorname{rang}_K(M)\leq\operatorname{rang}_k(M\otimes_A k)$이다.
이로써 주장이 증명된다.
\end{remark}

\subsection{고유 사상에 대한 응용: I. 교환 성질}
\label{subsection:III.7.7.fr}

이어지는 세 번호는 본질적으로 앞 번호들의 결과를 준스킴 사상의 언어로
옮긴 것이다.

\begin{env}[7.7.1]
\label{III.7.7.1-fr}
$f:X\to Y$를 준스킴들의 준콤팩트 분리 사상이라 하고,
$\mathscr P_\bullet$을 미분 연산자의 차수가 $-1$인 준연접
$\mathscr O_X$-가군들의 복합체라 하자. 더 나아가
$\mathscr O_X$-가군 $\mathscr P_i$들이 $Y$-평탄이라고 가정하자
($0_{\mathrm I}$, 6.7.1).
\oldpage[III]{198}
준연접 $\mathscr O_Y$-가군의 범주에서 같은 범주로 가는 $\partial$-함자
$\mathscr M\mapsto\mathscr T_\bullet(\mathscr M)$을 생각하자. 이 함자는
$\mathscr T_\bullet(\mathscr P_\bullet,\mathscr M)$으로도 쓰며,
(6.2.3)에 따라 다음과 같이 정의된다.
\[
  \mathscr T_n(\mathscr P_\bullet,\mathscr M)
  =\mathscr T_n(\mathscr M)
  =\mathscr H^{-n}
    \bigl(f,\mathscr P^\bullet\otimes_{\mathscr O_Y}\mathscr M\bigr)
  \qquad\text{모든 }n\in\mathbf Z,
  \tag{7.7.1.1}\label{III.7.7.1.1.fr}
\]
여기서 $\mathscr P^\bullet$은 차수 $j$의 항이 $\mathscr P_{-j}$인
복합체이고, 따라서 그 미분 연산자의 차수는 $+1$이다. 이렇게 정의한
$\mathscr T_\bullet$은 $\mathscr M$에 대한 \emph{호몰로지 함자}이다
(6.2.6).
\end{env}

\begin{env}[7.7.2]
\label{III.7.7.2-fr}
$g:Y'\to Y$를 사상이라 하고
\[
  X'=X_{(Y')}=X\times_Y Y',
  \qquad
  f'=f_{(Y')}:X'\longrightarrow Y',
\]
로 놓자. 이는 준콤팩트 분리 사상이다. 또
\[
  \mathscr P'_\bullet
  =\mathscr P_\bullet\otimes_{\mathscr O_Y}\mathscr O_{Y'};
\]
로 놓자. 이는 (I, 9.1.12)와 ($0_{\mathrm I}$, 6.2.1)에 따라
$Y'$-평탄인 준연접 $\mathscr O_{X'}$-가군들의 복합체이다. 같은 차수
규약을 사용하여
\[
  \mathscr T^{Y'}_\bullet(\mathscr M')
  =\mathscr H^\bullet
    \bigl(f',\mathscr P'^\bullet\otimes_{\mathscr O_{Y'}}\mathscr M'\bigr)
  =\mathscr H^\bullet
    \bigl(f',\mathscr P^\bullet\otimes_{\mathscr O_Y}\mathscr M'\bigr)
  \tag{7.7.2.1}\label{III.7.7.2.1.fr}
\]
로 놓는다. 이는 준연접 $\mathscr O_{Y'}$-가군 $\mathscr M'$에 대한
호몰로지 함자이다. $Y'$이 환 $A'$의 아핀 스킴일 때에는
$\mathscr T^{Y'}_\bullet$ 대신 $\mathscr T^{A'}_\bullet$이라 쓴다. 모든
$A'$-가군 $M'$에 대하여
\[
  \mathscr T^{A'}_\bullet(\widetilde{M'})
  =\bigl(\Gamma(Y',\mathscr T^{Y'}_\bullet(\widetilde{M'}))\bigr)^\sim;
\]
이고,
\[
  T^{A'}_\bullet(M')
  =\Gamma(Y',\mathscr T^{Y'}_\bullet(\widetilde{M'})),
\]
로 놓는다. 이는 $A'$-가군의 범주에 값을 갖는 $A'$-가군의 호몰로지
함자이다.

$Y=\operatorname{Spec}(A)$도 아핀이면 $A'$-가군의 함자
$T^{A'}_\bullet$은 $A$-가군의 호몰로지 함자 $T^A_\bullet$에서
$A$로부터 $A'$으로 스칼라를 확장하여 얻은 함자와 일치한다 (7.1.3).
실제로 $A\to A'$에 대응하는 사상을 $g:Y'\to Y$, 그에 대응하는 아핀
사상을 $g':X'\to X$라 하자 (II, 1.6.2). $\mathfrak U$가 $X$의 아핀
열린 덮개이면 $\mathfrak U'=g'^{-1}(\mathfrak U)$는 $X'$의 아핀 열린
덮개이다. (6.2.2)에 따라 다음 등식만 보이면 충분하다.
\[
  C^\bullet
    \bigl(\mathfrak U,
      \mathscr P_\bullet\otimes_{\mathscr O_Y}g_*(\mathscr M')\bigr)
  =C^\bullet
    \bigl(\mathfrak U',
      \mathscr P_\bullet\otimes_{\mathscr O_Y}\mathscr M'\bigr),
\]
이는 다시 $X$의 모든 아핀 열린집합 $U$에 대하여
$U'=g'^{-1}(U)$로 놓았을 때
\[
  \Gamma
    \bigl(U,\mathscr P_\bullet\otimes_{\mathscr O_Y}g_*(\mathscr M')\bigr)
  =\Gamma
    \bigl(U',\mathscr P_\bullet\otimes_{\mathscr O_Y}\mathscr M'\bigr),
\]
임을 보이는 것으로 귀착되며, 이는 자명하다 (I, 1.3과 3.2).

특히 $U$가 $Y$의 열린집합이면 모든 준연접 $\mathscr O_Y$-가군
$\mathscr M$에 대하여
\[
  \mathscr T^U_\bullet(\mathscr M|U)
  =\bigl(\mathscr T_\bullet(\mathscr M)\bigr)|U.
  \tag{7.7.2.2}\label{III.7.7.2.2.fr}
\]
이다.
\end{env}

\begin{env}[7.7.3]
\label{III.7.7.3-fr}
모든 준연접 $\mathscr O_Y$-가군 $\mathscr M$에 대하여 $\mathscr M$에
함자적인 표준 준동형
\[
  \mathscr T_p(\mathscr O_Y)\otimes_{\mathscr O_Y}\mathscr M
  \longrightarrow
  \mathscr T_p(\mathscr M).
  \tag{7.7.3.1}\label{III.7.7.3.1.fr}
\]
이 있다. 실제로 $Y$가 아핀이면 이 준동형은 (7.2.2)에서 정의했다.
일반적인 경우에도 어렵지 않게 확장된다. 이를 위해 $U$, $V$가
$V\subset U$인 $Y$의 두 아핀 열린집합이면 그림
\oldpage[III]{199}
\[
\xymatrix@C=18pt@R=32pt{
  {\begin{gathered}
    \bigl(\mathscr T_p(\mathscr O_Y)
      \otimes_{\mathscr O_Y}\mathscr M\bigr)|U\\
    =\mathscr T^U_p(\mathscr O_Y|U)
      \otimes_{\mathscr O_Y|U}(\mathscr M|U)
  \end{gathered}}
  \ar[r]\ar[d]
  &
  {\begin{gathered}
    \mathscr T^U_p(\mathscr M|U)\\
    =\bigl(\mathscr T_p(\mathscr M)\bigr)|U
  \end{gathered}}
  \ar[d]
  \\
  {\begin{gathered}
    \bigl(\mathscr T_p(\mathscr O_Y)
      \otimes_{\mathscr O_Y}\mathscr M\bigr)|V\\
    =\mathscr T^V_p(\mathscr O_Y|V)
      \otimes_{\mathscr O_Y|V}(\mathscr M|V)
  \end{gathered}}
  \ar[r]
  &
  {\begin{gathered}
    \mathscr T^V_p(\mathscr M|V)\\
    =\bigl(\mathscr T_p(\mathscr M)\bigr)|V
  \end{gathered}}
}
\]
이 (7.2.3.3)에 따라 가환함을 주의하면 된다.

모든 사상 $g:Y'\to Y$에 대하여 표준 준동형
\[
  \mathscr T_p(\mathscr O_Y)\otimes_{\mathscr O_Y}\mathscr O_{Y'}
  \longrightarrow
  \mathscr T^{Y'}_p(\mathscr O_{Y'}),
  \tag{7.7.3.2}\label{III.7.7.3.2.fr}
\]
이 있다. 이는 $S=Y$, $v_1=f$, $v_2=1_Y$이고
$\mathscr P^{(2)}_\bullet$이 차수 0의 유일한 항 $\mathscr M$으로
이루어진 경우에, 수렴대상에 대한 (6.7.11.2)의 특수한 경우이다.

$Y=\operatorname{Spec}(A)$와 $Y'=\operatorname{Spec}(A')$이 아핀이면,
(7.7.3.2)는 (7.2.2)에서 정의한 $A'$-가군의 표준 준동형
\[
  T^A_p(A)\otimes_A A'
  \longrightarrow
  T^{A'}_p(A')=T^A_p(A')
\]
에 대응하는 층의 준동형이다. 이는 (6.7.11)에서 쉽게 따르는데, 현재
경우에는 (6.7.11)에서
$\mathfrak U'^{(i)}=u_i^{-1}(\mathfrak U^{(i)})$로 잡을 수 있기 때문이다.
\end{env}

\begin{env}[7.7.4]
$f$가 \emph{고유} 사상이고 $Y=\operatorname{Spec}(A)$가 뇌터 아핀
스킴이며 $\mathscr P_\bullet$이 연접이고 $Y$-평탄인
$\mathscr O_X$-가군들의 \emph{아래로 유계인} 복합체일 때, (6.10.5)에서
동형을 제외하면 다음과 같이 쓸 수 있음을 보았다.
\[
  \mathscr T_p(\mathscr M)
  =\mathscr H_p
    \bigl(\mathscr L_\bullet\otimes_{\mathscr O_Y}\mathscr M\bigr),
  \qquad
  \mathscr L_\bullet=\widetilde{L_\bullet},
\]
여기서 $L_\bullet$은 \emph{유한 생성 자유 $A$-가군들의 아래로 유계인
복합체}이다. 따라서 함자 $\mathscr T_\bullet$은 (7.4)와 (7.6)에서
상세히 연구한 유형이다. 이제 그 연구의 결과를 옮기겠다.
\end{env}

\begin{theorem}[7.7.5]
$Y$를 국소 뇌터 준스킴, $(U_\alpha)$를 아핀 열린집합들로 이루어진
$Y$의 덮개, $f:X\to Y$를 고유 사상, $\mathscr P_\bullet$을 연접이고
$Y$-평탄인 $\mathscr O_X$-가군들의 아래로 유계인 복합체라 하자.
(7.7.1.1)로 정의한 호몰로지 함자 $\mathscr T_\bullet(\mathscr M)$은
다음 성질들을 갖는다.

\medskip
\noindent
\textnormal{I) (상반연속성).}\footnote{이 정리의 특수한 경우는 이미
Chow--Igusa의 논문 [3]에 있다. 상반연속성은 해석공간의 틀에서 (그리고
상당히 특수한 가정 아래) Kodaira--Spencer가 발견하였고
(« On the variations of almost-complex structures », \emph{Algebraic Geometry
and Topology, A Symposium in honor of S. Lefschetz}, Princeton Series
n\textsuperscript{o} 12, p. 139--150, Princeton, 1957), 일반적인 판본은
Grauert [5]가 증명하였다.}
\emph{함수}
\[
  y\longmapsto
  d_p(y)
  =
  \operatorname{rang}_{\kappa(y)}
    \mathscr T_p^{\kappa(y)}(\kappa(y))
  \tag{7.7.5.1}\label{III.7.7.5.1.fr}
\]
\emph{는 상반연속이다.}

\oldpage[III]{200}
\medskip
\noindent
\textnormal{II) (교환 성질).}
\emph{주어진 정수 $p$에 대하여 다음 조건들은 서로 동치이다.}
\begin{enumerate}
  \item[a)] $\mathscr T_p$는 오른쪽 완전하다.

  \item[a$'$)] $\mathscr T_p(\mathscr M)$은
    $\mathscr N\otimes_{\mathscr O_Y}\mathscr M$ 꼴의 함자와 동형이고,
    $\mathscr N$은 반드시 다음과 동형이다.
    \[
      \mathscr T_p(\mathscr O_Y)
      =
      \mathscr H^{-p}(f,\mathscr P^\bullet).
    \]

  \item[a$''$)] 함자적인 표준 준동형 (7.7.3.1)은 동형사상이다.

  \item[b)] $\mathscr T_{p-1}$은 왼쪽 완전하다.

  \item[b$'$)] $\mathscr O_Y$-가군 $\mathscr Q$와 함자들의 동형사상
    \[
      \mathscr T_{p-1}(\mathscr M)
      \xrightarrow{\sim}
      \mathscr{H}\!om_{\mathscr O_Y}(\mathscr Q,\mathscr M).
      \tag{7.7.5.2}\label{III.7.7.5.2.fr}
    \]
    이 존재한다. 여기서 $\mathscr Q$는 반드시 연접이고 유일한 동형까지
    결정된다.

  \item[c)] 아핀 열린집합 $U_\alpha$의 환을 $A_\alpha$라 할 때,
    모든 지표 $\alpha$에 대하여 $A_\alpha$-가군의 함자
    $T_p^{A_\alpha}$는 오른쪽 완전하다.

  \item[d)] 모든 사상 $g:Y'\to Y$에 대하여 표준 준동형
    \[
      \mathscr T_p(\mathscr O_Y)
      \otimes_{\mathscr O_Y}\mathscr O_{Y'}
      \longrightarrow
      \mathscr T_p^{Y'}(\mathscr O_{Y'})
      \tag{7.7.5.3}\label{III.7.7.5.3.fr}
    \]
    은 동형사상이다.
\end{enumerate}
\end{theorem}

\begin{proof}
상반연속성은 $Y$ 위에서 국소적인 성질이므로, (7.7.4)의 주의와
(7.6.9)에서 따른다. a$''$)가 a$'$)를 함의하고 a$'$)가 a)를 함의함은
분명하다. $Y$가 아핀일 때에는 (7.7.4)의 주의를 고려하면 a), a$''$),
b), b$'$)의 동치가 (7.3.1)과 (7.4.6)에서 증명되었다. 일반적인 경우로
넘어가기 위하여 먼저 a)와 c)가 동치임을 증명하자. 그러면 성질 a)가
$Y$ 위에서 국소적임이 증명된다. 같은 증명을 a$''$)와 b)의 국소성을
증명하는 데에도 적용할 수 있다. c)가 a)를 함의함은 분명하므로 역만
증명하면 된다. $Y$의 모든 아핀 열린집합 $U$와 준연접
$(\mathscr O_Y|U)$-가군들의 모든 완전열
\[
  0\longrightarrow\mathscr F'
  \longrightarrow\mathscr F
  \longrightarrow\mathscr F''
  \longrightarrow0
\]
에 대하여, 다음을 만족하는 준연접 $\mathscr O_Y$-가군들의 완전열
\[
  0\longrightarrow\mathscr G'
  \longrightarrow\mathscr G
  \longrightarrow\mathscr G''
  \longrightarrow0
\]
이 존재함을 보이면 충분함은 분명하다.
\[
  \mathscr F'=\mathscr G'|U,
  \qquad
  \mathscr F=\mathscr G|U,
  \qquad
  \mathscr F''=\mathscr G''|U;
\]
그런데 이는 $Y$가 국소 뇌터라는 가정과 (I, 9.4.2)에서 곧바로 따른다.
실제로 $\mathscr F$를 준연접 $\mathscr O_Y$-가군 $\mathscr G$로
연장하고, $\mathscr F'$을 $\mathscr G$의 부분 $\mathscr O_Y$-가군
$\mathscr G'$으로 연장한 뒤
$\mathscr G''=\mathscr G/\mathscr G'$으로 잡으면 충분하다.

일반적인 경우에 b)와 b$'$)의 동치를 증명하기 위하여, $Y$가 아핀이면
$\mathscr Q$가 유일한 동형까지 결정됨을 알고 있다는 점에 주의하자.
따라서 $U$가 아핀 스킴 $Y$의 아핀 열린집합이면 함자적인 동형사상
\[
  \mathscr T_{p-1}^{U}(\mathscr M|U)
  \xrightarrow{\sim}
  \mathscr{H}\!om_{\mathscr O_Y|U}
    (\mathscr Q|U,\mathscr M|U).
\]
이 존재한다. 일반적인 경우에는 $Y$의 모든 아핀 열린집합 $U$에 대하여
연접 $(\mathscr O_Y|U)$-가군 $\mathscr Q_U$와 함자적인 동형사상
\[
  \mathscr T_{p-1}^{U}(\mathscr M|U)
  \xrightarrow{\sim}
  \mathscr{H}\!om_{\mathscr O_Y|U}
    (\mathscr Q_U,\mathscr M|U);
\]
이 있다. 앞의 주의에 따르면 $V$가 $U$에 포함되는 아핀 열린집합이면
$\mathscr Q_U|V=\mathscr Q_V$이다. 따라서 (7.7.5.2)를 만족하는
$\mathscr O_Y$-가군 $\mathscr Q$가 존재하고 유일하다.

마지막으로 a)와 d)의 동치를 증명하면 된다. d)가 $Y$ 위에서 국소적인
성질임은 분명하고, 위에서 a)도 그러함을 보았다. 더 나아가 d)는 $Y'$
위에서도 국소적이다. 이제 $Y=\operatorname{Spec}(A)$,
$Y'=\operatorname{Spec}(A')$일 때 $T_\bullet^{A'}$은
$T_\bullet^A$에서 $A'$으로 스칼라를 확장하여 얻는 함자임을 보았다.

\oldpage[III]{201}
따라서 a$'$)이면 (7.7.5.3)이 동형사상임이 분명하다. 역으로 여전히
$Y=\operatorname{Spec}(A)$가 아핀이라고 가정하고, $M$을 임의의
$A$-가군이라 하자. 곱셈을
\[
  (a_1,m_1)(a_2,m_2)
  =
  (a_1a_2,a_1m_2+a_2m_1);
\]
로 정의한 $A$-대수 $A'=A\oplus M$을 취하면
\[
  T_p^{A'}(A')
  =
  T_p(A\oplus M)
  =
  T_p(A)\oplus T_p(M),
\]
이다. (7.7.5.3)이 전단사라는 가정은 표준 사상
\[
  T_p(A)\otimes_A M\longrightarrow T_p(M),
\]
도 전단사임을 함의한다. 달리 말해 d)는 a$''$)를 함의하며, 이로써
증명이 끝난다.
\end{proof}

\begin{theorem}[7.7.6]
$Y$를 국소 뇌터 준스킴, $f:X\to Y$를 고유 사상,
$\mathscr F$를 연접이고 $Y$-평탄인 $\mathscr O_X$-가군이라 하자.
그러면 연접 $\mathscr O_Y$-가군 $\mathscr Q$와 준연접
$\mathscr O_Y$-가군 $\mathscr M$에 관한 함자들의 동형사상
\[
  f_*(\mathscr F\otimes_{\mathscr O_Y}\mathscr M)
  \xrightarrow{\sim}
  \mathscr{H}\!om_{\mathscr O_Y}(\mathscr Q,\mathscr M)
  \tag{7.7.6.1}\label{III.7.7.6.1.fr}
\]
이 존재하며, $\mathscr Q$는 유일한 동형까지 결정된다. 따라서 함자들의
동형사상
\[
  \Gamma(X,\mathscr F\otimes_{\mathscr O_Y}\mathscr M)
  \xrightarrow{\sim}
  \operatorname{Hom}_{\mathscr O_Y}(\mathscr Q,\mathscr M).
  \tag{7.7.6.2}\label{III.7.7.6.2.fr}
\]
도 얻는다.
\end{theorem}

\begin{proof}
실제로 $\mathscr M\mapsto\mathscr F\otimes_{\mathscr O_Y}\mathscr M$은
완전하고 ($0_{\mathrm I}$, 6.7.4), $f_*$는 왼쪽 완전하므로 함자
$\mathscr M\mapsto f_*(\mathscr F\otimes_{\mathscr O_Y}\mathscr M)$은
왼쪽 완전하다. 이제 $p=1$에 대하여 (7.7.5, b))와 (7.7.5, b$'$))의
동치를 적용하면 충분하다.
\end{proof}

\begin{corollary}[7.7.7]
$Y$를 국소 뇌터 준스킴, $f:X\to Y$를 고유 사상,
$\mathscr F$, $\mathscr F'$을 연접이고 $Y$-평탄인 두
$\mathscr O_X$-가군, $u:\mathscr F\to\mathscr F'$을 준동형이라 하자.
준연접 $\mathscr O_Y$-가군 $\mathscr M$에 관한 두 함자
\begin{align*}
  \mathscr T(\mathscr M)
  &=
  \operatorname{Ker}
  \bigl(
    f_*(\mathscr F\otimes_{\mathscr O_Y}\mathscr M)
    \longrightarrow
    f_*(\mathscr F'\otimes_{\mathscr O_Y}\mathscr M)
  \bigr),\\
  T(\mathscr M)
  &=
  \Gamma(Y,\mathscr T(\mathscr M))\\
  &=
  \operatorname{Ker}
  \bigl(
    \Gamma(X,\mathscr F\otimes_{\mathscr O_Y}\mathscr M)
    \longrightarrow
    \Gamma(X,\mathscr F'\otimes_{\mathscr O_Y}\mathscr M)
  \bigr).
\end{align*}
을 생각하자. 그러면 연접 $\mathscr O_Y$-가군 $\mathscr R$과 함자들의
동형사상
\[
  \mathscr T(\mathscr M)
  \xrightarrow{\sim}
  \mathscr{H}\!om_{\mathscr O_Y}(\mathscr R,\mathscr M)
  \tag{7.7.7.1}\label{III.7.7.7.1.fr}
\]
및
\[
  T(\mathscr M)
  \xrightarrow{\sim}
  \operatorname{Hom}_{\mathscr O_Y}(\mathscr R,\mathscr M).
  \tag{7.7.7.2}\label{III.7.7.7.2.fr}
\]
이 존재하며, $\mathscr R$은 유일한 동형까지 결정된다.
\end{corollary}

\begin{proof}
(7.7.7.2)만 증명하면 충분하다. 실제로 이는 $Y$가 아핀인 경우에
(7.7.7.1)을 증명하며, 표현 가능한 함자의 표현 대상이 유일한 동형까지
유일하다는 사실 (0, 8.1.8)에 의하여 (7.7.5, b)와 b$'$))의 동치의
증명에서와 같이 논증하면 여기서 일반적인 경우로 넘어갈 수 있다.
(7.7.6)에 따라 함자적인 동형사상
\[
  \Gamma(X,\mathscr F\otimes_{\mathscr O_Y}\mathscr M)
  \xrightarrow{\sim}
  \operatorname{Hom}_{\mathscr O_Y}(\mathscr Q,\mathscr M),
  \qquad
  \Gamma(X,\mathscr F'\otimes_{\mathscr O_Y}\mathscr M)
  \xrightarrow{\sim}
  \operatorname{Hom}_{\mathscr O_Y}(\mathscr Q',\mathscr M).
\]
을 정의하는 두 연접 $\mathscr O_Y$-가군 $\mathscr Q$, $\mathscr Q'$이
존재한다. 한편 $u:\mathscr F\to\mathscr F'$은 표준적으로 함자들의 사상
\[
  \Gamma(X,\mathscr F\otimes_{\mathscr O_Y}\mathscr M)
  \longrightarrow
  \Gamma(X,\mathscr F'\otimes_{\mathscr O_Y}\mathscr M);
\]
을 정의한다.

\oldpage[III]{202}
이에 대응하여 다음 그림이 가환하도록 하는 $\mathscr O_Y$-가군의 유일한
준동형 $v:\mathscr Q'\to\mathscr Q$가 있다 (0, 8.1.4).
\[
\xymatrix@C=42pt@R=28pt{
  \Gamma(X,\mathscr F\otimes_{\mathscr O_Y}\mathscr M)
    \ar[r]\ar[d]_-{\sim}
  &
  \Gamma(X,\mathscr F'\otimes_{\mathscr O_Y}\mathscr M)
    \ar[d]^-{\sim}
  \\
  \operatorname{Hom}_{\mathscr O_Y}(\mathscr Q,\mathscr M)
    \ar[r]
  &
  \operatorname{Hom}_{\mathscr O_Y}(\mathscr Q',\mathscr M)
}
\]
반변 함자
$\mathscr N\mapsto\operatorname{Hom}_{\mathscr O_Y}(\mathscr N,\mathscr M)$은
$\mathscr O_Y$-가군의 범주에서 왼쪽 완전하므로,
$\mathscr R=\operatorname{Coker}(v)$로 잡으면 요구한 동형사상
(7.7.7.2)를 얻는다.
\end{proof}

\begin{corollary}[7.7.8]
$X$, $Y$, $f$에 관한 (7.7.6)의 가정 아래, $\mathscr F$와
$\mathscr G$를 다음 조건을 만족하는 두 연접 $\mathscr O_X$-가군이라
하자. (i) $\mathscr F$는 $Y$-평탄이다. (ii) $\mathscr G$는 유한 계수의
국소 자유 $\mathscr O_X$-가군들 사이의 준동형
$v:\mathscr E_1\to\mathscr E_0$의 여핵과 동형이다. 준연접
$\mathscr O_Y$-가군 $\mathscr M$에 관한 두 함자
\begin{align*}
  \mathscr T(\mathscr M)
  &=
  f_*
  \bigl(
    \mathscr{H}\!om_{\mathscr O_X}
      (\mathscr G,\mathscr F\otimes_{\mathscr O_Y}\mathscr M)
  \bigr),\\
  T(\mathscr M)
  &=
  \Gamma(Y,\mathscr T(\mathscr M))
  =
  \operatorname{Hom}_{\mathscr O_X}
    (\mathscr G,\mathscr F\otimes_{\mathscr O_Y}\mathscr M).
\end{align*}
을 생각하자. 그러면 연접 $\mathscr O_Y$-가군 $\mathscr N$과 함자들의
동형사상
\[
  \mathscr T(\mathscr M)
  \xrightarrow{\sim}
  \mathscr{H}\!om_{\mathscr O_Y}(\mathscr N,\mathscr M)
  \tag{7.7.8.1}\label{III.7.7.8.1.fr}
\]
및
\[
  T(\mathscr M)
  \xrightarrow{\sim}
  \operatorname{Hom}_{\mathscr O_Y}(\mathscr N,\mathscr M).
  \tag{7.7.8.2}\label{III.7.7.8.2.fr}
\]
이 존재하며, $\mathscr N$은 유일한 동형까지 결정된다.
\end{corollary}

\begin{proof}
함자적인 동형사상 ($0_{\mathrm I}$, 5.4.2.1)에 따라 $i=0,1$일 때
$\mathscr M$에 관하여 함자적인 동형사상
\begin{align*}
  \mathscr{H}\!om_{\mathscr O_X}
    (\mathscr E_i,\mathscr F\otimes_{\mathscr O_Y}\mathscr M)
  &\xrightarrow{\sim}
  \check{\mathscr E}_i
    \otimes_{\mathscr O_X}
    (\mathscr F\otimes_{\mathscr O_Y}\mathscr M)\\
  &\xrightarrow{\sim}
  (\check{\mathscr E}_i\otimes_{\mathscr O_X}\mathscr F)
    \otimes_{\mathscr O_Y}\mathscr M\\
  &\xrightarrow{\sim}
  \mathscr{H}\!om_{\mathscr O_X}(\mathscr E_i,\mathscr F)
    \otimes_{\mathscr O_Y}\mathscr M
\end{align*}
을 얻는다. $i=0,1$에 대하여
$\mathscr F_i=\mathscr{H}\!om_{\mathscr O_X}(\mathscr E_i,\mathscr F)$로
놓자. 이들은 연접 $\mathscr O_X$-가군이고 ($0_{\mathrm I}$, 5.3.5),
$Y$-평탄이다 ($0_{\mathrm I}$, 5.4.2). 또
$u=\mathscr{H}\!om(v,1_{\mathscr F}):\mathscr F_0\to\mathscr F_1$로 놓자.
함자
$\mathscr H\mapsto\mathscr{H}\!om_{\mathscr O_X}
(\mathscr H,\mathscr F\otimes_{\mathscr O_Y}\mathscr M)$의 왼쪽 완전성에
따라 $\mathscr M$에 관하여 함자적인 동형사상
\begin{align*}
  \mathscr{H}\!om_{\mathscr O_X}
    (\mathscr G,\mathscr F\otimes_{\mathscr O_Y}\mathscr M)
  &\xrightarrow{\sim}
  \operatorname{Ker}
  \bigl(
    \mathscr{H}\!om_{\mathscr O_X}
      (\mathscr E_0,\mathscr F\otimes_{\mathscr O_Y}\mathscr M)
    \longrightarrow
    \mathscr{H}\!om_{\mathscr O_X}
      (\mathscr E_1,\mathscr F\otimes_{\mathscr O_Y}\mathscr M)
  \bigr)\\
  &\xrightarrow{\sim}
  \operatorname{Ker}
  \bigl(
    \mathscr F_0\otimes_{\mathscr O_Y}\mathscr M
    \longrightarrow
    \mathscr F_1\otimes_{\mathscr O_Y}\mathscr M
  \bigr).
\end{align*}
을 얻는다. $f_*$가 왼쪽 완전하므로 여기서 함자적인 동형사상
\[
  f_*
  \bigl(
    \mathscr{H}\!om_{\mathscr O_X}
      (\mathscr G,\mathscr F\otimes_{\mathscr O_Y}\mathscr M)
  \bigr)
  \xrightarrow{\sim}
  \operatorname{Ker}
  \bigl(
    f_*(\mathscr F_0\otimes_{\mathscr O_Y}\mathscr M)
    \longrightarrow
    f_*(\mathscr F_1\otimes_{\mathscr O_Y}\mathscr M)
  \bigr)
\]
을 얻으며, 이제 (7.7.7)을 적용하면 충분하다.
\end{proof}

\oldpage[III]{203}
\begin{remarks}[7.7.9]
\textnormal{(i)} (7.7.6), (7.7.7), (7.7.8)에서
$\mathscr O_Y$-가군 $\mathscr Q$, $\mathscr R$, $\mathscr N$의 구성은
\emph{밑변경과 가환한다}. 예를 들어 (7.7.2)의 표기를 유지하면,
(7.7.6)의 경우 모든 준연접 $\mathscr O_{Y'}$-가군 $\mathscr M'$에
대하여 동형사상
\[
  f'_*(\mathscr F\otimes_{\mathscr O_Y}\mathscr M')
  \xrightarrow{\sim}
  \mathscr{H}\!om_{\mathscr O_{Y'}}(g^*(\mathscr Q),\mathscr M')
\]
이 있다. 실제로 (7.7.2)의 주의에 따라 이는 등식
\[
  \operatorname{Hom}_{\mathscr O_Y}(\mathscr Q,g_*(\mathscr M'))
  =
  \operatorname{Hom}_{\mathscr O_{Y'}}(g^*(\mathscr Q),\mathscr M')
\]
을 보이는 것으로 귀착되며, 이는 바로 ($0_{\mathrm I}$, 4.4.3.1)이다.
마찬가지로 (7.7.7)에서 $Y$, $f$, $\mathscr M$, $\mathscr F$,
$\mathscr F'$을 각각 $Y'$, $f'$, $\mathscr M'$,
$\mathscr F\otimes_{\mathscr O_X}\mathscr O_{X'}$,
$\mathscr F'\otimes_{\mathscr O_X}\mathscr O_{X'}$으로 바꾸면
$\mathscr R$을 $g^*(\mathscr R)$로 바꾸어야 한다. 마지막으로 (7.7.8)에서
$X$, $Y$, $f$, $\mathscr M$, $\mathscr F$를 각각 $X'$, $Y'$, $f'$,
$\mathscr M'$, $\mathscr F\otimes_{\mathscr O_X}\mathscr O_{X'}$으로
바꾸고 $\mathscr G$를
$\mathscr G'=\mathscr G\otimes_{\mathscr O_X}\mathscr O_{X'}$으로 바꾸면,
$\mathscr N$을 $g^*(\mathscr N)$로 바꾸어야 한다. 이는 여전히 완전열
\[
  \mathscr E'_1\longrightarrow
  \mathscr E'_0\longrightarrow
  \mathscr G'\longrightarrow0
  \qquad\text{여기서}\qquad
  \mathscr E'_i=\mathscr E_i\otimes_{\mathscr O_X}\mathscr O_{X'}
\]
이 있고
\[
  \check{\mathscr E}'_i
  \otimes_{\mathscr O_{X'}}
  (\mathscr F\otimes_{\mathscr O_Y}\mathscr M')
  =
  \check{\mathscr E}_i
  \otimes_{\mathscr O_X}
  (\mathscr F\otimes_{\mathscr O_Y}\mathscr M')
  \qquad(i=0,1).
\]
이기 때문이다.

\medskip
\textnormal{(ii)} $Y$에 관해 풍부한 가역 $\mathscr O_X$-가군이 존재하면,
임의의 연접 $\mathscr O_X$-가군 $\mathscr G$는 언제나 (7.7.8)의 조건
(ii)를 만족한다. 예를 들어 $Y$가 아핀이고 $f:X\to Y$가 사영 사상이면
그러하다. 실제로 (II, 5.5.1)을 고려하면, $\mathscr G$가 유한 계수의
국소 자유 $\mathscr O_X$-가군 $\mathscr E_0$의 몫과 동형이 되도록 하는
$\mathscr E_0$가 존재한다 (II, 2.7.10). $\mathscr E_0$와 $\mathscr G$가
연접이므로 $\mathscr E_0\to\mathscr G$의 핵 $\mathscr G_1$도 연접이다.
같은 결과를 $\mathscr G_1$에 적용하면, $\mathscr E_0$와
$\mathscr E_1$이 유한 계수의 국소 자유인 완전열
$\mathscr E_1\to\mathscr E_0\to\mathscr G\to0$을 얻는다.

\medskip
\textnormal{(iii)} 제V장에서 (7.7.8)의 제한적인 가정 (ii)가
\emph{불필요함}을 증명할 것이다.
\end{remarks}

\begin{proposition}[7.7.10]
\textnormal{(교환 성질의 국소 판정법.)} (7.7.5)의 일반적인 조건 아래,
$y$를 $Y$의 점, $p$를 정수라 하자. 다음 성질들은 서로 동치이다.
\begin{enumerate}
  \item[a)] 함자 $T_p^{\mathscr O_y}$는 오른쪽 완전하다.

  \item[b)] 표준 준동형
    \[
      T_p^{\mathscr O_y}(\mathscr O_y)
      \longrightarrow
      T_p^{\mathscr O_y}(\kappa(y))
    \]
    은 전사이다.

  \item[c)] 모든 정수 $n$에 대하여 표준 준동형
    \[
      T_p^{\mathscr O_y}
      (\mathscr O_y/\mathfrak m_y^{n+1})
      \longrightarrow
      T_p^{\mathscr O_y}(\kappa(y))
    \]
    은 전사이다.
\end{enumerate}
더 나아가 이 조건들을 만족하는 $y\in Y$들의 집합은
$\mathscr T_p^U$가 오른쪽 완전하도록 하는 $Y$의 가장 큰 열린집합
$U$이다.
\end{proposition}

\begin{proof}
(7.7.4)를 고려하면 a), b), c)의 동치는 (7.4.7)과 (7.5.2)에서 따른다.
$T_p^{\mathscr O_y}$가 오른쪽 완전인 점들의 집합 $U$가 열려 있음도
(7.4.4)의 결과이다. 역으로 $\mathscr T_p^V$가 오른쪽 완전이면,
(7.7.5)의 조건 c)와 (7.6.1)에 따라 모든 $y\in V$에 대하여
$T_p^{\mathscr O_y}$도 오른쪽 완전하다.
\end{proof}

\begin{corollary}[7.7.11]
$\mathscr T_p$가 오른쪽 완전(각각 왼쪽 완전)이면, 모든 사상
$g:Y'\to Y$에 대하여 $\mathscr T_p^{Y'}$도 오른쪽 완전(각각 왼쪽
완전)이다. $g$가 충실 평탄이면 그 역도 성립한다.
\end{corollary}

\oldpage[III]{204}
\begin{proof}
첫 번째 주장은 (7.6.1), 그리고 (7.7.5, c)와 b))에 따라 문제가 $Y$와
$Y'$ 위에서 국소적이라는 사실에서 곧바로 따른다. 두 번째 주장을
증명하려면 (7.7.10)에 따라 모든 $y\in Y$에 대하여
$T_p^{\mathscr O_y}$가 오른쪽 완전(각각 왼쪽 완전)임을 보이면 충분하다.
그런데 가정에 따라 $g(y')=y$인 $y'\in Y'$이 존재하고,
$\mathscr O_{y'}$은 충실 평탄 $\mathscr O_y$-가군이다. 따라서 결론은
가정과 (7.6.1)에서 따른다.
\end{proof}

\begin{remarks}[7.7.12]
\textnormal{(i)} (7.7.4)의 가정 아래 $\mathscr P_\bullet$이 더 나아가
\emph{유한} 복합체라고 하자. 그러면 $X$의 \emph{유한} 아핀 열린 덮개
$\mathfrak U$를 택할 수 있으므로, (7.7.1)에 따라 이중복합체
$C^\bullet(\mathfrak U,
\mathscr P^\bullet\otimes_{\mathscr O_Y}\mathscr M)$도 유한이다. 더
정확히 말하면, $\mathscr M$에 \emph{무관한} 순서쌍들의 유한집합 $E$가
존재하여
\[
  C^h(\mathfrak U,\mathscr P^k\otimes_{\mathscr O_Y}\mathscr M)=0
  \qquad\text{모든 }(h,k)\notin E.
\]
이다. 따라서 어떤 $i_1$이 존재하여 $i\geq i_1$이면 \emph{모든}
준연접 $\mathscr O_Y$-가군 $\mathscr M$에 대하여
$\mathscr T_i(\mathscr M)=0$이다. 특히 이 $i$들에 대하여
$\mathscr T_i$는 자명하게 $\mathscr M$에 관한 완전 함자이므로,
(7.4.1)에 따라 $Z'_i(L_\bullet)$은 유한 생성 \emph{평탄}
$A$-가군이며, 따라서 $A$가 뇌터이므로 유한 생성 \emph{사영}
$A$-가군이다. 이제 $A$-가군의 복합체 $L'_\bullet$을
$i<i_1$이면 $L'_i=L_i$, $L'_{i_1}=Z'_{i_1}(L_\bullet)$,
$i>i_1$이면 $L'_i=0$으로 정의하고,
$\mathscr L'_\bullet=\widetilde{L'_\bullet}$로 놓자. 그러면
\[
  \mathscr H_i(\mathscr L'_\bullet\otimes_{\mathscr O_Y}\mathscr M)
  =
  \mathscr H_i(\mathscr L_\bullet\otimes_{\mathscr O_Y}\mathscr M)
\]
임은 $i<i_1-1$일 때와 $i\geq i_1$일 때 분명하다. 뒤의 경우에는 양변이
영이다. 마지막으로 정의에 따라
$\operatorname{Im}(Z'_{i_1}\otimes_A M)
=\operatorname{Im}(L_{i_1}\otimes_A M)$이므로, $i=i_1-1$일 때에도
$\mathscr H_i(\mathscr L'_\bullet\otimes_{\mathscr O_Y}\mathscr M)
=\mathscr H_i(\mathscr L_\bullet\otimes_{\mathscr O_Y}\mathscr M)$이다.
따라서 $\mathscr L_i$들이 국소 자유 $\mathscr O_Y$-가군, 즉 유한
생성 사영 $A$-가군에 결부된 가군이기만 하면, (7.7.4)에서
$\mathscr L_\bullet$도 \emph{유한} 복합체라고 가정할 수 있다.

이 논증은 특히 $\mathscr P_\bullet$이 차수 $0$의 \emph{단 하나}의 항
$\mathscr F\ne0$으로 축소되는 경우에도 적용된다. 이때
$\mathscr T_n(\mathscr M)=\mathrm R^{-n}f_*
(\mathscr F\otimes_{\mathscr O_Y}\mathscr M)$이다. 그러면 $i>0$일 때
$\mathscr L_i=0$이라고 가정할 수 있다. 이 경우에는 코호몰로지적 표기를
우선 사용하여 $\mathscr T_p$ 대신 $\mathscr T^{-p}$라고 쓰겠다.

\medskip
\textnormal{(ii)} (7.7.5)의 진술에서 $\mathscr P_i$들이 $Y$-평탄이라고
더 이상 가정하지 않을 때에도, 이번에는
\[
  \mathscr T_p(\mathscr M)
  =
  \operatorname{Tor}^Y_p
    (f,1_Y;\mathscr P_\bullet,\mathscr M).
  \tag{7.7.12.1}\label{III.7.7.12.1.fr}
\]
로 놓으면 결론들은 여전히 성립한다. 실제로 (6.7.9)에 따라
$\operatorname{Tor}^Y_n(f,1_Y;\mathscr P_\bullet,\mathscr O_Y)$는 연접
$\mathscr O_Y$-가군이다. (6.10.1)을 고려하면 (6.10.5)의 증명은 바뀌지
않고 적용되며, $Y=\operatorname{Spec}(A)$가 아핀이면 여전히
\[
  \mathscr T_p(\mathscr M)
  =
  \mathscr H_p(\mathscr L_\bullet\otimes_{\mathscr O_Y}\mathscr M),
  \qquad
  \mathscr L_\bullet=\widetilde{L_\bullet},
\]
임을 보인다. 여기서 $L_\bullet$은 유한 생성 자유 $A$-가군들의
복합체이다. 이로써 주장이 증명된다.
\end{remarks}

\subsection{고유 사상에 대한 응용: II. 코호몰로지적 평탄성 판정법}
\label{subsection:III.7.8.fr}

\begin{definition}[7.8.1]
$X$, $Y$를 두 준스킴, $f:X\to Y$를 준콤팩트하고 분리인 사상,
$\mathscr P_\bullet$을 $Y$-평탄한 준연접 $\mathscr O_X$-가군들로 이루어진
복합체, $\mathscr T_\bullet$을 (7.7.1.1)로 정의한 준연접
$\mathscr O_Y$-가군의 호몰로지 함자, $y$를 $Y$의 한 점이라 하자.

\oldpage[III]{205}
$Y$ 안에서 $y$의 어떤 열린 근방 $U$에 대하여
$\mathscr T_p^U=\mathscr T_U^{-p}$가 완전하면, $\mathscr P_\bullet$이
\emph{점 $y$에서 차원 $p$에 관하여 $Y$ 위에 호몰로지적으로 평탄하다}
(또는 \emph{점 $y$에서 차원 $-p$에 관하여 $Y$ 위에 코호몰로지적으로
평탄하다})고 한다. $\mathscr P_\bullet$이 모든 점 $y\in Y$에서 차원
$p$에 관하여 $Y$ 위에 호몰로지적으로 평탄하면,
\emph{차원 $p$에서 $Y$ 위에 호몰로지적으로
평탄하다} (또는 \emph{차원 $-p$에서 $Y$ 위에 코호몰로지적으로
평탄하다})고 한다.

$\mathscr P_\bullet$이 \emph{모든 차원 $p$}에서 $Y$ 위에 (각각 점
$y$에서 $Y$ 위에) 호몰로지적으로 평탄하면, 간단히
$\mathscr P_\bullet$이 \emph{$Y$ 위에 호몰로지적으로 평탄하다}
(각각 \emph{점 $y$에서 $Y$ 위에 호몰로지적으로 평탄하다}), 또는
\emph{$Y$ 위에 코호몰로지적으로 평탄하다} (각각 \emph{점 $y$에서
$Y$ 위에 코호몰로지적으로 평탄하다})고 한다.
\end{definition}

\begin{env}[7.8.2]
정의에 따라 $Y$ 위의 호몰로지적 평탄성이라는 개념은 \emph{$Y$ 위에서
국소적}이다. $Y$가 국소 뇌터이거나 \emph{스킴}이면,
$\mathscr P_\bullet$이 차원 $p$에서 $Y$ 위에 호몰로지적으로 평탄하기
위한 필요충분조건은 함자 $\mathscr T_p$가 \emph{완전}한 것이다.
$Y$가 국소 뇌터인 경우의 증명은 (7.7.5)의 증명 도중에 이미
이루어졌다. $Y$가 스킴인 경우에도 논증은 같다. 곧 준콤팩트 스킴의
아핀 열린집합에 (I, 9.4.2)를 적용한다.
\end{env}

\begin{proposition}[7.8.3]
표기와 가정은 (7.8.1)과 같다고 하자. 다음 조건들은 동치이다.
\begin{enumerate}
  \item[a)] $\mathscr P_\bullet$은 점 $y$에서 차원 $p$에 관하여 $Y$
    위에 호몰로지적으로 평탄하다.
  \item[b)] $Y$ 안에서 $y$의 어떤 열린 근방 $U$에 대하여
    $\mathscr T_p^U$와 $\mathscr T_{p+1}^U$가 오른쪽 완전하다.
  \item[c)] $Y$ 안에서 $y$의 어떤 열린 근방 $U$에 대하여
    $\mathscr T_p^U$와 $\mathscr T_{p-1}^U$가 왼쪽 완전하다.
  \item[d)] $Y$ 안에서 $y$의 어떤 열린 근방 $U$에 대하여
    $\mathscr T_{p+1}^U$는 오른쪽 완전하고 $\mathscr T_{p-1}^U$는
    왼쪽 완전하다.
\end{enumerate}
\end{proposition}

\begin{proof}
$Y$가 아핀일 때의 $\mathscr T_p$의 해석을 고려하면, 이는 (7.3.3)의
일부를 옮겨 쓴 것에 지나지 않는다.
\end{proof}

\begin{proposition}[7.8.4]
$Y$를 국소 뇌터 준스킴, $f:X\to Y$를 고유 사상,
$\mathscr P_\bullet$을 $Y$-평탄한 연접 $\mathscr O_X$-가군들로 이루어진
아래 유계 복합체, $\mathscr T_\bullet$을 (7.7.1.1)로 정의한 함자라 하자.
모든 $y\in Y$에 대하여 다음 조건들은 동치이다.
\begin{enumerate}
  \item[a)] $\mathscr P_\bullet$은 점 $y$에서 차원 $p$에 관하여 $Y$
    위에 호몰로지적으로 평탄하다.
  \item[b)] 함자 $T_p^{\mathscr O_y}$는 완전하다.
  \item[c)] 어떤 정수 $n_0$가 존재하여 $n\geq n_0$이면
    \[
      \operatorname{long}T_p^{\mathscr O_y}
        (\mathscr O_y/\mathfrak m_y^{n+1})
      =
      \operatorname{long}T_p^{\mathscr O_y}(\kappa(y))
      \cdot
      \operatorname{long}(\mathscr O_y/\mathfrak m_y^{n+1})
      \tag{7.8.4.1}\label{III.7.8.4.1.fr}
    \]
    이다. 여기서 길이는 $\mathscr O_y$-가군의 길이이다.
  \item[d)] $y$의 어떤 열린 근방 $U$가 존재하여
    $(\mathscr H^{-p}(f,\mathscr P^\bullet))|U$는
    $(\mathscr O_Y|U)^m$ 꼴의 $(\mathscr O_Y|U)$-가군과 동형이고, 모든
    준연접 $(\mathscr O_Y|U)$-가군 $\mathscr M$에 대하여 표준 준동형
    \[
      ((\mathscr H^{-p}(f,\mathscr P^\bullet))|U)
      \otimes_{\mathscr O_Y|U}\mathscr M
      \longrightarrow
      \mathscr H^{-p}
      \bigl(
        f,(\mathscr P^\bullet|U)
          \otimes_{\mathscr O_Y|U}\mathscr M
      \bigr)
      \tag{7.8.4.2}\label{III.7.8.4.2.fr}
    \]
    은 전단사이다.
\end{enumerate}
이 조건들이 성립하면 다음 성질도 성립한다.
\begin{enumerate}
  \item[e)] $y$의 어떤 근방에서 함수 $z\mapsto d_p(z)$
    ((7.7.5.1)에서 정의됨)는 상수이다.
\end{enumerate}
더 나아가 $Y$가 점 $y$에서 축약이면 ($0_{\mathrm I}$, 4.1.4), e)는
다른 조건들과 동치이다.
\end{proposition}

\oldpage[III]{206}
\begin{proof}
실제로 조건 b)는 $T_p^{\mathscr O_y}$와 $T_{p+1}^{\mathscr O_y}$가
오른쪽 완전하다는 것과 동치이다 (7.3.3). 따라서 a)와 b)의 동치는
(7.7.10)과 (7.8.3)에서 따른다.
$\mathscr O_y/\mathfrak m_y^{n+1}$은 아르틴환이고,
$T_p^{\mathscr O_y}(\mathscr O_y/\mathfrak m_y^{n+1})$와
$T_p^{\mathscr O_y}(\kappa(y))$는 유한 생성
$(\mathscr O_y/\mathfrak m_y^{n+1})$-가군이며 (7.7.4), 따라서 유한
길이이다. 그러므로 b)와 c)의 동치도 (7.7.10)과 (7.3.5.7)에서 따른다.
a)가 e)를 함의하며 $\mathscr O_y$가 축약일 때에는 e)와 동치라는 사실은
(7.6.9)의 결과이다. 끝으로 정의 (7.8.1)과 (7.7.5)에 의하여 a)는
(7.8.4.2)가 전단사임을 함의한다. 또한 a)에 의하여
$(\mathscr H^{-p}(f,\mathscr P^\bullet))_y$는 평탄
$\mathscr O_y$-가군이고 (7.3.3, f)), (7.7.4)에 따라 유한 생성
$\mathscr O_y$-가군이므로
자유이다 ($0_{\mathrm I}$, 10.1.3). 한편
$\mathscr H^{-p}(f,\mathscr P^\bullet)$은 연접 $\mathscr O_Y$-가군이므로
(7.7.4), $y$의 어떤 근방에서 국소 자유이다
($0_{\mathrm I}$, 5.2.7). 역으로 d)가 a)를 함의함은 함자
$\mathscr T_p^U$의 정의 (7.7.2.2)에서 분명하다.
\end{proof}

\begin{proposition}[7.8.5]
(7.8.4)의 가정 아래 다음 조건들은 동치이다.
\begin{enumerate}
  \item[a)] $\mathscr P_\bullet$은 모든 차원 $i\leq p$에서 $Y$ 위에
    호몰로지적으로 평탄하다.
  \item[b)] $i\leq p+1$이면 함자 $\mathscr T_i$는 오른쪽 완전하다.
  \item[c)] $i\geq-p$이면 $\mathscr O_Y$-가군
    $\mathscr H^i(f,\mathscr P^\bullet)$은 국소 자유이다.
\end{enumerate}
\end{proposition}

\begin{proof}
a)와 b)의 동치는 (7.8.3)에 따라 자명하고, a)는 (7.8.4)에 따라 c)를
함의한다. 역으로 c)가 성립한다고 하자. 한편 (7.7.4)에 따라 $i<i_0$이면
$\mathscr L_i=0$이고, 따라서 $i\leq i_0$이면 $\mathscr T_i=0$임을
유의하자. 그러므로 모든 점 $y\in Y$에는 다음 성질을 갖는 아핀 근방
$U=\operatorname{Spec}(A)$가 있다. $i\leq p$이면 $T_i^A(A)$는 자유
$A$-가군이다. 따라서 (7.3.7)에 의하여 $i\leq p$이면
$\mathscr T_i^U=T_i^A$는 완전하다.
\end{proof}

이제 코호몰로지적 평탄성 판정법들을 주로 다음 경우에 적용하겠다.
복합체 $\mathscr P_\bullet$은 $Y$ 위에서 \emph{평탄}한 \emph{단 하나}의
연접 $\mathscr O_X$-가군 $\mathscr F$로 축소되고,
이를 $\mathscr P_0$로 택한다. 이때
$\mathscr T_p(\mathscr M)=\mathrm R^{-p}f_*
(\mathscr F\otimes_{\mathscr O_Y}\mathscr M)$임을 상기하자.

\begin{proposition}[7.8.6]
$Y$를 국소 뇌터 준스킴, $f:X\to Y$를 고유이고 평탄한 사상, $y$를
$Y$의 한 점이라 하자. 올 $f^{-1}(y)=X\otimes_Y\kappa(y)$를 $X_y$라
쓰자. $\Gamma(X_y,\mathscr O_{X_y})=R$가 분리가능 $\kappa(y)$-대수라고
가정하자 (Bourbaki, \emph{Alg.}, chap. VIII, § 7,
n\textsuperscript{o} 5). 달리 말하면 이는 $\kappa(y)$의 유한 차수
분리가능 확대체 유한 개의 직접곱이다. 그러면 $\mathscr O_X$는 점
$y$에서 차원 0에 관하여 $Y$ 위에 코호몰로지적으로 평탄하다.
\end{proposition}

\begin{proof}
(7.8.4)에 의하여 $Y$가 국소환 $A=\mathscr O_y$의 스펙트럼인 경우로
한정할 수 있다. $f$가 평탄하다는 가정은 $\mathscr T_{-1}=0$을
함의하므로, $T_0^A$가 왼쪽 완전함은 이미 알 수 있고 이제 오른쪽
완전함만 보이면 된다. 더 나아가 (7.7.10, c))에 의하여
$A=\mathscr O_y$가 \emph{아르틴환}인 경우로 환원된다.
$R$의 분해체인 $\kappa(y)$의 유한 확대체 $k'$을 택하자. 그러면
$R\otimes_{\kappa(y)}k'$은 $k'$과 동형인 유한 개의 체의 직접곱이다.
$A$에서 어떤 국소환 $A'$으로 가는 국소 준동형이 존재하여, $A'$은
$A$ 위에서 유한 자유인 $A$-대수이고 $A'$의 잉여체는 $k'$과 동형이다
($0_{\mathrm I}$, 10.3.2). (7.6.1)에 의하여 $T_0^{A'}$이 오른쪽
완전함을 증명하는 경우로 환원된다. 달리 말하면 $R$가 $\kappa(y)$와
동형인 $m$개의 체의 직접곱이라고 가정할 수 있다. 이제 다음 초등적인
보조정리를 유의하자.

\begin{lemma}[7.8.6.1]
$Z$를 국소환 달린 공간이라 하자. $Z$가 연결이기 위한 필요충분조건은

\oldpage[III]{207}
환 $\Gamma(Z,\mathscr O_Z)$가 영환이 아닌 두 환의 곱이 아닌 것이다.
\end{lemma}

\begin{proof}
실제로 $Z$가 서로소인 두 비어 있지 않은 열린집합의
합집합이면, $\Gamma(Z,\mathscr O_Z)$는 영환이 아닌 두 환
$\Gamma(Z_1,\mathscr O_Z)$와 $\Gamma(Z_2,\mathscr O_Z)$의 곱과
동형임이 분명하다. 역으로 $\Gamma(Z,\mathscr O_Z)$가 이러한 곱이라는
것은 $\Gamma(Z,\mathscr O_Z)$ 안에 0과 1과 다른 멱등원 $s$가
존재한다는 것과 동치이다.
모든 $z\in Z$에 대하여 $s_z$는 국소환 $\mathscr O_z$의 멱등원이므로
0 또는 1이다. 그런데 $s_z=0$인 $z$들의 집합은 열려 있음이
분명하다. 한편 $s_z=1$이면 정의에 따라 $s(z)\ne0$이므로, $s_z=1$인
$z$들의 집합도 열려 있다 ($0_{\mathrm I}$, 5.5.2). 이로부터 결론이
따른다.
\end{proof}

이 보조정리에 따르면 $X_y$는 정확히 $m$개의 연결 성분 $X'_i$를
가지며, 모든 $i$에 대하여
$\Gamma(X'_i,\mathscr O_{X'_i})=\kappa(y)$이다. $A$가 국소
아르틴환이라고 가정했으므로 그 스펙트럼은 한 점으로 이루어지고,
따라서 $X$와 $X_y$의 바탕공간은 같다. 그러므로 $X$는
$X'_i=X_i\otimes_Y\kappa(y)$인 $m$개의 연결 성분 $X_i$를 갖는다.
이로써 마지막으로 $R=\kappa(y)$인 경우로 환원된다. (7.7.10, b))에
따라 표준 준동형
\[
  \Gamma(X,\mathscr O_X)
  \longrightarrow
  \Gamma(X_y,\mathscr O_{X_y})
\]
이 전사임을 보이면 충분하다. 그런데 합성
\[
  \Gamma(Y,\mathscr O_Y)=A
  \longrightarrow
  \Gamma(X,\mathscr O_X)
  \longrightarrow
  \Gamma(X_y,\mathscr O_{X_y})=\kappa(y)
\]
이 이미 전사이므로 이는 자명하다.
\end{proof}

\begin{corollary}[7.8.7]
(7.8.6)의 가정 아래 다음 조건들을 만족하는 $y$의 열린 근방 $U$가
존재한다.
\begin{enumerate}
  \item[(i)] $f_*(\mathscr O_X)|U$는 $(\mathscr O_Y|U)^m$ 꼴의
    $(\mathscr O_Y|U)$-가군과 동형이다.
  \item[(ii)] 모든 $z\in U$에 대하여 표준 준동형
    \[
      (f_*(\mathscr O_X))_z
      \otimes_{\mathscr O_z}\kappa(z)
      \longrightarrow
      \Gamma(X_z,\mathscr O_{X_z})
    \]
    은 전단사이다.
\end{enumerate}
\end{corollary}

\begin{proof}
(i)은 (7.8.6)과 (7.8.4)에서 따른다.

(ii)는 $U$를 알맞게 택하면 $\mathscr T_0^U$가 완전하다는 사실과
(7.7.5.3)에서 따른다.
\end{proof}

\begin{corollary}[7.8.8]
(7.8.6)의 조건들이 성립하고 더 나아가
$\Gamma(X_y,\mathscr O_{X_y})=\kappa(y)$라고 하자. 그러면 표준
준동형 $\mathscr O_Y|U\to f_*(\mathscr O_X)|U$가 전단사인 $y$의 열린
근방 $U$가 존재한다.
\end{corollary}

\begin{proof}
실제로 (7.8.7, (ii))에 따르면 (7.8.7, (i))에 나타나는 정수 $m$은
반드시 1이다.
\end{proof}

\begin{corollary}[7.8.9]
(7.8.6)의 가정 아래 $y$의 열린 근방 $U$, 연접 $\mathscr O_U$-가군
$\mathscr Q$ (유일한 동형까지 결정됨), 그리고 준연접
$\mathscr O_U$-가군 $\mathscr M$을 변수로 하는 함자들의 동형사상
\[
  \mathrm R^1f_*(f^*(\mathscr M))
  \xrightarrow{\sim}
  \mathscr{H}\!om_{\mathscr O_U}(\mathscr Q,\mathscr M).
  \tag{7.8.9.1}\label{III.7.8.9.1.fr}
\]
이 존재한다.
\end{corollary}

\begin{proof}
실제로 가정에 따라 알맞은 $U$에 대하여 $\mathscr T_0^U$가 완전하다.
따라서 $p=0$인 경우에 (7.7.5, a))와 (7.7.5, b$'$))의 동치를 적용하되,
$\mathscr P_\bullet$은 차수 0의 항이 $\mathscr O_X$인 한 항 복합체로
택하면 충분하다.
\end{proof}

\begin{remarks}[7.8.10]
\textnormal{(i)} (7.8.6)의 조건 아래 $f$의 스테인 인자분해 (4.3.3)
\[
\xymatrix{
  X \ar[r]^{f'} & Y' \ar[r]^{g} & Y .
}
\]

\oldpage[III]{208}
$Y'=\operatorname{Spec}(f_*(\mathscr O_X))$를 생각하자. 이때 유한 사상
$g$는 $g_*(\mathscr O_{Y'})=f_*(\mathscr O_X)$가 $y$의 근방에서 국소
자유이고, $y$에서의 올이 분리가능 $\kappa(y)$-대수의 스펙트럼이라는
성질을 갖는다 (II, 1.5.1). 제IV장에서 이로부터 다음 사실을 이끌어 낼
것이다. $Y$ 안에서 $y$의 열린 근방 $U$가 존재하여 $g$의 제한
$g^{-1}(U)\to U$의 \emph{모든} 올 $g^{-1}(z)$ ($z\in U$)가 분리가능
$\kappa(z)$-대수의 스펙트럼이다 (이를 $U$의 \emph{에탈 덮개}라 하겠다).
그러면 (7.8.7, (ii))에 따라 (7.8.6)에서 점 $y$에 둔 가정은 $y$의 어떤
근방의 모든 점에서도 성립한다.

\medskip
\textnormal{(ii)} 제V장에서 다음 사실을 보겠다. $X$가 $Y$ 위에서
사영이고 (더 나아가 제IV장에서 정의할 성질인 $Y$ 위에서 ``단순''한
경우에도), (7.8.9)의 $\mathscr O_U$-가군 $\mathscr Q$는 반드시 국소
자유인 것은 아니다. 달리 말하면, 이러한 조건 아래에서도
$\mathscr O_X$는 점 $y$에서 \emph{차원 1}에 관하여 $Y$ 위에 반드시
코호몰로지적으로 평탄한 것은 아니다. 제V장에서 $\mathscr Q$를 $Y$에
대한 $X$의 피카르 스킴의 단위원 단면을 따른 1차 미분형식의 층으로
해석할 것이다.
\end{remarks}

\subsection{고유 사상에 대한 응용: III. 오일러--푸앵카레 특성과
힐베르트 다항식의 불변성}
\label{subsection:III.7.9.fr}
\label{III.7.9.fr}

\begin{env}[7.9.1]
\label{III.7.9.1.fr}
$A$를 환, $M$을 유한 생성 사영 $A$-가군이라 하자. $X=\operatorname{Spec}(A)$
위에서 이에 대응하는 $\mathscr O_X$-가군 $\widetilde M$이 \emph{유한
계수의 국소 자유}라는 것과 이는 동치임을 상기하자 (Bourbaki,
\emph{Alg. comm.}, chap. II, \S~5, n\textsuperscript{o}~2).
모든 $\mathfrak p\in\operatorname{Spec}(A)$에 대하여 자유
$A_{\mathfrak p}$-가군 $M_{\mathfrak p}$의 계수를 \emph{$\mathfrak p$에서
$M$의 계수}라 하고 $\operatorname{rang}_{\mathfrak p}(M)$로 나타낸다
(또는 이는 국소 자유 $\mathscr O_X$-가군 $\widetilde M$의
$\mathfrak p$에서의 계수이다). 따라서
\[
  \operatorname{rang}_{\mathfrak p}M
  =\operatorname{rang}_{\mathfrak p}(M_{\mathfrak p})
  =\operatorname{rang}_{\kappa(\mathfrak p)}
  \bigl(M\otimes_A\kappa(\mathfrak p)\bigr).
  \tag{7.9.1.1}\label{III.7.9.1.1.fr}
\]
\end{env}

\begin{proposition}[7.9.2]
\label{III.7.9.2.fr}
$P_\bullet$을 유한 생성 사영 $A$-가군들로 이루어진 유한 복합체라 하고,
모든 $A$-가군 $M$에 대하여
$T_\bullet(M)=H_\bullet(P_\bullet\otimes_A M)$으로 놓자. 그러면 모든
$\mathfrak p\in\operatorname{Spec}(A)$에 대하여
\[
  \sum_i(-1)^i\operatorname{rang}_{\kappa(\mathfrak p)}
  T_i\bigl(\kappa(\mathfrak p)\bigr)
  =\sum_i(-1)^i\operatorname{rang}_{\mathfrak p}(P_i).
  \tag{7.9.2.1}\label{III.7.9.2.1.fr}
\]
\end{proposition}

\begin{proof}
실제로 정의에 따라
$T_i(\kappa(\mathfrak p))=H_i(P_\bullet\otimes_A\kappa(\mathfrak p))$이고,
(7.9.1.1)을 고려하면 공식 (7.9.2.1)은 유한 차원 벡터공간들의 유한
복합체에서 호몰로지로 이행할 때 오일러--푸앵카레 특성이 불변한다는
사실에 지나지 않는다 ($0_{\mathrm I}$, 11.10.2).
\end{proof}

\begin{corollary}[7.9.3]
\label{III.7.9.3.fr}
함수
\[
  \mathfrak p\longmapsto
  \sum_i(-1)^i\operatorname{rang}_{\kappa(\mathfrak p)}
  T_i\bigl(\kappa(\mathfrak p)\bigr)
\]
는 $\operatorname{Spec}(A)$에서 국소 상수이다.
\end{corollary}

\begin{theorem}[7.9.4]
\label{III.7.9.4.fr}
$Y$를 국소 뇌터 준스킴, $f:X\to Y$를 고유 사상,
$\mathscr P_\bullet$을 연접이고 $Y$-평탄한 $\mathscr O_X$-가군들로
이루어진 유한 복합체라 하자.
\[
  \mathscr T_\bullet(\mathscr M)
  =\mathscr H^{-\bullet}
  \bigl(f,\mathscr P^\bullet\otimes_{\mathscr O_Y}\mathscr M\bigr)
  \qquad\text{((7.7.1.1) 참조)}
\]
로 놓으면 함수

\oldpage[III]{209}
\[
  y\longmapsto
  \sum_i(-1)^i\operatorname{rang}_{\kappa(y)}
  \mathscr T_i\bigl(\kappa(y)\bigr)
  \tag{7.9.4.1}\label{III.7.9.4.1.fr}
\]
는 $Y$에서 국소 상수이다.
\end{theorem}

\begin{proof}
$Y=\operatorname{Spec}(A)$가 뇌터 환 $A$의 스펙트럼인 아핀 경우로
한정할 수 있다. 복합체 $\mathscr P_\bullet$은 유한하므로 (7.7.12,
(i))에 따라
\[
  \mathscr T_p(\mathscr M)
  =\mathscr H_p(\mathscr L_\bullet\otimes_{\mathscr O_Y}\mathscr M),
  \qquad \mathscr L_\bullet=\widetilde{L_\bullet},
\]
이다. 여기서 $L_\bullet$은 유한 생성 사영 $A$-가군들로 이루어진
\emph{유한} 복합체이다. 이제 정리는 (7.9.3)에서 따른다.
\end{proof}

\begin{env}[7.9.5]
\label{III.7.9.5.fr}
(7.9.4)의 조건 아래 $Y$가 \emph{연결}이면 함수 (7.9.4.1)은
\emph{상수}이다. $Y$가 연결이고 공집합이 아니면 (7.9.4.1)의 유일한
(정수) 값을
\[
  \operatorname{EP}(f,\mathscr P_\bullet)
  \quad\text{또는}\quad
  \operatorname{EP}(Y,\mathscr P_\bullet),
  \quad\text{또는 단순히}\quad
  \operatorname{EP}(\mathscr P_\bullet)
\]
로 나타내되, 마지막 표기는 혼동할 염려가 없을 때 사용한다. 이 정수를
$f$(또는 $Y$)에 대한 \emph{$\mathscr P_\bullet$의
오일러--푸앵카레 특성}이라 한다. 일반적인 경우에는
\[
  \operatorname{EP}(f,\mathscr P_\bullet;y)
  \quad\text{또는}\quad
  \operatorname{EP}(Y,\mathscr P_\bullet;y)
  \quad\text{또는}\quad
  \operatorname{EP}(\mathscr P_\bullet;y)
\]
로 (7.9.4.1)의 오른쪽 항을 나타내기도 한다.
\end{env}

\begin{env}[7.9.6]
\label{III.7.9.6.fr}
$X$, $Y$, $f$에 관하여 (7.9.4)의 가정을 두고,
\[
  0\longrightarrow\mathscr P'_\bullet
  \xrightarrow{u}\mathscr P_\bullet
  \xrightarrow{v}\mathscr P''_\bullet
  \longrightarrow0
\]
를 연접이고 $Y$-평탄한 $\mathscr O_X$-가군들의 유한 복합체로 이루어진
완전열이라 하자. 준동형 $u$와 $v$의 차수는 각각 \emph{짝수} $2d$와
$2d'$이라고 하자. $\mathscr T_\bullet$은 호몰로지 함자이므로 (7.7.1),
호몰로지 완전열
\[
  \cdots\longrightarrow
  \mathscr T_i(\mathscr P'_\bullet,\kappa(y))
  \longrightarrow
  \mathscr T_{i+2d}(\mathscr P_\bullet,\kappa(y))
  \longrightarrow
  \mathscr T_{i+2d+2d'}(\mathscr P''_\bullet,\kappa(y))
  \longrightarrow
  \mathscr T_{i-1}(\mathscr P'_\bullet,\kappa(y))
  \longrightarrow\cdots
\]
을 얻으며, 이 열은 비영인 항이 유한 개뿐이다. 이 복합체의
오일러--푸앵카레 특성이 0임을 쓰면 ($0_{\mathrm I}$, 11.10.1), 곧바로
\[
  \operatorname{EP}(\mathscr P_\bullet;y)
  =\operatorname{EP}(\mathscr P'_\bullet;y)
  +\operatorname{EP}(\mathscr P''_\bullet;y)
  \tag{7.9.6.1}\label{III.7.9.6.1.fr}
\]
을 모든 $y\in Y$에 대하여 얻는다. 한편 예를 들어
$\mathscr P_\bullet=(\mathscr P_i)$이고 $i<0$이면 $\mathscr P_i=0$인
경우에는 복합체의 완전열
\[
\xymatrix@C=2.4em{
  \cdots \ar[r] & 0 \ar[r]\ar[d] & 0 \ar[r]\ar[d]
    & 0 \ar[r]\ar[d] & 0 \ar[r] & \cdots \\
  \cdots \ar[r] & 0 \ar[r]\ar[d] & 0 \ar[r]\ar[d]
    & \mathscr P_1 \ar[r]\ar[d] & \mathscr P_2 \ar[r] & \cdots \\
  \cdots \ar[r] & 0 \ar[r]\ar[d] & \mathscr P_0 \ar[r]\ar[d]
    & \mathscr P_1 \ar[r]\ar[d] & \mathscr P_2 \ar[r] & \cdots \\
  \cdots \ar[r] & 0 \ar[r]\ar[d] & \mathscr P_0 \ar[r]\ar[d]
    & 0 \ar[r]\ar[d] & 0 \ar[r] & \cdots \\
  \cdots \ar[r] & 0 \ar[r] & 0 \ar[r]
    & 0 \ar[r] & 0 \ar[r] & \cdots
}
\]
을 얻는다. 여기서 영이 아닌 세로 화살표들은 항등 자기동형이다.
(7.9.6.1)을 이 완전열에 적용하고 $\mathscr P_\bullet$의 길이에 관하여
귀납하면 공식
\[
  \operatorname{EP}(\mathscr P_\bullet;y)
  =\sum_i(-1)^i\operatorname{EP}(\mathscr P_i;y).
  \tag{7.9.6.2}\label{III.7.9.6.2.fr}
\]
을 얻는다.

\oldpage[III]{210}
여기서 $Y$ 위에 평탄한 모든 연접 $\mathscr O_X$-가군 $\mathscr F$에
대하여 $\operatorname{EP}(\mathscr F;y)$ (또는
$\operatorname{EP}(f,\mathscr F;y)$ 또는
$\operatorname{EP}(Y,\mathscr F;y)$)는 $\mathscr Q_\bullet$의 함수
$\operatorname{EP}(\mathscr Q_\bullet;y)$를 뜻한다. 이 복합체에서
$\ne0$인 유일한 항은 차수 $0$이고 $\mathscr F$와 같다. 따라서 한
항으로 축소된 복합체의 오일러--푸앵카레 특성만 연구하는 경우로
환원할 수 있다.
\end{env}

\begin{proposition}[7.9.7]
\label{III.7.9.7.fr}
(7.9.4)의 가정 아래 $Y'$을 국소 뇌터 준스킴, $g:Y'\to Y$를 사상이라
하고,
\[
  X'=X\times_Y Y',
  \qquad f'=f_{(Y')}:X'\longrightarrow Y',
\]
으로 놓자. $\mathscr P'_\bullet$을 $\mathscr O_{X'}$-가군들의 유한
복합체
\[
  \mathscr P_\bullet\otimes_{\mathscr O_Y}\mathscr O_{Y'}
\]
라 하자. 그러면 $\mathscr P'_\bullet$은 연접이고 $Y'$-평탄한
$\mathscr O_{X'}$-가군들로 이루어지고, 모든 $y'\in Y'$에 대하여
\[
  \operatorname{EP}(\mathscr P'_\bullet;y')
  =\operatorname{EP}(\mathscr P_\bullet;g(y')).
  \tag{7.9.7.1}\label{III.7.9.7.1.fr}
\]
이다.
\end{proposition}

\begin{proof}
$\mathscr O_{X'}$-가군 $\mathscr P'_i$들은 사영 $X'\to X$에 의한
$\mathscr P_i$들의 역상이므로 연접이다. 또한 문제는 $X$, $Y$, $Y'$에서
국소적이므로 ($0_{\mathrm I}$, 6.2.1)과 (I, 4.14.5)에 따라
$Y'$-평탄하다. 끝으로 $f'$은 고유임이 알려져 있으므로 (II, 5.4.2),
(7.9.7.1)의 왼쪽 항은 정의된다. 이제 언제나 할 수 있듯 $Y$와 $Y'$이
아핀인 경우로 환원하면, 공식 (7.9.7.1)은 (6.10.4.2), (7.7.2)와 보조정리
(7.6.7)에서 따른다.
\end{proof}
\begin{proposition}[7.9.8]
\label{III.7.9.8.fr}
(7.9.4)의 가정이 성립한다고 하고, 또한 어떤 정수 $i_0$가 존재하여
\[
  \mathscr T_i(\kappa(y))=0
  \qquad\text{$i\ne i_0$이고 모든 $y\in Y$에 대하여}
\]
라고 하자. 그러면
\[
  \mathscr T_{i_0}(\mathscr O_Y)
  =\mathscr H^{-i_0}(f,\mathscr P_\bullet)
\]
는 국소 자유 $\mathscr O_Y$-가군이고, $y\in Y$에서의 계수는
$(-1)^{i_0}\operatorname{EP}(f,\mathscr P_\bullet;y)$와 같다.
\end{proposition}

\begin{proof}
먼저 $\mathscr T_j^{\mathscr O_y}$들이 (7.4.4)의 가정을 만족하므로
(7.4.7)을 적용할 수 있음을 주목하자. 가정과 (7.5.3)에 따라
$i\ne i_0$일 때 $\mathscr T_i^{\mathscr O_y}$는 \emph{영}이다. 따라서
(7.3.3)에 따라 $\mathscr T_{i_0}$도 완전하고, 이어 (7.8.4)에 따라
$\mathscr H^{-i_0}(f,\mathscr P_\bullet)$은 국소 자유이며 한 점
$y\in Y$에서의 계수는
\[
  \operatorname{rang}_{\kappa(y)}
  \mathscr T_{i_0}(\kappa(y))
  =(-1)^{i_0}\operatorname{EP}(f,\mathscr P_\bullet;y)
\]
와 같다. 실제로 $i\ne i_0$일 때 $\mathscr T_i(\kappa(y))=0$이므로 이는
정의에서 따른다.
\end{proof}

\begin{corollary}[7.9.9]
\label{III.7.9.9.fr}
$Y$를 국소 뇌터 준스킴, $f:X\to Y$를 고유 사상, $\mathscr F$를
연접이고 $Y$-평탄한 $\mathscr O_X$-가군이라 하자. 어떤 정수 $i_0$가
존재하여
\[
  H^i\bigl(f^{-1}(y),
  \mathscr F\otimes_{\mathscr O_Y}\kappa(y)\bigr)=0
  \qquad\text{모든 $i\ne i_0$와 모든 $y\in Y$에 대하여}
\]
라고 가정하자. 그러면 $\mathrm R^{i_0}f_*(\mathscr F)$은 국소 자유
$\mathscr O_Y$-가군이고, $y$에서의 계수는
$(-1)^{i_0}\operatorname{EP}(f,\mathscr F;y)$와 같다.
\end{corollary}

특히 다음을 얻는다.

\begin{corollary}[7.9.10]
\label{III.7.9.10.fr}
$X$, $Y$, $\mathscr F$에 관하여 (7.9.9)의 예비 조건을 두고, 모든
$i>0$에 대하여 $\mathrm R^if_*(\mathscr F)=0$이라고 가정하자. 그러면
$f_*(\mathscr F)$는 국소 자유 $\mathscr O_Y$-가군이고, $y$에서의
계수는 $\operatorname{EP}(f,\mathscr F;y)$와 같다.
\end{corollary}

\begin{proof}
(7.9.9)에 따라 다음 보조정리를 증명하면 충분하다.
\end{proof}

\begin{lemma}[7.9.10.1]
\label{III.7.9.10.1.fr}
(7.9.10)의 가정 아래
\[
  H^i\bigl(f^{-1}(y),
  \mathscr F\otimes_{\mathscr O_Y}\kappa(y)\bigr)=0
\]
이 모든 $i>0$과 모든 $y\in Y$에 대하여 성립한다.
\end{lemma}

\begin{proof}
실제로 $Y=\operatorname{Spec}(A)$가 아핀인 경우로 한정할 수 있다.
(7.9.4)의 표기를 쓰고 $\mathscr P_\bullet$이 차수 $0$의 항
$\mathscr F$로 축소된 경우, 가정에 따라 $p<0$이면
$\mathscr T_p(\mathscr O_Y)=0$이다. (7.3.7)에 따라 $p<0$이면
$\mathscr T_p$가 완전하고, 이제 보조정리는 (7.7.5, a))와 (7.7.5, d))의
동치에서 따른다.
\end{proof}


\oldpage[III]{211}
\begin{proposition}[7.9.11]
\label{III.7.9.11.fr}
(7.9.4)의 가정을 두고, $\mathscr L$을 $Y$에 대하여 매우 풍부한 가역
$\mathscr O_X$-가군이라 하며, 모든 $n\in\mathbf Z$에 대하여
$\mathscr P_\bullet(n)=\mathscr P_\bullet\otimes_{\mathscr O_X}
\mathscr L^{\otimes n}$으로 놓자. 그러면 모든 $y\in Y$에 대하여 함수
\[
  n\longmapsto\operatorname{EP}(f,\mathscr P_\bullet(n);y)
  \tag{7.9.11.1}\label{III.7.9.11.1.fr}
\]
는 $\mathbf Q$-계수 다항식이고, $Y$의 같은 연결 성분에 속하는 모든
점에서 동일하다.
\end{proposition}

\begin{proof}
$\mathscr P_\bullet(n)$이 $Y$-평탄한 $\mathscr O_X$-가군들의
복합체임은 자명하다. (7.9.6.2)에 따라 $\mathscr P_\bullet$이 차수
0의 단 하나의 항 $\mathscr F\ne0$으로 축소된 경우로 한정할 수
있다. 또한 문제는 $Y$에서 국소적이므로 $Y$가 아핀이고 $f$가 사영
사상이라고 가정할 수 있다 (II, 5.5.3). $X_y=f^{-1}(y)$로 놓고,
$\mathscr L_y=\mathscr L\otimes_{\mathscr O_Y}\kappa(y)$라 하자. 이는 매우
풍부한 $\mathscr O_{X_y}$-가군이다 (II, 4.4.10).
(7.7.2)에 따라 복합체 $\mathscr P_\bullet(n)$에 관한 함자
$\mathscr T_\bullet$에 대하여
\[
  \mathscr T_i(\kappa(y))
  =H^{-i}(X_y,\mathscr F_y\otimes\mathscr L_y^{\otimes n})
  \qquad
  \bigl(\text{여기서 }\mathscr F_y=\mathscr F\otimes_{\mathscr O_Y}\kappa(y)\bigr)
\]
이다. 따라서 $\operatorname{EP}(f,\mathscr F(n);y)$는 (2.5.1)에서
정의한 오일러--푸앵카레 특성
$\chi_{\kappa(y)}(\mathscr F_y(n))$에 지나지 않는다. 이제 (7.9.11.1)이
다항식이라는 사실은 (2.5.3)에서 따른다. 또한 각 $n$에 대하여 그 값은
$Y$의 한 연결 성분에서 상수이므로 (7.9.4), 증명이 끝난다.
\end{proof}

\begin{env}[7.9.12]
\label{III.7.9.12.fr}
(7.9.11.1)의 유리수 계수 다항식을
$\operatorname{PH}(f,\mathscr P_\bullet;y)$ 또는
$\operatorname{PH}(\mathscr P_\bullet;y)$로 나타내고, 이를
\emph{$\mathscr P_\bullet$, $f$, $\mathscr L$에 대한 $y$에서의
힐베르트 다항식}이라 한다 (또는 혼동할 염려가 없으면 단순히
\emph{$\mathscr P_\bullet$의 $y$에서의 힐베르트 다항식}, 혹은 $f$의
힐베르트 다항식이라 한다). $Y$가 연결이고 공집합이 아니면 표기와
용어에서 $y$를 생략한다. 이렇게 얻은 불변량은 주어진 연접층의 연접
몫층들에 관한 “모듈라이 이론”에서 제V장의 핵심적 역할을 할 것이다.

(7.9.6)과 (7.9.11)의 표기 아래
\[
  \operatorname{PH}(\mathscr P_\bullet;y)
  =\operatorname{PH}(\mathscr P'_\bullet;y)
  +\operatorname{PH}(\mathscr P''_\bullet;y)
  \tag{7.9.12.1}\label{III.7.9.12.1.fr}
\]
이고, 특히
\[
  \operatorname{PH}(\mathscr P_\bullet;y)
  =\sum_i(-1)^i\operatorname{PH}(\mathscr P_i;y) ;
  \tag{7.9.12.2}\label{III.7.9.12.2.fr}
\]
이는 (7.9.6.1)과 (7.9.6.2)에서 자명하게 따른다. 마찬가지로 (7.9.7)의
표기와 가정 아래
\[
  \operatorname{PH}(\mathscr P'_\bullet;y')
  =\operatorname{PH}(\mathscr P_\bullet;g(y')).
  \tag{7.9.12.3}\label{III.7.9.12.3.fr}
\]
이다.
\end{env}

공식 (7.9.12.2)는 복합체의 힐베르트 다항식에 관한 연구를 하나의
$Y$-평탄한 $\mathscr O_X$-가군의 힐베르트 다항식에 관한 연구로
환원한다. 후자는 호몰로지적 고찰과 무관한 주목할 만한 해석을 갖는다.

\begin{corollary}[7.9.13]
\label{III.7.9.13.fr}
$Y$를 뇌터 준스킴, $f:X\to Y$를 고유 사상, $\mathscr L$을 $Y$에
대하여 매우 풍부한 가역 $\mathscr O_X$-가군, $\mathscr F$를 연접이고
$Y$-평탄한 $\mathscr O_X$-가군이라 하자. 어떤 정수 $n_0$가 존재하여
$n\geq n_0$이면 $f_*(\mathscr F(n))$은 국소 자유
$\mathscr O_Y$-가군이고, $y\in Y$에서의 계수는
$\operatorname{PH}(f,\mathscr F;y)(n)$과 같다.
\end{corollary}

\begin{proof}
사상 $f$는 사영 사상이므로 (II, 5.5.3), 어떤 $n_0$가 존재하여
$n\geq n_0$이면

\oldpage[III]{212}
모든 $i>0$에 대하여 $\mathrm R^if_*(\mathscr F(n))=0$이다 (2.2.1).
따라서 결론은 (7.9.10)에서 따른다.
\end{proof}


다음 평탄성 판정법은 제V장의 연접층 모듈라이 이론에서 중요할 것이다.

\begin{proposition}[7.9.14]
\label{III.7.9.14.fr}
$Y$를 뇌터 준스킴, $f:X\to Y$를 사영 사상, $\mathscr L$을 $f$에 대해
풍부한 가역 $\mathscr O_X$-가군이라 하고, 모든 $\mathscr O_X$-가군
$\mathscr F$와 모든 $n\in\mathbf Z$에 대하여
$\mathscr F(n)=\mathscr F\otimes\mathscr L^{\otimes n}$으로 놓자. 연접
$\mathscr O_X$-가군 $\mathscr F$가 $Y$-평탄이기 위한 필요충분조건은
어떤 정수 $n_0$가 존재하여 모든 $n\geq n_0$에 대해
$f_*(\mathscr F(n))$이 국소 자유 $\mathscr O_Y$-가군인 것이다.
\end{proposition}

\begin{proof}
조건의 필요성은 (7.9.13)과 같이 증명된다 ($f$가 사영 사상이므로
(2.2.1)의 결과를 풍부한 층 $\mathscr L$에 적용할 수 있다). 역을
증명하기 위해 $Y$가 환 $A$의 스펙트럼인 아핀 경우로 한정할 수 있다.
가정과 (2.2.2, (i))에 따라 $A$-가군
$\Gamma(X,\mathscr F(n))$들은 유한 생성 사영 가군이다 (Bourbaki,
\emph{Alg. comm.}, chap. II, \S~5, n\textsuperscript{o}~2, th. 1).
등급환 $S$를 생각하자.
\[
  S=\bigoplus_{n\geq0}\Gamma(X,\mathscr L^{\otimes n})
\]
$X$는 $\operatorname{Proj}(S)$와 표준적으로 동일시됨이
알려져 있다 (II, 4.5.2, (b) 및 5.4.4). 이제
\[
  M=\bigoplus_{n\geq n_0}\Gamma(X,\mathscr F(n))
\]
으로 놓자. 필요하면 $\mathscr L$을 거듭제곱
$\mathscr L^{\otimes d}$로 바꾸어, $S$가 차수 1의 유한 개 원소로
생성된다고 가정할 수 있다 (2.3.5.1). 그러면 (II, 2.7.5 및 2.7.2)에
따라 $\mathscr F$는 $\mathscr P\!roj_0(M)$과 동일시된다. 차수가
$>0$인 모든 동차 원소 $g\in S$에 대하여
$\Gamma(X_g,\mathscr F)=M_{(g)}$이다. 그런데 $M$은 사영
$A$-가군들의 직합이므로 평탄 $A$-가군이고, 따라서 $M_g$도 평탄하다
($0_{\mathrm I}$, 6.3.2). 나아가 $M_{(g)}$는 $M_g$의 직합인자이므로
역시 평탄하다 ($0_{\mathrm I}$, 6.1.2). 그러므로 (I, 4.14.5)에 따라
$\mathscr F$는 $X_g$의 모든 점에서 $Y$-평탄하다. $X_g$들이 $X$를
덮으므로 명제가 증명되었다.
\end{proof}

\bigskip
\noindent\emph{(계속.)}


\clearpage
\oldpage[III]{213}
\section*{표기 색인}

\begin{flushleft}
$\mathbf{Tor}_n^A(P_\bullet,Q_\bullet)$ : 6.3.1.

\medskip
$\mathscr{T}\!or_n^{\mathscr O_S}(\mathscr P_\bullet,\mathscr Q_\bullet)$,
$\mathscr{T}\!or_n^S(\mathscr P_\bullet,\mathscr Q_\bullet)$ : 6.4.1 및 6.5.1.

\medskip
$\mathscr{T}\!or_n^S(\mathscr F,\mathscr G)$ : 6.5.1.

\medskip
$\mathscr{T}\!or_q^S(\mathscr P_\bullet^{(1)},\mathscr P_\bullet^{(2)},\ldots,
\mathscr P_\bullet^{(m)})$ : 6.5.15.

\medskip
$\mathbf H^{-n}(\mathfrak U^{(i)},\mathscr P_\bullet^{(i)})$,
$\mathbf H^{-n}(X^{(i)},\mathscr P_\bullet^{(i)})$ : 6.6.1.

\medskip
$\mathbf{Tor}_n^A(L_{\bullet\bullet}^{(1)},L_{\bullet\bullet}^{(2)})$,
$\mathbf{Tor}_n^S(\mathfrak U^{(1)},\mathfrak U^{(2)};
\mathscr P_\bullet^{(1)},\mathscr P_\bullet^{(2)})$ : 6.6.2.

\medskip
$\mathscr{T}\!or_n^S(\mathfrak U^{(1)},\mathfrak U^{(2)};
\mathscr F^{(1)},\mathscr F^{(2)})$,
$\mathscr{T}\!or_n^S(\mathfrak U^{(1)},\mathfrak U^{(2)};
\mathscr P_\bullet^{(1)},\mathscr P_\bullet^{(2)})$ : 6.6.2.

\medskip
$\mathscr{T}\!or_n^S(f_1,f_2;\mathscr P_\bullet^{(1)},
\mathscr P_\bullet^{(2)})$ : 6.6.8, 6.7.1 및 6.7.3.

\medskip
$\mathscr{T}\!or_n^S((f_i)_{i\in I};(\mathscr P_\bullet^{(i)})_{i\in I})$,
$\mathscr{T}\!or_n^S(f_1,\ldots,f_m;\mathscr P_\bullet^{(1)},\ldots,
\mathscr P_\bullet^{(m)})$ : 6.7.11.

\medskip
$\mathbf{Ab}$, $\mathbf{Ab}_A$ : 7.1.1.

\medskip
$T_{(B)}$, $T^{(B)}$, $T\otimes_A B$
($T$는 $\mathbf{Ab}_A$에서 $\mathbf{Ab}$로 가는 함자) : 7.1.3.

\medskip
$t_M$ : 7.2.2.

\medskip
$T_{\mathfrak p}$ : 7.1.4.

\medskip
$\mathscr T_\bullet^{Y'}(\mathscr M')$,
$\mathscr T_\bullet^U(\mathscr M|U)$ : 7.7.2.

\medskip
$\operatorname{EP}(f,\mathscr P_\bullet;y)$,
$\operatorname{EP}(Y,\mathscr P_\bullet;y)$,
$\operatorname{EP}(\mathscr P_\bullet;y)$ : 7.9.5.

\medskip
$\operatorname{PH}(f,\mathscr P_\bullet;y)$,
$\operatorname{PH}(\mathscr P_\bullet;y)$ : 7.9.12.
\end{flushleft}

\clearpage
\oldpage[III]{214}
\section*{용어 색인}

\begin{flushleft}
가군들의 복합체의 오일러--푸앵카레 특성 : 7.9.5.

\medskip
점 $y\in Y$에서 차원 $p$에 관하여 코호몰로지적으로 평탄,
차원 $p$에서 $Y$ 위에 코호몰로지적으로 평탄, 점 $y$에서 $Y$ 위에
코호몰로지적으로 평탄, $Y$ 위에 코호몰로지적으로 평탄
($Y$-평탄한 $\mathscr O_X$-가군들의 복합체) : 7.8.1.

\medskip
호모토픽한 복합체들 : 6.1.4.

\medskip
코호몰로지 차원이 $\leq n$인 환 달린 공간 : 6.5.5.

\medskip
공변 가법 함자에서의 스칼라 확대 : 7.1.3.

\medskip
코호몰로지 차원이 $\leq n$인 환의 층 : 6.5.5.

\medskip
퀸네트 공식 : 6.7.8.

\medskip
점 $y\in Y$에서 차원 $p$에 관하여 호몰로지적으로 평탄,
차원 $p$에서 $Y$ 위에 호몰로지적으로 평탄, 점 $y$에서 $Y$ 위에
호몰로지적으로 평탄, $Y$ 위에 호몰로지적으로 평탄
($Y$-평탄한 $\mathscr O_X$-가군들의 복합체) : 7.8.1.

\medskip
호모토피 : 6.1.4.

\medskip
두 $A$-가군 복합체의 초Tor : 6.3.1.

\medskip
두 가군 복합체의 국소 초Tor : 6.4.1.

\medskip
두 가군 복합체의 전역 초Tor : 6.6.2, 6.6.8 및 6.7.3.

\medskip
가군들의 복합체에 대한 상대 힐베르트 다항식 : 7.9.12.

\medskip
퀸네트 스펙트럼 열들 : 6.7.3.
\end{flushleft}


\clearpage
\oldpage[III]{215}
\section*{목차}

\begin{flushleft}
\textsc{제III장.} --- \textbf{연접층의 코호몰로지적 연구 (계속)} \dotfill 5

\medskip
\S~6. 국소 및 전역 Tor 함자; 퀸네트 공식 \dotfill 5

6.1. 서론 \dotfill 5

6.2. 준스킴 위의 가군 복합체의 초코호몰로지 \dotfill 6

6.3. 두 가군 복합체의 초Tor \dotfill 9

6.4. 준연접 가군 복합체의 국소 초Tor 함자: 아핀 스킴의 경우 \dotfill 14

6.5. 준연접 가군 복합체의 국소 초Tor 함자: 일반적인 경우 \dotfill 16

6.6. 준연접 가군 복합체의 전역 초Tor 함자와 퀸네트 스펙트럼 열:
아핀 밑의 경우 \dotfill 21

6.7. 준연접 가군 복합체의 전역 초Tor 함자와 퀸네트 스펙트럼 열:
일반적인 경우 \dotfill 25

6.8. 전역 초Tor의 결합법칙 스펙트럼 열 \dotfill 32

6.9. 전역 초Tor의 밑변경 스펙트럼 열 \dotfill 34

6.10. 일부 코호몰로지 함자의 국소 구조 \dotfill 39

\medskip
\S~7. 가군의 공변 호몰로지 함자에서의 밑변경 연구 \dotfill 43

7.1. $A$-가군의 함자 \dotfill 43

7.2. 텐서곱 함자의 특징짓기 \dotfill 44

7.3. 가군의 호몰로지 함자에 대한 완전성 판정 \dotfill 48

7.4. 함자 $\mathrm H_i(P\otimes_A M)$의 완전성 판정 \dotfill 53

7.5. 뇌터 국소환의 경우 \dotfill 58

7.6. 완전성 성질들의 내림. 반연속성 정리와 Grauert의 완전성
판정법 \dotfill 60

7.7. 고유 사상에 대한 응용: I. 교환 성질 \dotfill 65

7.8. 고유 사상에 대한 응용: II. 코호몰로지적 평탄성 판정법 \dotfill 72

7.9. 고유 사상에 대한 응용: III. 오일러--푸앵카레 특성과
힐베르트 다항식의 불변성 \dotfill 76

\begin{flushright}
\emph{1962년 12월 15일 접수.}
\end{flushright}
\end{flushleft}

% 인쇄된 216쪽은 백지이다.
\clearpage
\thispagestyle{empty}
\null
\clearpage


\oldpage[III]{217}
\section*{정오 및 추가}
\begin{center}
(목록 2)

\medskip
A) \emph{오자}
\end{center}

\noindent\textbf{$(0_{\mathrm I},\,3.2.1)$} 27쪽 2행에서 $X$를 $U$로
바꾼다.

\noindent\textbf{$(0_{\mathrm I},\,3.2.6)$} 27쪽 아래에서 1행과 28쪽
1행에서 (세 곳의) $\mathfrak S_\lambda$를 $\mathscr H_\lambda$로
바꾼다.

\noindent\textbf{$(0_{\mathrm I},\,3.4.5)$} 3.4.5) 앞에 여는 괄호를
추가한다.

\noindent\textbf{$(0_{\mathrm I},\,4.2.1)$} 39쪽 10행에서
$\psi_*(A)$를 $\psi_*(\mathscr A)$로 바꾼다.

\noindent\textbf{$(0_{\mathrm I},\,5.1.3)$} 45쪽 16행에서
$\mathscr F|V$를 $\mathscr F|U$로 바꾼다.

\noindent\textbf{$(0_{\mathrm I},\,5.3.9)$} 48쪽 12행의
« telle que » 앞에 다음을 추가한다: 점 $x$의 어떤 열린 근방 $U$
위에서.

\noindent\textbf{$(0_{\mathrm I},\,7.6.9)$} 73쪽 아래에서 3행의
$I_\lambda$를 $J_\lambda$로 바꾼다.

\noindent\textbf{(I, 1.1.15)} 83쪽 3행에서 $\mathfrak a$를
$\mathfrak a\ne A$로 바꾼다.

\noindent\textbf{(I, 1.4.1)} 90쪽 아래에서 2행의 $\widetilde N$을
$N$으로 바꾼다.

\noindent\textbf{(I, 2.3.1)} 101쪽 아래에서 3행의 (0, 4.1.6)을
(0, 4.1.7)로 바꾼다.

\noindent\textbf{(I, 2.3.2)} 101쪽 11행에서 $B$를 $B_s$로, $C$를
$C_t$로 바꾼다.

\noindent\textbf{(I, 2.5.5)} 104쪽 19행에서 $S$-사상을
$S$-준스킴으로 바꾼다.

\noindent\textbf{(I, 3.3.7)} 109쪽 17행에서 $Y_{(S)}$를
$Y_{(S')}$로 바꾼다.

\noindent\textbf{(I, 3.4.5)} 113쪽 9행에서 $\mathrm s$를 $s$로
바꾼다.

\noindent\textbf{(I, 3.4.8)} 114쪽 4행에서 $p(x)$를 $f(x)$로
바꾼다.

\noindent\textbf{(I, 3.5.1)} 115쪽 3행에서 $1_X\times g$를
$1_{X'}\times g$로 바꾼다.

\noindent\textbf{(I, 3.7.2)} 119쪽 8행에서 첫 번째 $X$를 $X'$로
바꾼다.

\noindent\textbf{(I, 4.2.2)} 123쪽 8행의 도식에서 왼쪽 세로 화살표의
$\alpha_{\psi(y)}$를 $\varphi_{\psi(y)}$로 바꾼다.

\noindent\textbf{(I, 4.4.3)} 126쪽 16행에서 isomorphée를
isomorphes로 바꾼다.

\noindent\textbf{(I, 4.5.5)} 127쪽 아래에서 17행의 (4.2.4)를
(4.2.5)로 바꾼다.

\noindent\textbf{(I, 5.1.4)} 128쪽 아래에서 14행의 (2.1.7)을
(2.1.8)로 바꾼다.

\noindent\textbf{(I, 5.3.13)} 134쪽 20행에서 (4.2.4)를 (4.2.5)로
바꾼다.

\noindent\textbf{(I, 6.4.2)} 147쪽 아래에서 1행의 덮개를 유한 덮개로
바꾼다.

\noindent\textbf{(I, 9.1.13)} 171쪽 15행에서
$p^{-1}(\mathscr F)$를 $p^*(\mathscr F)$로 바꾼다.

\noindent\textbf{(I, 9.5.11)} 179쪽 7행에서 $Y$를 $Y'$로 바꾼다.

\noindent\textbf{(I, 9.6.5)} 180쪽 아래에서 17행의
$(0,\,4.1.4)$를 $(0,\,4.1.3)$으로 바꾼다.

\noindent\textbf{(I, 10.12.2)} 206쪽 아래에서 14행의
$(X_n)_{(S_n)}$을 $(X_n)_{(S_m)}$으로 바꾼다.

\clearpage
\oldpage[III]{218}
\noindent\textbf{(I), 용어 색인:} 219쪽 아래에서 6행의 0, 4.1.4를
0, 4.1.5로 바꾼다. 222쪽 아래에서 16, 17, 18행의 I, 2.1.7을
I, 2.1.8로 바꾸고, 222쪽 아래에서 19행의 I, 2.1.8을 I, 2.1.7로
바꾼다.

\noindent\textbf{(II, 1.5.2)} 12쪽 17행의 도식에서 왼쪽 세로 화살표
옆에 $f$를 하나 삽입한다.

\noindent\textbf{(II, 4.2.7)} 75쪽 아래에서 16행의 (III, 2.1.14)를
(III, 2.1.13)으로 바꾼다.

\noindent\textbf{(II, 5.2.1)} 97쪽 아래에서 16행의 $X-U$를 $U$로
바꾼다. 97쪽 아래에서 8행의 $X'$를 $X_f$로 바꾼다.

\noindent\textbf{(II, 5.3.6)} 100쪽 16행에서 (4.6.18)을 (4.6.17)로
바꾼다.

\noindent\textbf{(II, 6.2.7)} 116쪽 아래에서 9행의 제V장을 제IV장으로
바꾼다.

\noindent\textbf{(II, 6.6.5)} 132쪽 아래에서 15행의 제V장을
제IV장으로 바꾼다.

\noindent\textbf{(II, 7.4.12)} 151쪽 1행에서 제V장을 제III장으로
바꾼다.

\noindent\textbf{(II, 8.1.4)} 154쪽 15행에서 (III, 2.3.8)을
(III, 2.3.7)로 바꾼다.

\noindent\textbf{(II, 8.10.5)} 187쪽 아래에서 17행의
$g_{(j)x}$를 $g_{j(x)}$로 바꾼다.


\noindent\textbf{$(0_{\mathrm{III}},\,10.2.7)$} 20쪽 1행에서
$u$-아딕을 $\mathfrak m$-아딕으로 바꾼다. 20쪽 아래에서 13행의
$k$를 $k_i$로 바꾼다.

\noindent\textbf{$(0_{\mathrm{III}},\,11.1.1)$} 23쪽 18행에서
$\operatorname{Im}(d_r^{p+r,q-r+1})$을
$\operatorname{Im}(d_r^{p-r,q+r-1})$로 바꾼다.

\noindent\textbf{$(0_{\mathrm{III}},\,11.8.6)$} 45쪽 아래에서 14행의
화살표 $\to$를 $\rightsquigarrow$로 바꾼다 (네 곳).

\noindent\textbf{$(0_{\mathrm{III}},\,12.3.4)$} 61쪽 4행에서
$\rightsquigarrow$를 $\to$로 바꾼다.

\noindent\textbf{$(0_{\mathrm{III}},\,13.2.4)$} 67쪽 아래에서 2행의
$\varinjlim$을 $\varprojlim$으로 바꾼다 (두 곳).

\noindent\textbf{(III, 1.1.7)} 85쪽 15행에서 $A^r$을
$\bigwedge(A^r)$로 바꾼다 (두 곳).

\noindent\textbf{(III, 1.2.4)} 87쪽 7행에서
$\Gamma(U,\mathscr F)$를 $\Gamma(U,\mathscr O_X)$로 바꾸고,
$\mathscr G$를 $\mathscr F$로 바꾼다. 87쪽 8행에서
$\mathrm H^p(\mathfrak U,\mathscr G)$를
$\mathrm H^p(\mathfrak U,\mathscr F)$로 바꾼다.

\noindent\textbf{(III, 2.5.3.1)} 111쪽 3행에서 $n>0$을
$n\geq0$으로 바꾼다. 111쪽 8행에서 $-r\leq n\leq0$을
$-r\leq n<0$으로 바꾼다.

\noindent\textbf{(III, 3.1.2)} 116쪽 13행에서 $\mathscr G$를
$\mathscr G\in\mathbf K'$로 바꾼다.

\noindent\textbf{(III, 4.3.6)} 133쪽 7행에서
$K'=K\otimes_A\widehat A$를 $K'\supset K\otimes_A\widehat A$로
바꾼다. 133쪽 16행에서 $R(X)$를 $R(Y)$로 바꾼다.

\noindent\textbf{(III, 4.4.12)} 138쪽 10행에서 제V장을 제IV장으로
바꾼다.

\noindent\textbf{(III, 4.6.6)} 142쪽 19행에서 « 제V장 »을
« 뒤의 한 절에서 »로 바꾼다.

\begin{center}
B) \emph{본문 수정}\footnote{참조를 쉽게 찾을 수 있도록, 이제부터
제N장에 삽입된 정오표 목록에 속하는 수정은 $\mathbf{Err}_{\mathrm N}$
뒤에 번호를 붙여 나타낸다.}
\end{center}

\noindent\textbf{$(\mathbf{Err}_{\mathrm{III}},\,1)$}
$(0_{\mathrm I},\,5.1.3)$의 45쪽 18행에서 « 직합 »을 « 유한 직합 »으로
바꾼다.

\noindent\textbf{$(\mathbf{Err}_{\mathrm{III}},\,2)$}
$(0_{\mathrm I},\,5.4.3)$의 49쪽 아래에서 8행부터 4행까지
« 일반적인 경우에는\ldots{} » 이후를 다음 본문으로 바꾼다.

일반적으로 $X$가 모든
$x\in X$에 대하여 $\mathscr O_x$가 \emph{국소환}인 환 달린 공간이고,
\oldpage[III]{219}
$\mathscr L$이 어떤 $\mathscr O_X$-가군 $\mathscr F$에 대하여
$\mathscr L\otimes_{\mathscr O_X}\mathscr F$가
$\mathscr O_X$와 \emph{동형}이 되는 \emph{유한 생성}
$\mathscr O_X$-가군이면, 앞의 논증은 모든 $x\in X$에 대하여
$\mathscr L_x$가 $\mathscr O_x$와 동형임을 보여 준다. 여기서
$\mathscr L$이 \emph{가역}임을 얻는다. 실제로 모든 $x\in X$에 대하여
$U$를 $x$의 열린 근방으로 택하여 $\mathscr L|U$가 $U$ 위의
$n$개 절단 $s_i\ (1\leq i\leq n)$로 생성되게 하자 (5.2.1). 예컨대
$s=s_1$이고 $s_x\ne0$이라고 가정할 수 있다. $\mathscr L_x$가
$\mathscr O_x$와 동형이므로 각 $i$에 대하여 $x$의 어떤 열린 근방
$V_i\subset U$ 위의 $\mathscr O_X$의 절단 $t_i$가 존재하여
$(t_i)_x s_x=(s_i)_x$이다. 따라서 $x$의 어떤 열린 근방
$V\subset\bigcap_{i=1}^nV_i$를 택할 수 있고, $V$에서
$t_i s=s_i$이다. 다시 말해
$\mathscr L|V$는 \emph{단 하나의} 절단 $s|V$로 생성된다. 또한
$s$가 정의하는 준동형 $\mathscr O_X|V\to\mathscr L|V$의 핵 (5.1.1)의
한 절단 $z$를 열린집합 $W\subset V$ 위에서 택하면, 모든 $y\in W$에
대하여 $z_y$는 $\mathscr L_y$를 소멸시킨다. 따라서 가정에 의해
$z_y=0$이고, 그 결과 $z=0$이다. 이로써 주장이 증명된다. 더욱이
텐서곱 $\mathscr L^{-1}\otimes\mathscr L\otimes\mathscr F$를
고려하면 $\mathscr F$가 $\mathscr L^{-1}$과 동형임을 곧바로 알 수 있다.

\noindent\textbf{$(\mathbf{Err}_{\mathrm{III}},\,3)$}
$(0_{\mathrm I},\,7.2.4)$에서 명제가 성립하려면 허용환 $A$ 안의
아이디얼 $\mathfrak J^n$들이 \emph{닫혀 있다}는 조건을 부과해야 한다.
주어진 증명은 옳지 않으며, 수정된 명제는 Bourbaki, \emph{Top. gén.},
chap. III, 3\textsuperscript{e} éd., \S~3, n\textsuperscript{o}~5,
prop. 9의 cor. 1에서 따른다. 마찬가지로 $(0_{\mathrm I},\,7.2.5)$와
$(0_{\mathrm I},\,7.2.6)$의 명제들에서도 $\mathfrak J^n$들이 $A$ 안에서
\emph{닫혀 있다}고 가정해야 한다.

\noindent\textbf{$(\mathbf{Err}_{\mathrm{III}},\,4)$}
(I, 2.2.5) 뒤에 다음을 덧붙인다. (I, 2.2.5)의 표기에서
$\mathfrak J$가 $A$의 아이디얼이면, $(0,\,4.3.5)$에서 정의한
부분-$\mathscr O_X$-가군 $\mathfrak J\widetilde{\mathscr F}$를
$\widetilde{\mathfrak J\mathscr F}$로 나타낸다.

\noindent\textbf{$(\mathbf{Err}_{\mathrm{III}},\,5)$}
(I, 3.4.5)의 112쪽 아래에서 1행의 « \emph{체} $K$ »를
« \emph{대수적으로 닫힌 체} $K$ »로 바꾼다. 113쪽 1행의
« \emph{확대체} »를 « \emph{대수적으로 닫힌 확대체} »로 바꾼다.
113쪽 1행부터 4행까지 « $K$를 \ldots{} » 이후를 다음 본문으로 바꾼다.

체 $K$에 값을 갖는 $X$의 임의의 점에 대하여 $K$를 그에 대응하는
점의 \emph{값체}라 하고, $x$가 이 점의 소재점이면 이 점이
\emph{$x$에 위치한다}고도 한다. 이로써 $K$에 값을 갖는 점을 그
소재점에 대응시키는 사상 $X(K)\to X$를 정의한다.

113쪽 7행과 16행의 « 기하적 »을 삭제하고, 113쪽 10행과 13행의
« 기하적 »을 삭제한다. 115쪽 아래에서 10행과 11행의
« \emph{기하적} »을 삭제한다.

\noindent\textbf{$(\mathbf{Err}_{\mathrm{III}},\,6)$}
(I, 3.6.1)의 증명 끝, 117쪽 아래에서 12행에 다음을 덧붙인다.
$X'=X\times_Y\operatorname{Spec}(\mathscr O_y/\mathfrak a_y)$라 놓으면,
$p$에 의해 $X$의 점과 동일시되는 모든 점 $x\in X'$에 대하여
$\mathscr O_{X',x}=\mathscr O_{X,x}/\mathfrak a_y\mathscr O_{X,x}$이다.
실제로 문제는 $X$와 $Y$에 관하여 국소적이므로
$X=\operatorname{Spec}(B)$, $Y=\operatorname{Spec}(A)$라고 가정할 수 있고,
평탄성 $(0,\,1.3.2)$에 의해
$(B/\mathfrak a_yB)_x=B_x/\mathfrak a_yB_x$이다.

\noindent\textbf{$(\mathbf{Err}_{\mathrm{III}},\,7)$}
(I, 5.1.1)의 127쪽 아래에서 3행의 « 준연접 $\mathscr O_X$-가군 »을
« $\mathscr B$의 \emph{준연접 아이디얼} »로 바꾼다.

\noindent\textbf{$(\mathbf{Err}_{\mathrm{III}},\,8)$}
(I, 5.1.10) 뒤, 131쪽 16행에 다음을 덧붙인다.

\emph{주석 (5.1.11).} --- (5.1.9)의 논증을 약간 바꾸면, $X$가
준스킴이라고 미리 가정하지 않고 다음 조건만을 가정해도 결론이 여전히
성립함을 알 수 있다. $(X,\mathscr O_X)$가 국소환 달린 공간이고,
$\mathscr J$가 $\mathscr J^n=0$인 $\mathscr O_X$의 아이디얼이며,

\oldpage[III]{220}
환 달린 공간 $X_0=(X,\mathscr O_X/\mathscr J)$가 아핀 스킴이고, 끝으로
아이디얼 $\mathscr J^k/\mathscr J^{k+1}$들이 준연접
$\mathscr O_{X_0}$-가군이라고 가정한다. 실제로 그 논증에서 (2.2.4)
대신 (1.8.1) (cf. $(\mathbf{Err}_{\mathrm{II}})$)을 사용하면 충분하다.

\noindent\textbf{$(\mathbf{Err}_{\mathrm{III}},\,9)$}
(I, 5.3.1)의 132쪽 11행에서 « 또는 $\Delta_X$ » 뒤에
« 또는 $\varphi:X\to S$가 구조 사상이면 $\Delta_\varphi$ »를 삽입한다.

\noindent\textbf{$(\mathbf{Err}_{\mathrm{III}},\,10)$}
(I, 5.3.9)의 증명은 충분하지 않다. $\Delta_X(X)$가
$X\times_SX$ 안에서 국소적으로 닫혀 있음을 증명하지 않기 때문이다.
올바른 증명을 얻으려면 (4.2.4, $a$))를 사용하면 충분하다. 모든
$x\in X$와 $X$에서 $x$의 모든 아핀 근방 $U$에 대하여
$U\times_SU$는 $\Delta_X(x)$의 아핀 근방이다. (그 증명에서는 정의
(5.3.1.1)만을 사용하는) (5.3.16)을 고려하면,
$S=\operatorname{Spec}(B)$와 $X=\operatorname{Spec}(A)$가 아핀 스킴인
경우의 (5.3.9)를 증명하는 것으로 환원된다. 이때 (5.3.1.1)에 의해
$\Delta_X$는 $x\otimes y$를 $xy$로 보내는 표준 준동형
$A\otimes_BA\to A$에 대응함이 분명하다. 이 준동형은 전사이므로,
이 경우 $\Delta_X$는 닫힌 몰입 사상이다 (4.2.3). 이로써 증명이
완료된다.

\noindent\textbf{$(\mathbf{Err}_{\mathrm{III}},\,11)$}
(I, 7.1.15)와 (I, 7.1.16)의 158쪽 아래에서 7행과 15행의
« \emph{기하적} »을 삭제한다.

\noindent\textbf{$(\mathbf{Err}_{\mathrm{III}},\,12)$}
(I, 7.4.7)을 다음으로 바꾼다. $X$가 \emph{축소} 준스킴이고 모든 점이
기약 성분을 \emph{유한} 개만 갖는 열린 근방을 가질 때에도
(관용적으로) (7.4.1)의 정의들을 확장한다. 이때 (7.3.4)와 (7.4.6)에
따르면, 유한형 준연접 $\mathscr O_X$-가군 $\mathscr F$에 대하여
$\mathscr F$가 \emph{꼬임층}이라는 것은
$\operatorname{Supp}(\mathscr F)$가 \emph{$X$의 어떤 기약 성분도
포함하지 않는다}는 것과 동치이다.

\noindent\textbf{$(\mathbf{Err}_{\mathrm{III}},\,13)$}
(I, 10.11.7)의 205쪽 17행부터 23행에서 « 이제 증명할 것은\ldots{} »
이후의 증명 끝부분은 불필요하다. 고려하는 층들이 $(0,\,5.3.5)$에 의해
연접이므로 (10.11.6)에서 바로 따르기 때문이다.

\noindent\textbf{$(\mathbf{Err}_{\mathrm{III}},\,14)$}
(II, 1.7.8)의 16쪽 아래에서 10행 뒤에 다음을 덧붙인다.
$\mathscr E=\mathscr O_S^n$이면 $\mathbf V(\mathscr E)$ 대신
$\mathbf V_S^n$이라고도 쓴다. 더 나아가 $S=\operatorname{Spec}(A)$가
아핀이면 $\mathbf V_S^n$ 대신 $\mathbf V_A^n$이라고 쓴다. 이때
$T_i$들이 부정원인 경우
$\mathbf V_A^n=\operatorname{Spec}(A[T_1,\ldots,T_n])$이다.

\noindent\textbf{$(\mathbf{Err}_{\mathrm{III}},\,15)$}
(II, 1.7.10)의 17쪽 10행에서 « \emph{기하적} »을 삭제하고,
17쪽 12행과 16--17행에서 « 기하적 »을 삭제한다.

\noindent\textbf{$(\mathbf{Err}_{\mathrm{III}},\,16)$}
(II, 1.7.12)에 다음을 덧붙인다. 마찬가지로 다음과 같이 놓는다.
\[
  \mathbf V(\mathscr O_S^n)=\mathbf V_S^n=S[T_1,\ldots,T_n].
\]

\noindent\textbf{$(\mathbf{Err}_{\mathrm{III}},\,17)$}
(II, 4.2.6)의 75쪽 7행에서 « \emph{기하적} »을 삭제하고,
75쪽 8행에서 « \emph{기하적} »을 삭제한다.

\noindent\textbf{$(\mathbf{Err}_{\mathrm{III}},\,18)$}
(II, 4.6.13)의 92쪽 7행에서는 논증을 더 정밀하게 해야 한다. (4.4.10)의
표기에서 (i bis)를 적용하려면 여기서 몰입 사상 $\Gamma_f$가
\emph{준콤팩트}임을 증명해야 하기 때문이다. (4.6.4)에 따라
$\mathscr L$이 $f$에 대해 풍부함을 증명할 때에는 $Y$가 아핀인 경우로
한정할 수 있다. 한편 (v)의 각 가정에서 $f$는 준콤팩트이다
(I, 6.6.4). 또한 $g$가 분리 사상이면 $\Gamma_f$는 닫힌 몰입
사상이므로 (I, 5.4.3) 준콤팩트이다 (I, 6.6.4). 반대로 $X$가 국소
뇌터이면, $Y$가 아핀이고 $f$가 준콤팩트이므로 $X$의 바탕공간은

\oldpage[III]{221}
준콤팩트이고 따라서 뇌터이다. 그러므로 다시 (I, 6.6.4)를 적용하여
$\Gamma_f$가 준콤팩트임을 증명할 수 있다.

\noindent\textbf{$(\mathbf{Err}_{\mathrm{III}},\,19)$}
(II, 5.2.2)의 명제는 옳지 않다. $Y$의 열린집합 $U$에 대하여
$f^{-1}(U)$ 위의 준연접 $(\mathscr O_X|f^{-1}(U))$-가군이 어떤 준연접
$\mathscr O_X$-가군의 $f^{-1}(U)$로의 제한인지 알 수 없을 때에는 조건
$b)$, $c)$, $c')$가 $Y$ 위에서 국소적이지 않기 때문이다. 따라서
98쪽 16행의 « \emph{$f:X\to Y$를 분리 준콤팩트 사상이라 하자} »를
다음으로 바꾸어야 한다. « \emph{$X$, $Y$를 두 준스킴이라 하자.
$X$가 스킴이거나 $X$의 바탕공간이 국소 뇌터라고 가정하고,
$f:X\to Y$를 준콤팩트 사상이라 하자.} » 증명에서는 $X$가 스킴이라고
가정한 경우 그 가정들이 (I, 5.5.5와 5.5.8)에 의해 $f$가
\emph{분리 사상}임을 함의함에 유의한다. 따라서 두 경우 모두
(I, 9.2.2)를 적용할 수 있다.

\noindent\textbf{$(\mathbf{Err}_{\mathrm{III}},\,20)$}
(II, 6.2.3)의 115쪽 아래에서 18행 뒤에 다음을 덧붙인다.

\emph{$y=f(x)$의 아핀 열린 근방 $V$와 $x$의 아핀 열린 근방 $U$가
존재하여 $f(U)\subset V$이고 $f$의 제한인 사상 $U\to V$가 준유한이면,
사상 $f:X\to Y$가 점 $x\in X$에서 준유한이라고 한다. 사상
$f:X\to Y$가 $X$의 모든 점에서 준유한이면 $f$가 국소 준유한이라고
한다.}

\noindent\textbf{$(\mathbf{Err}_{\mathrm{III}},\,21)$}
(II, 6.4.3)의 증명은 옳지 않다. $K$의 유한형 부분-$A$-가군의 원소들이
반드시 $A$ 위에서 정수적인 것은 아니기 때문이다. 121쪽의 마지막
9행을 다음 본문으로 바꾼다. $E$가 꼬임이 없어서 충실하다고 가정할 수
있으므로, $E$는 충실한 $A[u]$-가군이기도 하다 ($A$는 $E$의
자기준동형환에 매입되어 있다). $E$가 유한형 $A$-가군이므로, 여기서
$u$는 $A$ 위에서 \emph{정수적}임을 얻는다 (Bourbaki,
\emph{Alg. comm.}, chap. V, \S~1, n\textsuperscript{o}~1, lemme 1).
따라서 $u\otimes1$의 고윳값들은 ($K$의 한 대수적 폐포 안에서) $A$
위에서 \emph{정수적인} 원소들이고, 그 결과 $\sigma_i(u)$들도 그러함을
곧바로 알 수 있다.

\noindent\textbf{$(\mathbf{Err}_{\mathrm{III}},\,22)$}
$(0_{\mathrm{III}},\,9.1.1)$의 12쪽 아래에서 18행의 « 열린 »을
삭제한다.

\noindent\textbf{$(\mathbf{Err}_{\mathrm{III}},\,23)$}
$(0_{\mathrm{III}},\,11.7.3)$의 43쪽 10행에서 $\mathsf C''$ 뒤에
다음을 덧붙인다. $\mathsf C'$의 모든 사영 대상 $P$에 대하여 (각각
$\mathsf C''$의 모든 사영 대상 $P'$에 대하여) 함자
$A'\rightsquigarrow T(P,A')$ (각각 $A\rightsquigarrow T(A,P')$)가
$\mathsf C'$에서 (각각 $\mathsf C$에서) 완전하다고 가정한다.

\noindent\textbf{$(\mathbf{Err}_{\mathrm{III}},\,24)$}
$(0_{\mathrm{III}},\,13.7.7)$의 명제는 정확하지 않으며 다음과 같이
고쳐야 한다. 78쪽 6행에서 « 및
$(\mathrm R^{n+1}T(A_k))_{k\in\mathbf Z}$ »를 삭제한다. 78쪽 7행의
« $\mathrm R'^nT(A)$는 유한형 $S$-가군이다 »를
« $\mathrm R'^nT(A)$와 $\mathrm R'^{n+1}T(A)$는 유한형 $S$-가군이다 »로
바꾼다. 78쪽 12--13행에서 « 및 $p+q=n+1$ »을 삭제한다. 증명의
78쪽 23행에서 « 및 $n+1$에 대하여 »를 삭제한다.

\noindent\textbf{$(\mathbf{Err}_{\mathrm{III}},\,25)$}
(III, 1.4.15)의 92쪽 아래에서 6행의 « \emph{유한형} »을
« \emph{준콤팩트} »로 바꾼다.

\noindent\textbf{$(\mathbf{Err}_{\mathrm{III}},\,26)$}
(III, 2.2.4)의 102쪽 1행에서 « \emph{$\mathscr F$와 $\mathscr H$의
지지집합들이 $Y$ 위에서 고유라고 가정한다} »를
« \emph{$\operatorname{Im}(\mathscr F\to\mathscr G)$의 지지집합이
$Y$ 위에서 고유라고 가정한다} »로 바꾼다. 증명에서는 10행부터
16행까지 « 그러면\ldots{} » 이후의 본문을 다음으로 바꾼다.

\oldpage[III]{222}
$\mathscr G_1=\operatorname{Im}(\mathscr F\to\mathscr G)
=\operatorname{Ker}(\mathscr G\to\mathscr H)$라 놓자. 이는 연접이다.
완전열
$0\to\mathscr G_1(n)\to\mathscr G(n)\to\mathscr H(n)$이 있고 함자
$f_*$가 왼쪽 완전이므로, 모든 $n$에 대하여 열
$0\to f_*(\mathscr G_1(n))\to f_*(\mathscr G(n))
\to f_*(\mathscr H(n))$이 완전하다. 따라서 $n$이 충분히 클 때 열
$f_*(\mathscr F(n))\to f_*(\mathscr G_1(n))\to0$이 완전함을 보이면
충분하다. 다시 말해 $\mathscr H=0$이고 $\mathscr G$의 지지집합이
$Y$ 위에서 고유인 경우로 한정할 수 있다. 따라서 그 지지집합은
\emph{$Z$에서 닫혀 있다} (II, 5.4.10). 그러므로
$\mathscr G'=i_*(\mathscr G)$는 $\mathscr G'|X=\mathscr G$인 연접
$\mathscr O_Z$-가군이다. 또한 $\mathscr F'|X=\mathscr F$인 연접
$\mathscr O_Z$-가군 $\mathscr F'$이 존재함이 알려져 있다 (I, 9.4.3).
더 나아가 $X$가 $Z$에서 열려 있고
$\operatorname{Supp}(\mathscr G)\subset X$가 $Z$에서 닫혀 있으므로,
$u'|X$가 주어진 전사 준동형 $u:\mathscr F\to\mathscr G$가 되도록 하는
층들의 전사 준동형 $u':\mathscr F'\to\mathscr G'$을 다음과 같이
정의할 수 있음이 명백하다. $\operatorname{Supp}(\mathscr G)$와 만나지
않는 $Z$의 모든 열린집합 $U$에서는 $u'|U=0$으로 놓고, 모든 열린집합
$U\subset X$에서는 $u'|U=u|U$로 놓는다. 이제 (2.2.2)의 증명 첫머리와
같이 $Y$가 아핀인 경우로 한정할 수 있다. 따라서 $n$이 충분히 클 때
준동형
$\Gamma(u):\Gamma(X,\mathscr F(n))\to\Gamma(X,\mathscr G(n))$이
전사임을 보이면 된다. 그런데 다음 가환도식이 있다.
\[
\xymatrix@C=35mm@R=12mm{
\Gamma(Z,\mathscr F'(n)) \ar[r]^{\Gamma(u')} \ar[d]_v &
\Gamma(Z,\mathscr G'(n)) \ar[d]^w \\
\Gamma(X,\mathscr F(n)) \ar[r]_{\Gamma(u)} & \Gamma(X,\mathscr G(n))
}
\]
여기서 $v$와 $w$는 제한 준동형들이다. 또한 $\mathscr G'$의 정의에
따르면 $w$는 \emph{전단사}이다. 한편 (2.2.3)에 따라 $n$이 충분히
클 때 $\Gamma(u')$는 \emph{전사}이고, 따라서 $\Gamma(u)$도 전사이다.

\noindent\textbf{$(\mathbf{Err}_{\mathrm{III}},\,27)$}
(III, 2.2.5)의 102쪽 아래에서 7행 뒤에 다음을 덧붙인다.

(iii) $X$, $Y$, $f$, $\mathscr L$에 관한 (2.2.4)의 가정 아래에서,
$\mathscr H$가 지지집합이 \emph{$Y$ 위에서 고유한} 연접
$\mathscr O_X$-가군이면 (2.2.4)와 비슷한 논증으로 어떤 정수 $N$이
존재하여 $n\geq N$일 때 모든 $i>0$에 대하여
$\mathrm R^if_*(\mathscr H(n))=0$임을 곧바로 알 수 있다. 코호몰로지
완전열로부터 다음을 얻는다. $u:\mathscr F\to\mathscr G$가 연접
$\mathscr O_X$-가군들의 준동형이고 $\operatorname{Ker}(u)$와
$\operatorname{Coker}(u)$의 지지집합들이 \emph{$Y$ 위에서 고유}하면,
어떤 $N$이 존재하여 $n\geq N$일 때 모든 $i>0$에 대하여 대응하는
준동형
$\mathrm R^if_*(\mathscr F(n))\to\mathrm R^if_*(\mathscr G(n))$이
\emph{전단사}이다.

\noindent\textbf{$(\mathbf{Err}_{\mathrm{III}},\,28)$}
(III, 4.4.9)의 137쪽 아래에서 17행의 « \emph{$Y$를 국소 뇌터
정역 준스킴이라 하자} »를 « \emph{$X$와 $Y$를 두 국소 뇌터 정역
준스킴이라 하자} »로 바꾼다. 137쪽 아래에서 16행의
« \emph{유한형} »을 « \emph{국소 유한형} »으로 바꾼다. 증명은
본질적으로 바뀌지 않는다. $f^{-1}(y)$가 $k(y)$ 위에서 국소 유한형이고,
따라서 여전히 이산이기 때문이다.

\noindent\textbf{$(\mathbf{Err}_{\mathrm{III}},\,29)$}
(I, 9.3.4)의 173쪽 아래에서 19행의 « \emph{뇌터} »를
« \emph{준콤팩트} »로 바꾸고, 173쪽 아래에서 18행과 19행의
« \emph{연접} »을 « \emph{유한형 준연접} »으로 바꾼다. 증명에서는

\oldpage[III]{223}
$X=\operatorname{Spec}(A)$, $\mathscr F=\widetilde M$이고 $M$이 유한형
$A$-가군이며, $\mathscr J=\widetilde{\mathfrak J}$이고
$\mathfrak J$가 $A$의 유한형 아이디얼인 경우로 환원한다. 그러면
나머지 논증은 바뀌지 않는다.

\noindent\textbf{$(\mathbf{Err}_{\mathrm{III}},\,30)$}
(I, 9.3.5)의 173쪽 아래에서 1행부터 7행까지를 다음 본문으로 바꾼다.

\emph{명제 (9.3.5). --- $X$를 준스킴, $\mathscr F$를 유한형 준연접
$\mathscr O_X$-가군이라 하자. 그러면 바탕공간이
$\operatorname{Supp}(\mathscr F)$와 같은 $X$의 닫힌 부분준스킴 $Y$와
유한형 준연접 $\mathscr O_Y$-가군 $\mathscr G$가 존재하여, 표준 포함
사상을 $j:Y\to X$라 하면 $\mathscr F$는 $j_*(\mathscr G)$와 동형이다.}

$\mathscr F$의 \emph{소멸자}인 $\mathscr O_X$의 아이디얼 $\mathscr J$가
\emph{준연접}임을 보이면 충분하다. 그러면 $Y$를 $\mathscr J$로 정의되는
$X$의 닫힌 부분준스킴으로 택하고 (I, 4.1.2),
$\mathscr J\mathscr F=0$이므로 $\mathscr F$는
$(\mathscr O_X/\mathscr J)$-가군이며 $\mathscr G=j^*(\mathscr F)$로
택하면 된다. $\mathscr J$가 준연접임을 보이기 위해, 이 문제는
국소적이므로 $X=\operatorname{Spec}(A)$,
$\mathscr F=\widetilde M$이고 $M$이 유한 개의 원소
$x_i\ (1\leq i\leq r)$로 생성되는 $A$-가군인 경우로 한정할 수 있다.
그러면 아이디얼 $\mathscr J$는 $x_i$들의 소멸자들의 교집합이다. 그런데
$x_i$의 소멸자는 $A$에서 $M$으로 가는 준동형 $s\mapsto sx_i$에
대응하는 준동형 $\mathscr O_X\to\mathscr F$의 핵이다. 따라서 이는
준연접 아이디얼이고 (I, 4.1.1), 이러한 아이디얼들의 모든 유한
교집합도 준연접 아이디얼이다 (I, 1.3.10).
